\documentclass[a4paper,12pt]{report}

% --------------------------------------------------
% Codifica e lingua
% --------------------------------------------------
\usepackage[T1]{fontenc}
\usepackage[utf8]{inputenc}
\usepackage[italian]{babel}
\usepackage{lmodern}

% --------------------------------------------------
% Impaginazione e tipografia
% --------------------------------------------------
\usepackage[a4paper,margin=2.5cm]{geometry}
\usepackage{microtype}
\usepackage{setspace}
\usepackage{parskip}

% --------------------------------------------------
% Matematica
% --------------------------------------------------
\usepackage{amsmath,amssymb,mathtools}
\usepackage{amsthm}
\usepackage{bm}
\usepackage{physics}

% --------------------------------------------------
% Figure, grafici, colori
% --------------------------------------------------
\usepackage{graphicx}
\usepackage{xcolor}
\usepackage{float}
\usepackage{caption}
\usepackage{subcaption}
\usepackage{booktabs}
\usepackage{array}
\usepackage{multirow}

% --------------------------------------------------
% TikZ e PGFPlots per diagrammi, blocchi, grafici
% --------------------------------------------------
\usepackage{tikz}
\usepackage{pgfplots}
\pgfplotsset{compat=1.18}

\usetikzlibrary{
    arrows.meta,
    calc,
    positioning,
    fit,
    matrix,
    backgrounds,
    decorations.pathreplacing,
    shapes.geometric,
    shapes.symbols,
    shapes.misc,
    patterns
}

% Eventuali circuiti / schemi tecnici
\usepackage{circuitikz}

% --------------------------------------------------
% Link ipertestuali
% --------------------------------------------------
\usepackage[
    colorlinks=true,
    linkcolor=blue!50!black,
    urlcolor=blue!50!black,
    citecolor=blue!50!black
]{hyperref}

% --------------------------------------------------
% Comandi utili
% --------------------------------------------------
\newcommand{\R}{\mathbb{R}}
\newcommand{\C}{\mathbb{C}}

% --------------------------------------------------
% Ambiente teoremi (facoltativo ma utile)
% --------------------------------------------------
\newtheorem{teorema}{Teorema}[chapter]
\newtheorem{proposizione}[teorema]{Proposizione}
\newtheorem{lemma}[teorema]{Lemma}
\theoremstyle{definition}
\newtheorem{definizione}[teorema]{Definizione}
\newtheorem{osservazione}[teorema]{Osservazione}
\newtheorem{esempio}[teorema]{Esempio}

\begin{document}

% --------------------------------------------------
% Frontespizio
% --------------------------------------------------
\begin{titlepage}
    \centering
    \vspace*{2cm}

    {\Huge\bfseries Teoria dei Sistemi \par}
    \vspace{0.8cm}

    {\Large Prof.ssa Pallottino \par}
    \vspace{0.5cm}

    {\large Università di Pisa\par}
    \vspace{0.3cm}

    {\large Ingegneria Robotica e dell'Automazione\par}

    \vfill

    {\Large\bfseries LaTex rielaborato - Appunti Ingegneri Robotizzati\par}
    \vspace{0.8cm}

    {\large\textbf{Materiale a cura di:}\par}
    \vspace{0.2cm}
    {\large Marcello, Alessio, Giovanni, Gabriele, Kevin\par}

    \vfill
\end{titlepage}

% --------------------------------------------------
% Disclaimer
% --------------------------------------------------
\chapter*{Disclaimer}
\addcontentsline{toc}{chapter}{Disclaimer}

Questo documento è stato realizzato a partire dagli appunti presi in aula durante le lezioni del corso di \emph{Teoria dei Sistemi e del Controllo}.  
Il materiale è stato successivamente rielaborato, ordinato e ampliato con l'obiettivo di renderlo più chiaro, coerente e utile allo studio.

Il testo ha esclusivamente finalità didattiche e di supporto personale allo studio. Non sostituisce in alcun modo le lezioni, il materiale ufficiale del corso, né i testi consigliati dai docenti.

Eventuali imprecisioni, refusi o semplificazioni sono possibili e vanno sempre confrontati con le dispense ufficiali e con quanto svolto a lezione.

Questo documento non è pensato per la vendita, la distribuzione commerciale o qualunque altra forma di utilizzo a scopo di lucro.

\tableofcontents

% --------------------------------------------------
% Capitoli
% --------------------------------------------------
\chapter{L1 -- Sistemi dinamici: ingredienti, funzioni di ingresso, stato, uscite e classificazione}

\section{Ingredienti di un sistema dinamico}

Per descrivere un sistema dinamico servono alcuni insiemi fondamentali.

\subsection{Insieme dei tempi}

Si indica con
\[
T
\]
l'insieme dei tempi. Nei casi più comuni:
\[
T=\mathbb{R} \qquad \text{(tempo continuo)},
\]
oppure
\[
T=\mathbb{Z} \qquad \text{(tempo discreto)}.
\]

L'ipotesi essenziale è che $T$ sia un insieme \emph{ordinato}: deve avere senso stabilire quale istante venga prima e quale dopo.

\subsection{Valori istantanei di ingresso, uscita e stato}

Si introducono poi tre insiemi:
\[
U, \qquad Y, \qquad X,
\]
dove:
\begin{itemize}
    \item $U$ è l'insieme dei possibili valori istantanei dell'ingresso;
    \item $Y$ è l'insieme dei possibili valori istantanei dell'uscita;
    \item $X$ è l'insieme dei possibili stati del sistema.
\end{itemize}

\paragraph{Osservazione.}
Qui stiamo parlando dei \emph{valori istantanei}.  
Per esempio, dire che
\[
u(t)\in U
\]
significa che, fissato un certo istante $t$, il valore assunto dall'ingresso appartiene all'insieme $U$.

\subsection{Funzioni di ingresso e di uscita}

Oltre agli insiemi dei valori istantanei, servono gli insiemi delle \emph{funzioni nel tempo}:

\[
\mathcal{U} \qquad \text{insieme delle funzioni di ingresso},
\]
\[
\mathcal{Y} \qquad \text{insieme delle funzioni di uscita}.
\]

Quindi:
\[
u(\cdot)\in \mathcal{U}, \qquad y(\cdot)\in \mathcal{Y},
\]
mentre per ogni istante $t$ vale:
\[
u(t)\in U, \qquad y(t)\in Y.
\]

La distinzione è molto importante:
\begin{itemize}
    \item $U$ e $Y$ raccolgono i valori istantanei;
    \item $\mathcal{U}$ e $\mathcal{Y}$ raccolgono le storie temporali complete.
\end{itemize}

\section{Funzioni di ingresso ammissibili}

Nello studio dei sistemi non si considera un insieme arbitrario di segnali, ma una classe di ingressi \emph{ammissibili}, cioè fisicamente o matematicamente utilizzabili.  
Gli appunti della lezione evidenziano tre proprietà naturali di $\mathcal{U}$.

\subsection{Esistenza di un istante iniziale}

Per ogni ingresso ammissibile $u(\cdot)\in \mathcal{U}$ si assume che esista un istante iniziale $t_0\in T$ a partire dal quale il segnale sia effettivamente definito.

In modo informale: il segnale \emph{inizia} a un certo istante; non è necessario assegnarlo per tempi arbitrariamente precedenti.

Una formulazione pratica consiste nel pensare che ogni ingresso abbia dominio del tipo
\[
[t_0,+\infty)\cap T.
\]

\paragraph{Interpretazione.}
Questa ipotesi è coerente con l'idea fisica di esperimento: scelgo un istante iniziale e da lì in poi applico l'ingresso al sistema.

\subsection{Chiusura per troncamento}

Se un ingresso $u(\cdot)$ è ammissibile, allora deve essere ammissibile anche il suo \emph{troncamento} in un istante successivo $t_1$.

Definiamo il segnale troncato:
\[
u^{[t_1]}(t)=u(t), \qquad t\ge t_1.
\]

L'ipotesi di chiusura per troncamento richiede che:
\[
u(\cdot)\in\mathcal{U}
\quad \Longrightarrow \quad
u^{[t_1]}(\cdot)\in\mathcal{U}.
\]

\paragraph{Significato.}
Se un ingresso è ammissibile, allora posso sempre decidere di ``cominciare a guardarlo'' da un certo istante in poi, senza uscire dalla classe degli ingressi ammissibili.

\begin{figure}[H]
\centering
\begin{tikzpicture}[scale=0.95]
    \draw[->] (0,0) -- (9,0) node[right] {$t$};
    \draw[->] (0,0) -- (0,3.2) node[above] {$u(t)$};

    \draw[dashed] (2,0) -- (2,2.6) node[above] {$t_0$};
    \draw[dashed] (5.2,0) -- (5.2,2.6) node[above] {$t_1$};

    \draw[thick,blue]
        plot[smooth] coordinates {(2,0.6) (2.6,1.8) (3.1,1.5) (3.8,2.0) (4.5,1.7) (5.2,1.9) (6.1,1.6) (7.0,1.9) (8.1,2.5)};

    \draw[very thick,red]
        plot[smooth] coordinates {(5.2,1.9) (6.1,1.6) (7.0,1.9) (8.1,2.5)};

    \node[blue] at (7.7,2.8) {$u(t)$};
    \node[red] at (7.6,1.1) {$u^{[t_1]}(t)$};
\end{tikzpicture}
\caption{Troncamento di un ingresso ammissibile. La parte successiva a $t_1$ deve restare un ingresso ammissibile.}
\end{figure}

\subsection{Chiusura per concatenazione}

Dati due ingressi ammissibili $u(\cdot),v(\cdot)\in\mathcal{U}$ e fissato un istante $\tau\in T$, si può costruire il segnale concatenato
\[
w(t)=
\begin{cases}
u(t), & t<\tau,\\[4pt]
v(t), & t\ge \tau.
\end{cases}
\]

La chiusura per concatenazione richiede che anche $w(\cdot)$ appartenga a $\mathcal{U}$.

\[
u(\cdot),v(\cdot)\in\mathcal{U}
\quad \Longrightarrow \quad
w(\cdot)\in\mathcal{U}.
\]

\paragraph{Significato.}
Posso prendere un tratto di un ingresso e, a partire da un certo istante, sostituirlo con un altro ingresso, ottenendo ancora un segnale ammissibile.

\begin{figure}[H]
\centering
\begin{tikzpicture}[scale=0.95]
    \draw[->] (0,0) -- (9,0) node[right] {$t$};
    \draw[->] (0,0) -- (0,3.2) node[above] {$w(t)$};

    \draw[dashed] (4.8,0) -- (4.8,2.8) node[above] {$\tau$};

    \draw[thick,blue]
        plot[smooth] coordinates {(0.8,0.8) (1.6,1.9) (2.3,1.4) (3.1,2.1) (4.0,1.6) (4.8,1.8)};

    \draw[thick,orange!90!black]
        plot[smooth] coordinates {(4.8,1.8) (5.4,1.4) (6.2,1.5) (7.0,2.0) (8.2,2.4)};

    \node[blue] at (2.1,2.5) {$u(t)$};
    \node[orange!90!black] at (7.1,2.8) {$v(t)$};
\end{tikzpicture}
\caption{Concatenazione di due ingressi ammissibili. Il segnale risultante deve appartenere ancora a $\mathcal{U}$.}
\end{figure}

\section{Funzione di transizione di stato}

Il cuore della descrizione dinamica è la \emph{funzione di transizione di stato}:
\[
\phi:\{(t,t_0)\in T^2:\, t\ge t_0\}\times X\times \mathcal{U}\to X.
\]

Essa permette di scrivere lo stato al tempo $t$ nella forma
\[
x(t)=\phi\bigl(t,t_0,x(t_0),u(\cdot)\bigr).
\]

Questa formula dice che lo stato al tempo corrente dipende da:
\begin{itemize}
    \item istante corrente $t$;
    \item istante iniziale $t_0$;
    \item stato iniziale $x(t_0)$;
    \item ingresso applicato.
\end{itemize}

\paragraph{Nota.}
Quando si scrive $\phi(t,t_0,x(t_0),u(\cdot))$, la notazione $u(\cdot)$ indica l'intera funzione di ingresso, non soltanto il suo valore istantaneo.

\subsection{Proprietà fondamentali di \texorpdfstring{$\phi$}{phi}}

\subsubsection{Consistenza}

All'istante iniziale il sistema deve restituire esattamente lo stato iniziale:
\[
\phi\bigl(t_0,t_0,x(t_0),u(\cdot)\bigr)=x(t_0).
\]

Questa proprietà è naturale: se non è ancora trascorso tempo, lo stato non può essere cambiato.

\subsubsection{Causalità}

La funzione di transizione non può dipendere dal futuro dell'ingresso.

Se due ingressi $u_1(\cdot)$ e $u_2(\cdot)$ coincidono fino all'istante $t_1$, allora devono produrre lo stesso stato a $t_1$:
\[
u_1(\tau)=u_2(\tau)\quad \forall \tau\in [t_0,t_1]
\quad \Longrightarrow \quad
\phi(t_1,t_0,\bar x,u_1)=\phi(t_1,t_0,\bar x,u_2).
\]

\paragraph{Interpretazione.}
Per sapere in quale stato si trovi il sistema all'istante $t_1$, basta conoscere l'ingresso fino a $t_1$; ciò che accadrà dopo non può influenzare il presente.

\subsubsection{Separazione rispetto al futuro}

Una formulazione equivalente e molto utile consiste nel dire che, per calcolare $x(t)$, è sufficiente il tratto di ingresso compreso tra $t_0$ e $t$:
\[
x(t)=\phi\bigl(t,t_0,x(t_0),u|_{[t_0,t]}\bigr).
\]

In altre parole, tutta la parte futura di $u(\cdot)$ è irrilevante per determinare lo stato attuale.

\subsubsection{Composizione}

L'evoluzione da $t_0$ a $t_2$ può essere spezzata in due evoluzioni successive:
\[
\phi(t_2,t_0,\bar x,u)
=
\phi\Bigl(t_2,t_1,\phi(t_1,t_0,\bar x,u),u\Bigr),
\qquad t_0\le t_1\le t_2.
\]

\paragraph{Significato.}
Prima porto il sistema da $t_0$ a $t_1$, ottenendo lo stato intermedio; poi uso quello stato come nuova condizione iniziale per evolvere da $t_1$ a $t_2$.

\section{Il significato dello stato}

Dalle proprietà precedenti emerge una delle idee centrali del corso:

\begin{quote}
Lo stato corrente contiene tutta l'informazione necessaria, insieme all'ingresso futuro, per determinare l'evoluzione futura del sistema.
\end{quote}

Per questo si dice che lo stato è una \emph{memoria compressa del passato}: non è necessario conservare tutta la storia passata, ma solo la parte essenziale riassunta in $x(t)$.

\section{Trasformazione d'uscita}

Per descrivere l'uscita si può introdurre una trasformazione globale
\[
G:\mathcal{U}\times X\to \mathcal{Y},
\]
che associa alla coppia ``ingresso + stato iniziale'' la funzione di uscita completa:
\[
(u(\cdot),x(t_0)) \mapsto y(\cdot).
\]

A livello istantaneo, in una rappresentazione locale, si scrive spesso
\[
y(t)=g(x(t),u(t),t),
\]
dove
\[
g:X\times U\times T \to Y.
\]

\paragraph{Osservazione.}
La scrittura con $G$ è una descrizione ``globale'' della relazione ingresso--uscita, mentre la scrittura con $g$ è una descrizione ``istantanea''.

\begin{figure}[H]
\centering
\begin{tikzpicture}[>=Latex]
    \node[draw, rounded corners, minimum width=3.3cm, minimum height=1.4cm] (sys) at (0,0) {Sistema dinamico};
    \draw[->, thick] (-4,0) -- (sys.west) node[midway, above] {$u(\cdot)$};
    \draw[->, thick] (sys.east) -- (4,0) node[midway, above] {$y(\cdot)$};
    \draw[->, thick] (0,-2) -- (sys.south) node[midway, right] {$x(t_0)$};
\end{tikzpicture}
\caption{Visione globale della trasformazione d'uscita: ingresso e stato iniziale determinano l'uscita.}
\end{figure}

\section{Classificazione dei sistemi dinamici}

\subsection{Sistemi propri e strettamente propri}

Nella forma locale
\[
y(t)=g(x(t),u(t),t),
\]
si distinguono due casi.

\paragraph{Sistema proprio.}
L'uscita può dipendere direttamente anche dall'ingresso corrente $u(t)$.

\paragraph{Sistema strettamente proprio.}
L'uscita non dipende istantaneamente da $u(t)$, ma solo dallo stato:
\[
y(t)=g(x(t),t).
\]

\paragraph{Interpretazione fisica.}
In un sistema strettamente proprio non esiste un passaggio istantaneo ingresso--uscita: l'ingresso deve prima modificare lo stato, e solo attraverso lo stato si riflette sull'uscita.

\paragraph{Caso lineare.}
Per un sistema lineare continuo o discreto:
\[
\begin{cases}
\dot x = Ax + Bu,\\
y = Cx + Du,
\end{cases}
\qquad \text{oppure} \qquad
\begin{cases}
x(k+1)=Ax(k)+Bu(k),\\
y(k)=Cx(k)+Du(k),
\end{cases}
\]
il sistema è strettamente proprio se e solo se
\[
D=0.
\]

Nel dominio della trasformata, una funzione di trasferimento strettamente propria ha denominatore di grado \emph{maggiore} del numeratore.

\subsection{Sistemi tempo-varianti e tempo-invarianti}

Un sistema è \emph{tempo-invariante} se il suo comportamento non cambia semplicemente traslando il riferimento temporale.

Definiamo, per un qualunque shift temporale $\tau$, il segnale traslato:
\[
u_\tau(t)=u(t-\tau).
\]

L'idea di tempo-invarianza è: se applico oggi un certo ingresso oppure applico domani lo stesso profilo, semplicemente traslato nel tempo, il sistema deve reagire nello stesso modo, a meno della stessa traslazione temporale.

Formalmente, per la dinamica, si richiede che
\[
\phi(t,t_0,\bar x,u)=\phi(t+\tau,t_0+\tau,\bar x,u_\tau).
\]

Per l'uscita, in forma locale, la tempo-invarianza si manifesta come assenza di dipendenza esplicita dal tempo:
\[
y(t)=g(x(t),u(t))
\]
invece di
\[
y(t)=g(x(t),u(t),t).
\]

\paragraph{Caso lineare.}
Un sistema lineare è tempo-invariante quando le matrici non dipendono dal tempo:
\[
A,B,C,D \quad \text{costanti}.
\]

\begin{figure}[H]
\centering
\begin{tikzpicture}[scale=0.95]
    \draw[->] (0,0) -- (10,0) node[right] {$t$};
    \draw[->] (0,0) -- (0,3.6) node[above] {$u(t)$};

    \draw[thick,blue]
        plot[smooth] coordinates {(0.8,0.7) (1.5,1.9) (2.3,1.2) (3.1,2.0) (4.0,1.3) (4.8,2.4)};

    \draw[thick,red]
        plot[smooth] coordinates {(3.0,0.7) (3.7,1.9) (4.5,1.2) (5.3,2.0) (6.2,1.3) (7.0,2.4)};

    \draw[dashed] (0.8,-0.15) -- (0.8,0.15);
    \draw[dashed] (3.0,-0.15) -- (3.0,0.15);

    \node[blue] at (4.7,3.0) {$u(t)$};
    \node[red] at (7.5,3.0) {$u_\tau(t)$};
    \node at (1.0,-0.45) {$t_0-\tau$};
    \node at (3.0,-0.45) {$t_0$};
\end{tikzpicture}
\caption{Tempo-invarianza: lo stesso profilo di ingresso, traslato nel tempo, produce una risposta semplicemente traslata.}
\end{figure}

\subsection{Scelta dell'istante iniziale}

Nei sistemi tempo-invarianti si può spesso scegliere senza perdita di generalità
\[
t_0=0.
\]

Questo non significa che il sistema ``nasca'' in $0$, ma che l'origine dei tempi è arbitraria: conta il tempo trascorso dall'inizio dell'osservazione, non il calendario assoluto.

\section{Sintesi della lezione}

Gli elementi fondamentali emersi in questa prima lezione sono:
\begin{enumerate}
    \item un sistema dinamico richiede insiemi di tempi, ingressi, uscite e stati;
    \item la dinamica si descrive tramite una funzione di transizione di stato $\phi$;
    \item l'uscita si descrive tramite una trasformazione globale $G$ oppure una legge locale $g$;
    \item le proprietà cruciali sono consistenza, causalità e composizione;
    \item una prima classificazione distingue sistemi propri, strettamente propri, tempo-varianti e tempo-invarianti.
\end{enumerate}

Questi concetti costituiranno la base per la rappresentazione a stato e, più avanti, per lo studio di linearità, raggiungibilità, osservabilità e stabilità.
\chapter{L2 -- Sistemi a stato vettore, rappresentazioni, moto, equilibrio, raggiungibilità, osservabilità e sistemi lineari}

\section{Sistemi a stato vettore}

In questa lezione si passa dalla descrizione astratta del sistema dinamico alla formulazione \emph{a stato vettore}.  
L'idea è che, in moltissimi problemi di interesse ingegneristico, i valori istantanei di ingresso, uscita e stato possano essere rappresentati come vettori.

Si assumono quindi:
\[
U,\quad Y,\quad X
\]
spazi vettoriali, con
\[
\dim U = m,\qquad \dim Y = \ell,\qquad \dim X = n.
\]

Qui:
\begin{itemize}
    \item $m$ è il numero di ingressi;
    \item $\ell$ è il numero di uscite;
    \item $n$ è il numero di variabili di stato.
\end{itemize}

\paragraph{Caso SISO e MIMO.}
Se
\[
m=\ell=1,
\]
il sistema è SISO (\emph{Single Input Single Output}); altrimenti, in generale, è MIMO.

\subsection{Spazi dei valori istantanei e spazi dei segnali}

Anche qui è fondamentale distinguere:
\[
U \ \text{da}\ \mathcal{U},
\qquad
Y \ \text{da}\ \mathcal{Y}.
\]

Infatti:
\begin{itemize}
    \item $U$ e $Y$ sono spazi vettoriali di valori \emph{istantanei};
    \item $\mathcal{U}$ e $\mathcal{Y}$ sono insiemi di \emph{funzioni nel tempo}.
\end{itemize}

Se $u_1(\cdot),u_2(\cdot)\in\mathcal{U}$ e $\alpha\in\mathbb{R}$, le operazioni naturali sono definite punto per punto:
\[
(u_1+u_2)(t)=u_1(t)+u_2(t),
\qquad
(\alpha u_1)(t)=\alpha\,u_1(t).
\]

Quando $\mathcal{U}$ è chiuso rispetto a tali operazioni, esso è uno spazio vettoriale. Lo stesso vale per $\mathcal{Y}$.

\paragraph{Osservazione importante.}
In generale:
\begin{itemize}
    \item $U$, $X$, $Y$ sono spazi vettoriali di dimensione finita;
    \item $\mathcal{U}$ e $\mathcal{Y}$ sono invece, tipicamente, spazi di dimensione infinita.
\end{itemize}

Questo diventerà cruciale quando si parlerà di rappresentazioni matriciali: una mappa lineare tra spazi finito-dimensionali si rappresenta con una matrice; una mappa lineare definita su uno spazio di funzioni, in generale, no.

\section{Rappresentazione esplicita del sistema}

Nel linguaggio introdotto nella lezione precedente, un sistema dinamico può essere descritto in forma esplicita come
\begin{align}
x(t) &= \phi\bigl(t,t_0,x(t_0),u_{[t_0,t]}\bigr), \label{eq:L2_phi}\\
y(t) &= g\bigl(t,t_0,x(t_0),u_{[t_0,t]}\bigr). \label{eq:L2_gglob}
\end{align}

Questa scrittura è detta \emph{esplicita} perché fornisce direttamente lo stato e l'uscita al tempo $t$.

\paragraph{Interpretazione.}
Se conosco:
\begin{itemize}
    \item lo stato iniziale $x(t_0)$,
    \item l'ingresso applicato nell'intervallo $[t_0,t]$,
    \item la funzione di transizione $\phi$,
\end{itemize}
allora conosco esattamente lo stato al tempo $t$.

\section{Rappresentazione implicita}

Spesso il sistema non viene descritto direttamente tramite $\phi$, ma tramite una legge locale di evoluzione.

\subsection{Caso a tempo continuo}

Nel tempo continuo, sotto opportune ipotesi di regolarità, la descrizione diventa:
\begin{equation}
\dot x(t)=f\bigl(x(t),u(t),t\bigr),
\label{eq:L2_ct_state}
\end{equation}
con condizione iniziale
\[
x(t_0)=x_0.
\]

L'uscita è descritta localmente da
\begin{equation}
y(t)=g\bigl(x(t),u(t),t\bigr).
\label{eq:L2_ct_output}
\end{equation}

Questa è la \emph{rappresentazione implicita}: non mi dice direttamente quale sia lo stato a un certo tempo futuro, ma mi dice \emph{come varia} lo stato.

\subsection{Da \texorpdfstring{$\phi$}{phi} a \texorpdfstring{$f$}{f}: idea del passaggio}

Dalla proprietà di composizione si ha, per $\Delta t>0$ piccolo,
\[
x(t+\Delta t)=\phi\bigl(t+\Delta t,t,x(t),u_{[t,t+\Delta t]}\bigr).
\]

Allora
\[
\frac{x(t+\Delta t)-x(t)}{\Delta t}
=
\frac{\phi\bigl(t+\Delta t,t,x(t),u_{[t,t+\Delta t]}\bigr)-x(t)}{\Delta t}.
\]

Se il limite esiste per $\Delta t\to 0$, si ottiene
\[
\dot x(t)=f\bigl(x(t),u(t),t\bigr).
\]

\paragraph{Lettura concettuale.}
\begin{itemize}
    \item La forma esplicita descrive il risultato finale dell'evoluzione.
    \item La forma implicita descrive la legge istantanea che genera quell'evoluzione.
\end{itemize}

\subsection{Caso a tempo discreto}

Nel tempo discreto non si usa la derivata, ma una relazione ricorsiva:
\begin{equation}
x(k+1)=f\bigl(x(k),u(k),k\bigr),
\label{eq:L2_dt_state}
\end{equation}
e
\begin{equation}
y(k)=g\bigl(x(k),u(k),k\bigr).
\label{eq:L2_dt_output}
\end{equation}

Qui la rappresentazione implicita coincide di fatto con un aggiornamento passo-passo dello stato.

\paragraph{Osservazione.}
Nel continuo passare da $f$ a $\phi$ significa risolvere un sistema di equazioni differenziali.  
Nel discreto passare da $f$ a $\phi$ significa iterare la dinamica.

\section{Eventi, moti, traiettorie e punti di equilibrio}

\subsection{Evento}

Un \emph{evento} del sistema è una coppia
\[
(x(t),t)\in X\times T.
\]

Si tratta quindi di uno stato osservato in un preciso istante.

\begin{figure}[H]
\centering
\begin{tikzpicture}[scale=1.0]
    \draw[->] (0,0) -- (7.5,0) node[right] {$t$};
    \draw[->] (0,0) -- (0,4.0) node[above] {$x(t)$};

    \draw[thick]
        plot[smooth] coordinates {(0.6,1.0) (1.3,1.8) (2.2,1.5) (3.2,2.4) (4.1,2.0) (5.3,3.0) (6.5,3.4)};

    \draw[dashed] (3.2,0) -- (3.2,2.4);
    \draw[dashed] (0,2.4) -- (3.2,2.4);
    \filldraw (3.2,2.4) circle (2pt);

    \node at (3.55,2.75) {evento $(x(t),t)$};
\end{tikzpicture}
\caption{Rappresentazione grafica di un evento del sistema.}
\end{figure}

\subsection{Moto del sistema}

Fissati:
\begin{itemize}
    \item uno stato iniziale $x(t_0)$,
    \item un ingresso $u(\cdot)$,
    \item un istante iniziale $t_0$,
\end{itemize}
il \emph{moto} del sistema è l'insieme
\[
\Bigl\{(x(t),t)\in X\times T:\ x(t)=\phi(t,t_0,x(t_0),u),\ t\ge t_0\Bigr\}.
\]

Il moto vive quindi nello spazio stato--tempo.

\subsection{Traiettoria}

La \emph{traiettoria} è la proiezione del moto sul solo spazio degli stati:
\[
\Bigl\{x(t)\in X:\ x(t)=\phi(t,t_0,x(t_0),u),\ t\ge t_0\Bigr\}.
\]

\paragraph{Differenza tra moto e traiettoria.}
\begin{itemize}
    \item Il moto tiene conto anche del tempo.
    \item La traiettoria descrive solo il percorso dello stato nello spazio $X$.
\end{itemize}

\subsection{Punto di equilibrio}

Uno stato $\bar x\in X$ è un \emph{punto di equilibrio} per un ingresso $\bar u(\cdot)$ se
\[
\bar x=\phi(t,t_0,\bar x,\bar u)\qquad \forall t\ge t_0.
\]

Questo significa che, se il sistema parte da $\bar x$ ed è soggetto all'ingresso $\bar u$, allora resta per sempre nello stesso stato.

\paragraph{Tempo continuo.}
Se il sistema è descritto da
\[
\dot x(t)=f(x(t),u(t),t),
\]
allora la condizione di equilibrio è
\[
0=f(\bar x,\bar u(t),t)\qquad \forall t\ge t_0.
\]

Nel caso tempo-invariante con ingresso costante $\bar u$, si riduce a
\[
0=f(\bar x,\bar u).
\]

\paragraph{Tempo discreto.}
Se
\[
x(k+1)=f(x(k),u(k),k),
\]
allora un equilibrio soddisfa
\[
\bar x=f(\bar x,\bar u(k),k)\qquad \forall k\ge k_0.
\]

Nel caso tempo-invariante con ingresso costante:
\[
\bar x=f(\bar x,\bar u).
\]

\section{Ciclo limite}

Un moto è periodico se esiste $T_p>0$ tale che
\[
x(t+T_p)=x(t)\qquad \forall t.
\]

Un \emph{ciclo limite} è una traiettoria periodica isolata nello spazio di stato.

\paragraph{Interpretazione.}
Il sistema non converge a un punto di equilibrio, ma evolve ripetutamente lungo una stessa orbita chiusa.

\begin{figure}[H]
\centering
\begin{tikzpicture}[scale=1.0]
    \draw[->] (-3.2,0) -- (3.2,0) node[right] {$x_1$};
    \draw[->] (0,-2.7) -- (0,2.7) node[above] {$x_2$};

    \draw[thick]
        plot[smooth cycle, tension=0.9]
        coordinates {(1.8,0.2) (1.0,1.6) (-0.8,1.7) (-1.8,0.2) (-1.2,-1.5) (0.8,-1.7)};
\end{tikzpicture}
\caption{Schema qualitativo di un ciclo limite nel piano di stato.}
\end{figure}

\paragraph{Nota concettuale.}
Il ciclo limite è tipico dei sistemi non lineari. Nei sistemi lineari possono comparire moti periodici, ma non nel senso strutturale forte tipico dei cicli limite isolati.

\section{Raggiungibilità e controllabilità}

L'ingresso modifica il moto del sistema. Da qui nascono due nozioni fondamentali.

Siano dati due eventi:
\[
(x_1,t_1),\qquad (x_2,t_2),\qquad t_2>t_1.
\]

Si dice che $(x_2,t_2)$ è \emph{raggiungibile} da $(x_1,t_1)$ se esiste un ingresso $u(\cdot)\in\mathcal{U}$ tale che
\[
x_2=\phi(t_2,t_1,x_1,u).
\]

In modo duale, si dice che $(x_1,t_1)$ è \emph{controllabile} da $(x_2,t_2)$ se esiste un ingresso che porta il sistema da $x_1$ al tempo $t_1$ a $x_2$ al tempo $t_2$.

\paragraph{Idea intuitiva.}
\begin{itemize}
    \item Raggiungibilità: posso arrivare nello stato desiderato?
    \item Controllabilità: posso forzare il sistema a seguire l'evoluzione voluta?
\end{itemize}

\section{Osservabilità e ricostruibilità}

La domanda di fondo è:

\begin{quote}
Osservando ingresso e uscita, riesco a capire da quale stato iniziale è partito il sistema?
\end{quote}

\subsection{Stati indistinguibili}

Due stati iniziali $x_0$ e $\bar x_0$ si dicono \emph{indistinguibili} nell'intervallo $[t_0,t_1]$ se, applicando lo stesso ingresso, producono la stessa uscita su tutto l'intervallo:
\[
g(\tau,t_0,x_0,u)=g(\tau,t_0,\bar x_0,u)
\qquad
\forall \tau\in [t_0,t_1],\ \forall u\in\mathcal U.
\]

Se esistono due stati distinti indistinguibili, il sistema non è osservabile.

\paragraph{Interpretazione fisica.}
Il sensore non riesce a distinguere quei due stati: l'uscita ``copre'' troppo poco dello stato.

\section{Esempio discreto di mancata raggiungibilità e mancata osservabilità}

Consideriamo il sistema a tempo discreto
\[
\begin{cases}
x(k+1)=Ax(k)+Bu(k),\\
y(k)=Cx(k),
\end{cases}
\]
con
\[
A=
\begin{bmatrix}
0 & 1\\
0 & 0
\end{bmatrix},
\qquad
B=
\begin{bmatrix}
1\\
0
\end{bmatrix},
\qquad
C=
\begin{bmatrix}
0 & 1
\end{bmatrix}.
\]

\subsection{Verifica della raggiungibilità}

La matrice di raggiungibilità è
\[
\mathcal{R}=\begin{bmatrix}B & AB\end{bmatrix}.
\]

Calcoliamo:
\[
AB=
\begin{bmatrix}
0 & 1\\
0 & 0
\end{bmatrix}
\begin{bmatrix}
1\\
0
\end{bmatrix}
=
\begin{bmatrix}
0\\
0
\end{bmatrix}.
\]

Quindi
\[
\mathcal{R}=
\begin{bmatrix}
1 & 0\\
0 & 0
\end{bmatrix},
\qquad
\operatorname{rank}(\mathcal{R})=1<2.
\]

Il sistema non è raggiungibile.

\paragraph{Lettura diretta.}
Se parto da
\[
x(0)=\begin{bmatrix}0\\0\end{bmatrix},
\]
allora
\[
x(1)=Ax(0)+Bu(0)=\begin{bmatrix}u(0)\\0\end{bmatrix},
\]
e poi
\[
x(2)=Ax(1)+Bu(1)=\begin{bmatrix}u(1)\\0\end{bmatrix}.
\]

La seconda componente resta sempre nulla.  
Quindi non posso raggiungere uno stato del tipo
\[
\begin{bmatrix}
1\\5
\end{bmatrix}.
\]

\subsection{Verifica dell'osservabilità}

La matrice di osservabilità è
\[
\mathcal{O}=
\begin{bmatrix}
C\\
CA
\end{bmatrix}.
\]

Calcoliamo:
\[
CA=
\begin{bmatrix}
0 & 1
\end{bmatrix}
\begin{bmatrix}
0 & 1\\
0 & 0
\end{bmatrix}
=
\begin{bmatrix}
0 & 0
\end{bmatrix}.
\]

Quindi
\[
\mathcal{O}=
\begin{bmatrix}
0 & 1\\
0 & 0
\end{bmatrix},
\qquad
\operatorname{rank}(\mathcal{O})=1<2.
\]

Il sistema non è osservabile.

\paragraph{Interpretazione.}
L'uscita vale
\[
y(k)=Cx(k)=
\begin{bmatrix}
0 & 1
\end{bmatrix}
\begin{bmatrix}
x_1(k)\\
x_2(k)
\end{bmatrix}
=x_2(k).
\]

Il sensore legge solo la seconda componente dello stato.  
La prima componente può cambiare senza essere direttamente rilevata: per questo parte dello stato rimane nascosta all'osservazione.

\section{Sistemi lineari}

Un sistema si dice lineare se la funzione di transizione $\phi$ e la trasformazione d'uscita $g$ sono lineari rispetto a stato iniziale e ingresso.

Formalmente, per ogni $x_1,x_2\in X$, $u_1,u_2\in\mathcal{U}$ e $\alpha_1,\alpha_2\in\mathbb{R}$,
\begin{multline}
\phi\bigl(t,t_0,\alpha_1x_1+\alpha_2x_2,\alpha_1u_1+\alpha_2u_2\bigr)
\\
=
\alpha_1\phi(t,t_0,x_1,u_1)+\alpha_2\phi(t,t_0,x_2,u_2),
\label{eq:L2_linearity_phi}
\end{multline}
e analogamente
\begin{multline}
g\bigl(t,t_0,\alpha_1x_1+\alpha_2x_2,\alpha_1u_1+\alpha_2u_2\bigr)
\\
=
\alpha_1g(t,t_0,x_1,u_1)+\alpha_2g(t,t_0,x_2,u_2).
\label{eq:L2_linearity_g}
\end{multline}

\subsection{Proprietà di decomposizione}

La linearità permette di separare l'effetto dello stato iniziale dall'effetto dell'ingresso.

Ponendo in \eqref{eq:L2_linearity_phi}
\[
x_2=0,\qquad u_1=0,\qquad \alpha_1=\alpha_2=1,
\]
si ottiene
\[
\phi(t,t_0,x,u)=\phi(t,t_0,x,0)+\phi(t,t_0,0,u).
\]

Quindi l'evoluzione dello stato si decompone in:
\[
\phi(t,t_0,x,u)=\phi_\ell(t,t_0,x)+\phi_f(t,t_0,u),
\]
dove
\[
\phi_\ell(t,t_0,x):=\phi(t,t_0,x,0)
\]
è la \emph{risposta libera}, e
\[
\phi_f(t,t_0,u):=\phi(t,t_0,0,u)
\]
è la \emph{risposta forzata}.

Allo stesso modo:
\[
g(t,t_0,x,u)=g_\ell(t,t_0,x)+g_f(t,t_0,u).
\]

\paragraph{Interpretazione.}
\begin{itemize}
    \item La risposta libera è quella dovuta unicamente allo stato iniziale.
    \item La risposta forzata è quella dovuta unicamente all'ingresso.
\end{itemize}

\begin{figure}[H]
\centering
\begin{tikzpicture}[>=Latex, scale=1.0]
    \node[draw, rounded corners, minimum width=2.7cm, minimum height=1.1cm] (sum) at (0,0) {$\phi$};
    \draw[->, thick] (-3,0.7) -- (sum.west) node[midway, above] {$x(t_0)$};
    \draw[->, thick] (-3,-0.7) -- (sum.west) node[midway, below] {$u(\cdot)$};
    \draw[->, thick] (sum.east) -- (3.5,0) node[midway, above] {$x(t)$};

    \node at (0,-2.0) {$x(t)=\underbrace{\phi_\ell(t,t_0,x(t_0))}_{\text{risposta libera}}
    +\underbrace{\phi_f(t,t_0,u)}_{\text{risposta forzata}}$};
\end{tikzpicture}
\caption{Idea di decomposizione della risposta nei sistemi lineari.}
\end{figure}

\subsection{Caso lineare tempo-invariante continuo}

Per un sistema LTI continuo
\[
\dot x(t)=Ax(t)+Bu(t),
\qquad
y(t)=Cx(t)+Du(t),
\]
la soluzione dello stato è
\[
x(t)=e^{A(t-t_0)}x(t_0)+\int_{t_0}^{t} e^{A(t-\tau)}Bu(\tau)\,d\tau.
\]

Qui:
\begin{itemize}
    \item
    \[
    e^{A(t-t_0)}x(t_0)
    \]
    è la risposta libera;
    \item
    \[
    \int_{t_0}^{t} e^{A(t-\tau)}Bu(\tau)\,d\tau
    \]
    è la risposta forzata.
\end{itemize}

\section{Nota finale sulle mappe lineari}

La decomposizione mette in evidenza una differenza importante.

\subsection{Parte libera}

La mappa
\[
\phi_\ell: X\to X
\]
è una mappa lineare tra spazi finito-dimensionali; quindi si può rappresentare con una matrice.

\subsection{Parte forzata}

La mappa
\[
\phi_f:\mathcal{U}\to X
\]
è una mappa lineare da uno spazio di funzioni, in generale infinito-dimensionale, verso lo spazio di stato.

Per questo motivo, in generale, $\phi_f$ \emph{non} si rappresenta con una matrice finita.

\paragraph{Conseguenza.}
Nei sistemi lineari:
\begin{itemize}
    \item la dinamica interna si presta naturalmente a una descrizione matriciale;
    \item l'effetto dell'ingresso richiede invece operatori più generali, oppure formule integrali/sommatorie.
\end{itemize}

\section{Sintesi della lezione}

In questa seconda lezione sono emersi i seguenti punti chiave:
\begin{enumerate}
    \item nei sistemi a stato vettore, ingresso, uscita e stato vivono in spazi vettoriali;
    \item la rappresentazione esplicita usa $\phi$, quella implicita usa $f$;
    \item si definiscono evento, moto, traiettoria, equilibrio e ciclo limite;
    \item si introducono raggiungibilità e osservabilità;
    \item nei sistemi lineari vale la decomposizione in risposta libera e risposta forzata.
\end{enumerate}

Questi concetti saranno alla base di tutto il seguito del corso: esponenziale di matrice, stabilità, raggiungibilità, osservabilità e progetto del controllo.
\chapter{L3 -- Sistemi lineari: matrice di transizione, evoluzione forzata, rappresentazioni esplicite e implicite}

\section{Forma lineare delle mappe di stato e di uscita}

Nella lezione precedente è emersa la proprietà di decomposizione dei sistemi lineari:
\[
\phi(t,t_0,x_0,u)=\phi_\ell(t,t_0,x_0)+\phi_f(t,t_0,u),
\]
\[
g(t,t_0,x_0,u)=g_\ell(t,t_0,x_0)+g_f(t,t_0,u).
\]

Quando il sistema è lineare e gli spazi
\[
X=\mathbb{R}^n,\qquad U=\mathbb{R}^m,\qquad Y=\mathbb{R}^{\ell}
\]
sono finito-dimensionali, la parte \emph{libera} e la parte \emph{istantanea di uscita} si possono rappresentare mediante matrici.

\subsection{Trasformazione d'uscita lineare}

La parte libera dell'uscita dipende linearmente dallo stato:
\[
g_\ell(x(t),t)=C(t)x(t),
\qquad
C(t)\in\mathbb{R}^{\ell\times n}.
\]

La parte forzata dell'uscita dipende linearmente dall'ingresso:
\[
g_f(u(t),t)=D(t)u(t),
\qquad
D(t)\in\mathbb{R}^{\ell\times m}.
\]

Quindi, in forma locale,
\begin{equation}
y(t)=C(t)x(t)+D(t)u(t).
\label{eq:L3_output_local}
\end{equation}

\paragraph{Interpretazione delle dimensioni.}
\begin{itemize}
    \item $x(t)\in \mathbb{R}^n$ è il vettore di stato;
    \item $u(t)\in \mathbb{R}^m$ è il vettore di ingresso;
    \item $y(t)\in \mathbb{R}^{\ell}$ è il vettore di uscita;
    \item $C(t)$ trasforma lo stato in uscita;
    \item $D(t)$ descrive l'eventuale passaggio istantaneo ingresso--uscita.
\end{itemize}

\subsection{Parte libera dello stato e matrice di transizione}

Poiché $\phi_\ell$ è una mappa lineare da $X$ in $X$, si può scrivere
\[
\phi_\ell(t,t_0,x(t_0))=\Phi(t,t_0)x(t_0),
\]
dove
\[
\Phi(t,t_0)\in\mathbb{R}^{n\times n}
\]
è la \emph{matrice di transizione di stato}.

\paragraph{Significato.}
La matrice $\Phi(t,t_0)$ trasforma lo stato iniziale $x(t_0)$ nello stato che il sistema avrebbe al tempo $t$ in assenza di ingresso.

\paragraph{Caso tempo-invariante.}
Se il sistema è lineare tempo-invariante, allora la matrice di transizione dipende solo dal tempo trascorso:
\[
\Phi(t,t_0)=\Phi(t-t_0).
\]

Questa dipendenza solo dalla differenza $t-t_0$ esprime precisamente il fatto che, in un sistema tempo-invariante, non conta \emph{quando} inizio l'evoluzione, ma solo \emph{quanto tempo} è trascorso dall'istante iniziale.

\section{Rappresentazioni esplicite di stato nei sistemi lineari}

\subsection{Tempo continuo}

Nel caso continuo, la rappresentazione esplicita dello stato assume la forma
\begin{equation}
x(t)=\Phi(t,t_0)x(t_0)+\int_{t_0}^{t} H(t,\tau)u(\tau)\,d\tau,
\label{eq:L3_explicit_TC}
\end{equation}
dove
\[
H(t,\tau)\in\mathbb{R}^{n\times m}
\]
è un opportuno \emph{nucleo di evoluzione forzata}.

L'uscita si scrive invece come
\begin{equation}
y(t)=C(t)x(t)+D(t)u(t).
\label{eq:L3_TC_output}
\end{equation}

\paragraph{Lettura della formula.}
La \eqref{eq:L3_explicit_TC} si divide in due contributi:
\begin{itemize}
    \item \emph{risposta libera}
    \[
    \Phi(t,t_0)x(t_0),
    \]
    dovuta solo allo stato iniziale;
    \item \emph{risposta forzata}
    \[
    \int_{t_0}^{t} H(t,\tau)u(\tau)\,d\tau,
    \]
    dovuta solo all'ingresso.
\end{itemize}

\subsection{Tempo discreto}

Nel caso discreto, l'integrale viene sostituito da una somma:
\begin{equation}
x(K)=\Phi(K,K_0)x(K_0)+\sum_{j=K_0}^{K-1} H(K,j)\,u(j),
\label{eq:L3_explicit_TD}
\end{equation}
con
\[
H(K,j)\in\mathbb{R}^{n\times m}.
\]

L'uscita si scrive
\begin{equation}
y(K)=C(K)x(K)+D(K)u(K).
\label{eq:L3_TD_output}
\end{equation}

\paragraph{Interpretazione.}
Per calcolare lo stato al tempo $K$:
\begin{itemize}
    \item porto avanti lo stato iniziale tramite $\Phi(K,K_0)$;
    \item sommo poi tutti i contributi degli ingressi applicati agli istanti
    \[
    j=K_0,K_0+1,\dots,K-1.
    \]
\end{itemize}

\section{Perché nel tempo discreto compare la somma}

Negli appunti questa parte è centrale: si vuole capire da dove nasce in modo naturale la formula
\[
x(K)=\Phi(K,K_0)x(K_0)+\sum_{j=K_0}^{K-1}H(K,j)u(j).
\]

L'idea consiste nel decomporre un ingresso discreto generico in una somma di \emph{ingressi elementari}.

\subsection{Un ingresso discreto su un intervallo finito}

Consideriamo l'ingresso ristretto all'intervallo
\[
[K_0,K)\cap \mathbb{Z}.
\]
Lo indichiamo con
\[
u_{[K_0,K)}(\cdot):T\to U.
\]

Esso è completamente individuato dai valori
\[
u(K_0),\ u(K_0+1),\ \dots,\ u(K-1).
\]

\paragraph{Osservazione importante.}
Anche se $u(\cdot)$ appartiene a uno spazio di funzioni, quando lo osservo solo su un intervallo discreto finito esso viene descritto da un numero finito di campioni. È proprio questo che permette di ricondurre l'evoluzione forzata a una somma finita di contributi elementari.

\begin{figure}[H]
\centering
\begin{tikzpicture}[scale=0.95]
    \draw[->] (0,0) -- (10.2,0) node[right] {$k$};
    \draw[->] (0,0) -- (0,4.2) node[above] {$U$};

    \foreach \x in {1,2,3,4,5,6,7,8,9} {
        \draw (\x,0.12) -- (\x,-0.12);
    }

    \draw[thick,red] (2,0)--(2,1.0);
    \fill[red] (2,1.0) circle (2.3pt);
    \node[red] at (2,-0.45) {$K_0$};
    \node[red] at (0.9,1.15) {$u(K_0)$};

    \draw[thick,green!60!black] (4,0)--(4,2.5);
    \fill[green!60!black] (4,2.5) circle (2.3pt);
    \node[green!60!black] at (4,-0.45) {$K_0+1$};
    \node[green!60!black] at (2.95,2.7) {$u(K_0+1)$};

    \draw[thick,violet] (5,0)--(5,1.8);
    \fill[violet] (5,1.8) circle (2.3pt);
    \node[violet] at (5,-0.45) {$K_0+2$};
    \node[violet] at (3.9,1.95) {$u(K_0+2)$};

    \draw[thick,orange!85!black] (9,0)--(9,3.3);
    \fill[orange!85!black] (9,3.3) circle (2.3pt);
    \node[orange!85!black] at (9,-0.45) {$K-1$};
    \node[orange!85!black] at (7.8,3.5) {$u(K-1)$};

    \fill (6,1.4) circle (2pt);
    \fill (7,3.0) circle (2pt);
    \fill (8,2.4) circle (2pt);
\end{tikzpicture}
\caption{Un ingresso discreto generico è completamente determinato dai suoi valori negli istanti $K_0,\dots,K-1$.}
\end{figure}

\subsection{Decomposizione in ingressi elementari}

Per ogni indice $j\in\{K_0,\dots,K-1\}$ definiamo il segnale elementare
\[
u_j(k)=
\begin{cases}
u(j), & k=j,\\[4pt]
0, & k\neq j.
\end{cases}
\]

Allora si ha la decomposizione
\begin{equation}
u_{[K_0,K)}(\cdot)=\sum_{j=K_0}^{K-1}u_j(\cdot).
\label{eq:L3_input_decomposition}
\end{equation}

\paragraph{Perché funziona.}
Per ogni istante $k$:
\begin{itemize}
    \item solo uno dei segnali $u_j$ può essere diverso da zero, precisamente quello con $j=k$;
    \item quindi la somma dei segnali elementari restituisce esattamente il valore originario $u(k)$.
\end{itemize}

\begin{figure}[H]
\centering
\begin{tikzpicture}[scale=0.9]
    % asse 1
    \begin{scope}[yshift=3.8cm]
        \draw[->] (0,0) -- (7.2,0) node[right] {$k$};
        \draw[->] (0,0) -- (0,2.4) node[above] {$U$};
        \foreach \x in {1,2,3,4,5,6} {\draw (\x,0.12)--(\x,-0.12);}
        \draw[thick,red] (2,0)--(2,1.2);
        \fill[red] (2,1.2) circle (2.1pt);
        \node[red] at (2,-0.35) {$K_0$};
        \node[red] at (0.9,1.45) {$u_{K_0}$};
        \foreach \x in {1,3,4,5,6} {\fill[red] (\x,0) circle (2pt);}
    \end{scope}

    % asse 2
    \begin{scope}[yshift=1.2cm]
        \draw[->] (0,0) -- (7.2,0) node[right] {$k$};
        \draw[->] (0,0) -- (0,2.4) node[above] {$U$};
        \foreach \x in {1,2,3,4,5,6} {\draw (\x,0.12)--(\x,-0.12);}
        \draw[thick,green!60!black] (3,0)--(3,1.8);
        \fill[green!60!black] (3,1.8) circle (2.1pt);
        \node[green!60!black] at (3,-0.35) {$K_0+1$};
        \node[green!60!black] at (1.6,2.05) {$u_{K_0+1}$};
        \foreach \x in {1,2,4,5,6} {\fill[green!60!black] (\x,0) circle (2pt);}
    \end{scope}

    % asse 3
    \begin{scope}[yshift=-1.4cm]
        \draw[->] (0,0) -- (7.2,0) node[right] {$k$};
        \draw[->] (0,0) -- (0,2.4) node[above] {$U$};
        \foreach \x in {1,2,3,4,5,6} {\draw (\x,0.12)--(\x,-0.12);}
        \draw[thick,violet] (4,0)--(4,1.4);
        \fill[violet] (4,1.4) circle (2.1pt);
        \node[violet] at (4,-0.35) {$K_0+2$};
        \node[violet] at (2.5,1.65) {$u_{K_0+2}$};
        \foreach \x in {1,2,3,5,6} {\fill[violet] (\x,0) circle (2pt);}
    \end{scope}
\end{tikzpicture}
\caption{Esempi di ingressi elementari: ciascuno è nullo ovunque tranne che in un solo istante.}
\end{figure}

\section{Dalla decomposizione dell'ingresso alla risposta forzata nel tempo discreto}

La risposta forzata è lineare rispetto all'ingresso. Applicando la decomposizione \eqref{eq:L3_input_decomposition} si ottiene
\[
\phi_f(K,K_0,u)=\sum_{j=K_0}^{K-1}\phi_f(K,K_0,u_j).
\]

Ora osserviamo che il segnale $u_j$ è nullo fino all'istante $j-1$ e si attiva solo all'istante $j$.  
Per causalità e composizione, il suo contributo allo stato a tempo $K$ dipende solo da ciò che accade a partire da $j$, quindi
\[
\phi_f(K,K_0,u_j)=\phi_f(K,j,u_j).
\]

Poiché $u_j$ è determinato dal solo vettore $u(j)\in U$, esiste una matrice
\[
H(K,j)\in\mathbb{R}^{n\times m}
\]
tale che
\[
\phi_f(K,j,u_j)=H(K,j)\,u(j).
\]

Sostituendo:
\[
\phi_f(K,K_0,u)=\sum_{j=K_0}^{K-1}H(K,j)\,u(j).
\]

Infine, sommando risposta libera e risposta forzata:
\[
x(K)=\Phi(K,K_0)x(K_0)+\sum_{j=K_0}^{K-1}H(K,j)\,u(j).
\]

\paragraph{Significato della matrice $H(K,j)$.}
La matrice $H(K,j)$ descrive come un ingresso applicato all'istante $j$ contribuisce allo stato osservato all'istante finale $K$.

\section{Rappresentazione esplicita discreta e passaggio alla forma ricorsiva}

La formula
\[
x(K)=\Phi(K,K_0)x(K_0)+\sum_{j=K_0}^{K-1}H(K,j)\,u(j)
\]
è una \emph{rappresentazione esplicita} del sistema a tempo discreto.

Se ora scegliamo
\[
K=j+1,
\]
otteniamo
\[
x(j+1)=\Phi(j+1,j)x(j)+H(j+1,j)u(j).
\]

Definendo
\[
A(j):=\Phi(j+1,j),
\qquad
B(j):=H(j+1,j),
\]
si ricava la forma ricorsiva standard
\begin{equation}
x(j+1)=A(j)x(j)+B(j)u(j).
\label{eq:L3_TD_implicit}
\end{equation}

L'uscita rimane
\begin{equation}
y(j)=C(j)x(j)+D(j)u(j).
\label{eq:L3_TD_output_local}
\end{equation}

\paragraph{Conclusione.}
Nel tempo discreto:
\begin{itemize}
    \item la rappresentazione con $\Phi$ e $H$ è la forma \emph{esplicita};
    \item la rappresentazione con $A(j)$ e $B(j)$ è la forma \emph{implicita} o ricorsiva.
\end{itemize}

\section{Proprietà di consistenza e composizione di \texorpdfstring{$\Phi$}{Phi}}

Le proprietà della funzione di transizione $\phi$ si traducono in proprietà matriciali di $\Phi$ e $H$.

\subsection{Consistenza}

Per definizione di consistenza:
\[
\phi(t_0,t_0,x(t_0),u)=x(t_0).
\]

Nel caso lineare e libero:
\[
x(t_0)=\Phi(t_0,t_0)x(t_0)\qquad \forall x(t_0).
\]

Poiché questa relazione deve valere per ogni stato iniziale, segue che
\begin{equation}
\Phi(t_0,t_0)=I_n.
\label{eq:L3_consistency_Phi}
\end{equation}

La stessa proprietà vale sia nel tempo continuo sia nel tempo discreto:
\[
\Phi(t,t)=I_n,
\qquad
\Phi(K,K)=I_n.
\]

\subsection{Composizione per \texorpdfstring{$\Phi$}{Phi}}

Dalla proprietà generale di composizione della dinamica:
\[
\phi(t_2,t_0,x_0,u)=\phi\bigl(t_2,t_1,\phi(t_1,t_0,x_0,u),u\bigr),
\qquad t_0\le t_1\le t_2,
\]
e scegliendo ingresso nullo $u=0$, si ottiene
\[
\Phi(t_2,t_0)x_0=\Phi(t_2,t_1)\Phi(t_1,t_0)x_0.
\]

Essendo vero per ogni $x_0$, segue:
\begin{equation}
\Phi(t_2,t_0)=\Phi(t_2,t_1)\Phi(t_1,t_0),
\qquad t_0\le t_1\le t_2.
\label{eq:L3_composition_Phi_TC}
\end{equation}

Nel discreto:
\begin{equation}
\Phi(K_2,K_0)=\Phi(K_2,K_1)\Phi(K_1,K_0),
\qquad K_0\le K_1\le K_2.
\label{eq:L3_composition_Phi_TD}
\end{equation}

\paragraph{Interpretazione.}
La matrice di transizione da $t_0$ a $t_2$ si ottiene componendo la transizione da $t_0$ a $t_1$ con quella da $t_1$ a $t_2$.

\section{Proprietà di composizione di \texorpdfstring{$H$}{H}}

Usiamo ora la rappresentazione esplicita discreta in due modi diversi.

Da un lato:
\[
x(K_2)=\Phi(K_2,K_0)x(K_0)+\sum_{j=K_0}^{K_2-1}H(K_2,j)u(j).
\]

Dall'altro, applicando prima la dinamica da $K_0$ a $K_1$ e poi da $K_1$ a $K_2$:
\[
x(K_2)=\Phi(K_2,K_1)x(K_1)+\sum_{j=K_1}^{K_2-1}H(K_2,j)u(j).
\]

Sostituiamo nella seconda formula
\[
x(K_1)=\Phi(K_1,K_0)x(K_0)+\sum_{j=K_0}^{K_1-1}H(K_1,j)u(j),
\]
ottenendo
\[
x(K_2)=
\Phi(K_2,K_1)\Phi(K_1,K_0)x(K_0)
+
\sum_{j=K_0}^{K_1-1}\Phi(K_2,K_1)H(K_1,j)u(j)
+
\sum_{j=K_1}^{K_2-1}H(K_2,j)u(j).
\]

Confrontando con la formula diretta e usando l'unicità dei coefficienti, si ottengono due risultati:

\begin{enumerate}
    \item la proprietà di composizione di $\Phi$ già vista;
    \item la proprietà di composizione di $H$:
    \begin{equation}
    \Phi(K_2,K_1)H(K_1,j)=H(K_2,j),
    \qquad K_0\le j<K_1\le K_2.
    \label{eq:L3_composition_H_TD}
    \end{equation}
\end{enumerate}

Nel tempo continuo l'analogo è
\begin{equation}
\Phi(t_2,t_1)H(t_1,\tau)=H(t_2,\tau),
\qquad t_0\le \tau \le t_1\le t_2.
\label{eq:L3_composition_H_TC}
\end{equation}

\paragraph{Interpretazione.}
Per propagare fino al tempo finale l'effetto di un ingresso applicato in un istante precedente, posso:
\begin{itemize}
    \item calcolarne prima l'effetto fino a un tempo intermedio;
    \item poi trasportare quell'effetto dal tempo intermedio a quello finale con $\Phi$.
\end{itemize}

\section{Riassunto strutturale delle rappresentazioni esplicite}

Le proprietà appena ricavate permettono di scrivere in modo ordinato le rappresentazioni esplicite dei sistemi lineari.

\subsection{Tempo discreto}

\begin{equation}
\boxed{
\begin{aligned}
x(K) &= \Phi(K,K_0)x(K_0)+\sum_{j=K_0}^{K-1}H(K,j)u(j),\\
y(K) &= C(K)x(K)+D(K)u(K).
\end{aligned}
}
\label{eq:L3_explicit_discrete_boxed}
\end{equation}

\subsection{Tempo continuo}

\begin{equation}
\boxed{
\begin{aligned}
x(t) &= \Phi(t,t_0)x(t_0)+\int_{t_0}^{t}H(t,\tau)u(\tau)\,d\tau,\\
y(t) &= C(t)x(t)+D(t)u(t).
\end{aligned}
}
\label{eq:L3_explicit_continuous_boxed}
\end{equation}

\section{Come si passa nel continuo alla forma implicita}

Negli appunti vengono poi ricavate due relazioni fondamentali, che permettono di passare dalla rappresentazione esplicita continua a quella implicita.

\subsection{Primo lemma: derivata della matrice di transizione}

Usiamo la proprietà di composizione:
\[
\Phi(t+\Delta t,t_0)=\Phi(t+\Delta t,t)\,\Phi(t,t_0).
\]

Sottraendo $\Phi(t,t_0)$ da entrambi i membri:
\[
\Phi(t+\Delta t,t_0)-\Phi(t,t_0)
=
\bigl[\Phi(t+\Delta t,t)-I_n\bigr]\Phi(t,t_0).
\]

Dividendo per $\Delta t$ e passando al limite per $\Delta t\to 0$:
\[
\frac{\partial \Phi(t,t_0)}{\partial t}
=
\left[
\frac{\partial \Phi(\zeta,t)}{\partial \zeta}
\right]_{\zeta=t}
\Phi(t,t_0).
\]

Definiamo allora
\begin{equation}
A(t):=
\left[
\frac{\partial \Phi(\zeta,t)}{\partial \zeta}
\right]_{\zeta=t}.
\label{eq:L3_def_A}
\end{equation}

Quindi il primo lemma diventa
\begin{equation}
\frac{\partial \Phi(t,t_0)}{\partial t}=A(t)\Phi(t,t_0).
\label{eq:L3_lemma1}
\end{equation}

\subsection{Secondo lemma: derivata del nucleo forzato}

Usando la proprietà di composizione di $H$,
\[
H(t+\Delta t,\tau)=\Phi(t+\Delta t,t)\,H(t,\tau),
\]
si ottiene in modo analogo:
\[
\frac{H(t+\Delta t,\tau)-H(t,\tau)}{\Delta t}
=
\frac{\Phi(t+\Delta t,t)-I_n}{\Delta t}\,H(t,\tau).
\]

Passando al limite:
\begin{equation}
\frac{\partial H(t,\tau)}{\partial t}=A(t)H(t,\tau).
\label{eq:L3_lemma2}
\end{equation}

\section{Derivazione della forma implicita nel tempo continuo}

Partiamo dalla rappresentazione esplicita
\[
x(t)=\Phi(t,t_0)x(t_0)+\int_{t_0}^{t}H(t,\tau)u(\tau)\,d\tau.
\]

Deriviamo rispetto a $t$.

\subsection{Derivata del primo termine}

Usando il lemma \eqref{eq:L3_lemma1}:
\[
\frac{d}{dt}\Bigl(\Phi(t,t_0)x(t_0)\Bigr)=A(t)\Phi(t,t_0)x(t_0).
\]

\subsection{Derivata del termine integrale}

Applichiamo la regola di Leibniz:
\[
\frac{d}{dt}\left(
\int_{t_0}^{t}H(t,\tau)u(\tau)\,d\tau
\right)
=
H(t,t)u(t)+\int_{t_0}^{t}\frac{\partial H(t,\tau)}{\partial t}u(\tau)\,d\tau.
\]

Usando il lemma \eqref{eq:L3_lemma2}:
\[
=
H(t,t)u(t)+\int_{t_0}^{t}A(t)H(t,\tau)u(\tau)\,d\tau.
\]

Poiché $A(t)$ non dipende dalla variabile di integrazione $\tau$, si può portare fuori dall'integrale:
\[
=
H(t,t)u(t)+A(t)\int_{t_0}^{t}H(t,\tau)u(\tau)\,d\tau.
\]

\subsection{Ricomposizione finale}

Sommiamo i due contributi:
\[
\dot x(t)
=
A(t)\Phi(t,t_0)x(t_0)
+
A(t)\int_{t_0}^{t}H(t,\tau)u(\tau)\,d\tau
+
H(t,t)u(t).
\]

Raccogliendo $A(t)$:
\[
\dot x(t)
=
A(t)\left[
\Phi(t,t_0)x(t_0)+\int_{t_0}^{t}H(t,\tau)u(\tau)\,d\tau
\right]
+
H(t,t)u(t).
\]

Ma la quantità tra parentesi quadre è proprio $x(t)$, quindi
\[
\dot x(t)=A(t)x(t)+H(t,t)u(t).
\]

Definendo
\begin{equation}
B(t):=H(t,t),
\label{eq:L3_def_B}
\end{equation}
si ottiene finalmente la forma implicita del sistema lineare continuo:
\begin{equation}
\boxed{
\dot x(t)=A(t)x(t)+B(t)u(t).
}
\label{eq:L3_TC_implicit}
\end{equation}

L'uscita è
\begin{equation}
\boxed{
y(t)=C(t)x(t)+D(t)u(t).
}
\label{eq:L3_TC_output_implicit}
\end{equation}

\section{Relazione tra forma esplicita e forma implicita}

A questo punto il quadro è completo.

\subsection{Tempo discreto}

\paragraph{Forma esplicita}
\[
x(K)=\Phi(K,K_0)x(K_0)+\sum_{j=K_0}^{K-1}H(K,j)u(j).
\]

\paragraph{Forma implicita}
\[
x(k+1)=A(k)x(k)+B(k)u(k),
\]
con
\[
A(k)=\Phi(k+1,k),\qquad B(k)=H(k+1,k).
\]

\subsection{Tempo continuo}

\paragraph{Forma esplicita}
\[
x(t)=\Phi(t,t_0)x(t_0)+\int_{t_0}^{t}H(t,\tau)u(\tau)\,d\tau.
\]

\paragraph{Forma implicita}
\[
\dot x(t)=A(t)x(t)+B(t)u(t),
\]
con
\[
A(t)=\left[\frac{\partial \Phi(\zeta,t)}{\partial \zeta}\right]_{\zeta=t},
\qquad
B(t)=H(t,t).
\]

\paragraph{Messaggio concettuale della lezione.}
Le matrici $\Phi$ e $H$ descrivono l'evoluzione \emph{globale} del sistema:
\begin{itemize}
    \item $\Phi$ trasporta nel tempo lo stato iniziale;
    \item $H$ pesa i contributi dell'ingresso applicato ai vari istanti.
\end{itemize}

Le matrici $A$ e $B$, invece, descrivono la legge \emph{locale} di evoluzione:
\begin{itemize}
    \item $A$ governa la dinamica interna;
    \item $B$ governa il modo in cui l'ingresso agisce istante per istante sullo stato.
\end{itemize}

\section{Sintesi finale della lezione}

In questa lezione sono stati introdotti e chiariti i seguenti punti:
\begin{enumerate}
    \item nei sistemi lineari l'uscita si scrive come
    \[
    y=Cx+Du;
    \]
    \item la risposta libera dello stato è rappresentata dalla matrice di transizione $\Phi$;
    \item la risposta forzata è descritta da un nucleo $H$;
    \item nel tempo discreto l'evoluzione forzata diventa una somma finita grazie alla decomposizione dell'ingresso in segnali elementari;
    \item $\Phi$ e $H$ soddisfano proprietà di consistenza e composizione;
    \item nel continuo, derivando la forma esplicita, si ottiene la forma implicita
    \[
    \dot x=A(t)x+B(t)u(t).
    \]
\end{enumerate}

Questa struttura è fondamentale perché nei capitoli successivi permetterà di studiare:
\begin{itemize}
    \item esponenziale di matrice;
    \item evoluzione libera e forzata;
    \item stabilità;
    \item raggiungibilità e osservabilità dei sistemi lineari.
\end{itemize}
\chapter{L4 -- Soluzione dei sistemi lineari: caso discreto generale, caso stazionario, successione di Peano--Baker ed esponenziale di matrice}

\section{Collegamento tra forma esplicita e forma implicita}

Nella lezione precedente abbiamo ottenuto le relazioni che collegano le quantità ``globali'' $\Phi$ e $H$ alle quantità ``locali'' $A$ e $B$.

\subsection{Tempo continuo}

Per il tempo continuo valgono:
\begin{equation}
A(t)=\left.\frac{\partial \Phi(\xi,t)}{\partial \xi}\right|_{\xi=t},
\qquad
B(t)=H(t,t).
\label{eq:L4_def_A_B_TC}
\end{equation}

Quindi la rappresentazione implicita locale del sistema è
\begin{equation}
\dot x(t)=A(t)x(t)+B(t)u(t),
\qquad
y(t)=C(t)x(t)+D(t)u(t).
\label{eq:L4_TC_implicita}
\end{equation}

\subsection{Tempo discreto}

Nel tempo discreto, invece, si ha:
\begin{equation}
A(j)=\Phi(j+1,j),
\qquad
B(j)=H(j+1,j),
\label{eq:L4_def_A_B_TD}
\end{equation}
e dunque la forma implicita diventa
\begin{equation}
x(k+1)=A(k)x(k)+B(k)u(k),
\qquad
y(k)=C(k)x(k)+D(k)u(k).
\label{eq:L4_TD_implicita}
\end{equation}

L'obiettivo di questa lezione è capire come, a partire da \eqref{eq:L4_TD_implicita} e \eqref{eq:L4_TC_implicita}, si ricostruiscono le forme esplicite e come si presentano le matrici $\Phi$ e $H$.

\section{Sistemi lineari discreti generici}

Consideriamo il sistema discreto non necessariamente stazionario
\begin{equation}
x(k+1)=A(k)x(k)+B(k)u(k),
\qquad
x(K_0)\ \text{assegnato}.
\label{eq:L4_sys_TD}
\end{equation}

\subsection{Sviluppo iterativo dei primi passi}

Iteriamo la dinamica a partire da $K_0$.

Per $k=K_0$:
\[
x(K_0+1)=A(K_0)x(K_0)+B(K_0)u(K_0).
\]

Per $k=K_0+1$:
\begin{align*}
x(K_0+2)
&=A(K_0+1)x(K_0+1)+B(K_0+1)u(K_0+1)\\
&=A(K_0+1)\bigl[A(K_0)x(K_0)+B(K_0)u(K_0)\bigr]+B(K_0+1)u(K_0+1)\\
&=A(K_0+1)A(K_0)x(K_0)+A(K_0+1)B(K_0)u(K_0)+B(K_0+1)u(K_0+1).
\end{align*}

Per $k=K_0+2$:
\begin{align*}
x(K_0+3)
&=A(K_0+2)x(K_0+2)+B(K_0+2)u(K_0+2)\\
&=A(K_0+2)A(K_0+1)A(K_0)x(K_0)\\
&\quad +A(K_0+2)A(K_0+1)B(K_0)u(K_0)\\
&\quad +A(K_0+2)B(K_0+1)u(K_0+1)\\
&\quad +B(K_0+2)u(K_0+2).
\end{align*}

A questo punto il pattern è chiaro:
\begin{itemize}
    \item compare un termine libero, cioè la propagazione dello stato iniziale;
    \item compare una somma di termini forzati, uno per ogni ingresso passato, ciascuno trasportato fino all'istante finale.
\end{itemize}

\subsection{Forma esplicita generale}

Introduciamo la convenzione di prodotto ordinato
\[
\prod_{i=K_0}^{K-1}A(i):=A(K-1)A(K-2)\cdots A(K_0),
\]
cioè gli indici crescono da destra verso sinistra, come richiesto dalla composizione delle matrici di stato.

Con questa convenzione si ottiene:
\begin{equation}
x(K)=\Phi(K,K_0)x(K_0)+\sum_{j=K_0}^{K-1}H(K,j)u(j),
\label{eq:L4_explicit_TD}
\end{equation}
dove
\begin{equation}
\Phi(K,K_0)=\prod_{i=K_0}^{K-1}A(i),
\label{eq:L4_Phi_TD_generic}
\end{equation}
e
\begin{equation}
H(K,j)=\left(\prod_{i=j+1}^{K-1}A(i)\right)B(j).
\label{eq:L4_H_TD_generic}
\end{equation}

\paragraph{Convenzione importante.}
Se il prodotto è vuoto, esso vale l'identità:
\[
\prod_{i=K}^{K-1}A(i)=I.
\]
Per esempio, per $j=K-1$ si ha
\[
H(K,K-1)=B(K-1).
\]

\paragraph{Interpretazione di $H(K,j)$.}
Il termine $H(K,j)$ dice come un ingresso applicato all'istante $j$ viene trasportato fino all'istante finale $K$.  
In altre parole:
\begin{itemize}
    \item $B(j)$ trasforma l'ingresso in variazione immediata dello stato;
    \item il prodotto di matrici $A$ successive trasporta quella variazione fino al tempo finale.
\end{itemize}

\begin{figure}[H]
\centering
\begin{tikzpicture}[>=Latex, node distance=1.45cm]
    \node[draw, rounded corners, minimum width=1.2cm, minimum height=0.9cm] (x0) {$x(K_0)$};
    \node[draw, rounded corners, right=of x0, minimum width=1.2cm, minimum height=0.9cm] (x1) {$x(K_0+1)$};
    \node[draw, rounded corners, right=of x1, minimum width=1.2cm, minimum height=0.9cm] (x2) {$x(K_0+2)$};
    \node[right=0.7cm of x2] (dots) {$\cdots$};
    \node[draw, rounded corners, right=0.7cm of dots, minimum width=1.2cm, minimum height=0.9cm] (xk) {$x(K)$};

    \draw[->, thick] (x0) -- (x1) node[midway, above] {$A(K_0)$};
    \draw[->, thick] (x1) -- (x2) node[midway, above] {$A(K_0+1)$};
    \draw[->, thick] (x2) -- (dots);
    \draw[->, thick] (dots) -- (xk) node[midway, above] {$A(K-1)$};

    \draw[->, thick] ($(x0)+(0,-1.4)$) -- (x1.south) node[midway, left] {$B(K_0)u(K_0)$};
    \draw[->, thick] ($(x1)+(0,-1.4)$) -- (x2.south) node[midway, left] {$B(K_0+1)u(K_0+1)$};
\end{tikzpicture}
\caption{Interpretazione della dinamica discreta: lo stato iniziale si propaga tramite i prodotti di $A(k)$, mentre ciascun ingresso entra attraverso $B(k)$ e poi viene trasportato in avanti.}
\end{figure}

\section{Caso discreto stazionario}

Consideriamo ora il caso in cui il sistema sia tempo-invariante:
\[
A(k)=A,\qquad B(k)=B,\qquad C(k)=C,\qquad D(k)=D
\qquad \forall k.
\]

Allora le formule precedenti si semplificano molto.

\subsection{Matrice di transizione}

Da \eqref{eq:L4_Phi_TD_generic} segue
\begin{equation}
\Phi(K,K_0)=A^{K-K_0}.
\label{eq:L4_Phi_TD_LTI}
\end{equation}

\subsection{Nucleo forzato}

Da \eqref{eq:L4_H_TD_generic} segue
\begin{equation}
H(K,j)=A^{K-j-1}B.
\label{eq:L4_H_TD_LTI}
\end{equation}

\subsection{Soluzione esplicita}

Quindi la soluzione del sistema discreto stazionario è
\begin{equation}
x(K)=A^{K-K_0}x(K_0)+\sum_{j=K_0}^{K-1}A^{K-j-1}Bu(j).
\label{eq:L4_sol_TD_LTI}
\end{equation}

L'uscita vale
\begin{equation}
y(K)=Cx(K)+Du(K).
\label{eq:L4_out_TD_LTI}
\end{equation}

\paragraph{Messaggio chiave.}
Nel caso discreto stazionario, tutto il comportamento del sistema dipende dallo studio delle potenze di $A$.  
Per questo motivo il corso dedicherà grande attenzione alle proprietà spettrali di $A$.

\section{Caso continuo: dinamica libera e matrice di transizione}

Passiamo ora al tempo continuo:
\begin{equation}
\dot x(t)=A(t)x(t)+B(t)u(t).
\label{eq:L4_sys_TC}
\end{equation}

Per capire come è fatta la matrice di transizione, si parte dall'evoluzione libera, cioè dal caso
\[
u(t)=0.
\]

In questo caso la dinamica diventa
\begin{equation}
\dot x(t)=A(t)x(t),
\qquad x(t_0)=x_0.
\label{eq:L4_free_TC}
\end{equation}

Poiché l'evoluzione libera è descritta da $\Phi(t,t_0)$, deve valere
\[
x(t)=\Phi(t,t_0)x(t_0).
\]

Derivando rispetto a $t$:
\[
\dot x(t)=\frac{\partial \Phi(t,t_0)}{\partial t}x(t_0).
\]

Ma dalla dinamica libera \eqref{eq:L4_free_TC} si ha anche
\[
\dot x(t)=A(t)x(t)=A(t)\Phi(t,t_0)x(t_0).
\]

Poiché la relazione vale per ogni $x(t_0)$, si ottiene l'\emph{equazione differenziale matriciale}:
\begin{equation}
\boxed{
\frac{\partial \Phi(t,t_0)}{\partial t}=A(t)\Phi(t,t_0),
\qquad
\Phi(t_0,t_0)=I.
}
\label{eq:L4_matrix_ode}
\end{equation}

\paragraph{Questa è la caratterizzazione fondamentale di $\Phi$.}
Essa dice che trovare la matrice di transizione equivale a risolvere un problema di Cauchy matriciale.

\section{Successione di Peano}

Per risolvere \eqref{eq:L4_matrix_ode} in modo costruttivo si introduce la \emph{successione di Peano}.

Si cerca $\Phi(t,t_0)$ sotto forma di serie
\begin{equation}
\Phi(t,t_0)=\sum_{k=0}^{\infty} S_k(t,t_0),
\label{eq:L4_peano_series}
\end{equation}
dove i termini sono definiti ricorsivamente da
\begin{equation}
S_0(t,t_0)=I,
\label{eq:L4_S0}
\end{equation}
\begin{equation}
S_1(t,t_0)=\int_{t_0}^{t}A(\tau)S_0(\tau,t_0)\,d\tau
=
\int_{t_0}^{t}A(\tau)\,d\tau,
\label{eq:L4_S1}
\end{equation}
e, più in generale,
\begin{equation}
S_{k+1}(t,t_0)=\int_{t_0}^{t}A(\tau)S_k(\tau,t_0)\,d\tau.
\label{eq:L4_Sk_rec}
\end{equation}

\paragraph{Interpretazione.}
Ogni termine della successione rappresenta un ordine successivo di ``accumulo'' dell'effetto della matrice $A(\tau)$ lungo l'intervallo di integrazione.

\subsection{Primi termini}

Scriviamo i primi termini per capire la struttura:
\[
S_0(t,t_0)=I,
\]
\[
S_1(t,t_0)=\int_{t_0}^{t}A(\tau_1)\,d\tau_1,
\]
\[
S_2(t,t_0)=\int_{t_0}^{t}A(\tau_1)\left(\int_{t_0}^{\tau_1}A(\tau_2)\,d\tau_2\right)d\tau_1,
\]
e così via.

In modo informale, la successione di Peano è l'analogo matriciale della serie iterativa che si ottiene riscrivendo il problema differenziale come equazione integrale.

\section{Stima dei termini della successione di Peano}

Per dimostrare che la serie \eqref{eq:L4_peano_series} converge, serve una stima uniforme sui termini $S_k$.

Sia
\[
M=\sup_{\tau\in[t_0,t]}\norm{A(\tau)}.
\]
Supponiamo che $A(\tau)$ sia limitata sull'intervallo considerato. Allora vale il seguente risultato.

\begin{proposizione}
Per ogni $k\ge 0$,
\begin{equation}
\norm{S_k(t,t_0)}\le \frac{M^k(t-t_0)^k}{k!}.
\label{eq:L4_bound_Sk}
\end{equation}
\end{proposizione}

\begin{proof}
Procediamo per induzione.

Per $k=0$:
\[
S_0(t,t_0)=I,
\]
e la stima vale banalmente, perché
\[
\norm{S_0(t,t_0)}=\norm{I}\le 1.
\]
Nel seguito si assume una norma matriciale compatibile e, se necessario, si assorbe il valore di $\norm{I}$ nella costante.

Per $k=1$:
\[
S_1(t,t_0)=\int_{t_0}^{t}A(\tau)\,d\tau.
\]
Usando disuguaglianza triangolare e definizione di $M$:
\[
\norm{S_1(t,t_0)}
\le
\int_{t_0}^{t}\norm{A(\tau)}\,d\tau
\le
\int_{t_0}^{t}M\,d\tau
=
M(t-t_0).
\]

Supponiamo ora vera la stima per un certo $k$ e dimostriamola per $k+1$:
\[
S_{k+1}(t,t_0)=\int_{t_0}^{t}A(\tau)S_k(\tau,t_0)\,d\tau.
\]
Allora
\begin{align*}
\norm{S_{k+1}(t,t_0)}
&\le
\int_{t_0}^{t}\norm{A(\tau)}\,\norm{S_k(\tau,t_0)}\,d\tau\\
&\le
\int_{t_0}^{t}
M\,
\frac{M^k(\tau-t_0)^k}{k!}
\,d\tau\\
&=
\frac{M^{k+1}}{k!}\int_{t_0}^{t}(\tau-t_0)^k\,d\tau\\
&=
\frac{M^{k+1}}{k!}\cdot \frac{(t-t_0)^{k+1}}{k+1}\\
&=
\frac{M^{k+1}(t-t_0)^{k+1}}{(k+1)!}.
\end{align*}
La tesi è dimostrata.
\end{proof}

\paragraph{Conseguenza.}
La serie di Peano è dominata dalla serie esponenziale
\[
\sum_{k=0}^{\infty}\frac{M^k(t-t_0)^k}{k!}=e^{M(t-t_0)},
\]
quindi converge assolutamente e uniformemente su ogni intervallo limitato.

\section{Verifica che la serie di Peano risolve il problema matriciale}

Una volta definita
\[
\Phi(t,t_0)=\sum_{k=0}^{\infty}S_k(t,t_0),
\]
bisogna verificare che soddisfi davvero \eqref{eq:L4_matrix_ode}.

\subsection{Condizione iniziale}

Valutiamo la serie per $t=t_0$:
\[
\Phi(t_0,t_0)=\sum_{k=0}^{\infty}S_k(t_0,t_0).
\]

Ora:
\[
S_0(t_0,t_0)=I,
\]
mentre per ogni $k\ge 1$ i termini sono integrali su intervalli degeneri, quindi
\[
S_k(t_0,t_0)=0.
\]

Ne segue
\begin{equation}
\Phi(t_0,t_0)=I.
\label{eq:L4_Phi_initial}
\end{equation}

\subsection{Equazione differenziale}

Dalla definizione ricorsiva:
\[
S_k(t,t_0)=\int_{t_0}^{t}A(\tau)S_{k-1}(\tau,t_0)\,d\tau,
\qquad k\ge 1.
\]

Derivando termine a termine:
\[
\frac{\partial S_k(t,t_0)}{\partial t}=A(t)S_{k-1}(t,t_0),\qquad k\ge 1.
\]

Allora
\begin{align*}
\frac{\partial \Phi(t,t_0)}{\partial t}
&=
\sum_{k=1}^{\infty}\frac{\partial S_k(t,t_0)}{\partial t}\\
&=
\sum_{k=1}^{\infty}A(t)S_{k-1}(t,t_0)\\
&=
A(t)\sum_{k=1}^{\infty}S_{k-1}(t,t_0)\\
&=
A(t)\sum_{j=0}^{\infty}S_j(t,t_0)\\
&=
A(t)\Phi(t,t_0).
\end{align*}

Quindi
\begin{equation}
\frac{\partial \Phi(t,t_0)}{\partial t}=A(t)\Phi(t,t_0),
\end{equation}
come volevamo dimostrare.

\section{Espressione di \texorpdfstring{$H(t,\tau)$}{H(t,tau)} nel continuo}

Nella lezione precedente era emersa la proprietà di composizione
\[
\Phi(t,\tau)H(\tau,t_0)=H(t,t_0).
\]

Scegliendo come istante iniziale il tempo $\tau$, si ottiene
\[
H(t,\tau)=\Phi(t,\tau)H(\tau,\tau).
\]

Ma per definizione
\[
H(\tau,\tau)=B(\tau).
\]

Quindi
\begin{equation}
H(t,\tau)=\Phi(t,\tau)B(\tau).
\label{eq:L4_H_TC_from_Phi}
\end{equation}

Sostituendo nella rappresentazione esplicita dello stato:
\begin{equation}
x(t)=\Phi(t,t_0)x(t_0)+\int_{t_0}^{t}\Phi(t,\tau)B(\tau)u(\tau)\,d\tau.
\label{eq:L4_explicit_TC_general}
\end{equation}

Questa è la formula generale della soluzione dei sistemi lineari a tempo continuo.

\section{Caso continuo stazionario}

Supponiamo ora che il sistema sia tempo-invariante:
\[
A(t)=A,\qquad B(t)=B,\qquad C(t)=C,\qquad D(t)=D.
\]

In questo caso la ricorsione di Peano si semplifica moltissimo.

\subsection{Termini della successione di Peano}

Poiché $A(\tau)=A$ è costante, i primi termini diventano:
\[
S_0(t,t_0)=I,
\]
\[
S_1(t,t_0)=\int_{t_0}^{t}A\,d\tau=A(t-t_0),
\]
\[
S_2(t,t_0)=\int_{t_0}^{t}A\,S_1(\tau,t_0)\,d\tau
=
\int_{t_0}^{t}A\cdot A(\tau-t_0)\,d\tau
=
A^2\frac{(t-t_0)^2}{2!}.
\]

Procedendo per induzione si ottiene
\begin{equation}
S_k(t,t_0)=\frac{A^k(t-t_0)^k}{k!}.
\label{eq:L4_Peano_TI_term}
\end{equation}

\subsection{Matrice di transizione}

Allora la serie di Peano diventa
\[
\Phi(t,t_0)=\sum_{k=0}^{\infty}\frac{A^k(t-t_0)^k}{k!}.
\]

Questa è precisamente la definizione di \emph{esponenziale di matrice}:
\begin{equation}
\boxed{
\Phi(t,t_0)=e^{A(t-t_0)}.
}
\label{eq:L4_Phi_exp_matrix}
\end{equation}

\paragraph{Osservazione fondamentale.}
L'esponenziale di matrice gioca, per i sistemi lineari continui, lo stesso ruolo che le potenze di $A$ giocano per i sistemi lineari discreti stazionari.

\subsection{Nucleo forzato}

Da \eqref{eq:L4_H_TC_from_Phi} segue
\[
H(t,\tau)=\Phi(t,\tau)B=e^{A(t-\tau)}B.
\]

\subsection{Soluzione completa}

La soluzione del sistema lineare continuo stazionario è quindi
\begin{equation}
x(t)=e^{A(t-t_0)}x(t_0)+\int_{t_0}^{t}e^{A(t-\tau)}Bu(\tau)\,d\tau.
\label{eq:L4_sol_TC_LTI}
\end{equation}

L'uscita vale
\begin{equation}
y(t)=Cx(t)+Du(t).
\label{eq:L4_out_TC_LTI}
\end{equation}

\begin{figure}[H]
\centering
\begin{tikzpicture}[>=Latex]
    \draw[->] (0,0) -- (9.8,0) node[right] {$\tau$};
    \draw[->] (0,0) -- (0,2.6) node[above] {};
    \draw[dashed] (1.5,0) -- (1.5,2.1) node[above] {$t_0$};
    \draw[dashed] (8.2,0) -- (8.2,2.1) node[above] {$t$};

    \draw[thick,blue] (2.3,0.6) -- (2.3,1.3);
    \draw[thick,blue] (4.2,0.6) -- (4.2,1.6);
    \draw[thick,blue] (6.1,0.6) -- (6.1,1.1);

    \node[blue] at (2.3,1.55) {$u(\tau_1)$};
    \node[blue] at (4.2,1.85) {$u(\tau_2)$};
    \node[blue] at (6.1,1.35) {$u(\tau_3)$};

    \draw[->, thick,red] (2.3,1.3) to[bend left=15] (8.0,2.2);
    \draw[->, thick,red] (4.2,1.6) to[bend left=10] (8.0,2.0);
    \draw[->, thick,red] (6.1,1.1) to[bend left=5]  (8.0,1.8);

    \node[red] at (8.8,2.25) {$x(t)$};
\end{tikzpicture}
\caption{Interpretazione della convoluzione nel continuo: ogni ingresso passato $u(\tau)$ contribuisce allo stato finale $x(t)$ attraverso il peso $\Phi(t,\tau)B$.}
\end{figure}

\section{Confronto finale tra continuo e discreto}

A questo punto il parallelo tra i due casi è completo.

\subsection{Tempo discreto stazionario}

\[
x(K)=A^{K-K_0}x(K_0)+\sum_{j=K_0}^{K-1}A^{K-j-1}Bu(j),
\]
\[
y(K)=Cx(K)+Du(K).
\]

\subsection{Tempo continuo stazionario}

\[
x(t)=e^{A(t-t_0)}x(t_0)+\int_{t_0}^{t}e^{A(t-\tau)}Bu(\tau)\,d\tau,
\]
\[
y(t)=Cx(t)+Du(t).
\]

\paragraph{Analogia concettuale.}
\begin{itemize}
    \item Nel discreto l'evoluzione libera è governata da $A^k$.
    \item Nel continuo l'evoluzione libera è governata da $e^{At}$.
    \item Nel discreto la risposta forzata è una somma.
    \item Nel continuo la risposta forzata è un integrale di convoluzione.
\end{itemize}

\section{Sintesi della lezione}

In questa lezione abbiamo visto che:

\begin{enumerate}
    \item nel tempo discreto, iterando la dinamica
    \[
    x(k+1)=A(k)x(k)+B(k)u(k),
    \]
    si ottengono esplicitamente
    \[
    \Phi(K,K_0)=\prod_{i=K_0}^{K-1}A(i),
    \qquad
    H(K,j)=\left(\prod_{i=j+1}^{K-1}A(i)\right)B(j);
    \]

    \item nel caso discreto stazionario queste formule si riducono a
    \[
    \Phi(K,K_0)=A^{K-K_0},
    \qquad
    H(K,j)=A^{K-j-1}B;
    \]

    \item nel tempo continuo la matrice di transizione è la soluzione del problema
    \[
    \dot \Phi(t,t_0)=A(t)\Phi(t,t_0),
    \qquad
    \Phi(t_0,t_0)=I;
    \]

    \item tale soluzione si costruisce tramite la successione di Peano;

    \item nel caso tempo-invariante la serie di Peano coincide con l'esponenziale di matrice:
    \[
    \Phi(t,t_0)=e^{A(t-t_0)};
    \]

    \item la soluzione completa del sistema continuo stazionario è
    \[
    x(t)=e^{A(t-t_0)}x(t_0)+\int_{t_0}^{t}e^{A(t-\tau)}Bu(\tau)\,d\tau.
    \]
\end{enumerate}

Questi risultati sono fondamentali: da qui in avanti tutto lo studio dei sistemi lineari ruoterà intorno a
\begin{itemize}
    \item potenze di matrici nel discreto;
    \item esponenziale di matrice nel continuo;
    \item struttura spettrale della matrice $A$.
\end{itemize}
\chapter{L5 -- Proprietà dell'esponenziale di matrice, funzioni di matrice, sequenze di Krylov, sottospazi ciclici e forma di Jordan}

\section{Esponenziale di matrice}

Nel caso dei sistemi lineari tempo-invarianti continui, la soluzione libera è
\[
x(t)=e^{A(t-t_0)}x(t_0).
\]
Per questo motivo è fondamentale capire bene che cosa sia l'esponenziale di matrice e quali proprietà possieda.

\subsection{Definizione}

Data una matrice quadrata
\[
A\in\mathbb{R}^{n\times n}
\qquad
(\text{o più in generale }A\in\mathbb{C}^{n\times n}),
\]
si definisce l'esponenziale di matrice come la serie
\begin{equation}
e^{At}
:=
\sum_{k=0}^{\infty}\frac{A^k t^k}{k!}.
\label{eq:L5_def_exp}
\end{equation}

Più in generale, fissato un istante iniziale $t_0$,
\begin{equation}
\Phi(t,t_0)=e^{A(t-t_0)}
=
\sum_{k=0}^{\infty}\frac{A^k (t-t_0)^k}{k!}.
\label{eq:L5_phi_exp}
\end{equation}

\paragraph{Osservazione.}
Questa definizione è perfettamente analoga a quella dell'esponenziale scalare:
\[
e^z=\sum_{k=0}^{\infty}\frac{z^k}{k!}.
\]
La differenza è che ora, al posto di una variabile scalare, compare una matrice.

\section{Proprietà fondamentali dell'esponenziale di matrice}

Le proprietà che seguono sono quelle emerse negli appunti e sono centrali per tutto il resto del corso.

\begin{proposizione}
Per ogni matrice quadrata $A$ valgono le seguenti proprietà.
\end{proposizione}

\subsection{Valore all'origine}

Ponendo $t=t_0$ in \eqref{eq:L5_phi_exp}, si ha
\[
\Phi(t_0,t_0)=e^{A(t_0-t_0)}=e^{A\cdot 0}=e^{0}.
\]
Poiché nella serie tutti i termini con $k\ge 1$ si annullano, resta solo il termine di ordine zero:
\[
e^{0}=I.
\]

Quindi
\begin{equation}
e^{A\cdot 0}=I,
\qquad
\Phi(t_0,t_0)=I.
\label{eq:L5_exp_identity}
\end{equation}

\subsection{Derivata rispetto al tempo}

Derivando termine a termine la serie:
\[
\frac{d}{dt}e^{A(t-t_0)}
=
\frac{d}{dt}\left(
\sum_{k=0}^{\infty}\frac{A^k (t-t_0)^k}{k!}
\right).
\]

Si ottiene
\[
=
\sum_{k=1}^{\infty}\frac{A^k k (t-t_0)^{k-1}}{k!}
=
\sum_{k=1}^{\infty}\frac{A^k (t-t_0)^{k-1}}{(k-1)!}.
\]

Raccogliendo una $A$:
\[
=
A\sum_{k=1}^{\infty}\frac{A^{k-1}(t-t_0)^{k-1}}{(k-1)!}
=
A\sum_{j=0}^{\infty}\frac{A^j (t-t_0)^j}{j!}.
\]

Dunque
\begin{equation}
\frac{d}{dt}e^{A(t-t_0)}=A\,e^{A(t-t_0)}.
\label{eq:L5_exp_derivative}
\end{equation}

Poiché ogni potenza di $A$ commuta con $A$, si può anche scrivere
\[
\frac{d}{dt}e^{A(t-t_0)}=e^{A(t-t_0)}A.
\]

\subsection{Non singolarità}

La matrice esponenziale è sempre invertibile. Infatti
\[
e^{At}e^{-At}=e^{0}=I.
\]

Quindi
\begin{equation}
\left(e^{At}\right)^{-1}=e^{-At}.
\label{eq:L5_exp_inverse}
\end{equation}

In particolare, $e^{At}$ non è mai singolare, per qualunque matrice $A$ e per qualunque $t$.

\subsection{Proprietà di composizione}

Poiché la matrice $A$ commuta con sé stessa, si ha
\[
e^{A(t_2-t_1)}e^{A(t_1-t_0)}=e^{A[(t_2-t_1)+(t_1-t_0)]}=e^{A(t_2-t_0)}.
\]

Quindi
\begin{equation}
e^{A(t_2-t_0)}=e^{A(t_2-t_1)}e^{A(t_1-t_0)}.
\label{eq:L5_exp_composition}
\end{equation}

Questa è esattamente la proprietà di composizione della matrice di transizione:
\[
\Phi(t_2,t_0)=\Phi(t_2,t_1)\Phi(t_1,t_0).
\]

\paragraph{Osservazione importante.}
La formula
\[
e^{X+Y}=e^Xe^Y
\]
\emph{non} è vera in generale per matrici arbitrarie $X$ e $Y$. È vera qui perché stiamo sommando multipli della \emph{stessa} matrice $A$, e quindi i termini commutano.

\subsection{Invarianza per similitudine}

Se $S$ è invertibile, allora
\begin{equation}
e^{(S^{-1}AS)t}=S^{-1}e^{At}S.
\label{eq:L5_similarity_exp}
\end{equation}

\begin{proof}
Basta sviluppare la serie:
\[
e^{(S^{-1}AS)t}
=
\sum_{k=0}^{\infty}\frac{(S^{-1}AS)^k t^k}{k!}.
\]

Ma
\[
(S^{-1}AS)^k=S^{-1}A^kS,
\]
per ogni $k\ge 0$. Quindi
\[
e^{(S^{-1}AS)t}
=
\sum_{k=0}^{\infty}\frac{S^{-1}A^kS\,t^k}{k!}
=
S^{-1}\left(\sum_{k=0}^{\infty}\frac{A^k t^k}{k!}\right)S
=
S^{-1}e^{At}S.
\]
\end{proof}

\paragraph{Significato.}
Questa proprietà è essenziale perché, se riesco a portare una matrice in una forma più semplice mediante similitudine, posso calcolare l'esponenziale nella forma semplice e poi tornare alla matrice originale.

\section{Funzioni di matrice e matrici polinomiali}

L'esponenziale di matrice è un esempio di \emph{funzione di matrice}. Prima di affrontare i casi pratici di calcolo, è utile introdurre le funzioni polinomiali di una matrice.

\subsection{Polinomio e matrice polinomiale associata}

Sia
\[
p(\lambda)=a_m\lambda^m+a_{m-1}\lambda^{m-1}+\cdots+a_1\lambda+a_0
\]
un polinomio. Alla matrice $A\in\mathbb{R}^{n\times n}$ si associa la matrice polinomiale
\begin{equation}
p(A)=a_mA^m+a_{m-1}A^{m-1}+\cdots+a_1A+a_0I.
\label{eq:L5_poly_matrix}
\end{equation}

\paragraph{Caso monico.}
Spesso si suppone che il polinomio sia \emph{monico}, cioè
\[
a_m=1.
\]
In tal caso
\[
p(\lambda)=\lambda^m+a_{m-1}\lambda^{m-1}+\cdots+a_1\lambda+a_0.
\]

\subsection{Fattorizzazione del polinomio}

Se il polinomio si scinde su $\mathbb{C}$, allora esistono radici $\lambda_1,\dots,\lambda_m$ tali che
\[
p(\lambda)=\prod_{i=1}^{m}(\lambda-\lambda_i).
\]

Di conseguenza si può scrivere
\begin{equation}
p(A)=\prod_{i=1}^{m}(A-\lambda_i I).
\label{eq:L5_factorized_poly_matrix}
\end{equation}

\section{Autovalori di una matrice polinomiale}

Una proprietà molto utile è la seguente.

\begin{proposizione}
Se $\lambda$ è un autovalore di $A$ con autovettore $v\neq 0$, allora $p(\lambda)$ è un autovalore di $p(A)$ con lo stesso autovettore $v$.
\end{proposizione}

\begin{proof}
Se
\[
Av=\lambda v,
\]
allora
\[
A^2v=A(Av)=A(\lambda v)=\lambda Av=\lambda^2 v.
\]
Per induzione si ottiene
\[
A^kv=\lambda^k v\qquad \forall k\ge 0.
\]

Allora
\begin{align*}
p(A)v
&=
(a_mA^m+a_{m-1}A^{m-1}+\cdots+a_1A+a_0I)v\\
&=
a_mA^mv+a_{m-1}A^{m-1}v+\cdots+a_1Av+a_0v\\
&=
a_m\lambda^m v+a_{m-1}\lambda^{m-1}v+\cdots+a_1\lambda v+a_0v\\
&=
p(\lambda)v.
\end{align*}

Quindi $v$ è autovettore di $p(A)$ e l'autovalore associato è proprio $p(\lambda)$.
\end{proof}

\paragraph{Osservazione.}
Questa proprietà è una forma elementare della \emph{mappatura spettrale} per i polinomi.

\section{Sequenza di Krylov}

Fissiamo uno stato iniziale
\[
x_0\in X.
\]

Si chiama \emph{sequenza di Krylov generata da $x_0$ rispetto ad $A$} la successione di vettori
\begin{equation}
x_0,\ Ax_0,\ A^2x_0,\ A^3x_0,\ \dots
\label{eq:L5_krylov_sequence}
\end{equation}

\begin{figure}[H]
\centering
\begin{tikzpicture}[>=Latex, node distance=1.55cm]
    \node[draw, circle, minimum size=9mm] (x0) {$x_0$};
    \node[draw, circle, right=of x0, minimum size=9mm] (x1) {$Ax_0$};
    \node[draw, circle, right=of x1, minimum size=9mm] (x2) {$A^2x_0$};
    \node[right=of x2] (dots) {$\cdots$};
    \draw[->, thick] (x0) -- node[above] {$A$} (x1);
    \draw[->, thick] (x1) -- node[above] {$A$} (x2);
    \draw[->, thick] (x2) -- (dots);
\end{tikzpicture}
\caption{Schema della sequenza di Krylov generata da $x_0$.}
\end{figure}

Poiché lo spazio di stato $X$ ha dimensione finita, non è possibile avere infiniti vettori linearmente indipendenti.  
Dunque esiste un primo indice $p$ tale che
\[
A^p x_0
\]
sia combinazione lineare dei precedenti:
\[
x_0,\ Ax_0,\ \dots,\ A^{p-1}x_0.
\]

Di conseguenza, per ogni $q\ge p$, anche $A^q x_0$ sarà combinazione lineare dei primi $p$ vettori.

\section{Sottospazio ciclico}

Si definisce \emph{sottospazio ciclico generato da $x_0$} il sottospazio
\begin{equation}
S_{x_0}:=\operatorname{span}\{x_0,Ax_0,A^2x_0,\dots,A^{p-1}x_0\},
\label{eq:L5_cyclic_subspace}
\end{equation}
dove $p$ è il numero di vettori linearmente indipendenti nella sequenza di Krylov.

\paragraph{Interpretazione.}
Questo sottospazio è il ``luogo naturale'' in cui vive l'evoluzione libera che parte da $x_0$.

\subsection{Esempio 1: autovettore}

Se $x_0=v$ è autovettore di $A$:
\[
Av=\lambda v,
\]
allora
\[
A^k v=\lambda^k v,
\]
quindi tutti i vettori della sequenza di Krylov sono multipli di $v$.  
In tal caso
\[
S_v=\operatorname{span}\{v\},
\qquad
\dim S_v=1.
\]

\subsection{Esempio 2: dimensione maggiore di uno}

Sia
\[
A=
\begin{pmatrix}
1 & 1\\
0 & -1
\end{pmatrix},
\qquad
x_0=
\begin{pmatrix}
0\\
1
\end{pmatrix}.
\]

Allora
\[
Ax_0=
\begin{pmatrix}
1\\
-1
\end{pmatrix},
\]
che non è multiplo di $x_0$. Quindi
\[
S_{x_0}=\operatorname{span}\left\{
\begin{pmatrix}
0\\1
\end{pmatrix},
\begin{pmatrix}
1\\-1
\end{pmatrix}
\right\}
=\mathbb{R}^2.
\]

In questo caso il sottospazio ciclico coincide con tutto lo spazio di stato.

\section{Sottospazi \texorpdfstring{$A$}{A}-invarianti}

Un sottospazio $S\subseteq X$ si dice \emph{$A$-invariante} se
\begin{equation}
x\in S \quad \Longrightarrow \quad Ax\in S.
\label{eq:L5_A_invariant_def}
\end{equation}

\paragraph{Significato dinamico.}
Se l'evoluzione libera parte da uno stato contenuto in un sottospazio $A$-invariante, allora tutta l'evoluzione libera rimane in quel sottospazio.

\subsection{Esempio di sottospazio non invariante}

Sia
\[
A=
\begin{pmatrix}
1 & 1\\
0 & -1
\end{pmatrix},
\qquad
S=\operatorname{span}\left\{
\begin{pmatrix}
1\\
1
\end{pmatrix}
\right\}.
\]

Un vettore generico di $S$ è
\[
v=
\begin{pmatrix}
\alpha\\
\alpha
\end{pmatrix}.
\]

Allora
\[
Av=
\begin{pmatrix}
1 & 1\\
0 & -1
\end{pmatrix}
\begin{pmatrix}
\alpha\\
\alpha
\end{pmatrix}
=
\begin{pmatrix}
2\alpha\\
-\alpha
\end{pmatrix},
\]
che in generale non appartiene a $S$, perché non ha componenti uguali.

Quindi $S$ non è $A$-invariante.

\section{Il sottospazio ciclico è \texorpdfstring{$A$}{A}-invariante}

Questo è uno dei risultati centrali della lezione.

\begin{teorema}
Dato $A$ e fissato $x_0$, il sottospazio ciclico $S_{x_0}$ è $A$-invariante.
\end{teorema}

\begin{proof}
Sia
\[
S_{x_0}=\operatorname{span}\{x_0,Ax_0,\dots,A^{p-1}x_0\}.
\]
Prendiamo un vettore arbitrario $x\in S_{x_0}$. Allora esistono coefficienti $\alpha_0,\dots,\alpha_{p-1}$ tali che
\[
x=\sum_{i=0}^{p-1}\alpha_i A^i x_0.
\]

Moltiplicando per $A$:
\[
Ax=\sum_{i=0}^{p-1}\alpha_i A^{i+1}x_0
=
\alpha_0Ax_0+\alpha_1A^2x_0+\cdots+\alpha_{p-2}A^{p-1}x_0+\alpha_{p-1}A^p x_0.
\]

Ora, per definizione di $p$, il vettore $A^p x_0$ è combinazione lineare dei precedenti:
\[
A^p x_0 \in \operatorname{span}\{x_0,Ax_0,\dots,A^{p-1}x_0\}=S_{x_0}.
\]

Di conseguenza anche $Ax$ è combinazione lineare di vettori di $S_{x_0}$, quindi
\[
Ax\in S_{x_0}.
\]

Pertanto $S_{x_0}$ è $A$-invariante.
\end{proof}

\paragraph{Conseguenza dinamica.}
Se l'evoluzione libera parte da $x_0$, allora resta tutta dentro $S_{x_0}$. Questo riduce il problema dinamico a un sottospazio spesso di dimensione molto più piccola di $n$.

\section{Polinomi annullatori}

\subsection{Per la matrice}

Un polinomio $p(\lambda)$ si dice \emph{annullatore} della matrice $A$ se
\begin{equation}
p(A)=0.
\label{eq:L5_annihilating_poly}
\end{equation}

Questo significa che la matrice polinomiale associata è la matrice nulla.

\subsection{Per un vettore}

Più in generale, rispetto a un dato vettore $x_0$, si può chiedere che
\begin{equation}
p(A)x_0=0.
\label{eq:L5_annihilating_poly_vector}
\end{equation}

In questo caso $p$ annulla il vettore $x_0$ attraverso la dinamica indotta da $A$.

\paragraph{Osservazione.}
Il minimo grado di un polinomio monico che annulla $x_0$ è strettamente collegato alla dimensione del sottospazio ciclico generato da $x_0$.

\section{Teorema di Cayley--Hamilton}

Il teorema chiave è il seguente.

\begin{teorema}[Cayley--Hamilton]
Sia
\[
p_A(\lambda)=\det(\lambda I-A)
\]
il polinomio caratteristico di $A$. Allora
\begin{equation}
p_A(A)=0.
\label{eq:L5_CH}
\end{equation}
\end{teorema}

Se
\[
p_A(\lambda)=\lambda^n+a_{n-1}\lambda^{n-1}+\cdots+a_1\lambda+a_0,
\]
allora
\begin{equation}
A^n+a_{n-1}A^{n-1}+\cdots+a_1A+a_0I=0.
\label{eq:L5_CH_expanded}
\end{equation}

\paragraph{Conseguenza fondamentale.}
Ogni potenza $A^n$ si può riscrivere come combinazione lineare delle potenze precedenti:
\[
I,\ A,\ A^2,\ \dots,\ A^{n-1}.
\]

Questo è uno dei motivi per cui lo studio delle funzioni di matrice può essere ricondotto a uno spazio finito-dimensionale.

\subsection{Esempio 1}

Sia
\[
A=
\begin{pmatrix}
1 & 1\\
0 & -1
\end{pmatrix}.
\]

Poiché è triangolare superiore, gli autovalori sono gli elementi sulla diagonale:
\[
\lambda_1=1,\qquad \lambda_2=-1.
\]

Il polinomio caratteristico è
\[
p_A(\lambda)=(\lambda-1)(\lambda+1)=\lambda^2-1.
\]

Quindi Cayley--Hamilton dice che
\[
p_A(A)=A^2-I=0.
\]

Infatti
\[
A^2=
\begin{pmatrix}
1 & 1\\
0 & -1
\end{pmatrix}^2
=
\begin{pmatrix}
1 & 0\\
0 & 1
\end{pmatrix}
=I.
\]

\subsection{Esempio 2}

Sia
\[
A=
\begin{pmatrix}
1 & -1\\
1 & -1
\end{pmatrix}.
\]

Si ha
\[
\det(A)=0,
\qquad
\operatorname{tr}(A)=0.
\]
Quindi il polinomio caratteristico è
\[
p_A(\lambda)=\lambda^2.
\]

Pertanto
\[
p_A(A)=A^2=0.
\]

In questo caso il polinomio caratteristico è anche annullatore.

\paragraph{Conseguenza.}
Se una potenza di $A$ si annulla o si esprime come combinazione delle precedenti, allora anche le successive si esprimono nello stesso modo. Questo semplifica enormemente il calcolo di $e^{At}$.

\section{Evoluzione libera e sottospazi ciclici}

Nel tempo discreto, in assenza di ingresso, l'evoluzione libera è
\[
x(k)=A^k x_0.
\]

Nel tempo continuo, sempre in assenza di ingresso, è
\[
x(t)=e^{At}x_0.
\]

L'idea chiave è che, se conosco bene la struttura del sottospazio ciclico generato da $x_0$, posso comprendere dove vive l'evoluzione libera.

Infatti:
\begin{itemize}
    \item nel discreto, $A^k x_0$ appartiene per definizione alla sequenza di Krylov;
    \item nel continuo, espandendo in serie
    \[
    e^{At}x_0=\sum_{k=0}^{\infty}\frac{A^k t^k}{k!}x_0,
    \]
    si vede che anche qui l'evoluzione è costruita con combinazioni lineari dei vettori della sequenza di Krylov.
\end{itemize}

Quindi il sottospazio ciclico contiene l'intera evoluzione libera che parte da $x_0$.

\section{Casi in cui è facile calcolare l'esponenziale}

Dopo aver introdotto le proprietà generali, gli appunti si concentrano sui casi in cui l'esponenziale di matrice si calcola facilmente.

\subsection{Caso 1: matrice diagonale}

Sia
\[
A=
\begin{pmatrix}
\lambda_1 & 0 & \cdots & 0\\
0 & \lambda_2 & \cdots & 0\\
\vdots & \vdots & \ddots & \vdots\\
0 & 0 & \cdots & \lambda_n
\end{pmatrix}.
\]

Allora
\[
A^k=
\begin{pmatrix}
\lambda_1^k & 0 & \cdots & 0\\
0 & \lambda_2^k & \cdots & 0\\
\vdots & \vdots & \ddots & \vdots\\
0 & 0 & \cdots & \lambda_n^k
\end{pmatrix}.
\]

Quindi
\begin{align}
e^{At}
&=
\sum_{k=0}^{\infty}\frac{A^k t^k}{k!}\nonumber\\
&=
\begin{pmatrix}
\sum_{k=0}^{\infty}\dfrac{\lambda_1^k t^k}{k!} & 0 & \cdots & 0\\
0 & \sum_{k=0}^{\infty}\dfrac{\lambda_2^k t^k}{k!} & \cdots & 0\\
\vdots & \vdots & \ddots & \vdots\\
0 & 0 & \cdots & \sum_{k=0}^{\infty}\dfrac{\lambda_n^k t^k}{k!}
\end{pmatrix}\nonumber\\
&=
\begin{pmatrix}
e^{\lambda_1 t} & 0 & \cdots & 0\\
0 & e^{\lambda_2 t} & \cdots & 0\\
\vdots & \vdots & \ddots & \vdots\\
0 & 0 & \cdots & e^{\lambda_n t}
\end{pmatrix}.
\label{eq:L5_exp_diagonal}
\end{align}

\paragraph{Conclusione.}
Se la matrice è diagonale, basta fare l'esponenziale degli elementi diagonali.

\subsection{Caso 2: matrice diagonalizzabile}

Supponiamo che esista una matrice invertibile $S$ tale che
\[
SAS^{-1}=D,
\]
dove $D$ è diagonale.

Equivalentemente,
\[
A=S^{-1}DS.
\]

Allora, usando l'invarianza per similitudine,
\begin{equation}
e^{At}=S^{-1}e^{Dt}S.
\label{eq:L5_exp_diagonalizable}
\end{equation}

Poiché $e^{Dt}$ si calcola facilmente tramite \eqref{eq:L5_exp_diagonal}, si conclude che l'esponenziale di qualunque matrice diagonalizzabile si ottiene:
\begin{enumerate}
    \item diagonalizzando la matrice;
    \item facendo l'esponenziale della diagonale;
    \item tornando indietro mediante la trasformazione di similitudine.
\end{enumerate}

\paragraph{Osservazione.}
Questo metodo funziona molto bene, ma solo se la matrice è davvero diagonalizzabile.

\subsection{Caso 3: matrice simile a una matrice di Jordan}

Non tutte le matrici sono diagonalizzabili. Tuttavia, su $\mathbb{C}$, ogni matrice è simile a una matrice in forma di Jordan.

\begin{teorema}[Forma di Jordan]
Per ogni matrice quadrata $A$ esiste una matrice invertibile $S$ tale che
\[
SAS^{-1}=J,
\]
dove $J$ è una matrice a blocchi di Jordan.
\end{teorema}

Una matrice di Jordan ha la forma
\[
J=
\begin{pmatrix}
J_{\lambda_1}^{(p_1)} & & & 0\\
& J_{\lambda_2}^{(p_2)} & & \\
& & \ddots & \\
0 & & & J_{\lambda_r}^{(p_r)}
\end{pmatrix},
\]
dove ogni blocco di Jordan è del tipo
\begin{equation}
J_\lambda^{(p)}=
\begin{pmatrix}
\lambda & 1 & 0 & \cdots & 0\\
0 & \lambda & 1 & \cdots & 0\\
0 & 0 & \lambda & \ddots & \vdots\\
\vdots & \vdots & \ddots & \ddots & 1\\
0 & 0 & \cdots & 0 & \lambda
\end{pmatrix}.
\label{eq:L5_jordan_block}
\end{equation}

\begin{figure}[H]
\centering
\begin{tikzpicture}[scale=1]
    \draw (0,0) rectangle (2.7,2.2);
    \draw (3.2,0.4) rectangle (5.2,2.2);
    \draw (5.7,0.8) rectangle (8.7,2.2);

    \node at (1.35,1.1) {$J_{\lambda_1}^{(p_1)}$};
    \node at (4.2,1.3) {$J_{\lambda_2}^{(p_2)}$};
    \node at (7.2,1.5) {$J_{\lambda_r}^{(p_r)}$};

    \node at (9.4,1.1) {$\cdots$};
\end{tikzpicture}
\caption{Schema di una matrice di Jordan a blocchi diagonali.}
\end{figure}

\paragraph{Osservazione.}
Gli autovalori della matrice compaiono sulla diagonale dei blocchi, ripetuti secondo la molteplicità.

\section{Esponenziale di una matrice in forma di Jordan}

Se
\[
SAS^{-1}=J,
\]
allora
\[
e^{At}=S^{-1}e^{Jt}S.
\]

Poiché $J$ è a blocchi diagonali, si ha
\begin{equation}
e^{Jt}=
\begin{pmatrix}
e^{J_{\lambda_1}^{(p_1)}t} & & & 0\\
& e^{J_{\lambda_2}^{(p_2)}t} & & \\
& & \ddots & \\
0 & & & e^{J_{\lambda_r}^{(p_r)}t}
\end{pmatrix}.
\label{eq:L5_exp_Jordan_blocks}
\end{equation}

Quindi il problema si riduce a calcolare l'esponenziale di un singolo blocco di Jordan.

\subsection{Esponenziale di un blocco di Jordan}

Scriviamo
\[
J_\lambda^{(p)}=\lambda I + N,
\]
dove $N$ è nilpotente e ha la forma
\[
N=
\begin{pmatrix}
0 & 1 & 0 & \cdots & 0\\
0 & 0 & 1 & \cdots & 0\\
0 & 0 & 0 & \ddots & \vdots\\
\vdots & \vdots & \vdots & \ddots & 1\\
0 & 0 & 0 & \cdots & 0
\end{pmatrix},
\qquad
N^p=0.
\]

Poiché $\lambda I$ e $N$ commutano, si ha
\[
e^{J_\lambda^{(p)}t}
=
e^{(\lambda I+N)t}
=
e^{\lambda t}e^{Nt}.
\]

Ma, essendo $N$ nilpotente, la serie di $e^{Nt}$ si tronca:
\[
e^{Nt}=
I+Nt+\frac{N^2t^2}{2!}+\cdots+\frac{N^{p-1}t^{p-1}}{(p-1)!}.
\]

Quindi
\begin{equation}
e^{J_\lambda^{(p)}t}
=
e^{\lambda t}
\left(
I+Nt+\frac{N^2t^2}{2!}+\cdots+\frac{N^{p-1}t^{p-1}}{(p-1)!}
\right).
\label{eq:L5_exp_Jordan_block_compact}
\end{equation}

Esplicitamente:
\begin{equation}
e^{J_\lambda^{(p)}t}
=
e^{\lambda t}
\begin{pmatrix}
1 & t & \dfrac{t^2}{2!} & \cdots & \dfrac{t^{p-1}}{(p-1)!}\\
0 & 1 & t & \cdots & \dfrac{t^{p-2}}{(p-2)!}\\
0 & 0 & 1 & \ddots & \vdots\\
\vdots & \vdots & \ddots & \ddots & t\\
0 & 0 & \cdots & 0 & 1
\end{pmatrix}.
\label{eq:L5_exp_Jordan_block_explicit}
\end{equation}

\paragraph{Conclusione pratica.}
Per calcolare l'esponenziale di una matrice generica:
\begin{enumerate}
    \item la porto in forma di Jordan;
    \item calcolo l'esponenziale di ciascun blocco;
    \item ricompongo i blocchi;
    \item torno alla matrice originale tramite similitudine.
\end{enumerate}

\section{Collegamento con la dinamica dei sistemi lineari}

Per i sistemi lineari continui in evoluzione libera:
\[
\dot x(t)=Ax(t),\qquad u(t)=0,
\]
la soluzione è
\[
x(t)=e^{A(t-t_0)}x_0.
\]

Quindi tutto lo studio della dinamica libera si riduce al calcolo e alla comprensione di $e^{At}$.

Gli strumenti introdotti in questa lezione mostrano che:
\begin{itemize}
    \item se $A$ è diagonale, il problema è immediato;
    \item se $A$ è diagonalizzabile, il problema si riduce al caso diagonale;
    \item se $A$ non è diagonalizzabile, si passa alla forma di Jordan.
\end{itemize}

\section{Sintesi finale della lezione}

In questa lezione abbiamo visto che:

\begin{enumerate}
    \item l'esponenziale di matrice si definisce come serie
    \[
    e^{At}=\sum_{k=0}^{\infty}\frac{A^k t^k}{k!};
    \]

    \item esso soddisfa proprietà analoghe all'esponenziale scalare:
    \[
    e^{0}=I,\qquad
    \frac{d}{dt}e^{At}=Ae^{At},\qquad
    (e^{At})^{-1}=e^{-At};
    \]

    \item l'esponenziale è compatibile con la similitudine:
    \[
    e^{(S^{-1}AS)t}=S^{-1}e^{At}S;
    \]

    \item un polinomio $p$ applicato a una matrice si definisce come
    \[
    p(A)=a_mA^m+\cdots+a_1A+a_0I;
    \]

    \item se $\lambda$ è autovalore di $A$, allora $p(\lambda)$ è autovalore di $p(A)$;

    \item la sequenza di Krylov
    \[
    x_0,\ Ax_0,\ A^2x_0,\dots
    \]
    genera il sottospazio ciclico $S_{x_0}$, che è sempre $A$-invariante;

    \item il teorema di Cayley--Hamilton implica che le potenze alte di $A$ si possono scrivere come combinazioni lineari delle precedenti;

    \item l'esponenziale di una matrice diagonale si calcola diagonalmente, quello di una matrice diagonalizzabile si ricava per similitudine, e il caso generale si affronta tramite forma di Jordan.
\end{enumerate}

Questi risultati saranno la base per lo studio più fine dell'analisi modale, delle molteplicità algebriche e geometriche, e del legame tra struttura spettrale della matrice $A$ e andamento temporale delle soluzioni.
\chapter{L6 -- Blocchi di Jordan, autovalori complessi e autovettori generalizzati}

\section{Esponenziale di un blocco di Jordan nilpotente}

Nella lezione precedente abbiamo visto che, se una matrice è simile a una forma di Jordan, allora il calcolo di $e^{At}$ si riduce al calcolo dell'esponenziale dei singoli blocchi di Jordan.  
Conviene quindi partire dal caso più semplice: il blocco di Jordan \emph{nilpotente} associato all'autovalore nullo.

\subsection{Blocco nilpotente elementare}

Indichiamo con
\[
J_0^{(p)}=
\begin{pmatrix}
0 & 1 & 0 & \cdots & 0\\
0 & 0 & 1 & \cdots & 0\\
\vdots & \vdots & \vdots & \ddots & \vdots\\
0 & 0 & 0 & \cdots & 1\\
0 & 0 & 0 & \cdots & 0
\end{pmatrix}\in\mathbb{R}^{p\times p}
\]
il blocco di Jordan nilpotente di ordine $p$.

Per esempio, per $p=3$:
\[
J_0^{(3)}=
\begin{pmatrix}
0 & 1 & 0\\
0 & 0 & 1\\
0 & 0 & 0
\end{pmatrix},
\qquad
\left(J_0^{(3)}\right)^2=
\begin{pmatrix}
0 & 0 & 1\\
0 & 0 & 0\\
0 & 0 & 0
\end{pmatrix},
\qquad
\left(J_0^{(3)}\right)^3=0.
\]

In generale:
\begin{equation}
\left(J_0^{(p)}\right)^p=0.
\label{eq:L6_nilpotent_power_zero}
\end{equation}

Questa proprietà dice che $J_0^{(p)}$ è \emph{nilpotente} di indice $p$.

\subsection{Esponenziale del blocco nilpotente}

Dalla definizione di esponenziale di matrice:
\[
e^{J_0^{(p)}t}
=
\sum_{k=0}^{\infty}\frac{\left(J_0^{(p)}\right)^k t^k}{k!}.
\]

Ma, grazie a \eqref{eq:L6_nilpotent_power_zero}, tutti i termini con $k\ge p$ sono nulli. Quindi la serie si tronca:
\begin{equation}
e^{J_0^{(p)}t}
=
I+\frac{J_0^{(p)}t}{1!}+\frac{\left(J_0^{(p)}\right)^2 t^2}{2!}+\cdots+\frac{\left(J_0^{(p)}\right)^{p-1} t^{p-1}}{(p-1)!}.
\label{eq:L6_exp_nilpotent_jordan}
\end{equation}

Per $p=3$ si ottiene
\[
e^{J_0^{(3)}t}
=
I+J_0^{(3)}t+\frac{\left(J_0^{(3)}\right)^2t^2}{2!}
=
\begin{pmatrix}
1 & t & \dfrac{t^2}{2}\\
0 & 1 & t\\
0 & 0 & 1
\end{pmatrix}.
\]

In generale, l'esponenziale del blocco nilpotente ha la forma
\begin{equation}
e^{J_0^{(p)}t}
=
\begin{pmatrix}
1 & t & \dfrac{t^2}{2!} & \cdots & \dfrac{t^{p-1}}{(p-1)!}\\
0 & 1 & t & \cdots & \dfrac{t^{p-2}}{(p-2)!}\\
0 & 0 & 1 & \ddots & \vdots\\
\vdots & \vdots & \ddots & \ddots & t\\
0 & 0 & \cdots & 0 & 1
\end{pmatrix}.
\label{eq:L6_exp_nilpotent_general}
\end{equation}

\paragraph{Interpretazione dinamica.}
Questo risultato è molto importante: un blocco nilpotente genera fattori polinomiali in $t$.  
Più il blocco è grande, più compaiono potenze alte di $t$ nella soluzione.

\section{Esponenziale di un blocco di Jordan generico}

Consideriamo ora un blocco di Jordan associato a un autovalore reale $d$:
\begin{equation}
J_d^{(p)}=
\begin{pmatrix}
d & 1 & 0 & \cdots & 0\\
0 & d & 1 & \cdots & 0\\
\vdots & \vdots & \vdots & \ddots & \vdots\\
0 & 0 & 0 & \cdots & 1\\
0 & 0 & 0 & \cdots & d
\end{pmatrix}.
\label{eq:L6_jordan_block_real}
\end{equation}

Si può scrivere come
\[
J_d^{(p)}=dI+J_0^{(p)}.
\]

Poiché $dI$ commuta con $J_0^{(p)}$, vale
\[
e^{J_d^{(p)}t}=e^{(dI+J_0^{(p)})t}=e^{dt}e^{J_0^{(p)}t}.
\]

Usando \eqref{eq:L6_exp_nilpotent_general}:
\begin{equation}
e^{J_d^{(p)}t}
=
e^{dt}
\begin{pmatrix}
1 & t & \dfrac{t^2}{2!} & \cdots & \dfrac{t^{p-1}}{(p-1)!}\\
0 & 1 & t & \cdots & \dfrac{t^{p-2}}{(p-2)!}\\
0 & 0 & 1 & \ddots & \vdots\\
\vdots & \vdots & \ddots & \ddots & t\\
0 & 0 & \cdots & 0 & 1
\end{pmatrix}.
\label{eq:L6_exp_jordan_real_block}
\end{equation}

\paragraph{Conseguenza qualitativa.}
L'andamento temporale di un blocco di Jordan reale non è governato soltanto da $e^{dt}$, ma da
\[
e^{dt},\qquad te^{dt},\qquad \frac{t^2}{2!}e^{dt},\ \dots
\]
fino al grado $p-1$.  
Quindi la dimensione del blocco conta, non solo il valore dell'autovalore.

\section{Potenze di un blocco di Jordan nel tempo discreto}

Nel tempo discreto la dinamica libera è
\[
x(k)=A^{k-k_0}x(k_0).
\]
Se $A$ è simile a una matrice di Jordan, allora bisogna capire come è fatto $J^k$.

Per un blocco reale
\[
J_d^{(p)}=dI+J_0^{(p)},
\]
vale il binomio matriciale
\begin{equation}
\left(J_d^{(p)}\right)^k
=
\sum_{i=0}^{p-1}\binom{k}{i}d^{\,k-i}\left(J_0^{(p)}\right)^i,
\label{eq:L6_power_jordan_block}
\end{equation}
dove la somma si ferma a $p-1$ perché $\left(J_0^{(p)}\right)^p=0$.

\paragraph{Caso $p=2$.}
Se
\[
J_d^{(2)}=
\begin{pmatrix}
d & 1\\
0 & d
\end{pmatrix},
\]
allora
\[
\left(J_d^{(2)}\right)^k
=
\begin{pmatrix}
d^k & k d^{k-1}\\
0 & d^k
\end{pmatrix}.
\]

\paragraph{Caso $p=3$.}
Se
\[
J_d^{(3)}=
\begin{pmatrix}
d & 1 & 0\\
0 & d & 1\\
0 & 0 & d
\end{pmatrix},
\]
allora
\[
\left(J_d^{(3)}\right)^k
=
\begin{pmatrix}
d^k & \binom{k}{1}d^{k-1} & \binom{k}{2}d^{k-2}\\
0 & d^k & \binom{k}{1}d^{k-1}\\
0 & 0 & d^k
\end{pmatrix}.
\]

\paragraph{Interpretazione dinamica nel discreto.}
Qui compaiono fattori del tipo
\[
d^k,\qquad kd^{k-1},\qquad \binom{k}{2}d^{k-2},\dots
\]
quindi non ci sono esponenziali continui, ma potenze e fattori polinomiali in $k$.

\paragraph{Conseguenza asintotica.}
Nel tempo discreto un blocco di Jordan associato a un autovalore reale $d$ converge a zero se
\[
|d|<1.
\]
Se invece $|d|>1$, diverge.  
Il caso $|d|=1$ è delicato: la presenza di blocchi di Jordan di ordine superiore a $1$ può generare crescita polinomiale.

\section{Autovalori complessi e necessità di una forma reale}

Fino ad ora abbiamo trattato il caso di autovalori reali.  
Se però $A\in\mathbb{R}^{n\times n}$ ha autovalori complessi, bisogna fare attenzione: lo stato del sistema è reale, quindi la descrizione finale dell'evoluzione deve restare reale.

\subsection{Coppie coniugate}

Se $A$ è reale e possiede un autovalore complesso
\[
\lambda=\sigma+j\omega,
\]
allora anche il suo coniugato
\[
\bar\lambda=\sigma-j\omega
\]
è autovalore di $A$.

Analogamente, se $v$ è autovettore associato a $\lambda$, allora $\bar v$ è autovettore associato a $\bar\lambda$.

Scriviamo
\[
v=q_r+jq_i,
\qquad
\bar v=q_r-jq_i,
\]
dove $q_r,q_i\in\mathbb{R}^n$ sono parte reale e parte immaginaria dell'autovettore.

\subsection{Caso semplice \texorpdfstring{$2\times 2$}{2x2}}

Supponiamo che $A\in\mathbb{R}^{2\times 2}$ abbia autovalori
\[
\lambda=\sigma+j\omega,
\qquad
\bar\lambda=\sigma-j\omega.
\]

Se
\[
Q=[\,v\ \bar v\,],
\qquad
D=
\begin{pmatrix}
\lambda & 0\\
0 & \bar\lambda
\end{pmatrix},
\]
allora
\[
AQ=QD.
\]

Il problema è che $Q$ è complessa. Per costruire una base reale si introduce la matrice
\[
E=\frac12
\begin{pmatrix}
1 & -j\\
1 & j
\end{pmatrix},
\]
così che
\[
QE=[\,q_r\ q_i\,]\in\mathbb{R}^{2\times 2}.
\]

Facendo i conti si ottiene
\[
E^{-1}DE
=
\begin{pmatrix}
\sigma & \omega\\
-\omega & \sigma
\end{pmatrix}
=:M.
\]

Quindi
\[
AQE=QDE=QE\,(E^{-1}DE)=QE\,M,
\]
e pertanto
\begin{equation}
A=(QE)\,M\,(QE)^{-1},
\label{eq:L6_real_reduction_complex}
\end{equation}
dove $QE$ è reale e
\begin{equation}
M=
\begin{pmatrix}
\sigma & \omega\\
-\omega & \sigma
\end{pmatrix}.
\label{eq:L6_real_block_complex_pair}
\end{equation}

\paragraph{Messaggio importante.}
Una coppia di autovalori complessi coniugati viene sostituita, in una base reale, con un blocco reale $2\times 2$.

\section{Esponenziale del blocco reale associato a una coppia complessa}

Consideriamo
\[
M=
\begin{pmatrix}
\sigma & \omega\\
-\omega & \sigma
\end{pmatrix}
=
\sigma I+
\begin{pmatrix}
0 & \omega\\
-\omega & 0
\end{pmatrix}.
\]

Poiché la parte diagonale commuta con la parte antisimmetrica, si ottiene
\begin{equation}
e^{Mt}
=
e^{\sigma t}
\begin{pmatrix}
\cos(\omega t) & \sin(\omega t)\\
-\sin(\omega t) & \cos(\omega t)
\end{pmatrix}.
\label{eq:L6_exp_real_complex_block}
\end{equation}

\paragraph{Interpretazione.}
La soluzione è il prodotto di:
\begin{itemize}
    \item un fattore esponenziale reale $e^{\sigma t}$, che governa crescita o decadimento;
    \item una rotazione/oscillazione con pulsazione $\omega$.
\end{itemize}

Quindi:
\begin{itemize}
    \item se $\sigma<0$, si hanno oscillazioni smorzate;
    \item se $\sigma=0$, si hanno oscillazioni persistenti;
    \item se $\sigma>0$, si hanno oscillazioni amplificate.
\end{itemize}

\subsection{Caso discreto}

Nel tempo discreto, se $\lambda=\sigma+j\omega$, è spesso utile passare alla forma polare
\[
\lambda=\rho(\cos\theta+j\sin\theta),
\qquad
\rho=\sqrt{\sigma^2+\omega^2}.
\]

Per il blocco reale
\[
M=
\begin{pmatrix}
\sigma & \omega\\
-\omega & \sigma
\end{pmatrix}
\]
vale
\begin{equation}
M^k
=
\rho^k
\begin{pmatrix}
\cos(k\theta) & \sin(k\theta)\\
-\sin(k\theta) & \cos(k\theta)
\end{pmatrix}.
\label{eq:L6_power_real_complex_block}
\end{equation}

\paragraph{Interpretazione discreta.}
Nel discreto:
\begin{itemize}
    \item $\rho$ regola l'ampiezza;
    \item $\theta$ regola la rotazione a ogni passo.
\end{itemize}

Quindi:
\begin{itemize}
    \item se $\rho<1$, la traiettoria tende a zero;
    \item se $\rho=1$, rimane limitata e puramente rotante;
    \item se $\rho>1$, diverge.
\end{itemize}

\section{Esempio con autovalori complessi}

Consideriamo
\[
A=
\begin{pmatrix}
3 & -1\\
2 & 1
\end{pmatrix}.
\]

Il polinomio caratteristico è
\[
p_A(\lambda)=\det(\lambda I-A)
=
(\lambda-3)(\lambda-1)+2
=
\lambda^2-4\lambda+5.
\]

Gli autovalori sono
\[
\lambda_{1,2}=2\pm j.
\]

Per $\lambda=2+j$ risolviamo
\[
(A-\lambda I)v=0.
\]

Si trova un autovettore complesso, per esempio
\[
v=
\begin{pmatrix}
1\\
1-j
\end{pmatrix}.
\]

Scrivendo
\[
v=q_r+jq_i
\]
si ha
\[
q_r=
\begin{pmatrix}
1\\
1
\end{pmatrix},
\qquad
q_i=
\begin{pmatrix}
0\\
-1
\end{pmatrix}.
\]

Quindi la matrice reale di cambio di base è
\[
Q_R=[\,q_r\ q_i\,]
=
\begin{pmatrix}
1 & 0\\
1 & -1
\end{pmatrix},
\]
e il blocco reale associato è
\[
M=
\begin{pmatrix}
2 & 1\\
-1 & 2
\end{pmatrix}.
\]

Ne segue che
\[
A=Q_R\,M\,Q_R^{-1},
\]
e quindi
\[
e^{At}=Q_R\,e^{Mt}\,Q_R^{-1}
=
Q_R\left[
e^{2t}
\begin{pmatrix}
\cos t & \sin t\\
-\sin t & \cos t
\end{pmatrix}
\right]Q_R^{-1}.
\]

\paragraph{Commento dinamico.}
Poiché la parte reale degli autovalori è $2>0$, il sistema ha andamento oscillatorio ma con ampiezza crescente: la dinamica libera esplode.

\section{Autovalore nullo e blocchi di Jordan}

Gli appunti insistono molto su un punto importante: l'autovalore da solo non basta, conta anche la dimensione del blocco di Jordan.

\subsection{Caso \texorpdfstring{$\lambda=0$}{lambda=0}}

\paragraph{Blocco semplice.}
Se
\[
J_0^{(1)}=(0),
\]
allora
\[
e^{J_0^{(1)}t}=1.
\]

Per la matrice nulla $2\times 2$:
\[
J=
\begin{pmatrix}
0 & 0\\
0 & 0
\end{pmatrix},
\qquad
e^{Jt}=I.
\]

Questa è una dinamica costante nel tempo: la soluzione non diverge ma neppure decade.

\paragraph{Blocco di ordine 2.}
Se
\[
J_0^{(2)}=
\begin{pmatrix}
0 & 1\\
0 & 0
\end{pmatrix},
\]
allora
\[
e^{J_0^{(2)}t}
=
\begin{pmatrix}
1 & t\\
0 & 1
\end{pmatrix}.
\]

Qui compare il termine $t$, che cresce senza limite.  
Quindi un autovalore nullo con blocco di Jordan non banale genera crescita polinomiale.

\paragraph{Messaggio.}
Nel continuo, per avere stabilità associata a un autovalore con parte reale nulla, è necessario che i blocchi di Jordan associati siano tutti di ordine $1$.

\subsection{Caso \texorpdfstring{$\lambda=-1$}{lambda=-1}}

Se
\[
J_{-1}^{(1)}=(-1),
\]
allora
\[
e^{J_{-1}^{(1)}t}=e^{-t},
\]
che tende a zero.

Se invece
\[
J_{-1}^{(2)}=
\begin{pmatrix}
-1 & 1\\
0 & -1
\end{pmatrix},
\]
allora
\[
e^{J_{-1}^{(2)}t}
=
e^{-t}
\begin{pmatrix}
1 & t\\
0 & 1
\end{pmatrix}
=
\begin{pmatrix}
e^{-t} & te^{-t}\\
0 & e^{-t}
\end{pmatrix}.
\]

Anche in questo caso la soluzione tende a zero, perché il decadimento esponenziale vince sul fattore polinomiale.

\section{Autovettori generalizzati}

Quando una matrice non è diagonalizzabile, gli autovettori non bastano per costruire una base.  
Per questo si introducono gli \emph{autovettori generalizzati}.

\subsection{Motivazione}

Consideriamo
\[
A=
\begin{pmatrix}
2 & 1 & 0\\
0 & 2 & 1\\
0 & 0 & 2
\end{pmatrix}
=
J_2^{(3)}.
\]

L'unico autovalore è $\lambda=2$, con molteplicità algebrica $3$.  
Calcoliamo
\[
A-2I=
\begin{pmatrix}
0 & 1 & 0\\
0 & 0 & 1\\
0 & 0 & 0
\end{pmatrix}.
\]

Il nucleo di $A-2I$ è formato dai vettori del tipo
\[
v=
\begin{pmatrix}
\alpha\\
0\\
0
\end{pmatrix},
\]
quindi
\[
\dim \ker(A-2I)=1.
\]

Dunque la molteplicità geometrica è $1$, minore della molteplicità algebrica $3$: la matrice non è diagonalizzabile.

\subsection{Definizione}

Un vettore $q$ si dice \emph{autovettore generalizzato di ordine $p$ associato all'autovalore $\lambda$} se
\begin{equation}
(A-\lambda I)^p q=0,
\qquad
(A-\lambda I)^{p-1}q\neq 0.
\label{eq:L6_generalized_eigenvector_def}
\end{equation}

Quindi:
\begin{itemize}
    \item per $p=1$ si ritrova il normale autovettore;
    \item per $p>1$ si ottengono i vettori necessari per completare le catene di Jordan.
\end{itemize}

\subsection{Autospazio generalizzato}

L'insieme
\[
\ker(A-\lambda I)^p
\]
è il sottospazio degli autovettori generalizzati di ordine al più $p$.

Si ha una catena crescente di sottospazi:
\begin{equation}
\ker(A-\lambda I)\subseteq \ker(A-\lambda I)^2 \subseteq \cdots \subseteq \ker(A-\lambda I)^p \subseteq \cdots
\label{eq:L6_chain_kernels}
\end{equation}

Questa catena si stabilizza quando si raggiunge la molteplicità algebrica dell'autovalore.

\section{Catene di autovettori generalizzati}

Per un blocco di Jordan di ordine $k$ associato all'autovalore $\lambda$, si costruisce una catena
\[
q^{(1)},q^{(2)},\dots,q^{(k)}
\]
tale che
\begin{equation}
(A-\lambda I)q^{(1)}=0,
\label{eq:L6_chain_1}
\end{equation}
\begin{equation}
(A-\lambda I)q^{(2)}=q^{(1)},
\label{eq:L6_chain_2}
\end{equation}
\[
\vdots
\]
\begin{equation}
(A-\lambda I)q^{(k)}=q^{(k-1)}.
\label{eq:L6_chain_k}
\end{equation}

Equivalentemente:
\[
Aq^{(1)}=\lambda q^{(1)},
\]
\[
Aq^{(2)}=q^{(1)}+\lambda q^{(2)},
\]
\[
\vdots
\]
\[
Aq^{(k)}=q^{(k-1)}+\lambda q^{(k)}.
\]

\begin{figure}[H]
\centering
\begin{tikzpicture}[>=Latex, node distance=1.65cm]
    \node[draw, rounded corners] (qk) {$q^{(k)}$};
    \node[draw, rounded corners, left=of qk] (qk1) {$q^{(k-1)}$};
    \node[left=of qk1] (dots) {$\cdots$};
    \node[draw, rounded corners, left=of dots] (q2) {$q^{(2)}$};
    \node[draw, rounded corners, left=of q2] (q1) {$q^{(1)}$};

    \draw[->, thick] (qk) -- node[above] {$A-\lambda I$} (qk1);
    \draw[->, thick] (qk1) -- (dots);
    \draw[->, thick] (dots) -- (q2);
    \draw[->, thick] (q2) -- node[above] {$A-\lambda I$} (q1);
    \draw[->, thick] (q1) -- ++(0,-1.2) node[below] {$0$};
\end{tikzpicture}
\caption{Catena di autovettori generalizzati associata a un blocco di Jordan.}
\end{figure}

\paragraph{Interpretazione.}
Una catena di ordine $k$ genera un blocco di Jordan di dimensione $k$.

\section{Significato della molteplicità geometrica}

Sia $\lambda_i$ un autovalore di $A$. Allora
\[
m_g(\lambda_i)=\dim\ker(A-\lambda_i I)
\]
è la molteplicità geometrica dell'autovalore.

\paragraph{Interpretazione strutturale.}
La molteplicità geometrica dice \emph{quanti blocchi di Jordan} sono associati a quell'autovalore.

In particolare:
\begin{itemize}
    \item se $m_g(\lambda_i)=m_a(\lambda_i)$, tutti i blocchi sono di ordine $1$ e la matrice è diagonalizzabile rispetto a quell'autovalore;
    \item se $m_g(\lambda_i)<m_a(\lambda_i)$, compaiono blocchi di Jordan di ordine superiore a $1$.
\end{itemize}

\section{Esempio di costruzione di una catena}

Riprendiamo
\[
A=
\begin{pmatrix}
2 & 1 & 0\\
0 & 2 & 1\\
0 & 0 & 2
\end{pmatrix}.
\]

Cerchiamo una catena per $\lambda=2$.

\subsection{Autovettore}

Serve
\[
(A-2I)q_1=0.
\]
Poiché
\[
A-2I=
\begin{pmatrix}
0 & 1 & 0\\
0 & 0 & 1\\
0 & 0 & 0
\end{pmatrix},
\]
si può scegliere
\[
q_1=
\begin{pmatrix}
1\\
0\\
0
\end{pmatrix}.
\]

\subsection{Autovettore generalizzato di ordine 2}

Cerchiamo $q_2$ tale che
\[
(A-2I)q_2=q_1.
\]
Una scelta naturale è
\[
q_2=
\begin{pmatrix}
0\\
1\\
0
\end{pmatrix},
\]
infatti
\[
(A-2I)q_2=
\begin{pmatrix}
1\\
0\\
0
\end{pmatrix}
=q_1.
\]

\subsection{Autovettore generalizzato di ordine 3}

Cerchiamo ora $q_3$ tale che
\[
(A-2I)q_3=q_2.
\]
Si può prendere
\[
q_3=
\begin{pmatrix}
0\\
0\\
1
\end{pmatrix},
\]
perché
\[
(A-2I)q_3=
\begin{pmatrix}
0\\
1\\
0
\end{pmatrix}
=q_2.
\]

\paragraph{Conclusione.}
La catena
\[
q_1,\ q_2,\ q_3
\]
costruisce esattamente il blocco $J_2^{(3)}$.

\section{Esempio con autovalore di molteplicità algebrica 3 e geometrica 2}

Consideriamo una matrice con autovalore $\lambda=5$ di molteplicità algebrica $3$ ma molteplicità geometrica $2$.  
Allora i possibili schemi di Jordan associati a $\lambda=5$ sono:
\begin{itemize}
    \item un blocco di ordine $3$;
    \item un blocco di ordine $2$ più un blocco di ordine $1$;
    \item tre blocchi di ordine $1$.
\end{itemize}

Se
\[
\dim\ker(A-5I)=2,
\]
allora i blocchi sono due, e quindi la struttura possibile è
\[
J=
\begin{pmatrix}
5 & 1 & 0\\
0 & 5 & 0\\
0 & 0 & 5
\end{pmatrix},
\]
cioè un blocco di ordine $2$ e uno di ordine $1$.

\paragraph{Procedura.}
\begin{enumerate}
    \item si trova una base di $\ker(A-5I)$, che fornisce gli autovettori di ordine $1$;
    \item si cerca poi un vettore in $\ker(A-5I)^2$ ma non in $\ker(A-5I)$;
    \item applicando $(A-5I)$ a questo vettore si ottiene uno degli autovettori del primo passo;
    \item questa costruisce una catena di lunghezza $2$.
\end{enumerate}

\section{Conteggio degli autovettori generalizzati}

Sia
\[
M_p:=\dim\ker(A-\lambda I)^p.
\]

Allora:
\begin{itemize}
    \item il numero di autovettori generalizzati di ordine almeno $1$ è $M_1$;
    \item il numero di autovettori generalizzati di ordine almeno $2$ è $M_2-M_1$;
    \item il numero di autovettori generalizzati di ordine almeno $3$ è $M_3-M_2$;
    \item e così via.
\end{itemize}

Più precisamente, il numero di catene di lunghezza almeno $p$ è
\[
M_p-M_{p-1}.
\]

\paragraph{Interpretazione.}
Questo permette di ricostruire la dimensione dei blocchi di Jordan semplicemente studiando la crescita delle dimensioni dei nuclei
\[
\ker(A-\lambda I),\ \ker(A-\lambda I)^2,\ \ker(A-\lambda I)^3,\dots
\]

\section{Schema pratico per costruire la forma di Jordan}

Se si vuole costruire una matrice $Q$ tale che
\[
Q^{-1}AQ=J,
\]
si procede così:

\begin{enumerate}
    \item si trovano gli autovalori di $A$ e le loro molteplicità algebriche;
    \item per ogni autovalore $\lambda_i$ si calcola
    \[
    \dim\ker(A-\lambda_i I),
    \quad
    \dim\ker(A-\lambda_i I)^2,
    \quad \dots
    \]
    fino a stabilizzazione;
    \item da queste dimensioni si deduce la struttura dei blocchi di Jordan;
    \item si costruiscono le catene di autovettori generalizzati;
    \item si mettono come colonne di $Q$ i vettori delle catene, ordinati in modo coerente con i blocchi di Jordan.
\end{enumerate}

\paragraph{Caso particolare: un solo autovalore e un solo blocco.}
Se la matrice ha un solo autovalore $\lambda$ con:
\[
m_a(\lambda)=n,\qquad m_g(\lambda)=1,
\]
allora esiste un unico blocco di Jordan di ordine $n$.  
In tal caso basta costruire una singola catena completa
\[
q_1,\ q_2,\ \dots,\ q_n.
\]

\section{Sintesi finale della lezione}

In questa lezione abbiamo visto che:

\begin{enumerate}
    \item l'esponenziale di un blocco nilpotente $J_0^{(p)}$ è un polinomio matriciale che si tronca al grado $p-1$;

    \item l'esponenziale di un blocco di Jordan reale è dato da
    \[
    e^{J_d^{(p)}t}=e^{dt}e^{J_0^{(p)}t},
    \]
    quindi compaiono fattori del tipo $t^k e^{dt}$;

    \item nel tempo discreto le potenze di un blocco di Jordan generano fattori del tipo
    \[
    \binom{k}{i}d^{k-i};
    \]

    \item per autovalori complessi coniugati, una coppia complessa si traduce in un blocco reale
    \[
    \begin{pmatrix}
    \sigma & \omega\\
    -\omega & \sigma
    \end{pmatrix},
    \]
    il cui esponenziale descrive oscillazioni smorzate, persistenti o amplificate;

    \item quando la matrice non è diagonalizzabile si introducono gli autovettori generalizzati;

    \item le catene di autovettori generalizzati permettono di costruire i blocchi di Jordan;

    \item la molteplicità geometrica di un autovalore coincide con il numero di blocchi di Jordan associati a quell'autovalore;

    \item studiando i nuclei di
    \[
    (A-\lambda I)^p
    \]
    si può ricostruire operativamente la forma di Jordan.
\end{enumerate}

Questa parte è decisiva per comprendere il comportamento dinamico dei sistemi lineari, soprattutto nei casi non diagonalizzabili e nei casi oscillatori con autovalori complessi.
\chapter{L7 -- Costruzione pratica della forma di Jordan e catene di autovettori generalizzati}

\section{Obiettivo della lezione}

Dopo aver introdotto nella lezione precedente gli autovettori generalizzati, in questa lezione vogliamo passare dalla teoria alla \emph{costruzione pratica} della forma di Jordan di una matrice generica.

Il problema è il seguente: data una matrice
\[
A\in\mathbb{R}^{n\times n},
\]
si vuole costruire una matrice invertibile $Q$ tale che
\[
Q^{-1}AQ=J,
\]
dove $J$ è la forma di Jordan associata ad $A$.

Equivalentemente, si può scrivere
\begin{equation}
AQ=QJ.
\label{eq:L7_AQ_QJ}
\end{equation}

Il punto chiave è capire come devono essere scelte le colonne di $Q$ affinché questa relazione produca esattamente i blocchi di Jordan desiderati.

\section{Caso fondamentale: un solo blocco di Jordan}

Conviene cominciare dal caso più semplice ma più istruttivo: supponiamo che $A$ sia simile a un unico blocco di Jordan di dimensione $m$ associato all'autovalore $\lambda$:
\[
J=
J_\lambda^{(m)}
=
\begin{pmatrix}
\lambda & 1 & 0 & \cdots & 0\\
0 & \lambda & 1 & \cdots & 0\\
0 & 0 & \lambda & \ddots & \vdots\\
\vdots & \vdots & \ddots & \ddots & 1\\
0 & 0 & \cdots & 0 & \lambda
\end{pmatrix}.
\]

Sia
\[
Q=\begin{pmatrix} q_1 & q_2 & \cdots & q_m \end{pmatrix},
\]
dove $q_1,\dots,q_m$ sono le colonne di $Q$.

Dalla relazione \eqref{eq:L7_AQ_QJ}, confrontando colonna per colonna, si ottiene:

\subsection{Prima colonna}

La prima colonna di $J$ è
\[
\begin{pmatrix}
\lambda\\
0\\
\vdots\\
0
\end{pmatrix},
\]
quindi
\[
Aq_1=Q
\begin{pmatrix}
\lambda\\
0\\
\vdots\\
0
\end{pmatrix}
=\lambda q_1.
\]

Dunque
\begin{equation}
(A-\lambda I)q_1=0.
\label{eq:L7_chain_first}
\end{equation}

Quindi $q_1$ è un \emph{autovettore} associato a $\lambda$.

\subsection{Colonne successive}

Per $j=2,\dots,m$, la $j$-esima colonna di $J$ ha un $1$ nella riga $j-1$, un $\lambda$ sulla diagonale e zeri altrove.  
Quindi:
\[
Aq_j=q_{j-1}+\lambda q_j,
\]
ossia
\begin{equation}
(A-\lambda I)q_j=q_{j-1},
\qquad j=2,\dots,m.
\label{eq:L7_chain_relation}
\end{equation}

Queste relazioni sono la struttura fondamentale di una \emph{catena di autovettori generalizzati}.

\begin{figure}[H]
\centering
\begin{tikzpicture}[>=Latex, node distance=1.6cm]
    \node[draw, rounded corners] (q1) {$q_1$};
    \node[draw, rounded corners, right=of q1] (q2) {$q_2$};
    \node[right=of q2] (dots) {$\cdots$};
    \node[draw, rounded corners, right=of dots] (qm1) {$q_{m-1}$};
    \node[draw, rounded corners, right=of qm1] (qm) {$q_m$};

    \draw[<- , thick] (q1) -- node[above] {$A-\lambda I$} (q2);
    \draw[<- , thick] (q2) -- (dots);
    \draw[<- , thick] (dots) -- (qm1);
    \draw[<- , thick] (qm1) -- node[above] {$A-\lambda I$} (qm);

    \draw[->, thick] (q1) -- +(0,-1.2) node[below] {$0$};
\end{tikzpicture}
\caption{Catena di autovettori generalizzati associata a un blocco di Jordan.}
\end{figure}

\section{Ordine degli autovettori generalizzati}

Dalle relazioni \eqref{eq:L7_chain_first}--\eqref{eq:L7_chain_relation} segue immediatamente che:
\[
q_1\in\ker(A-\lambda I),
\]
\[
q_2\in\ker(A-\lambda I)^2,\qquad q_2\notin\ker(A-\lambda I),
\]
\[
q_3\in\ker(A-\lambda I)^3,\qquad q_3\notin\ker(A-\lambda I)^2,
\]
e così via.

In generale:
\begin{equation}
q_j\in\ker(A-\lambda I)^j
\qquad\text{e}\qquad
q_j\notin\ker(A-\lambda I)^{j-1}.
\label{eq:L7_order_generalized}
\end{equation}

Questo significa che $q_j$ è un autovettore generalizzato di ordine $j$.

\paragraph{Osservazione.}
Per costruire una catena conviene spesso partire \emph{dall'ultimo vettore}, cioè da un elemento di ordine massimo, e poi risalire applicando $A-\lambda I$:
\[
q_m \xmapsto{\,A-\lambda I\,} q_{m-1} \xmapsto{\,A-\lambda I\,} \cdots \xmapsto{\,A-\lambda I\,} q_1 \xmapsto{\,A-\lambda I\,} 0.
\]

\section{Costruzione delle catene a partire dai nuclei successivi}

Sia $\lambda$ un autovalore di $A$. Consideriamo la catena crescente di sottospazi
\begin{equation}
\ker(A-\lambda I)\subseteq \ker(A-\lambda I)^2 \subseteq \cdots \subseteq \ker(A-\lambda I)^p=\ker(A-\lambda I)^{p+1},
\label{eq:L7_kernel_chain}
\end{equation}
dove $p$ è il primo indice di stabilizzazione.

Definiamo
\begin{equation}
\mu_j:=\dim\ker(A-\lambda I)^j,
\label{eq:L7_mu_j}
\end{equation}
e poniamo convenzionalmente
\[
\mu_0:=0.
\]

Allora la quantità
\begin{equation}
m_j:=\mu_j-\mu_{j-1}
\label{eq:L7_m_j}
\end{equation}
misura quanti \emph{nuovi} vettori compaiono passando dall'ordine $j-1$ all'ordine $j$.

\paragraph{Interpretazione pratica.}
Nella costruzione adattata a Jordan:
\begin{itemize}
    \item $m_1$ è il numero di autovettori indipendenti, cioè la molteplicità geometrica;
    \item $m_2$ dice quante catene raggiungono almeno l'ordine $2$;
    \item $m_3$ dice quante catene raggiungono almeno l'ordine $3$;
    \item e così via.
\end{itemize}

\paragraph{Procedura costruttiva per un autovalore fissato $\lambda$.}
\begin{enumerate}
    \item Si calcolano i nuclei successivi $\ker(A-\lambda I)^j$ fino a stabilizzazione.
    \item Si determinano le dimensioni $\mu_j$ e le differenze $m_j$.
    \item Si parte dal massimo ordine $p$ e si scelgono $m_p$ vettori
    \[
    q_1^{(p)},\dots,q_{m_p}^{(p)}\in \ker(A-\lambda I)^p\setminus \ker(A-\lambda I)^{p-1}.
    \]
    \item Da ciascuno di essi si genera una catena ponendo
    \[
    q_i^{(p-1)}=(A-\lambda I)q_i^{(p)},\qquad
    q_i^{(p-2)}=(A-\lambda I)q_i^{(p-1)},\quad \dots
    \]
    fino ad arrivare a un autovettore.
    \item Si passa poi all'ordine $p-1$, scegliendo \emph{nuovi} vettori che non siano già coinvolti nelle catene costruite.
\end{enumerate}

\paragraph{Idea di fondo.}
Si parte dall'alto per evitare di scegliere vettori che finirebbero per essere già contenuti nelle catene costruite in precedenza.

\section{Come si dispongono le colonne di \texorpdfstring{$Q$}{Q}}

Una volta costruite le catene, le colonne di $Q$ devono essere ordinate \emph{catena per catena}.

Se, per esempio, abbiamo due catene:
\[
q_1^{(1)},q_1^{(2)},\dots,q_1^{(p)},
\qquad
q_2^{(1)},q_2^{(2)},\dots,q_2^{(r)},
\]
allora una scelta compatibile è
\begin{equation}
Q=
\begin{pmatrix}
q_1^{(1)} & q_1^{(2)} & \cdots & q_1^{(p)} &
q_2^{(1)} & q_2^{(2)} & \cdots & q_2^{(r)}
\end{pmatrix}.
\label{eq:L7_Q_ordered_by_chains}
\end{equation}

\paragraph{Nota importante.}
Si possono permutare i blocchi di Jordan, ma non si può spezzare una catena internamente, altrimenti la relazione $AQ=QJ$ non assume più la forma desiderata.

\section{Procedura generale per costruire la forma di Jordan}

Riassumiamo adesso l'algoritmo completo.

\subsection*{Passo 1: autovalori e molteplicità algebriche}

Si calcolano gli autovalori $\lambda_i$ di $A$ e le loro molteplicità algebriche $m_a(\lambda_i)$.

\subsection*{Passo 2: nuclei successivi}

Per ogni autovalore $\lambda_i$ si calcolano i nuclei
\[
\ker(A-\lambda_i I),\ \ker(A-\lambda_i I)^2,\ \dots,\ \ker(A-\lambda_i I)^p,
\]
fino al primo indice $p$ per cui
\[
\ker(A-\lambda_i I)^p=\ker(A-\lambda_i I)^{p+1}.
\]

\subsection*{Passo 3: dimensioni e differenze}

Si definiscono
\[
\mu_j=\dim\ker(A-\lambda_i I)^j,
\qquad
m_j=\mu_j-\mu_{j-1}.
\]

\subsection*{Passo 4: costruzione delle catene}

Si parte dall'ordine massimo $p$ e si scelgono i vettori generalizzati di ordine massimo; poi si fa scendere la catena applicando $A-\lambda_i I$.

\subsection*{Passo 5: matrice \texorpdfstring{$Q$}{Q}}

Si mettono come colonne di $Q$ tutti i vettori delle catene, ordinati blocco per blocco.

\subsection*{Passo 6: forma di Jordan}

La matrice $J$ si legge direttamente dalle lunghezze delle catene:
\begin{itemize}
    \item una catena di lunghezza $1$ produce un blocco di Jordan di ordine $1$;
    \item una catena di lunghezza $2$ produce un blocco di ordine $2$;
    \item ecc.
\end{itemize}

\section{Esempio completo}

Negli appunti compare un esempio molto istruttivo. Consideriamo la matrice
\begin{equation}
A=
\begin{pmatrix}
0 & 1 & 0 & 0 & -1 & 0\\
0 & 0 & 0 & 0 & 0 & -1\\
0 & -1 & -1 & 1 & 0 & -1\\
0 & 0 & 0 & -1 & 0 & 0\\
0 & 0 & 0 & 0 & -1 & 1\\
0 & 0 & 0 & 0 & 0 & -1
\end{pmatrix}.
\label{eq:L7_example_A}
\end{equation}

La matrice è triangolare superiore a blocchi, quindi gli autovalori si leggono sulla diagonale:
\[
\lambda=0 \quad \text{con molteplicità algebrica }2,
\]
\[
\lambda=-1 \quad \text{con molteplicità algebrica }4.
\]

Dunque:
\[
m_a(0)=2,
\qquad
m_a(-1)=4.
\]

\subsection{Autovalore \texorpdfstring{$\lambda=0$}{lambda=0}}

Dobbiamo studiare i nuclei di $A^j$.

\subsubsection*{Primo passo}

Poiché la prima colonna di $A$ è nulla, si ha subito
\[
Ae_1=0,
\]
quindi
\[
\ker(A)=\operatorname{span}\{e_1\},
\qquad
\mu_1=\dim\ker(A)=1.
\]

\subsubsection*{Secondo passo}

Calcolando $A^2$ si ottiene
\[
\mu_2=\dim\ker(A^2)=2.
\]
Poiché la molteplicità algebrica di $0$ è già pari a $2$, la catena si stabilizza qui:
\[
\ker(A^3)=\ker(A^2).
\]

Quindi
\[
m_1=\mu_1=1,
\qquad
m_2=\mu_2-\mu_1=1.
\]

\paragraph{Conclusione.}
Esiste una sola catena per $\lambda=0$, di lunghezza $2$.

\subsubsection*{Scelta della catena}

Una scelta possibile è
\[
q_{0,1}^{(2)}=
\begin{pmatrix}
1\\
1\\
-1\\
0\\
0\\
0
\end{pmatrix},
\]
per cui
\[
q_{0,1}^{(1)}=Aq_{0,1}^{(2)}=
\begin{pmatrix}
1\\
0\\
0\\
0\\
0\\
0
\end{pmatrix}
=e_1.
\]

Questa produce il blocco di Jordan
\[
J_0^{(2)}=
\begin{pmatrix}
0 & 1\\
0 & 0
\end{pmatrix}.
\]

\subsection{Autovalore \texorpdfstring{$\lambda=-1$}{lambda=-1}}

Ora studiamo
\[
A+I.
\]

\subsubsection*{Primo passo}

Si trova
\[
\dim\ker(A+I)=2,
\]
quindi
\[
\mu_1=2.
\]

Una possibile base di $\ker(A+I)$ è
\[
q_{-1,1}^{(1)}=
\begin{pmatrix}
0\\
0\\
1\\
0\\
0\\
0
\end{pmatrix},
\qquad
q_{-1,2}^{(1)}=
\begin{pmatrix}
1\\
0\\
0\\
0\\
1\\
0
\end{pmatrix}.
\]

\subsubsection*{Secondo passo}

Si calcola poi
\[
\mu_2=\dim\ker(A+I)^2=4.
\]

Quindi
\[
m_2=\mu_2-\mu_1=4-2=2.
\]

\subsubsection*{Terzo passo}

Si verifica che
\[
\mu_3=\dim\ker(A+I)^3=4,
\]
perciò la catena si stabilizza:
\[
\ker(A+I)^3=\ker(A+I)^2.
\]

Quindi
\[
m_3=\mu_3-\mu_2=0.
\]

\paragraph{Conclusione.}
Per l'autovalore $-1$ ci sono due catene di lunghezza $2$.

\subsubsection*{Scelta delle catene}

Una prima scelta naturale è
\[
q_{-1,1}^{(2)}=
\begin{pmatrix}
0\\
0\\
0\\
1\\
0\\
0
\end{pmatrix},
\qquad
(A+I)q_{-1,1}^{(2)}=
\begin{pmatrix}
0\\
0\\
1\\
0\\
0\\
0
\end{pmatrix}
=q_{-1,1}^{(1)}.
\]

Per la seconda catena, una scelta possibile è
\[
q_{-1,2}^{(2)}=
\begin{pmatrix}
0\\
1\\
0\\
2\\
0\\
1
\end{pmatrix},
\]
infatti
\[
(A+I)q_{-1,2}^{(2)}=
\begin{pmatrix}
1\\
0\\
0\\
0\\
1\\
0
\end{pmatrix}
=q_{-1,2}^{(1)}.
\]

\subsection{Costruzione di \texorpdfstring{$Q$}{Q}}

Disponiamo ora le colonne per catene:
\[
Q=
\begin{pmatrix}
q_{0,1}^{(1)} & q_{0,1}^{(2)} &
q_{-1,1}^{(1)} & q_{-1,1}^{(2)} &
q_{-1,2}^{(1)} & q_{-1,2}^{(2)}
\end{pmatrix}.
\]

Esplicitamente, una possibile matrice di cambio di base è
\begin{equation}
Q=
\begin{pmatrix}
1 & 1 & 0 & 0 & 1 & 0\\
0 & 1 & 0 & 0 & 0 & 1\\
0 & -1 & 1 & 0 & 0 & 0\\
0 & 0 & 0 & 1 & 0 & 2\\
0 & 0 & 0 & 0 & 1 & 0\\
0 & 0 & 0 & 0 & 0 & 1
\end{pmatrix}.
\label{eq:L7_example_Q}
\end{equation}

\subsection{Forma di Jordan finale}

La forma di Jordan corrispondente è
\begin{equation}
J=
\operatorname{diag}\!\left(
J_0^{(2)},
J_{-1}^{(2)},
J_{-1}^{(2)}
\right)
=
\begin{pmatrix}
0 & 1 & 0 & 0 & 0 & 0\\
0 & 0 & 0 & 0 & 0 & 0\\
0 & 0 & -1 & 1 & 0 & 0\\
0 & 0 & 0 & -1 & 0 & 0\\
0 & 0 & 0 & 0 & -1 & 1\\
0 & 0 & 0 & 0 & 0 & -1
\end{pmatrix}.
\label{eq:L7_example_J}
\end{equation}

Dunque
\[
Q^{-1}AQ=J.
\]

\paragraph{Commento.}
L'esempio mostra bene che:
\begin{itemize}
    \item l'autovalore $0$ genera un unico blocco di ordine $2$;
    \item l'autovalore $-1$ genera due blocchi di ordine $2$;
    \item l'informazione cruciale non è solo la molteplicità algebrica, ma la crescita delle dimensioni dei nuclei successivi.
\end{itemize}

\section{Proprietà degli autovettori generalizzati}

Gli appunti sottolineano alcune proprietà che è utile fissare.

\subsection{Linearità lungo una catena}

Se
\[
(A-\lambda I)q^{(k)}=q^{(k-1)},
\]
allora i vettori
\[
q^{(k)},\,q^{(k-1)},\,\dots,\,q^{(1)}
\]
sono linearmente indipendenti.

\paragraph{Idea della dimostrazione.}
Se fossero linearmente dipendenti, applicando successivamente $A-\lambda I$ si arriverebbe a una dipendenza anche tra vettori di ordine inferiore, fino a contraddire il fatto che $q^{(1)}\neq 0$ è un autovettore.

\subsection{Connessione con i sottospazi \texorpdfstring{$A$}{A}-invarianti}

Ricordiamo che il sottospazio ciclico generato da $x_0$,
\[
S_{x_0}=\operatorname{span}\{x_0,Ax_0,\dots,A^{p-1}x_0\},
\]
è $A$-invariante.

Se $x_0$ è un autovettore, allora
\[
S_{x_0}=\operatorname{span}\{x_0\},
\]
cioè una retta.

Se invece $x_0$ è un autovettore generalizzato di ordine $k$, allora la catena associata genera un sottospazio $A$-invariante di dimensione $k$, che coincide con il blocco di Jordan corrispondente.

\section{Interpretazione strutturale}

Tutta la costruzione fatta fin qui serve a capire la struttura dinamica del sistema.

\paragraph{Per un autovalore reale $\lambda$:}
\begin{itemize}
    \item se il blocco è di ordine $1$, compare solo il modo $e^{\lambda t}$ nel continuo o $\lambda^k$ nel discreto;
    \item se il blocco è di ordine $r>1$, compaiono anche fattori polinomiali:
    \[
    e^{\lambda t},\ te^{\lambda t},\ \dots,\ t^{r-1}e^{\lambda t},
    \]
    nel continuo, oppure
    \[
    \lambda^k,\ k\lambda^{k-1},\ \binom{k}{2}\lambda^{k-2},\dots
    \]
    nel discreto.
\end{itemize}

\paragraph{Per questo motivo:}
non basta conoscere gli autovalori; serve conoscere anche la forma di Jordan, cioè la dimensione dei blocchi.

\section{Schema riassuntivo finale}

Per un autovalore fissato $\lambda$:

\begin{enumerate}
    \item Calcolo le dimensioni
    \[
    \mu_j=\dim\ker(A-\lambda I)^j.
    \]
    \item Calcolo le differenze
    \[
    m_j=\mu_j-\mu_{j-1}.
    \]
    \item Interpreto $m_j$ come numero di catene che arrivano almeno all'ordine $j$.
    \item Scelgo i vettori di ordine massimo.
    \item Li faccio scendere con l'operatore $A-\lambda I$.
    \item Dispongo le colonne di $Q$ per catene.
    \item Ricavo $J$.
\end{enumerate}

\section{Sintesi della lezione}

In questa lezione abbiamo visto che:

\begin{enumerate}
    \item la relazione
    \[
    AQ=QJ
    \]
    letta colonna per colonna produce naturalmente le catene di autovettori generalizzati;

    \item un vettore $q_j$ di una catena soddisfa
    \[
    (A-\lambda I)q_j=q_{j-1},
    \]
    e quindi appartiene a
    \[
    \ker(A-\lambda I)^j\setminus \ker(A-\lambda I)^{j-1};
    \]

    \item per costruire la forma di Jordan bisogna studiare la crescita dei nuclei
    \[
    \ker(A-\lambda I)^j;
    \]

    \item le quantità
    \[
    \mu_j=\dim\ker(A-\lambda I)^j,
    \qquad
    m_j=\mu_j-\mu_{j-1}
    \]
    permettono di capire quante catene e di quali lunghezze servono;

    \item le colonne della matrice di cambio di base $Q$ vanno disposte seguendo le catene, senza spezzarle;

    \item nell'esempio proposto, la matrice ha:
    \[
    J=\operatorname{diag}\!\left(J_0^{(2)},J_{-1}^{(2)},J_{-1}^{(2)}\right).
    \]
\end{enumerate}

Questa procedura completa il passaggio dalla teoria degli autovalori generalizzati alla costruzione operativa della forma di Jordan, che sarà poi usata per studiare in modo sistematico l'evoluzione libera dei sistemi lineari.
\chapter{L8 -- Proprietà degli autovettori generalizzati e modi associati ai blocchi di Jordan}

\section{Obiettivo della lezione}

Nelle lezioni precedenti abbiamo introdotto:
\begin{itemize}
    \item la forma di Jordan;
    \item gli autovettori generalizzati;
    \item la costruzione delle catene associate a ciascun autovalore.
\end{itemize}

In questa lezione vogliamo chiarire tre idee fondamentali:

\begin{enumerate}
    \item quali proprietà hanno gli autovettori generalizzati;
    \item come sono collegati ai sottospazi ciclici e ai sottospazi $A$-invarianti;
    \item perché la struttura dei blocchi di Jordan determina direttamente i \emph{modi} che compaiono nell'evoluzione temporale del sistema.
\end{enumerate}

La morale della lezione è molto importante: per capire davvero l'andamento nel tempo di un sistema lineare non basta conoscere gli autovalori; bisogna conoscere anche la \emph{dimensione dei blocchi di Jordan} associati a ciascun autovalore.

\section{Richiamo: catene di autovettori generalizzati}

Sia $\lambda$ un autovalore di $A \in \R^{n\times n}$.  
Una catena di autovettori generalizzati associata a $\lambda$ è una famiglia di vettori
\[
q^{(1)}, q^{(2)}, \dots, q^{(k)}
\]
tale che
\begin{equation}
(A-\lambda I)q^{(1)} = 0,
\label{eq:L8_chain_1}
\end{equation}
\begin{equation}
(A-\lambda I)q^{(2)} = q^{(1)},
\label{eq:L8_chain_2}
\end{equation}
\[
\vdots
\]
\begin{equation}
(A-\lambda I)q^{(k)} = q^{(k-1)}.
\label{eq:L8_chain_k}
\end{equation}

Il vettore $q^{(1)}$ è un normale autovettore, mentre $q^{(j)}$ con $j\ge 2$ è un autovettore generalizzato di ordine $j$.

Equivalentemente si può scrivere:
\[
Aq^{(1)} = \lambda q^{(1)},
\]
\[
Aq^{(2)} = q^{(1)} + \lambda q^{(2)},
\]
\[
\vdots
\]
\[
Aq^{(k)} = q^{(k-1)} + \lambda q^{(k)}.
\]

\begin{figure}[H]
\centering
\begin{tikzpicture}[>=Latex, node distance=1.7cm]
    \node[draw, rounded corners] (q1) {$q^{(1)}$};
    \node[draw, rounded corners, right=of q1] (q2) {$q^{(2)}$};
    \node[right=of q2] (dots) {$\cdots$};
    \node[draw, rounded corners, right=of dots] (qk1) {$q^{(k-1)}$};
    \node[draw, rounded corners, right=of qk1] (qk) {$q^{(k)}$};

    \draw[<-, thick] (q1) -- node[above] {$A-\lambda I$} (q2);
    \draw[<-, thick] (q2) -- (dots);
    \draw[<-, thick] (dots) -- (qk1);
    \draw[<-, thick] (qk1) -- node[above] {$A-\lambda I$} (qk);

    \draw[->, thick] (q1) -- +(0,-1.1) node[below] {$0$};
\end{tikzpicture}
\caption{Catena di autovettori generalizzati associata all'autovalore $\lambda$.}
\end{figure}

\section{Ordine degli autovettori generalizzati}

Dalle relazioni precedenti segue subito che:
\[
q^{(1)} \in \ker(A-\lambda I),
\]
\[
q^{(2)} \in \ker(A-\lambda I)^2 \setminus \ker(A-\lambda I),
\]
\[
q^{(3)} \in \ker(A-\lambda I)^3 \setminus \ker(A-\lambda I)^2,
\]
e così via.

In generale:
\begin{equation}
q^{(j)} \in \ker(A-\lambda I)^j
\qquad \text{e} \qquad
q^{(j)} \notin \ker(A-\lambda I)^{j-1}.
\label{eq:L8_order}
\end{equation}

Questo giustifica la terminologia: $q^{(j)}$ è un autovettore generalizzato di ordine $j$.

\section{Conteggio tramite i nuclei successivi}

Per un autovalore fissato $\lambda$, definiamo
\begin{equation}
\mu_j := \dim \ker(A-\lambda I)^j.
\label{eq:L8_mu_j}
\end{equation}

Poniamo anche, per comodità,
\[
\mu_0 := 0.
\]

Allora la quantità
\begin{equation}
n_j := \mu_j - \mu_{j-1}
\label{eq:L8_n_j}
\end{equation}
rappresenta il numero di \emph{nuovi} autovettori generalizzati che compaiono all'ordine $j$.

In pratica:
\begin{itemize}
    \item $n_1$ è il numero di autovettori indipendenti, cioè la molteplicità geometrica;
    \item $n_2$ dice quanti autovettori generalizzati di ordine $2$ servono;
    \item $n_3$ dice quanti autovettori generalizzati di ordine $3$ servono;
    \item ecc.
\end{itemize}

\paragraph{Osservazione importante.}
Questi numeri permettono di ricostruire la struttura dei blocchi di Jordan senza doverla indovinare.

\section{Prima proprietà: indipendenza lineare lungo una catena}

Una catena di autovettori generalizzati
\[
q^{(1)}, q^{(2)}, \dots, q^{(k)}
\]
è costituita da vettori linearmente indipendenti.

\subsection*{Dimostrazione}

Supponiamo per assurdo che esista una combinazione lineare nulla:
\[
\alpha_1 q^{(1)} + \alpha_2 q^{(2)} + \cdots + \alpha_k q^{(k)} = 0.
\]

Applichiamo $(A-\lambda I)^{k-1}$ a entrambi i membri.  
Poiché:
\[
(A-\lambda I)^{k-1} q^{(k)} = q^{(1)},
\]
mentre
\[
(A-\lambda I)^{k-1} q^{(j)} = 0 \qquad \text{per ogni } j<k,
\]
si ottiene
\[
\alpha_k q^{(1)} = 0.
\]

Essendo $q^{(1)} \neq 0$, segue $\alpha_k = 0$.

Ripetendo lo stesso ragionamento con $(A-\lambda I)^{k-2}$, poi con $(A-\lambda I)^{k-3}$, e così via, si trova successivamente:
\[
\alpha_{k-1}=0,\quad \alpha_{k-2}=0,\quad \dots,\quad \alpha_1=0.
\]

Quindi i vettori sono linearmente indipendenti.

\paragraph{Conclusione.}
Una catena di lunghezza $k$ produce sempre uno spazio vettoriale di dimensione $k$.

\section{Seconda proprietà: sottospazio ciclico generato da un autovettore generalizzato}

Sia $x_0 \in \R^n$. Ricordiamo che il sottospazio ciclico generato da $x_0$ è
\begin{equation}
S_{x_0} := \operatorname{span}\{x_0, Ax_0, A^2x_0, \dots, A^{p-1}x_0\},
\label{eq:L8_cyclic_subspace}
\end{equation}
dove $p$ è il primo indice per cui i vettori diventano linearmente dipendenti.

Questo sottospazio è $A$-invariante.

\subsection{Caso particolare: autovettore}

Se $x_0 = v$ è un autovettore, cioè
\[
Av = \lambda v,
\]
allora
\[
A^2v = \lambda^2 v,\qquad A^3v = \lambda^3 v,\qquad \dots
\]
quindi tutti i vettori $A^j v$ sono multipli di $v$.

Perciò:
\[
S_v = \operatorname{span}\{v\},
\]
cioè una retta.

\subsection{Caso generale: autovettore generalizzato di ordine \texorpdfstring{$k$}{k}}

Sia ora $q^{(k)}$ un autovettore generalizzato di ordine $k$.  
Allora vale
\[
(A-\lambda I)q^{(k)} = q^{(k-1)},
\]
cioè
\[
Aq^{(k)} = q^{(k-1)} + \lambda q^{(k)}.
\]

Quindi $Aq^{(k)}$ appartiene allo span di $q^{(k)}$ e $q^{(k-1)}$.

Vediamo cosa succede continuando:

\[
(A-\lambda I)^2 q^{(k)} = q^{(k-2)}.
\]

Sviluppando:
\[
(A-\lambda I)^2 q^{(k)} = \bigl(A^2 - 2\lambda A + \lambda^2 I \bigr)q^{(k)},
\]
per cui
\[
q^{(k-2)} = A^2 q^{(k)} - 2\lambda A q^{(k)} + \lambda^2 q^{(k)}.
\]

Da qui si ricava
\[
A^2 q^{(k)} \in \operatorname{span}\{q^{(k)}, q^{(k-1)}, q^{(k-2)}\}.
\]

Proseguendo nello stesso modo, si ottiene che ogni potenza $A^r q^{(k)}$ appartiene allo span dei vettori della catena:
\[
q^{(1)}, q^{(2)}, \dots, q^{(k)}.
\]

Viceversa, ogni $q^{(j)}$ si può ottenere da $q^{(k)}$ applicando una potenza di $(A-\lambda I)$, quindi è anch'esso nel sottospazio ciclico generato da $q^{(k)}$.

\paragraph{Conclusione fondamentale.}
Vale
\begin{equation}
S_{q^{(k)}} = \operatorname{span}\{q^{(1)}, q^{(2)}, \dots, q^{(k)}\}.
\label{eq:L8_cyclic_chain}
\end{equation}

Poiché i vettori della catena sono linearmente indipendenti, si ha anche
\[
\dim S_{q^{(k)}} = k.
\]

\paragraph{Interpretazione.}
Una catena di lunghezza $k$ genera un sottospazio ciclico $A$-invariante di dimensione $k$: è esattamente lo spazio ``visto'' da quel blocco di Jordan.

\section{Esempio 1: una matrice simile a un unico blocco di Jordan}

Consideriamo la matrice
\begin{equation}
A_2 =
\begin{pmatrix}
5 & 0 & 1\\
1 & 5 & \dfrac12\\
0 & 0 & 5
\end{pmatrix}.
\label{eq:L8_A2}
\end{equation}

L'unico autovalore è
\[
\lambda = 5,
\qquad
m_a(5)=3.
\]

Studiamo i nuclei di $A_2 - 5I$.

\subsection*{Primo nucleo}

Si ha
\[
A_2 - 5I =
\begin{pmatrix}
0 & 0 & 1\\
1 & 0 & \dfrac12\\
0 & 0 & 0
\end{pmatrix}.
\]

Risolvendo
\[
(A_2-5I)x=0
\]
si trova
\[
\ker(A_2-5I)=\operatorname{span}\left\{
\begin{pmatrix}
0\\
1\\
0
\end{pmatrix}
\right\},
\]
quindi
\[
\mu_1 = 1.
\]

Dunque c'è un solo blocco di Jordan associato a $\lambda=5$.

\subsection*{Secondo nucleo}

Calcolando il quadrato:
\[
(A_2-5I)^2 =
\begin{pmatrix}
0 & 0 & 0\\
0 & 0 & 1\\
0 & 0 & 0
\end{pmatrix}.
\]

Quindi
\[
\ker(A_2-5I)^2
=
\operatorname{span}\left\{
\begin{pmatrix}
0\\
1\\
0
\end{pmatrix},
\begin{pmatrix}
1\\
0\\
0
\end{pmatrix}
\right\},
\]
e dunque
\[
\mu_2 = 2.
\]

\subsection*{Terzo nucleo}

Infine
\[
(A_2-5I)^3 = 0,
\]
per cui
\[
\ker(A_2-5I)^3 = \R^3,
\qquad
\mu_3=3.
\]

Quindi:
\[
n_1 = \mu_1 = 1,\qquad
n_2 = \mu_2 - \mu_1 = 1,\qquad
n_3 = \mu_3 - \mu_2 = 1.
\]

\paragraph{Conclusione.}
Esiste una sola catena, di lunghezza $3$.  
Quindi la forma di Jordan deve essere
\[
J_2=
\begin{pmatrix}
5 & 1 & 0\\
0 & 5 & 1\\
0 & 0 & 5
\end{pmatrix}.
\]

\subsection*{Costruzione della catena}

Scegliamo un vettore in $\ker(A_2-5I)^3 \setminus \ker(A_2-5I)^2$:
\[
q^{(3)}=
\begin{pmatrix}
0\\
0\\
1
\end{pmatrix}.
\]

Allora
\[
q^{(2)} = (A_2-5I)q^{(3)}
=
\begin{pmatrix}
1\\
\dfrac12\\
0
\end{pmatrix},
\]
e poi
\[
q^{(1)} = (A_2-5I)q^{(2)}
=
\begin{pmatrix}
0\\
1\\
0
\end{pmatrix}.
\]

Costruiamo
\[
Q_2=
\begin{pmatrix}
q^{(1)} & q^{(2)} & q^{(3)}
\end{pmatrix}
=
\begin{pmatrix}
0 & 1 & 0\\
1 & \dfrac12 & 0\\
0 & 0 & 1
\end{pmatrix}.
\]

Allora
\[
Q_2^{-1}A_2Q_2 = J_2.
\]

\paragraph{Osservazione strutturale.}
Una matrice come
\[
A_1=
\begin{pmatrix}
5 & 1 & 0\\
0 & 5 & 1\\
0 & 0 & 5
\end{pmatrix}
\]
è già in forma di Jordan; $A_2$ invece non lo è, ma è \emph{simile} ad essa.  
Dal punto di vista dinamico, le due matrici hanno gli stessi modi temporali.

\section{Significato strutturale dei blocchi di Jordan}

Se
\[
A = QJQ^{-1},
\]
allora
\[
e^{At} = Qe^{Jt}Q^{-1}.
\]

Questo significa che le matrici $Q$ e $Q^{-1}$ mescolano solo linearmente le componenti dello stato: gli \emph{andamenti nel tempo} veri e propri compaiono dentro $e^{Jt}$.

Quindi, per capire i modi temporali di un sistema, bisogna guardare i blocchi di Jordan.

\subsection{Un solo blocco di ordine 3}

Per esempio, se
\[
J=
\begin{pmatrix}
5 & 1 & 0\\
0 & 5 & 1\\
0 & 0 & 5
\end{pmatrix},
\]
allora
\[
e^{Jt}
=
e^{5t}
\begin{pmatrix}
1 & t & \dfrac{t^2}{2}\\
0 & 1 & t\\
0 & 0 & 1
\end{pmatrix}.
\]

Quindi i modi che compaiono sono:
\[
e^{5t},\qquad te^{5t},\qquad t^2 e^{5t}.
\]

Più precisamente il fattore $\frac{t^2}{2}$ cambia solo il coefficiente, non il tipo di andamento: il modo strutturalmente nuovo è il termine polinomiale in $t^2$ moltiplicato per $e^{5t}$.

\section{Esempio 2: lettura dei blocchi dai numeri \texorpdfstring{$\mu_j$}{muj}}

Supponiamo ora di avere un autovalore $\lambda=0$ e di aver trovato:
\[
\mu_1=2,\qquad \mu_2=4,\qquad \mu_3=5.
\]

Allora
\[
n_1 = \mu_1 = 2,
\qquad
n_2 = \mu_2-\mu_1 = 2,
\qquad
n_3 = \mu_3-\mu_2 = 1.
\]

Come si interpreta questo risultato?

\begin{itemize}
    \item $n_1=2$ significa che ci sono due catene, quindi due blocchi di Jordan.
    \item $n_2=2$ significa che entrambe le catene arrivano almeno all'ordine $2$.
    \item $n_3=1$ significa che una sola delle due catene arriva all'ordine $3$.
\end{itemize}

Quindi la struttura dei blocchi è necessariamente:
\[
J = \operatorname{diag}\!\left(J_0^{(2)}, J_0^{(3)}\right)
\]
oppure la stessa a blocchi permutati.

\begin{figure}[H]
\centering
\begin{tikzpicture}[scale=0.85]
    % blocco da 3
    \draw[fill=blue!15] (0,0) rectangle (1,1);
    \draw[fill=blue!15] (0,1) rectangle (1,2);
    \draw[fill=blue!15] (0,2) rectangle (1,3);

    % blocco da 2
    \draw[fill=green!15] (1.8,0) rectangle (2.8,1);
    \draw[fill=green!15] (1.8,1) rectangle (2.8,2);

    \node at (0.5,3.35) {$3$};
    \node at (2.3,2.35) {$2$};

    \node at (0.5,-0.4) {\small una catena di lunghezza $3$};
    \node at (2.3,-0.4) {\small una catena di lunghezza $2$};
\end{tikzpicture}
\caption{Lettura delle lunghezze delle catene dai valori $\mu_1=2$, $\mu_2=4$, $\mu_3=5$.}
\end{figure}

\subsection{Andamento temporale associato}

Nel continuo, poiché $\lambda=0$, si ha
\[
e^{J_0^{(2)}t}
=
\begin{pmatrix}
1 & t\\
0 & 1
\end{pmatrix},
\qquad
e^{J_0^{(3)}t}
=
\begin{pmatrix}
1 & t & \dfrac{t^2}{2}\\
0 & 1 & t\\
0 & 0 & 1
\end{pmatrix}.
\]

Dunque
\[
e^{Jt}
=
\begin{pmatrix}
1 & t & 0 & 0 & 0\\
0 & 1 & 0 & 0 & 0\\
0 & 0 & 1 & t & \dfrac{t^2}{2}\\
0 & 0 & 0 & 1 & t\\
0 & 0 & 0 & 0 & 1
\end{pmatrix}.
\]

I modi presenti sono quindi del tipo:
\[
1,\qquad t,\qquad t^2.
\]

\paragraph{Interpretazione.}
Questo è un punto molto importante:
\begin{itemize}
    \item il modo costante $1$ corrisponde a una componente che non decade;
    \item il termine $t$ introduce crescita polinomiale;
    \item il termine $t^2$ introduce una crescita ancora più forte.
\end{itemize}

Quindi il sistema non è asintoticamente stabile.

\section{I modi del sistema e il ruolo dei blocchi di Jordan}

Possiamo ora formulare con precisione il messaggio strutturale della lezione.

\subsection{Nel tempo continuo}

Se un blocco di Jordan reale di ordine $r$ è associato all'autovalore $\lambda$, allora i modi che compaiono sono:
\begin{equation}
e^{\lambda t},\quad
te^{\lambda t},\quad
t^2 e^{\lambda t},\quad
\dots,\quad
t^{r-1}e^{\lambda t}.
\label{eq:L8_modes_CT}
\end{equation}

\subsection{Nel tempo discreto}

In modo analogo, nel discreto compaiono andamenti del tipo
\begin{equation}
\lambda^k,\quad
k\lambda^k,\quad
k^2\lambda^k,\quad
\dots
\label{eq:L8_modes_DT}
\end{equation}
a meno di coefficienti costanti o equivalenze del tipo $\binom{k}{j}\lambda^{k-j}$.

\paragraph{Osservazione.}
Dal punto di vista qualitativo, ciò che conta è la presenza del fattore polinomiale in $t$ oppure in $k$.

\section{Quando i blocchi di Jordan diventano decisivi}

Per autovalori ``non critici'', il segno della parte reale (nel continuo) o il modulo (nel discreto) basta spesso a descrivere il comportamento asintotico.

Ma per autovalori \emph{critici} bisogna guardare il blocco di Jordan massimo associato.

\subsection{Tempo continuo}

Nel continuo, il caso critico è
\[
\Re(\lambda)=0.
\]

Infatti:
\begin{itemize}
    \item se $\Re(\lambda)<0$, anche eventuali fattori polinomiali vengono dominati dal decadimento esponenziale;
    \item se $\Re(\lambda)>0$, il sistema diverge comunque;
    \item se $\Re(\lambda)=0$, la presenza di blocchi di ordine superiore a $1$ introduce termini tipo $t$, $t^2$, \dots, e può far perdere la stabilità.
\end{itemize}

\subsection{Tempo discreto}

Nel discreto, il caso critico è
\[
|\lambda|=1.
\]

Infatti:
\begin{itemize}
    \item se $|\lambda|<1$, il sistema tende a zero;
    \item se $|\lambda|>1$, il sistema diverge;
    \item se $|\lambda|=1$, la presenza di blocchi di Jordan di ordine superiore a $1$ introduce fattori tipo $k$, $k^2$, \dots, e può rendere il sistema instabile.
\end{itemize}

\paragraph{Morale pratica.}
Per ogni autovalore critico bisogna sempre controllare la dimensione del blocco di Jordan più grande a esso associato.

\section{Osservazione conclusiva sul significato dei modi}

Quando si porta una matrice in forma di Jordan, non lo si fa per un gusto puramente algebrico, ma perché nella matrice esponenziale dei blocchi di Jordan si leggono direttamente i modi temporali del sistema.

In altre parole:
\[
A = QJQ^{-1}
\quad \Longrightarrow \quad
e^{At} = Qe^{Jt}Q^{-1}.
\]

La matrice $Q$ cambia coordinate, ma non cambia la natura degli andamenti.  
Gli andamenti elementari li troviamo in $e^{Jt}$.

Per questo, i modi del sistema dipendono da:
\begin{itemize}
    \item gli autovalori;
    \item il loro valore numerico;
    \item la loro molteplicità algebrica;
    \item la dimensione dei blocchi di Jordan associati.
\end{itemize}

\section{Sintesi finale della lezione}

In questa lezione abbiamo visto che:

\begin{enumerate}
    \item gli autovettori generalizzati di una stessa catena sono linearmente indipendenti;
    \item una catena di lunghezza $k$ genera un sottospazio ciclico $A$-invariante di dimensione $k$;
    \item i numeri
    \[
    \mu_j=\dim\ker(A-\lambda I)^j
    \quad\text{e}\quad
    n_j=\mu_j-\mu_{j-1}
    \]
    permettono di ricostruire la struttura dei blocchi di Jordan;
    \item un blocco di Jordan di ordine $r$ associato a $\lambda$ produce nel continuo modi del tipo
    \[
    e^{\lambda t},\ te^{\lambda t},\ \dots,\ t^{r-1}e^{\lambda t};
    \]
    \item nel discreto compaiono analogamente termini del tipo
    \[
    \lambda^k,\ k\lambda^k,\ k^2\lambda^k,\dots;
    \]
    \item nei casi critici, cioè $\Re(\lambda)=0$ nel continuo oppure $|\lambda|=1$ nel discreto, la dimensione del blocco di Jordan è decisiva per stabilire il comportamento del sistema.
\end{enumerate}

Questa è la vera ragione per cui si studia la forma di Jordan: essa rende leggibili, in modo immediato, i modi strutturali del sistema.
\chapter{L9 -- Analisi modale e stabilità a partire dalla forma di Jordan}

\section{Obiettivo della lezione}

Dopo aver introdotto la forma di Jordan e le catene di autovettori generalizzati, il passo naturale successivo è capire \emph{che cosa si vede nel tempo}.

L'idea fondamentale è questa:

\begin{quote}
se porto la matrice di stato in forma di Jordan, allora nell'evoluzione libera del sistema compaiono in modo esplicito tutti i possibili andamenti temporali.
\end{quote}

Questi andamenti temporali elementari vengono chiamati \textbf{modi del sistema}.  
La struttura dei modi dipende da:
\begin{itemize}
    \item il valore degli autovalori;
    \item il fatto che siano reali oppure complessi;
    \item la dimensione del blocco di Jordan più grande associato a ciascun autovalore.
\end{itemize}

Questa lezione serve quindi a collegare in modo chiaro:
\[
\text{autovalori e blocchi di Jordan}
\qquad \Longleftrightarrow \qquad
\text{modi temporali del sistema}
\qquad \Longleftrightarrow \qquad
\text{stabilità / instabilità}.
\]

\section{Perché i blocchi di Jordan sono importanti}

Nel tempo continuo, l'evoluzione libera è
\[
x(t)=e^{At}x(0),
\]
mentre nel tempo discreto è
\[
x(k)=A^k x(0).
\]

Se $A$ è simile a una forma di Jordan $J$, cioè
\[
A=QJQ^{-1},
\]
allora
\[
e^{At}=Qe^{Jt}Q^{-1}
\qquad \text{(tempo continuo)},
\]
e
\[
A^k=QJ^kQ^{-1}
\qquad \text{(tempo discreto)}.
\]

Questo significa che:
\begin{itemize}
    \item le matrici $Q$ e $Q^{-1}$ fanno solo un cambio di coordinate;
    \item i veri andamenti nel tempo compaiono dentro $e^{Jt}$ oppure dentro $J^k$.
\end{itemize}

Di conseguenza, per fare analisi modale, bisogna guardare i blocchi di Jordan.

\section{Modi critici: perché compaiono i termini polinomiali}

Supponiamo di avere, nel continuo, un blocco di Jordan reale associato a un autovalore critico.  
Se ad esempio compare il blocco
\[
J_0^{(3)}=
\begin{pmatrix}
0 & 1 & 0\\
0 & 0 & 1\\
0 & 0 & 0
\end{pmatrix},
\]
allora
\[
e^{J_0^{(3)}t}
=
I + J_0^{(3)}t + \frac{(J_0^{(3)})^2 t^2}{2!}
=
\begin{pmatrix}
1 & t & \dfrac{t^2}{2}\\
0 & 1 & t\\
0 & 0 & 1
\end{pmatrix}.
\]

Quindi i modi che compaiono sono:
\[
1,\qquad t,\qquad t^2.
\]

\paragraph{Interpretazione.}
\begin{itemize}
    \item Il modo $1$ è un andamento costante: non esplode, ma nemmeno converge a zero.
    \item Il modo $t$ diverge lentamente, in modo polinomiale.
    \item Il modo $t^2$ diverge ancora più rapidamente.
\end{itemize}

Per questo motivo, quando un autovalore è \emph{critico} (nel continuo $\Re(\lambda)=0$, nel discreto $|\lambda|=1$), non basta conoscere il suo valore: bisogna sapere anche la dimensione del blocco di Jordan più grande a esso associato.

\begin{figure}[H]
\centering
\begin{tikzpicture}[>=Latex, node distance=1.4cm]
    \node[draw, rounded corners, align=center] (b1) {$J^{(1)}_0$};
    \node[draw, rounded corners, align=center, right=of b1] (m1) {$1$};

    \node[draw, rounded corners, align=center, below=1.1cm of b1] (b2) {$J^{(2)}_0$};
    \node[draw, rounded corners, align=center, right=of b2] (m2) {$1,\; t$};

    \node[draw, rounded corners, align=center, below=1.1cm of b2] (b3) {$J^{(3)}_0$};
    \node[draw, rounded corners, align=center, right=of b3] (m3) {$1,\; t,\; t^2$};

    \draw[->, thick] (b1) -- (m1);
    \draw[->, thick] (b2) -- (m2);
    \draw[->, thick] (b3) -- (m3);
\end{tikzpicture}
\caption{Un blocco di Jordan critico di ordine crescente genera modi polinomiali di grado crescente.}
\end{figure}

\section{Esempio strutturale: lettura dei modi da una matrice a blocchi}

Consideriamo una matrice a blocchi triangolare superiore del tipo
\begin{equation}
A=
\begin{pmatrix}
A_1^\star & B_{12}\\
0 & A_2
\end{pmatrix},
\label{eq:L9_Ablock}
\end{equation}
con
\[
A_1^\star=
\begin{pmatrix}
1 & 0 & 0\\
-1 & -1 & 1\\
-3 & -1 & -1
\end{pmatrix},
\qquad
A_2=
\begin{pmatrix}
1 & 0 & 0\\
-1 & -2 & 0\\
0 & 1 & -2
\end{pmatrix}.
\]

Il blocco di accoppiamento $B_{12}$ può essere non nullo, ma gli autovalori della matrice triangolare a blocchi sono comunque l'unione degli autovalori dei blocchi diagonali.

\subsection{Autovalori dei blocchi diagonali}

Per $A_1^\star$ si riconoscono:
\begin{itemize}
    \item un autovalore reale $\lambda=1$;
    \item una coppia complessa coniugata $\lambda=-1\pm j$.
\end{itemize}

Per $A_2$ si hanno:
\begin{itemize}
    \item un autovalore reale $\lambda=1$;
    \item l'autovalore $\lambda=-2$ con molteplicità algebrica $2$.
\end{itemize}

Quindi, per l'intera matrice $A$, gli autovalori sono:
\[
\lambda_1 = 1 \quad \text{con } m_a(1)=2,
\]
\[
\lambda_2 = -2 \quad \text{con } m_a(-2)=2,
\]
\[
\lambda_{3,4} = -1\pm j \quad \text{con molteplicità }1.
\]

\paragraph{Primo commento importante.}
Sapere gli autovalori non basta ancora.  
Dobbiamo capire \emph{quanto sono grandi i blocchi di Jordan} associati a ciascuno di essi.

\subsection{Autovalore \texorpdfstring{$\lambda=1$}{lambda=1}}

Dallo studio del rango di $A-I$ si trova
\[
\rank(A-I)=5,
\]
quindi
\[
\mu_1(1)=\dim\ker(A-I)=6-5=1.
\]

Poiché la molteplicità algebrica di $\lambda=1$ è $2$, ma la molteplicità geometrica è $1$, segue che per $\lambda=1$ c'è \emph{un solo blocco di Jordan}, necessariamente di ordine $2$:
\[
J^{(2)}_{1}=
\begin{pmatrix}
1 & 1\\
0 & 1
\end{pmatrix}.
\]

\subsection{Autovalore \texorpdfstring{$\lambda=-2$}{lambda=-2}}

In modo analogo, studiando $A+2I$ si trova
\[
\rank(A+2I)=5,
\]
quindi
\[
\mu_1(-2)=\dim\ker(A+2I)=1.
\]

Anche qui la molteplicità algebrica è $2$ ma la molteplicità geometrica è $1$, dunque per $\lambda=-2$ compare un solo blocco di Jordan di ordine $2$:
\[
J^{(2)}_{-2}=
\begin{pmatrix}
-2 & 1\\
0 & -2
\end{pmatrix}.
\]

\subsection{Coppia complessa \texorpdfstring{$-1\pm j$}{-1 +- j}}

La coppia complessa coniugata $-1\pm j$ compare una sola volta, quindi nel formalismo reale è associata a un unico blocco reale $2\times 2$:
\[
M=
\begin{pmatrix}
-1 & 1\\
-1 & -1
\end{pmatrix}.
\]

\paragraph{Conclusione.}
La forma di Jordan reale della matrice avrà la struttura
\begin{equation}
J=
\operatorname{diag}
\left(
J_1^{(2)},
J_{-2}^{(2)},
M
\right),
\label{eq:L9_J_example}
\end{equation}
a meno di una diversa permutazione dei blocchi sulla diagonale.

\section{Modi che compaiono nell'esempio}

Una volta trovata la struttura di Jordan \eqref{eq:L9_J_example}, i modi si leggono immediatamente.

\subsection{Contributo del blocco \texorpdfstring{$J_1^{(2)}$}{J1}}

Nel continuo:
\[
e^{J_1^{(2)}t}
=
e^t
\begin{pmatrix}
1 & t\\
0 & 1
\end{pmatrix},
\]
quindi i modi associati sono
\[
e^t,
\qquad
t e^t.
\]

\subsection{Contributo del blocco \texorpdfstring{$J_{-2}^{(2)}$}{J-2}}

Si ha
\[
e^{J_{-2}^{(2)}t}
=
e^{-2t}
\begin{pmatrix}
1 & t\\
0 & 1
\end{pmatrix},
\]
quindi i modi associati sono
\[
e^{-2t},
\qquad
t e^{-2t}.
\]

\subsection{Contributo del blocco complesso reale \texorpdfstring{$M$}{M}}

Per il blocco
\[
M=
\begin{pmatrix}
-1 & 1\\
-1 & -1
\end{pmatrix}
=
- I + 
\begin{pmatrix}
0 & 1\\
-1 & 0
\end{pmatrix},
\]
si ottiene
\[
e^{Mt}
=
e^{-t}
\begin{pmatrix}
\cos t & \sin t\\
-\sin t & \cos t
\end{pmatrix}.
\]

Dunque i modi associati alla coppia complessa sono:
\[
e^{-t}\cos t,
\qquad
e^{-t}\sin t.
\]

\paragraph{Riassumendo.}
I modi dell'esempio sono:
\begin{equation}
e^t,\qquad te^t,\qquad e^{-2t},\qquad te^{-2t},\qquad e^{-t}\cos t,\qquad e^{-t}\sin t.
\label{eq:L9_modes_example}
\end{equation}

\paragraph{Interpretazione.}
\begin{itemize}
    \item $e^t$ e $te^t$ divergono;
    \item $e^{-2t}$ e $te^{-2t}$ convergono a zero;
    \item $e^{-t}\cos t$ ed $e^{-t}\sin t$ sono oscillazioni smorzate.
\end{itemize}

Quindi il sistema dell'esempio è instabile, perché compaiono modi divergenti.

\section{Osservazione importante: non sempre serve costruire tutte le catene}

Se l'obiettivo è capire gli \emph{andamenti nel tempo}, spesso non è necessario costruire esplicitamente tutte le catene di autovettori generalizzati.

Infatti:
\begin{itemize}
    \item per sapere quali modi compaiono, basta conoscere gli autovalori e la dimensione del blocco di Jordan più grande a essi associato;
    \item la costruzione delle catene serve invece quando si vuole anche esplicitare il cambio di coordinate $Q$ e quindi scrivere in dettaglio le componenti dello stato nelle coordinate originali.
\end{itemize}

\section{Esempio: come ottenere una traiettoria costante}

Consideriamo il sistema continuo
\begin{equation}
\dot{x}(t)=Ax(t),
\qquad
A=
\begin{pmatrix}
6 & -3\\
2 & -1
\end{pmatrix}.
\label{eq:L9_Asmall}
\end{equation}

Vogliamo trovare $x(0)$ tale che la traiettoria sia costante, cioè
\[
x(t)=
\begin{pmatrix}
k_1\\
k_2
\end{pmatrix}
\qquad \forall t.
\]

\subsection*{Autovalori}

Si ha
\[
\det(A-\lambda I)=
\det\begin{pmatrix}
6-\lambda & -3\\
2 & -1-\lambda
\end{pmatrix}
=
(6-\lambda)(-1-\lambda)+6
=
\lambda(\lambda-5).
\]

Quindi gli autovalori sono
\[
\lambda_1=0,
\qquad
\lambda_2=5.
\]

\subsection*{Autovettori}

Per $\lambda=0$:
\[
Ax=0
\quad \Longrightarrow \quad
\ker A = \operatorname{span}\left\{
\begin{pmatrix}
1\\
2
\end{pmatrix}
\right\}.
\]

Per $\lambda=5$:
\[
(A-5I)x=0
\quad \Longrightarrow \quad
\ker(A-5I)=\operatorname{span}\left\{
\begin{pmatrix}
3\\
1
\end{pmatrix}
\right\}.
\]

Possiamo quindi scegliere
\[
Q=
\begin{pmatrix}
1 & 3\\
2 & 1
\end{pmatrix},
\qquad
J=
\begin{pmatrix}
0 & 0\\
0 & 5
\end{pmatrix}.
\]

Allora
\[
e^{At}=Qe^{Jt}Q^{-1}
=
Q
\begin{pmatrix}
1 & 0\\
0 & e^{5t}
\end{pmatrix}
Q^{-1}.
\]

\paragraph{Idea chiave.}
Per ottenere una traiettoria costante, devo eliminare il modo $e^{5t}$ e lasciare solo il modo associato all'autovalore nullo.

Questo accade se la condizione iniziale appartiene all'autospazio di $\lambda=0$, cioè se
\[
x(0)=\alpha
\begin{pmatrix}
1\\
2
\end{pmatrix}.
\]

\paragraph{Conclusione.}
Tutte e sole le condizioni iniziali del tipo
\begin{equation}
x(0)=\alpha
\begin{pmatrix}
1\\
2
\end{pmatrix},
\qquad \alpha\in\R,
\label{eq:L9_constant_traj}
\end{equation}
generano traiettorie costanti.

\paragraph{Interpretazione modale.}
Se scelgo una condizione iniziale nel sottospazio associato al modo $1=e^{0t}$, allora lo stato non vede il modo esponenziale $e^{5t}$ e rimane costante nel tempo.

\section{Analisi modale generale nel tempo continuo}

\subsection{Autovalore reale}

Sia $\lambda_i \in \R$ un autovalore e sia $p_i$ la dimensione del blocco di Jordan \emph{più grande} associato a $\lambda_i$.

Allora i modi associati a $\lambda_i$ sono del tipo
\begin{equation}
e^{\lambda_i t},\qquad
t e^{\lambda_i t},\qquad
t^2 e^{\lambda_i t},\qquad
\dots,\qquad
t^{p_i-1} e^{\lambda_i t}.
\label{eq:L9_modes_real_CT}
\end{equation}

\paragraph{Commento.}
\begin{itemize}
    \item Se $p_i=1$, compare solo il modo esponenziale puro.
    \item Se $p_i>1$, compaiono anche fattori polinomiali in $t$.
\end{itemize}

\subsection{Autovalore complesso}

Se invece
\[
\lambda_i = \sigma_i + j\omega_i,
\qquad
\overline{\lambda_i} = \sigma_i - j\omega_i,
\]
allora nel formalismo reale i modi compaiono a coppie:
\begin{equation}
e^{\sigma_i t}\cos(\omega_i t),
\qquad
e^{\sigma_i t}\sin(\omega_i t).
\label{eq:L9_modes_complex_CT_base}
\end{equation}

Se il blocco reale di Jordan associato ha dimensione $q_i$ (equivalentemente, se nel complesso il blocco associato ha ordine $q_i$), allora compaiono anche i fattori polinomiali:
\begin{equation}
e^{\sigma_i t}\cos(\omega_i t),\;
e^{\sigma_i t}\sin(\omega_i t),\;
t e^{\sigma_i t}\cos(\omega_i t),\;
t e^{\sigma_i t}\sin(\omega_i t),\;
\dots,\;
t^{q_i-1} e^{\sigma_i t}\cos(\omega_i t),\;
t^{q_i-1} e^{\sigma_i t}\sin(\omega_i t).
\label{eq:L9_modes_complex_CT}
\end{equation}

\paragraph{Interpretazione.}
Questi modi rappresentano oscillazioni:
\begin{itemize}
    \item smorzate se $\sigma_i<0$;
    \item amplificate se $\sigma_i>0$;
    \item non smorzate se $\sigma_i=0$.
\end{itemize}

\begin{figure}[H]
\centering
\begin{tikzpicture}[>=Latex, node distance=1.25cm]
    \node[draw, rounded corners, align=center] (r1) {$J_{\lambda}^{(1)}$\\$\lambda\in\R$};
    \node[draw, rounded corners, align=center, right=3.8cm of r1] (m1) {$e^{\lambda t}$};

    \node[draw, rounded corners, align=center, below=1.2cm of r1] (r2) {$J_{\lambda}^{(3)}$\\$\lambda\in\R$};
    \node[draw, rounded corners, align=center, right=3.8cm of r2] (m2) {$e^{\lambda t},\; t e^{\lambda t},\; t^2 e^{\lambda t}$};

    \node[draw, rounded corners, align=center, below=1.2cm of r2] (c1) {$\text{blocco reale }M$\\$\lambda=\sigma\pm j\omega$};
    \node[draw, rounded corners, align=center, right=3.8cm of c1] (m3) {$e^{\sigma t}\cos(\omega t),\; e^{\sigma t}\sin(\omega t)$};

    \draw[->, thick] (r1) -- (m1);
    \draw[->, thick] (r2) -- (m2);
    \draw[->, thick] (c1) -- (m3);
\end{tikzpicture}
\caption{Ogni blocco di Jordan genera una famiglia caratteristica di modi temporali.}
\end{figure}

\section{Analisi modale generale nel tempo discreto}

Nel discreto, se
\[
x(k+1)=Ax(k),
\]
allora
\[
x(k)=A^k x(0)=QJ^kQ^{-1}x(0).
\]

\subsection{Autovalore reale}

Se $\lambda_i\in\R$ e il blocco di Jordan più grande associato ha dimensione $p_i$, allora i modi sono del tipo
\begin{equation}
\lambda_i^k,\qquad
k\lambda_i^k,\qquad
k^2\lambda_i^k,\qquad
\dots,\qquad
k^{p_i-1}\lambda_i^k,
\label{eq:L9_modes_real_DT}
\end{equation}
a meno di costanti moltiplicative o equivalenze del tipo $\binom{k}{r}\lambda_i^{k-r}$.

\paragraph{Commento.}
Per l'analisi qualitativa, quello che conta davvero è la presenza del fattore polinomiale in $k$.

\subsection{Autovalore complesso}

Se
\[
\lambda_i=\sigma_i+j\omega_i=\rho_i e^{j\theta_i},
\qquad
\overline{\lambda_i}=\rho_i e^{-j\theta_i},
\]
allora i modi reali di base sono
\begin{equation}
\rho_i^k \cos(k\theta_i),
\qquad
\rho_i^k \sin(k\theta_i).
\label{eq:L9_modes_complex_DT_base}
\end{equation}

Se il blocco reale di Jordan più grande associato alla coppia ha ordine $q_i$, allora compaiono anche i fattori polinomiali:
\begin{equation}
\rho_i^k \cos(k\theta_i),\;
\rho_i^k \sin(k\theta_i),\;
k\rho_i^k \cos(k\theta_i),\;
k\rho_i^k \sin(k\theta_i),\;
\dots,\;
k^{q_i-1}\rho_i^k \cos(k\theta_i),\;
k^{q_i-1}\rho_i^k \sin(k\theta_i).
\label{eq:L9_modes_complex_DT}
\end{equation}

\paragraph{Osservazione.}
In formule più precise, nei blocchi discreti possono comparire anche potenze del tipo $\rho_i^{k-r}$; tuttavia, dal punto di vista asintotico, la struttura è quella di un termine oscillatorio-amplificato/smorzato moltiplicato per un polinomio in $k$.

\section{Tutti questi andamenti sono i modi del sistema}

Possiamo ora formulare chiaramente il concetto.

\paragraph{Definizione operativa.}
Le funzioni di tempo che compaiono nelle componenti di
\[
e^{Jt}
\qquad \text{oppure} \qquad
J^k
\]
sono i \textbf{modi propri del sistema}.

Nel continuo i modi sono funzioni di $t$; nel discreto sono successioni in $k$.

\paragraph{Messaggio chiave.}
Una volta nota la forma di Jordan, i modi si leggono direttamente e quindi si può prevedere qualitativamente l'evoluzione dello stato senza risolvere ogni volta il sistema da zero.

\section{Stabilità letta direttamente dai modi}

Adesso usiamo i modi per capire quando il sistema converge, diverge oppure resta limitato.

\subsection{Tempo continuo}

Sia $\lambda_i$ un autovalore di $A$.

\paragraph{Caso 1: \texorpdfstring{$\Re(\lambda_i)<0$}{Re(lambda)<0}.}
Tutti i modi associati hanno il fattore
\[
e^{\Re(\lambda_i)t},
\]
che tende a zero esponenzialmente. Anche se compaiono fattori polinomiali $t^r$, il decadimento esponenziale domina, quindi i modi convergono a zero.

\paragraph{Caso 2: \texorpdfstring{$\Re(\lambda_i)>0$}{Re(lambda)>0}.}
I modi divergono esponenzialmente. Eventuali fattori polinomiali rendono solo più marcata la divergenza, ma il comportamento qualitativo è già deciso dal segno positivo della parte reale.

\paragraph{Caso 3: \texorpdfstring{$\Re(\lambda_i)=0$}{Re(lambda)=0}.}
Questo è il caso critico.

\begin{itemize}
    \item Se il blocco di Jordan massimo associato a $\lambda_i$ ha dimensione $1$, allora compaiono solo modi costanti oppure oscillatori limitati, per esempio
    \[
    1,\qquad \cos(\omega t),\qquad \sin(\omega t).
    \]
    \item Se invece il blocco massimo ha dimensione maggiore di $1$, compaiono modi con fattori polinomiali:
    \[
    t,\; t^2,\; \dots,\; t^r,
    \qquad
    t\cos(\omega t),\; t\sin(\omega t),\; \dots
    \]
    e quindi il sistema diventa polinomialmente divergente.
\end{itemize}

\paragraph{Riassunto per il continuo.}
\begin{itemize}
    \item $\Re(\lambda_i)<0$ $\Rightarrow$ modi convergenti;
    \item $\Re(\lambda_i)>0$ $\Rightarrow$ modi divergenti;
    \item $\Re(\lambda_i)=0$ $\Rightarrow$ bisogna guardare la dimensione del blocco di Jordan massimo.
\end{itemize}

\subsection{Tempo discreto}

Nel discreto il ruolo di $\Re(\lambda_i)$ è svolto dal modulo $|\lambda_i|$.

\paragraph{Caso 1: \texorpdfstring{$|\lambda_i|<1$}{|lambda|<1}.}
Tutti i modi convergono a zero. Se $\lambda_i$ è complesso, possono esserci oscillazioni, ma con ampiezza decrescente.

\paragraph{Caso 2: \texorpdfstring{$|\lambda_i|>1$}{|lambda|>1}.}
I modi divergono esponenzialmente. Se $\lambda_i$ è complesso, la divergenza può essere oscillatoria.

\paragraph{Caso 3: \texorpdfstring{$|\lambda_i|=1$}{|lambda|=1}.}
Anche qui siamo nel caso critico.

\begin{itemize}
    \item Se il blocco di Jordan massimo ha dimensione $1$, i modi restano limitati:
    \[
    1,\qquad \cos(k\theta),\qquad \sin(k\theta).
    \]
    \item Se il blocco ha dimensione $>1$, compaiono modi con fattori $k$, $k^2$, \dots, quindi il sistema diventa polinomialmente divergente.
\end{itemize}

\paragraph{Riassunto per il discreto.}
\begin{itemize}
    \item $|\lambda_i|<1$ $\Rightarrow$ modi convergenti;
    \item $|\lambda_i|>1$ $\Rightarrow$ modi divergenti;
    \item $|\lambda_i|=1$ $\Rightarrow$ bisogna guardare la dimensione del blocco di Jordan massimo.
\end{itemize}

\section{Schema riassuntivo finale}

\subsection*{Tempo continuo}

\begin{center}
\renewcommand{\arraystretch}{1.4}
\begin{tabular}{>{\raggedright\arraybackslash}p{4.2cm} >{\raggedright\arraybackslash}p{9.5cm}}
\toprule
Condizione sugli autovalori & Comportamento dei modi \\
\midrule
$\Re(\lambda_i)<0$ &
Convergenza esponenziale a zero, indipendentemente dalla dimensione dei blocchi di Jordan associati. \\
\addlinespace
$\Re(\lambda_i)>0$ &
Divergenza esponenziale, indipendentemente dalla dimensione dei blocchi di Jordan associati. \\
\addlinespace
$\Re(\lambda_i)=0$ e blocco massimo di dimensione $1$ &
Modi costanti oppure oscillatori limitati e non smorzati. \\
\addlinespace
$\Re(\lambda_i)=0$ e blocco massimo di dimensione $>1$ &
Modi polinomialmente divergenti: compaiono termini del tipo $t, t^2, \dots$. \\
\bottomrule
\end{tabular}
\end{center}

\subsection*{Tempo discreto}

\begin{center}
\renewcommand{\arraystretch}{1.4}
\begin{tabular}{>{\raggedright\arraybackslash}p{4.2cm} >{\raggedright\arraybackslash}p{9.5cm}}
\toprule
Condizione sugli autovalori & Comportamento dei modi \\
\midrule
$|\lambda_i|<1$ &
Convergenza a zero; se $\lambda_i$ è complesso possono esserci oscillazioni smorzate. \\
\addlinespace
$|\lambda_i|>1$ &
Divergenza esponenziale; se $\lambda_i$ è complesso possono esserci oscillazioni amplificate. \\
\addlinespace
$|\lambda_i|=1$ e blocco massimo di dimensione $1$ &
Modi costanti oppure oscillatori a ampiezza costante. \\
\addlinespace
$|\lambda_i|=1$ e blocco massimo di dimensione $>1$ &
Modi polinomialmente divergenti: compaiono termini del tipo $k, k^2, \dots$. \\
\bottomrule
\end{tabular}
\end{center}

\section{Conclusione della lezione}

Il contenuto concettuale della lezione può essere riassunto in una frase:

\begin{quote}
gli andamenti temporali di un sistema lineare sono completamente determinati dagli autovalori e dalla struttura dei blocchi di Jordan associati.
\end{quote}

Più precisamente:
\begin{itemize}
    \item gli autovalori fissano il tipo base di andamento (esponenziale, oscillatorio, crescente, decrescente);
    \item la dimensione dei blocchi di Jordan introduce eventualmente fattori polinomiali in $t$ oppure in $k$;
    \item nei casi critici, questi fattori polinomiali sono decisivi per distinguere tra comportamento limitato e comportamento divergente.
\end{itemize}

Per questo motivo la forma di Jordan non è soltanto uno strumento algebrico: è lo strumento che rende visibile, in modo immediato, la struttura modale del sistema.

\vspace{0.5cm}

Nel seguito del corso, questa idea tornerà continuamente:
\[
\text{spettro + Jordan}
\quad \Longrightarrow \quad
\text{modi}
\quad \Longrightarrow \quad
\text{stabilità, risposta libera, comportamento asintotico}.
\]
\chapter{L10 -- Modi dominanti, spazio delle fasi ed equivalenza algebrica}

\section{Obiettivo della lezione}

In questa lezione si tirano le fila del lavoro fatto nelle lezioni precedenti su:
\begin{itemize}
    \item evoluzione libera dei sistemi lineari;
    \item forma di Jordan;
    \item modi del sistema;
    \item interpretazione geometrica nello spazio delle fasi.
\end{itemize}

L'idea centrale è la seguente:

\begin{quote}
conoscere gli autovalori e la struttura dei blocchi di Jordan permette di capire \emph{quali andamenti temporali} possono comparire nello stato e, quindi, di capire stabilità, instabilità, velocità di convergenza e dinamica dominante.
\end{quote}

Inoltre, si introduce una seconda idea fondamentale del corso:

\begin{quote}
due sistemi ottenuti l'uno dall'altro mediante un cambio invertibile di coordinate di stato sono \emph{algebricamente equivalenti}: cambiano le coordinate con cui descriviamo lo stato, ma non cambia la sostanza dinamica del sistema.
\end{quote}

\section{Richiamo: evoluzione libera nei casi continuo e discreto}

Per un sistema lineare stazionario a tempo continuo
\[
\dot x(t)=Ax(t)+Bu(t), \qquad x(0)=x_0,
\]
la \emph{evoluzione libera} si ottiene ponendo $u(t)\equiv 0$:
\[
\dot x(t)=Ax(t), \qquad x(0)=x_0,
\]
e vale
\begin{equation}
x(t)=e^{At}x_0.
\label{eq:L10_free_CT}
\end{equation}

Per un sistema lineare stazionario a tempo discreto
\[
x(k+1)=Ax(k)+Bu(k), \qquad x(0)=x_0,
\]
l'evoluzione libera si ottiene ponendo $u(k)\equiv 0$:
\[
x(k+1)=Ax(k), \qquad x(0)=x_0,
\]
e vale
\begin{equation}
x(k)=A^k x_0.
\label{eq:L10_free_DT}
\end{equation}

Queste due formule sono il punto di partenza di tutta l'analisi modale:
\begin{itemize}
    \item nel continuo studiamo $e^{At}$;
    \item nel discreto studiamo $A^k$.
\end{itemize}

Se $A$ è simile a una forma di Jordan $J$, cioè se
\[
A=QJQ^{-1},
\]
allora
\[
e^{At}=Qe^{Jt}Q^{-1},
\qquad
A^k=QJ^kQ^{-1}.
\]
Dunque i \emph{modi temporali} veri e propri compaiono dentro $e^{Jt}$ e $J^k$.

\section{Richiamo sui modi elementari}

Dalle lezioni precedenti sappiamo che:

\begin{itemize}
    \item se $\lambda\in\mathbb{R}$ è un autovalore reale, compaiono modi del tipo
    \[
    e^{\lambda t},\quad te^{\lambda t},\quad t^2 e^{\lambda t},\ \dots
    \qquad \text{(tempo continuo)},
    \]
    e
    \[
    \lambda^k,\quad k\lambda^k,\quad k^2\lambda^k,\ \dots
    \qquad \text{(tempo discreto)};
    \]
    \item se $\lambda=\sigma+j\omega$ è complesso, allora nel continuo compaiono modi del tipo
    \[
    e^{\sigma t}\cos(\omega t), \qquad e^{\sigma t}\sin(\omega t),
    \]
    ed eventualmente anche
    \[
    t e^{\sigma t}\cos(\omega t),\quad
    t e^{\sigma t}\sin(\omega t),\quad
    t^2 e^{\sigma t}\cos(\omega t),\quad
    t^2 e^{\sigma t}\sin(\omega t),\ \dots
    \]
    se il blocco di Jordan associato è di ordine maggiore di $1$;
    \item nel discreto, se $\lambda=\rho e^{j\theta}$, compaiono modi del tipo
    \[
    \rho^k\cos(k\theta),\qquad \rho^k\sin(k\theta),
    \]
    ed eventualmente anche con fattori polinomiali in $k$.
\end{itemize}

\subsection{Interpretazione immediata dei casi reali}

Per $\lambda\in\mathbb{R}$ nel continuo:
\begin{itemize}
    \item se $\lambda<0$, il modo si smorza;
    \item se $\lambda>0$, il modo diverge;
    \item se $\lambda=0$, il modo è critico: può essere costante, oppure polinomialmente divergente se compaiono fattori $t, t^2,\dots$
\end{itemize}

Nel discreto:
\begin{itemize}
    \item se $|\lambda|<1$, il modo converge a zero;
    \item se $|\lambda|>1$, il modo diverge;
    \item se $|\lambda|=1$, il modo è critico.
\end{itemize}

\section{Grafici qualitativi dei modi nel continuo}

I grafici seguenti non servono per fare calcoli numerici, ma per allenare l'intuizione: guardando questi andamenti bisogna imparare a riconoscere subito che tipo di autovalore o di blocco di Jordan c'è dietro.

\subsection{Modi reali non oscillatori}

\begin{figure}[H]
\centering

\begin{subfigure}{0.31\textwidth}
\centering
\begin{tikzpicture}
\begin{axis}[
    width=\textwidth,
    height=5cm,
    axis lines=middle,
    xmin=0, xmax=4,
    ymin=0, ymax=6,
    xtick=\empty, ytick=\empty,
    title={$e^{\sigma t}$}
]
\addplot[domain=0:4, samples=200, thick] {1};
\addplot[domain=0:4, samples=200, thick] {exp(x/1.6)};
\addplot[domain=0:4, samples=200, thick] {exp(-x/1.2)};
\end{axis}
\end{tikzpicture}
\caption{Caso base}
\end{subfigure}
\hfill
\begin{subfigure}{0.31\textwidth}
\centering
\begin{tikzpicture}
\begin{axis}[
    width=\textwidth,
    height=5cm,
    axis lines=middle,
    xmin=0, xmax=4,
    ymin=0, ymax=6,
    xtick=\empty, ytick=\empty,
    title={$t\,e^{\sigma t}$}
]
\addplot[domain=0:4, samples=200, thick] {x};
\addplot[domain=0:4, samples=200, thick] {x*exp(x/3)};
\addplot[domain=0:4, samples=200, thick] {x*exp(-x)};
\end{axis}
\end{tikzpicture}
\caption{Con fattore lineare}
\end{subfigure}
\hfill
\begin{subfigure}{0.31\textwidth}
\centering
\begin{tikzpicture}
\begin{axis}[
    width=\textwidth,
    height=5cm,
    axis lines=middle,
    xmin=0, xmax=4,
    ymin=0, ymax=8,
    xtick=\empty, ytick=\empty,
    title={$t^2e^{\sigma t}$}
]
\addplot[domain=0:4, samples=200, thick] {x^2};
\addplot[domain=0:4, samples=200, thick] {x^2*exp(x/3.5)};
\addplot[domain=0:4, samples=200, thick] {x^2*exp(-x)};
\end{axis}
\end{tikzpicture}
\caption{Con fattore quadratico}
\end{subfigure}

\caption{Andamenti qualitativi dei modi reali nel continuo: nello stesso grafico si leggono i casi $\sigma=0$, $\sigma>0$ e $\sigma<0$.}
\label{fig:L10_real_modes_CT}
\end{figure}

\paragraph{Lettura della figura.}
\begin{itemize}
    \item La curva costante $1$ corrisponde al caso $\sigma=0$ e blocco di Jordan di ordine $1$.
    \item Le curve $t$ e $t^2$ corrispondono ancora al caso $\sigma=0$, ma con blocchi di Jordan di ordine maggiore.
    \item Le curve crescenti corrispondono a $\sigma>0$.
    \item Le curve che, dopo un transitorio, vanno a zero corrispondono a $\sigma<0$.
\end{itemize}

Un punto molto importante è il seguente:

\begin{quote}
nel continuo, anche se compare un fattore polinomiale $t^r$, se $\sigma<0$ l'esponenziale vince comunque e il modo tende a zero.
\end{quote}

Ad esempio
\[
t e^{-t}\to 0,
\qquad
t^2 e^{-t}\to 0,
\qquad
t^{10}e^{-t}\to 0
\qquad \text{per } t\to +\infty.
\]

\subsection{Modi oscillatori nel continuo}

Se $\lambda=\sigma+j\omega$, i modi reali si ottengono prendendo parte reale e immaginaria di $e^{\lambda t}$, cioè
\[
e^{\sigma t}\cos(\omega t),
\qquad
e^{\sigma t}\sin(\omega t).
\]

Se il blocco di Jordan associato ha ordine maggiore di $1$, compaiono anche
\[
t e^{\sigma t}\cos(\omega t),\qquad
t e^{\sigma t}\sin(\omega t),
\]
oppure
\[
t^2 e^{\sigma t}\cos(\omega t),\qquad
t^2 e^{\sigma t}\sin(\omega t),
\]
e così via.

\begin{figure}[H]
\centering

\begin{subfigure}{0.31\textwidth}
\centering
\begin{tikzpicture}
\begin{axis}[
    width=\textwidth,
    height=5cm,
    axis lines=middle,
    xmin=0, xmax=10,
    ymin=-1.5, ymax=1.5,
    xtick=\empty, ytick=\empty,
    title={$e^{\sigma t}\sin(\omega t)$}
]
\addplot[domain=0:10, samples=300, thick] {sin(deg(2*x))};
\addplot[domain=0:10, samples=300, thick, dashed] {exp(-x/4)};
\addplot[domain=0:10, samples=300, thick, dashed] {-exp(-x/4)};
\addplot[domain=0:10, samples=300, thick, dashed] {exp(x/12)};
\addplot[domain=0:10, samples=300, thick, dashed] {-exp(x/12)};
\end{axis}
\end{tikzpicture}
\caption{Ordine $1$}
\end{subfigure}
\hfill
\begin{subfigure}{0.31\textwidth}
\centering
\begin{tikzpicture}
\begin{axis}[
    width=\textwidth,
    height=5cm,
    axis lines=middle,
    xmin=0, xmax=10,
    ymin=-3.5, ymax=3.5,
    xtick=\empty, ytick=\empty,
    title={$t\,e^{\sigma t}\sin(\omega t)$}
]
\addplot[domain=0:10, samples=300, thick] {x*sin(deg(2*x))/4};
\addplot[domain=0:10, samples=300, thick, dashed] {x*exp(-x/4)/2.7};
\addplot[domain=0:10, samples=300, thick, dashed] {-x*exp(-x/4)/2.7};
\addplot[domain=0:10, samples=300, thick, dashed] {x*exp(x/12)/7};
\addplot[domain=0:10, samples=300, thick, dashed] {-x*exp(x/12)/7};
\end{axis}
\end{tikzpicture}
\caption{Ordine $2$}
\end{subfigure}
\hfill
\begin{subfigure}{0.31\textwidth}
\centering
\begin{tikzpicture}
\begin{axis}[
    width=\textwidth,
    height=5cm,
    axis lines=middle,
    xmin=0, xmax=10,
    ymin=-5, ymax=5,
    xtick=\empty, ytick=\empty,
    title={$t^2e^{\sigma t}\sin(\omega t)$}
]
\addplot[domain=0:10, samples=300, thick] {x^2*sin(deg(2*x))/18};
\addplot[domain=0:10, samples=300, thick, dashed] {x^2*exp(-x/4)/14};
\addplot[domain=0:10, samples=300, thick, dashed] {-x^2*exp(-x/4)/14};
\addplot[domain=0:10, samples=300, thick, dashed] {x^2*exp(x/12)/42};
\addplot[domain=0:10, samples=300, thick, dashed] {-x^2*exp(x/12)/42};
\end{axis}
\end{tikzpicture}
\caption{Ordine $3$}
\end{subfigure}

\caption{Modi oscillatori nel continuo. Le linee tratteggiate sono gli inviluppi. Se $\sigma<0$ l'oscillazione è smorzata, se $\sigma=0$ è limitata, se $\sigma>0$ è amplificata.}
\label{fig:L10_complex_modes_CT}
\end{figure}

\paragraph{Osservazione.}
Se $\sigma=0$, l'inviluppo è costante: quindi compaiono oscillazioni non smorzate.
Se invece compaiono fattori $t$, $t^2$, \dots\ anche con $\sigma=0$, l'ampiezza cresce polinomialmente.

\section{Un primo messaggio pratico: i modi permettono di riconoscere il sistema}

Se conosco autovalori e dimensioni dei blocchi di Jordan, allora conosco i modi.
Viceversa, osservando i modi, posso risalire alla struttura spettrale qualitativa del sistema.

Per esempio:
\begin{itemize}
    \item se vedo un termine $e^{-t}$, penso a un autovalore reale negativo;
    \item se vedo un termine $te^{-t}$, penso a un autovalore reale negativo con blocco di Jordan di ordine almeno $2$;
    \item se vedo un termine $e^{-t}\sin t$, penso a una coppia complessa con parte reale negativa;
    \item se vedo un termine $t e^{-t}\sin t$, penso a una coppia complessa con blocco reale di Jordan di ordine almeno $2$;
    \item se vedo un termine costante oppure oscillatorio limitato, penso a un caso critico.
\end{itemize}

Questa capacità di lettura qualitativa è molto importante nella progettazione dei controllori: non interessa solo sapere \emph{quanto vale} un autovalore, ma soprattutto \emph{che andamento induce nel tempo}.

\section{Dominanza modale}

\subsection{Idea intuitiva}

Nella pratica, soprattutto nel controllo, capita spesso che il sistema abbia molti autovalori, ma solo alcuni siano davvero rilevanti per l'andamento a lungo termine.

Il concetto chiave è questo:

\begin{quote}
il modo che decade più lentamente è quello che, a lungo andare, resta visibile; quelli che decadono più velocemente diventano trascurabili.
\end{quote}

Per esempio, nel continuo:
\[
e^{-t}
\qquad \text{e} \qquad
e^{-10t}.
\]
Entrambi vanno a zero, ma $e^{-10t}$ ci va molto più velocemente.  
Infatti
\[
\frac{e^{-10t}}{e^{-t}}=e^{-9t}\to 0.
\]
Quindi, per tempi grandi, il comportamento del sistema è dominato da $e^{-t}$.

\subsection{Polo dominante nel continuo}

Supponiamo di avere autovalori
\[
\lambda_1,\lambda_2,\dots,\lambda_n
\]
e che uno di essi abbia parte reale più grande di tutti:
\[
\Re(\lambda_{\mathrm{dom}})=\max_i \Re(\lambda_i).
\]

Allora, \emph{a meno di cancellazioni dovute alla particolare condizione iniziale}, il comportamento a lungo termine è dominato dai modi associati a $\lambda_{\mathrm{dom}}$.

Se il blocco di Jordan più grande associato a $\lambda_{\mathrm{dom}}$ ha ordine $r$, allora l'andamento dominante sarà del tipo
\begin{equation}
t^{r-1}e^{\lambda_{\mathrm{dom}}t}.
\label{eq:L10_dom_CT}
\end{equation}

\paragraph{Perché compare proprio $t^{r-1}$?}
Perché, per un blocco di Jordan di ordine $r$, i modi associati sono
\[
e^{\lambda t},\quad
t e^{\lambda t},\quad
t^2 e^{\lambda t},\ \dots,\quad
t^{r-1}e^{\lambda t},
\]
e quello con grado polinomiale massimo è il più lento a scomparire rispetto agli altri associati allo stesso autovalore.

\paragraph{Osservazione importante.}
Quando due autovalori hanno parti reali diverse, vince quello con parte reale più grande.
Quando invece hanno la stessa parte reale, allora per capire chi domina bisogna guardare:
\begin{itemize}
    \item il grado del fattore polinomiale;
    \item la presenza di eventuali oscillazioni;
    \item la condizione iniziale.
\end{itemize}

\subsection{Polo dominante nel discreto}

Nel discreto il criterio analogo non è la parte reale, ma il modulo.

Se
\[
|\lambda_{\mathrm{dom}}|=\max_i |\lambda_i|,
\]
allora, a meno di cancellazioni, il comportamento a lungo termine è dominato dai modi associati a $\lambda_{\mathrm{dom}}$.

Se il blocco di Jordan massimo associato ha ordine $r$, allora il modo dominante è del tipo
\begin{equation}
k^{r-1}\lambda_{\mathrm{dom}}^k.
\label{eq:L10_dom_DT}
\end{equation}

Se $\lambda_{\mathrm{dom}}$ è complesso, si avranno oscillazioni moltiplicate da un inviluppo dominato da $|\lambda_{\mathrm{dom}}|^k$.

\subsection{Perché la dominanza modale è importante nel controllo}

Nella progettazione dei controllori spesso si impongono specifiche come:
\begin{itemize}
    \item tempo di assestamento;
    \item sovraelongazione;
    \item rapidità di convergenza;
    \item smorzamento delle oscillazioni.
\end{itemize}

Anche se il sistema a ciclo chiuso può essere di ordine elevato, spesso si cerca di far sì che la dinamica visibile sia dominata da uno o due modi principali, spesso assimilabili a un comportamento del secondo ordine.

Per questo si parla di \textbf{poli dominanti}: sono i poli che determinano la dinamica osservabile a lungo termine.

\section{Un esempio molto semplice di dominanza}

Supponiamo che nel continuo compaiano i due modi
\[
e^{-t}, \qquad e^{-10t}.
\]
Allora
\[
e^{-10t}=e^{-t}e^{-9t},
\]
e il fattore $e^{-9t}$ tende molto rapidamente a zero.

Quindi, per tempi grandi,
\[
c_1e^{-t}+c_2e^{-10t}\approx c_1e^{-t}.
\]

In altre parole:
\begin{quote}
il modo più lento è quello che resta.
\end{quote}

Allo stesso modo, nel discreto, tra
\[
(0.8)^k
\qquad \text{e} \qquad
(0.2)^k
\]
domina $(0.8)^k$, perché
\[
\frac{(0.2)^k}{(0.8)^k}=\left(\frac14\right)^k\to 0.
\]

\section{Spazio delle fasi}

\subsection{Significato geometrico}

Se il sistema ha due variabili di stato, per esempio $x_1$ e $x_2$, allora possiamo rappresentare lo stato come un punto del piano:
\[
x=
\begin{pmatrix}
x_1\\
x_2
\end{pmatrix}.
\]

Lo \textbf{spazio delle fasi} è proprio lo spazio delle variabili di stato.  
Ogni traiettoria del sistema diventa quindi una curva nel piano $(x_1,x_2)$.

Questo è estremamente utile perché ci permette di capire:
\begin{itemize}
    \item se le traiettorie si avvicinano all'origine;
    \item se si allontanano;
    \item se esistono direzioni privilegiate;
    \item se l'equilibrio è stabile, instabile o di tipo sella.
\end{itemize}

\subsection{Sistema omogeneo 2D}

Consideriamo il sistema omogeneo
\[
\dot x = Ax
\]
oppure
\[
x(k+1)=Ax(k).
\]

L'origine $x=0$ è sempre un punto di equilibrio.  
Il tipo di equilibrio dipende dagli autovalori e dagli autovettori (o, più in generale, dalla struttura di Jordan).

\paragraph{Messaggio importante.}
Non basta conoscere solo gli autovalori: per capire \emph{come} si dispongono le traiettorie nel piano, servono anche le direzioni eigenspaziali e quindi, in generale, la struttura completa di $A$.

\subsection{Ritratti di fase qualitativi}

\begin{figure}[H]
\centering

\begin{subfigure}{0.31\textwidth}
\centering
\begin{tikzpicture}[scale=0.9,>=Latex]
\draw[->] (-2.2,0)--(2.2,0) node[right] {$x_1$};
\draw[->] (0,-2.2)--(0,2.2) node[above] {$x_2$};
\fill (0,0) circle (2pt);

\draw[thick,->] (1.8,1.8) .. controls (1,1) and (0.4,0.4) .. (0.08,0.08);
\draw[thick,->] (-1.8,-1.8) .. controls (-1,-1) and (-0.4,-0.4) .. (-0.08,-0.08);
\draw[thick,->] (1.8,-1.5) .. controls (0.8,-1) and (0.3,-0.4) .. (0.06,-0.06);
\draw[thick,->] (-1.6,1.7) .. controls (-0.8,0.9) and (-0.25,0.3) .. (-0.05,0.05);
\end{tikzpicture}
\caption{Nodo stabile}
\end{subfigure}
\hfill
\begin{subfigure}{0.31\textwidth}
\centering
\begin{tikzpicture}[scale=0.9,>=Latex]
\draw[->] (-2.2,0)--(2.2,0) node[right] {$x_1$};
\draw[->] (0,-2.2)--(0,2.2) node[above] {$x_2$};
\fill (0,0) circle (2pt);

\draw[thick,->] (0.08,0.08) .. controls (0.5,0.5) and (1.0,1.0) .. (1.8,1.8);
\draw[thick,->] (-0.08,-0.08) .. controls (-0.5,-0.5) and (-1.0,-1.0) .. (-1.8,-1.8);
\draw[thick,->] (0.06,-0.06) .. controls (0.3,-0.4) and (0.8,-1.0) .. (1.8,-1.5);
\draw[thick,->] (-0.05,0.05) .. controls (-0.25,0.3) and (-0.8,0.9) .. (-1.6,1.7);
\end{tikzpicture}
\caption{Nodo instabile}
\end{subfigure}
\hfill
\begin{subfigure}{0.31\textwidth}
\centering
\begin{tikzpicture}[scale=0.9,>=Latex]
\draw[->] (-2.2,0)--(2.2,0) node[right] {$x_1$};
\draw[->] (0,-2.2)--(0,2.2) node[above] {$x_2$};
\fill (0,0) circle (2pt);

\draw[thick,->] (-1.9,0)--(-0.15,0);
\draw[thick,->] (1.9,0)--(0.15,0);
\draw[thick,->] (0,0.15)--(0,1.9);
\draw[thick,->] (0,-0.15)--(0,-1.9);

\draw[thick,->] (-1.6,1.6) .. controls (-0.7,1.1) and (-0.2,0.5) .. (0,0.18);
\draw[thick,->] (1.6,-1.6) .. controls (0.7,-1.1) and (0.2,-0.5) .. (0,-0.18);

\draw[thick,->] (0,0.18) .. controls (0.2,0.5) and (0.7,1.1) .. (1.6,1.6);
\draw[thick,->] (0,-0.18) .. controls (-0.2,-0.5) and (-0.7,-1.1) .. (-1.6,-1.6);
\end{tikzpicture}
\caption{Sella}
\end{subfigure}

\caption{Ritratti di fase qualitativi per sistemi omogenei bidimensionali.}
\label{fig:L10_phase_portraits}
\end{figure}

\paragraph{Interpretazione.}
\begin{itemize}
    \item Se tutti i modi convergono, l'origine è stabile asintoticamente.
    \item Se tutti i modi divergono, l'origine è instabile.
    \item Se alcuni modi convergono e altri divergono, l'origine è una sella.
\end{itemize}

\section{Equivalenza algebrica tra sistemi lineari stazionari}

\subsection{Definizione}

Consideriamo due sistemi lineari stazionari
\[
\Sigma:
\begin{cases}
x^{+}=Ax+Bu\\
y=Cx+Du
\end{cases}
\qquad
\widetilde{\Sigma}:
\begin{cases}
\tilde x^{+}=\tilde A \tilde x+\tilde B u\\
y=\tilde C \tilde x+\tilde D u
\end{cases}
\]
dove con $x^+$ si intende:
\begin{itemize}
    \item $\dot x$ nel continuo;
    \item $x(k+1)$ nel discreto.
\end{itemize}

I due sistemi si dicono \textbf{algebricamente equivalenti} se esiste una matrice invertibile $T$ tale che
\begin{equation}
x=T\tilde x.
\label{eq:L10_state_change}
\end{equation}

Sostituendo \eqref{eq:L10_state_change} nel sistema originale si ottiene
\[
T\tilde x^{+}=AT\tilde x+Bu
\]
e dunque
\[
\tilde x^{+}=T^{-1}AT\,\tilde x + T^{-1}B\,u.
\]
Inoltre, siccome
\[
y=Cx+Du=CT\tilde x+Du,
\]
si ottiene
\[
y=\tilde C \tilde x+\tilde D u
\]
con
\[
\tilde C=CT,
\qquad
\tilde D=D.
\]

Quindi le relazioni di trasformazione sono:
\begin{equation}
\boxed{
\tilde A=T^{-1}AT,\qquad
\tilde B=T^{-1}B,\qquad
\tilde C=CT,\qquad
\tilde D=D.
}
\label{eq:L10_similarity_system}
\end{equation}

\subsection{Interpretazione}

Questa trasformazione non cambia il sistema in senso fisico o ingresso-uscita: cambia solo il modo in cui descriviamo internamente lo stato.

In altre parole:
\begin{quote}
i due sistemi rappresentano la stessa dinamica, ma in coordinate di stato diverse.
\end{quote}

Per questo motivo:
\begin{itemize}
    \item hanno gli stessi autovalori;
    \item hanno la stessa forma di Jordan a meno dell'ordinamento dei blocchi;
    \item hanno la stessa funzione di trasferimento;
    \item hanno le stesse proprietà strutturali fondamentali.
\end{itemize}

\subsection{Perché questo è importante}

Portare $A$ in forma di Jordan non significa cambiare la sostanza del sistema: significa scegliere coordinate più comode per leggere i modi del sistema.

Perciò tutto quello che abbiamo detto su $J$ vale anche per $A$, semplicemente letto in una base diversa.

\section{Evoluzione libera ed evoluzione forzata}

Fin qui abbiamo ragionato soprattutto sull'evoluzione libera, che dipende solo da $A$:
\[
x(t)=e^{At}x(0),
\qquad
x(k)=A^k x(0).
\]

Ma per il sistema completo l'evoluzione è
\begin{equation}
x(t)=e^{At}x(0)+\int_0^t e^{A(t-\tau)}B\,u(\tau)\,d\tau
\qquad \text{(continuo)},
\label{eq:L10_forced_CT}
\end{equation}
e
\begin{equation}
x(k)=A^k x(0)+\sum_{i=0}^{k-1}A^{k-1-i}B\,u(i)
\qquad \text{(discreto)}.
\label{eq:L10_forced_DT}
\end{equation}

\paragraph{Osservazione concettuale.}
\begin{itemize}
    \item L'evoluzione libera dipende solo da $A$.
    \item L'evoluzione forzata dipende invece sia da $A$ sia da $B$.
\end{itemize}

Questo anticipa il problema della \textbf{raggiungibilità} e della \textbf{controllabilità}:
\begin{quote}
dato un ingresso $u$, quali stati posso raggiungere?
\end{quote}

\subsection{Esempio intuitivo}

Se
\[
B=
\begin{pmatrix}
1\\0\\0\\0
\end{pmatrix},
\]
l'ingresso entra direttamente solo nella prima variabile di stato.

Ma questo non vuol dire che influenzi solo quella: la matrice $A$ può trasferire l'effetto dell'ingresso da una componente dello stato alle altre.

Quindi la capacità di controllare il sistema non dipende solo da dove entra l'ingresso, ma da come \emph{A e B insieme} propagano tale ingresso nello spazio di stato.

\section{Spazio ciclico e struttura dell'evoluzione libera}

L'evoluzione libera
\[
x(t)=e^{At}x(0),
\qquad
x(k)=A^k x(0)
\]
resta confinata nel sottospazio generato dalle potenze di $A$ applicate alla condizione iniziale.

In particolare, nel discreto si introduce il sottospazio ciclico
\[
S_{x_0}=\operatorname{span}\{x_0,Ax_0,A^2x_0,\dots,A^{p-1}x_0\},
\]
dove $p=\dim S_{x_0}$.

Nel continuo il senso geometrico è lo stesso: l'evoluzione libera resta nel sottospazio $A$-invariante generato da $x_0$.

Questo è coerente con l'analisi modale:
\begin{itemize}
    \item se $x_0$ eccita certi modi, essi compaiono nella traiettoria;
    \item se $x_0$ non ha componente lungo un certo sottospazio modale, quel modo può non comparire.
\end{itemize}

\section{Riepilogo finale sugli andamenti nel tempo}

Possiamo riassumere gli andamenti possibili nel modo seguente.

\subsection*{Tempo continuo}

\begin{itemize}
    \item Se $\lambda_i\in\mathbb{R}$, i modi associati sono
    \[
    e^{\lambda_i t},\quad
    t e^{\lambda_i t},\quad
    \dots,\quad
    t^{p_i-1}e^{\lambda_i t},
    \]
    dove $p_i$ è la dimensione del blocco di Jordan più grande associato a $\lambda_i$.
    \item Se $\lambda_i=\sigma_i+j\omega_i\in\mathbb{C}$, i modi reali associati sono
    \[
    e^{\sigma_i t}\cos(\omega_i t),\quad
    e^{\sigma_i t}\sin(\omega_i t),
    \]
    ed eventualmente anche con fattori $t, t^2,\dots$
\end{itemize}

\subsection*{Tempo discreto}

\begin{itemize}
    \item Se $\lambda_i\in\mathbb{R}$, i modi associati sono
    \[
    \lambda_i^k,\quad
    k\lambda_i^k,\quad
    \dots,\quad
    k^{p_i-1}\lambda_i^k.
    \]
    \item Se $\lambda_i=\rho_i e^{j\theta_i}\in\mathbb{C}$, i modi reali associati sono
    \[
    \rho_i^k\cos(k\theta_i),\quad
    \rho_i^k\sin(k\theta_i),
    \]
    eventualmente moltiplicati per $k, k^2,\dots$
\end{itemize}

\section{Messaggio conclusivo della lezione}

Questa lezione si può sintetizzare così:

\begin{enumerate}
    \item l'evoluzione libera di un sistema lineare è governata da $e^{At}$ oppure da $A^k$;
    \item gli andamenti temporali si leggono dai blocchi di Jordan;
    \item il comportamento a lungo termine è determinato dai \textbf{modi dominanti};
    \item due sistemi simili sono algebricamente equivalenti, cioè descrivono la stessa dinamica in coordinate diverse;
    \item quando si passa all'evoluzione forzata entra in gioco anche $B$, e questo apre naturalmente ai concetti di raggiungibilità e controllabilità.
\end{enumerate}

In particolare, per l'analisi qualitativa, bisogna abituarsi a pensare così:
\[
\text{autovalori} + \text{blocchi di Jordan}
\quad \Longrightarrow \quad
\text{modi}
\quad \Longrightarrow \quad
\text{stabilità, rapidità, oscillazioni, dominanza}.
\]

Questo è il vero ponte tra algebra lineare e comportamento dinamico dei sistemi.
\chapter{L11 -- Raggiungibilità e controllabilità all'origine}

\section{Idea della lezione}

In questa lezione introduciamo uno dei problemi centrali della teoria del controllo:

\begin{quote}
dato un sistema dinamico lineare, esiste un ingresso di controllo che permetta di portare lo stato dove vogliamo?
\end{quote}

Questa domanda può essere letta in due modi leggermente diversi.

\begin{itemize}
    \item \textbf{Raggiungibilità}: partendo dall'origine, quali stati posso raggiungere in un tempo finito?
    \item \textbf{Controllabilità all'origine}: partendo da uno stato iniziale assegnato, posso riportare il sistema all'origine in un tempo finito?
\end{itemize}

Nel caso dei sistemi lineari tempo-invarianti questi concetti si descrivono in modo molto chiaro attraverso:
\begin{itemize}
    \item la formula di evoluzione dello stato;
    \item uno spazio di stati raggiungibili;
    \item la matrice di raggiungibilità;
    \item nel continuo, anche il gramiano di raggiungibilità.
\end{itemize}

L’idea fondamentale da tenere a mente è questa:

\begin{quote}
la possibilità di muovere il sistema nello spazio di stato non dipende da un singolo ingresso scelto a mano, ma dalla struttura della coppia \((A,B)\).
\end{quote}

\section{Sistema lineare stazionario a tempo continuo}

Consideriamo un sistema lineare stazionario a tempo continuo
\begin{equation}
\Sigma_c:
\begin{cases}
\dot x(t)=Ax(t)+Bu(t),\\
y(t)=Cx(t)+Du(t),
\end{cases}
\qquad x(t_0)=x_0,
\label{eq:L11_sistema_tc}
\end{equation}
dove
\[
x(t)\in\mathbb{R}^n,\qquad
u(t)\in\mathbb{R}^m,\qquad
y(t)\in\mathbb{R}^p.
\]

La soluzione del sistema è
\begin{equation}
x(t)=e^{A(t-t_0)}x_0+\int_{t_0}^{t} e^{A(t-\tau)}Bu(\tau)\,d\tau.
\label{eq:L11_soluzione_tc}
\end{equation}

Questa formula separa in modo netto due contributi:
\begin{itemize}
    \item il termine
    \[
    e^{A(t-t_0)}x_0,
    \]
    che è l'\textbf{evoluzione libera} e dipende solo dallo stato iniziale;
    \item il termine
    \[
    \int_{t_0}^{t} e^{A(t-\tau)}Bu(\tau)\,d\tau,
    \]
    che è l'\textbf{evoluzione forzata} e dipende dall’ingresso.
\end{itemize}

\subsection{Riduzione al caso con stato iniziale nullo}

Valutiamo \eqref{eq:L11_soluzione_tc} all’istante finale \(t_f\):
\begin{equation}
x(t_f)=e^{A(t_f-t_0)}x(t_0)+\int_{t_0}^{t_f} e^{A(t_f-\tau)}Bu(\tau)\,d\tau.
\label{eq:L11_x_tf_tc}
\end{equation}

Portando a sinistra il termine di evoluzione libera otteniamo
\begin{equation}
x(t_f)-e^{A(t_f-t_0)}x(t_0)
=
\int_{t_0}^{t_f} e^{A(t_f-\tau)}Bu(\tau)\,d\tau.
\label{eq:L11_traslazione_tc}
\end{equation}

Questa formula è molto importante, perché ci dice che il problema

\begin{quote}
\emph{partendo da \(x(t_0)\), posso arrivare in \(x(t_f)\)?}
\end{quote}

è equivalente al problema

\begin{quote}
\emph{partendo dall’origine, posso arrivare nel vettore}
\[
\bar x(t_f):=x(t_f)-e^{A(t_f-t_0)}x(t_0)\, ?
\]
\end{quote}

Quindi, per studiare la raggiungibilità, basta concentrarsi sul caso \(x(t_0)=0\).

\section{Insieme degli stati raggiungibili dall'origine nel continuo}

Se \(x(t_0)=0\), la formula dello stato diventa
\begin{equation}
x(t_f)=\int_{t_0}^{t_f} e^{A(t_f-\tau)}Bu(\tau)\,d\tau.
\label{eq:L11_raggiungibile_tc}
\end{equation}

Questo porta alla seguente definizione.

\subsection*{Definizione}
Fissato \(t_f>t_0\), l’insieme degli stati raggiungibili dall’origine è
\begin{equation}
\mathcal{R}^{t_f}
=
\left\{
x\in\mathbb{R}^n \;:\;
\exists\,u(\cdot)\ \text{tale che}\
x=
\int_{t_0}^{t_f} e^{A(t_f-\tau)}Bu(\tau)\,d\tau
\right\}.
\label{eq:L11_def_R_tf}
\end{equation}

In altre parole, \(\mathcal{R}^{t_f}\) è l’insieme di tutti gli stati in cui posso portare il sistema all’istante \(t_f\), partendo dall’origine.

\subsection{Operatore di raggiungibilità}

Definiamo l’operatore lineare
\begin{equation}
\varphi:\mathcal{U}\to\mathbb{R}^n,
\qquad
\varphi(u)=\int_{t_0}^{t_f} e^{A(t_f-\tau)}Bu(\tau)\,d\tau,
\label{eq:L11_operatore_phi}
\end{equation}
dove \(\mathcal{U}\) è uno spazio di ingressi, ad esempio \(L^2([t_0,t_f],\mathbb{R}^m)\).

Allora
\begin{equation}
\mathcal{R}^{t_f}=\operatorname{Im}(\varphi).
\label{eq:L11_Rtf_immagine_phi}
\end{equation}

Quindi lo studio della raggiungibilità coincide con lo studio dell’immagine di un operatore lineare.

\begin{figure}[H]
\centering
\begin{tikzpicture}[>=Latex, node distance=4cm]
    \node[draw, circle, minimum size=1.7cm] (U) {$\mathcal U$};
    \node[draw, circle, minimum size=1.7cm, right of=U] (X) {$\mathbb{R}^n$};
    \draw[->, thick] (U) -- node[above] {$\varphi$} (X);
    \node[below=0.7cm of X] {\(\operatorname{Im}(\varphi)=\mathcal R^{t_f}\)};
\end{tikzpicture}
\caption{L’insieme raggiungibile dall’origine nel continuo è l’immagine dell’operatore di raggiungibilità.}
\label{fig:L11_phi}
\end{figure}

\section{Operatore aggiunto e gramiano di raggiungibilità}

Per analizzare meglio \(\operatorname{Im}(\varphi)\), introduciamo l’operatore aggiunto \(\varphi^\ast\).

\subsection{Prodotti scalari}

Sullo spazio di stato usiamo il prodotto scalare euclideo:
\[
\langle x_1,x_2\rangle_X=x_1^\top x_2.
\]

Sullo spazio degli ingressi \(\mathcal U=L^2([t_0,t_f],\mathbb{R}^m)\) usiamo
\[
\langle u_1,u_2\rangle_{\mathcal U}
=
\int_{t_0}^{t_f} u_1^\top(\tau)u_2(\tau)\,d\tau.
\]

Per definizione, \(\varphi^\ast:\mathbb{R}^n\to\mathcal U\) è l’operatore tale che
\[
\langle \varphi(u),x\rangle_X
=
\langle u,\varphi^\ast(x)\rangle_{\mathcal U}
\qquad
\forall u\in\mathcal U,\ \forall x\in\mathbb{R}^n.
\]

\subsection{Calcolo esplicito di \(\varphi^\ast\)}

Partiamo da
\[
\varphi(u)=\int_{t_0}^{t_f} e^{A(t_f-\tau)}Bu(\tau)\,d\tau.
\]
Allora
\[
\langle \varphi(u),x\rangle_X
=
x^\top\int_{t_0}^{t_f} e^{A(t_f-\tau)}Bu(\tau)\,d\tau
=
\int_{t_0}^{t_f} x^\top e^{A(t_f-\tau)}Bu(\tau)\,d\tau.
\]
Usando la proprietà di trasposizione,
\[
x^\top e^{A(t_f-\tau)}Bu(\tau)
=
u^\top(\tau)B^\top e^{A^\top(t_f-\tau)}x.
\]
Quindi
\[
\langle \varphi(u),x\rangle_X
=
\int_{t_0}^{t_f}
u^\top(\tau)\,B^\top e^{A^\top(t_f-\tau)}x\,d\tau
=
\langle u,\varphi^\ast(x)\rangle_{\mathcal U}.
\]
Si riconosce allora
\begin{equation}
\varphi^\ast(x)(\tau)=B^\top e^{A^\top(t_f-\tau)}x.
\label{eq:L11_phi_star}
\end{equation}

\subsection{Gramiano di raggiungibilità}

Consideriamo ora la composizione \(\varphi\varphi^\ast\):
\[
(\varphi\varphi^\ast)(x)
=
\int_{t_0}^{t_f}
e^{A(t_f-\tau)}B\Big(B^\top e^{A^\top(t_f-\tau)}x\Big)\,d\tau.
\]
Raccogliendo \(x\), otteniamo
\[
(\varphi\varphi^\ast)(x)
=
\left(
\int_{t_0}^{t_f}
e^{A(t_f-\tau)}BB^\top e^{A^\top(t_f-\tau)}\,d\tau
\right)x.
\]

Questo suggerisce di definire il \textbf{gramiano di raggiungibilità}
\begin{equation}
G_{\mathcal R}(t_f)
=
\int_{t_0}^{t_f}
e^{A(t_f-\tau)}BB^\top e^{A^\top(t_f-\tau)}\,d\tau.
\label{eq:L11_gramiano}
\end{equation}

Allora
\begin{equation}
\varphi\varphi^\ast(x)=G_{\mathcal R}(t_f)x.
\label{eq:L11_phi_phi_star}
\end{equation}

Poiché in dimensione finita vale
\[
\operatorname{Im}(\varphi)=\operatorname{Im}(\varphi\varphi^\ast),
\]
concludiamo che
\begin{equation}
\mathcal{R}^{t_f}
=
\operatorname{Im}(G_{\mathcal R}(t_f)).
\label{eq:L11_Rtf_gramiano}
\end{equation}

Quindi nel continuo lo spazio raggiungibile può essere studiato anche tramite il gramiano.

\section{Matrice di raggiungibilità nel continuo}

Introduciamo ora la classica matrice di raggiungibilità:
\begin{equation}
R=
\begin{bmatrix}
B & AB & A^2B & \cdots & A^{n-1}B
\end{bmatrix}.
\label{eq:L11_matrice_raggiungibilita}
\end{equation}

Il risultato fondamentale è che, per sistemi lineari stazionari,
\begin{equation}
\operatorname{Im}(G_{\mathcal R}(t_f))=\operatorname{Im}(R).
\label{eq:L11_gramiano_matrice}
\end{equation}

Quindi, nel continuo,
\begin{equation}
\mathcal{R}^{t_f}=\operatorname{Im}(R).
\label{eq:L11_Rtf_uguale_R}
\end{equation}

\subsection{Commento sul significato}

Questo è un fatto molto importante.  
Anche se il gramiano contiene un integrale e dipende apparentemente dall’istante finale \(t_f\), il sottospazio degli stati raggiungibili è lo stesso che si ottiene guardando la matrice finita
\[
R=
\begin{bmatrix}
B & AB & \cdots & A^{n-1}B
\end{bmatrix}.
\]

Quindi, in pratica, il problema della raggiungibilità si riduce a un semplice problema di rango.

\section{Sistema completamente raggiungibile}

\subsection*{Definizione}
La coppia \((A,B)\) si dice \textbf{completamente raggiungibile} se
\begin{equation}
\operatorname{rank}(R)=n.
\label{eq:L11_rank_full}
\end{equation}

In tal caso
\[
\operatorname{Im}(R)=\mathbb{R}^n,
\]
e quindi ogni stato è raggiungibile dall’origine in tempo finito.

\subsection{Interpretazione geometrica}

Se \(\operatorname{rank}(R)=r<n\), allora gli stati raggiungibili non riempiono tutto lo spazio di stato, ma solo un sottospazio di dimensione \(r\).

Per esempio:
\begin{itemize}
    \item se \(r=1\), gli stati raggiungibili stanno su una retta;
    \item se \(r=2\) in \(\mathbb{R}^3\), stanno su un piano;
    \item se \(r=n\), si riempie tutto \(\mathbb{R}^n\).
\end{itemize}

\begin{figure}[H]
\centering
\begin{tikzpicture}[>=Latex, scale=1.0]
    \draw[->] (0,0) -- (4,0) node[right] {$x_1$};
    \draw[->] (0,0) -- (0,3.5) node[above] {$x_2$};
    \draw[->] (0,0) -- (-1.6,-1.2) node[below left] {$x_3$};

    \filldraw[fill=gray!20, draw=black]
    (0.4,0.4) -- (2.8,0.7) -- (1.7,2.3) -- (-0.6,2.0) -- cycle;

    \node at (2.3,2.1) {\(\operatorname{Im}(R)\)};
\end{tikzpicture}
\caption{Se \(\operatorname{rank}(R)=2\) in \(\mathbb{R}^3\), gli stati raggiungibili formano un piano.}
\label{fig:L11_piano_raggiungibile}
\end{figure}

\section{Spazio raggiungibile da uno stato iniziale generico}

Se \(x_0\neq 0\), allora dallo sviluppo dello stato segue che
\[
x(t_f)=e^{A(t_f-t_0)}x_0+r,
\qquad r\in\mathcal{R}^{t_f}.
\]
Quindi l’insieme degli stati raggiungibili a partire da \(x_0\) è
\begin{equation}
\mathcal{R}^{t_f}_{x_0}=e^{A(t_f-t_0)}x_0+\mathcal{R}^{t_f}.
\label{eq:L11_R_x0}
\end{equation}

Questo insieme non è, in generale, un sottospazio vettoriale: è una \textbf{traslazione affine} dello spazio raggiungibile dall’origine.

\section{Caso discreto}

Passiamo ora al sistema lineare stazionario a tempo discreto:
\begin{equation}
\Sigma_d:
\begin{cases}
x(k+1)=Ax(k)+Bu(k),\\
y(k)=Cx(k)+Du(k).
\end{cases}
\label{eq:L11_sistema_td}
\end{equation}

Iterando la dinamica si ottiene
\begin{equation}
x(t)=A^t x_0+\sum_{k=0}^{t-1}A^{t-k-1}Bu(k).
\label{eq:L11_soluzione_td}
\end{equation}

Se \(x_0=0\), allora
\begin{equation}
x(t)=\sum_{k=0}^{t-1}A^{t-k-1}Bu(k).
\label{eq:L11_td_zero}
\end{equation}

Raccogliendo i termini rispetto agli ingressi:
\begin{equation}
x(t)=
\begin{bmatrix}
B & AB & A^2B & \cdots & A^{t-1}B
\end{bmatrix}
\begin{bmatrix}
u(t-1)\\
u(t-2)\\
\vdots\\
u(0)
\end{bmatrix}.
\label{eq:L11_td_fattorizzata}
\end{equation}

Definiamo allora la \textbf{matrice di raggiungibilità in \(t\) passi}
\begin{equation}
R_t=
\begin{bmatrix}
B & AB & A^2B & \cdots & A^{t-1}B
\end{bmatrix}.
\label{eq:L11_Rt}
\end{equation}

e il corrispondente insieme degli stati raggiungibili dall’origine in \(t\) passi:
\begin{equation}
\mathcal{R}_t=\operatorname{Im}(R_t).
\label{eq:L11_spazio_Rt}
\end{equation}

\subsection{Monotonia degli spazi raggiungibili}

Vale la catena crescente
\begin{equation}
\mathcal{R}_1\subseteq \mathcal{R}_2\subseteq \cdots \subseteq \mathcal{R}_t\subseteq \mathcal{R}_{t+1}\subseteq \cdots
\label{eq:L11_catena_Rt}
\end{equation}
perché tutto ciò che posso raggiungere in \(t\) passi posso ancora raggiungerlo in \(t+1\) passi, eventualmente scegliendo l’ultimo ingresso nullo.

Poiché lo spazio di stato ha dimensione \(n\), questa catena si stabilizza al più dopo \(n\) passi.

\subsection*{Definizione}
Il sistema discreto è \textbf{completamente raggiungibile} se esiste \(t\) tale che
\begin{equation}
\dim \mathcal{R}_t=n.
\label{eq:L11_completa_raggiungibilita_td}
\end{equation}

Per un sistema lineare tempo-invariante basta controllare la matrice
\[
R=
\begin{bmatrix}
B & AB & \cdots & A^{n-1}B
\end{bmatrix}.
\]

\section{Controllabilità all'origine}

Finora abbiamo chiesto: \emph{quali stati posso raggiungere partendo dall’origine?}  
Ora ribaltiamo il punto di vista: \emph{da quali stati iniziali posso tornare all’origine?}

\subsection{Definizione nel discreto}

Fissato \(t\), definiamo
\begin{equation}
\mathcal{C}_t=
\left\{
x_0\in\mathbb{R}^n \;:\;
\exists\,u(0),\dots,u(t-1)\ \text{tali che}\ x(t)=0
\right\}.
\label{eq:L11_Ct}
\end{equation}

Usando \eqref{eq:L11_soluzione_td}, la condizione \(x(t)=0\) diventa
\[
A^t x_0+\sum_{k=0}^{t-1}A^{t-k-1}Bu(k)=0.
\]
Quindi
\begin{equation}
-A^t x_0\in \mathcal{R}_t.
\label{eq:L11_zero_condition}
\end{equation}

\subsection{Raggiungibilità implica controllabilità all'origine}

Se il sistema è completamente raggiungibile in \(t\) passi, allora
\[
\mathcal{R}_t=\mathbb{R}^n.
\]
Quindi per ogni \(x_0\) vale automaticamente
\[
-A^t x_0\in\mathcal{R}_t,
\]
e dunque ogni stato iniziale può essere portato all’origine.

Perciò nel discreto vale
\begin{equation}
\text{raggiungibilità} \;\Longrightarrow\; \text{controllabilità all'origine}.
\label{eq:L11_ragg_impl_ctrl0_td}
\end{equation}

\subsection{Nel discreto il viceversa non vale in generale}

Nel discreto può accadere che
\[
\operatorname{Im}(A^t)\subseteq \mathcal{R}_t
\]
pur senza avere \(\mathcal{R}_t=\mathbb{R}^n\).

In questo caso:
\begin{itemize}
    \item ogni stato iniziale può essere portato all’origine;
    \item ma non ogni stato è raggiungibile partendo dall’origine.
\end{itemize}

Quindi nel discreto:
\begin{equation}
\boxed{
\text{raggiungibilità} \Longrightarrow \text{controllabilità all'origine},
\qquad
\text{ma non viceversa.}
}
\label{eq:L11_no_equiv_td}
\end{equation}

\section{Nel continuo raggiungibilità e controllabilità all'origine coincidono}

Nel continuo la situazione è diversa.  
Infatti la condizione per arrivare a zero all’istante \(t_f\) è
\[
0=e^{A(t_f-t_0)}x_0+\int_{t_0}^{t_f} e^{A(t_f-\tau)}Bu(\tau)\,d\tau,
\]
cioè
\[
-e^{A(t_f-t_0)}x_0\in \mathcal{R}^{t_f}.
\]
Ma la matrice esponenziale è sempre invertibile:
\[
e^{A(t_f-t_0)} \ \text{è invertibile}\qquad \forall t_f>t_0.
\]
Quindi il cambio di variabile
\[
z=e^{A(t_f-t_0)}x_0
\]
è biiettivo, e perciò

\begin{equation}
\boxed{
\text{nel continuo, controllabilità all'origine e raggiungibilità coincidono.}
}
\label{eq:L11_equiv_tc}
\end{equation}

\subsection{Significato intuitivo}

Nel continuo l’evoluzione libera è governata da una matrice invertibile.  
Per questo motivo, il problema
\begin{quote}
\emph{portare \(x_0\) a zero}
\end{quote}
è equivalente al problema
\begin{quote}
\emph{raggiungere un certo stato partendo dall’origine.}
\end{quote}

\section{Lo spazio raggiungibile è \(A\)-invariante}

Una proprietà strutturale molto importante è che lo spazio raggiungibile è \(A\)-invariante.

\subsection*{Proposizione}
Lo spazio
\[
\operatorname{Im}(R)
\]
è \(A\)-invariante, cioè
\[
x\in\operatorname{Im}(R)\quad\Longrightarrow\quad Ax\in\operatorname{Im}(R).
\]

\subsection*{Dimostrazione}
Se \(x\in\operatorname{Im}(R)\), allora esiste \(y\) tale che
\[
x=Ry=
\begin{bmatrix}
B & AB & A^2B & \cdots & A^{n-1}B
\end{bmatrix}y.
\]
Allora
\[
Ax=
\begin{bmatrix}
AB & A^2B & A^3B & \cdots & A^nB
\end{bmatrix}y.
\]
Per il teorema di Cayley--Hamilton, \(A^nB\) è combinazione lineare delle colonne
\[
B,\ AB,\ \dots,\ A^{n-1}B.
\]
Di conseguenza anche \(Ax\) è combinazione lineare delle colonne di \(R\), e quindi
\[
Ax\in\operatorname{Im}(R).
\]
\hfill\(\square\)

\subsection{Commento}

Questa proprietà sarà decisiva nella lezione successiva, quando riscriveremo il sistema separando:
\begin{itemize}
    \item la parte raggiungibile;
    \item la parte non raggiungibile.
\end{itemize}

Infatti il fatto che \(\operatorname{Im}(R)\) sia \(A\)-invariante permette di costruire una base di coordinate adatta a questa decomposizione.

\section{Esempio semplice di non completa raggiungibilità}

Consideriamo il sistema discreto
\[
A=
\begin{bmatrix}
1 & 0 & 0\\
0 & 0 & 0\\
1 & 0 & 0
\end{bmatrix},
\qquad
B=
\begin{bmatrix}
1\\
0\\
1
\end{bmatrix}.
\]
Allora
\[
AB=
\begin{bmatrix}
1\\
0\\
1
\end{bmatrix}=B,
\qquad
A^2B=AB=B.
\]
Quindi
\[
R=
\begin{bmatrix}
B & AB & A^2B
\end{bmatrix}
=
\begin{bmatrix}
1 & 1 & 1\\
0 & 0 & 0\\
1 & 1 & 1
\end{bmatrix}.
\]
Il rango è
\[
\operatorname{rank}(R)=1<3.
\]
Dunque il sistema non è completamente raggiungibile.

Tuttavia
\[
\operatorname{Im}(A)=\operatorname{span}\!\left\{\begin{bmatrix}1\\0\\1\end{bmatrix}\right\}
=
\operatorname{Im}(R_1).
\]
Questo significa che ogni \(-Ax_0\) appartiene già a \(R_1\), quindi il sistema è controllabile all’origine in un passo, pur non essendo completamente raggiungibile.

Questo esempio mostra in modo molto chiaro la differenza tra i due concetti nel discreto.

\section{Schema riassuntivo finale}

\begin{center}
\fbox{
\parbox{0.93\textwidth}{
\textbf{Tempo continuo}
\begin{itemize}
    \item
    \[
    x(t)=e^{A(t-t_0)}x_0+\int_{t_0}^{t}e^{A(t-\tau)}Bu(\tau)\,d\tau
    \]
    \item
    \[
    \mathcal{R}^{t_f}=\operatorname{Im}(R),
    \qquad
    R=\begin{bmatrix}B & AB & \cdots & A^{n-1}B\end{bmatrix}
    \]
    \item
    \[
    \operatorname{rank}(R)=n
    \iff
    \text{sistema completamente raggiungibile}
    \]
    \item
    \[
    \text{raggiungibilità}
    \iff
    \text{controllabilità all'origine}
    \]
\end{itemize}

\vspace{0.2cm}

\textbf{Tempo discreto}
\begin{itemize}
    \item
    \[
    x(t)=A^t x_0+\sum_{k=0}^{t-1}A^{t-k-1}Bu(k)
    \]
    \item
    \[
    \mathcal{R}_t=\operatorname{Im}(R_t),
    \qquad
    R_t=\begin{bmatrix}B & AB & \cdots & A^{t-1}B\end{bmatrix}
    \]
    \item
    \[
    \mathcal{R}_1\subseteq \mathcal{R}_2\subseteq\cdots
    \]
    \item
    \[
    \text{raggiungibilità}
    \Longrightarrow
    \text{controllabilità all'origine}
    \]
    ma in generale il viceversa \emph{non} vale.
\end{itemize}
}}
\end{center}

\section{Conclusione}

La lezione ha introdotto il primo vero strumento strutturale della teoria del controllo: lo \textbf{spazio raggiungibile}.  
Da qui in avanti l’idea sarà sempre la stessa:

\begin{quote}
capire se l’ingresso riesce a generare tutto lo spazio di stato, oppure solo una sua parte.
\end{quote}

Questa domanda verrà sviluppata ulteriormente nella prossima lezione, dove useremo l’\(A\)-invarianza dello spazio raggiungibile per costruire la \textbf{forma standard di raggiungibilità}.
\chapter{L12 -- Sottospazio di raggiungibilità e forma standard di raggiungibilità}

\section{Idea della lezione}

In questa lezione vogliamo rispondere a una domanda strutturale molto importante:

\begin{quote}
\emph{È possibile riscrivere il sistema separando chiaramente la parte raggiungibile da quella non raggiungibile?}
\end{quote}

La risposta è sì.  
Il punto di partenza è capire bene che lo \textbf{spazio di raggiungibilità} non è solo l’insieme degli stati che si possono raggiungere con un opportuno ingresso, ma è anche un sottospazio con una struttura molto precisa: è un sottospazio \textbf{\(A\)-invariante} che contiene l’immagine di \(B\).

Questa proprietà permette di scegliere un cambio di coordinate adatto, con cui il sistema assume una forma a blocchi nella quale:
\begin{itemize}
    \item una parte dello stato è effettivamente influenzabile dall’ingresso;
    \item l’altra parte evolve da sola, senza poter essere modificata dal controllo.
\end{itemize}

Questa è la cosiddetta \textbf{forma standard di raggiungibilità}.

\section{Richiamo: spazio di raggiungibilità}

Consideriamo un sistema lineare stazionario.

Nel caso tempo continuo:
\begin{equation}
\dot x(t)=Ax(t)+Bu(t),
\qquad x(t)\in\R^n,\quad u(t)\in\R^m.
\end{equation}

Nel caso tempo discreto:
\begin{equation}
x(k+1)=Ax(k)+Bu(k),
\qquad x(k)\in\R^n,\quad u(k)\in\R^m.
\end{equation}

La matrice di raggiungibilità associata alla coppia \((A,B)\) è
\begin{equation}
R=\begin{bmatrix}
B & AB & A^2B & \cdots & A^{n-1}B
\end{bmatrix}.
\end{equation}

Lo \textbf{spazio di raggiungibilità} è definito come
\begin{equation}
\mathcal{R}=\operatorname{Im}(R).
\end{equation}

Nel discreto, se si vuole enfatizzare il numero di passi, si considera anche
\begin{equation}
R_t=\begin{bmatrix}
B & AB & \cdots & A^{t-1}B
\end{bmatrix},
\qquad
\mathcal{R}_t=\operatorname{Im}(R_t).
\end{equation}

Il significato è il seguente:
\begin{itemize}
    \item \(\mathcal{R}_t\) contiene gli stati raggiungibili dall’origine in \(t\) passi;
    \item \(\mathcal{R}\) contiene tutti gli stati raggiungibili dall’origine e coincide con il sottospazio generato dagli effetti successivi di \(B,AB,\dots,A^{n-1}B\).
\end{itemize}

\section{Teorema fondamentale sul sottospazio di raggiungibilità}

\begin{teorema}
Lo spazio di raggiungibilità
\[
\mathcal{R}=\operatorname{Im}(R)
\]
è il più piccolo sottospazio \(A\)-invariante che contiene \(\operatorname{Im}(B)\).
\end{teorema}

Questo teorema è fondamentale, perché caratterizza \(\mathcal{R}\) non solo come insieme di stati ottenibili con un ingresso, ma come oggetto geometrico ben preciso.

\subsection{Dimostrazione}

Dimostriamo il risultato in tre passaggi.

\subsubsection*{1. \(\mathcal{R}\) contiene \(\operatorname{Im}(B)\)}

Per definizione,
\[
R=\begin{bmatrix}
B & AB & A^2B & \cdots & A^{n-1}B
\end{bmatrix}.
\]
Le colonne di \(B\) sono tra le colonne di \(R\), quindi ogni vettore di \(\operatorname{Im}(B)\) appartiene a \(\operatorname{Im}(R)\). Dunque
\[
\operatorname{Im}(B)\subseteq \mathcal{R}.
\]

\subsubsection*{2. \(\mathcal{R}\) è \(A\)-invariante}

Prendiamo un qualunque vettore \(x\in\mathcal{R}\).  
Allora esiste \(y\) tale che
\[
x=Ry.
\]
Moltiplicando per \(A\),
\[
Ax=A(Ry)=
\begin{bmatrix}
AB & A^2B & \cdots & A^nB
\end{bmatrix}y.
\]

Il problema apparente è il termine \(A^nB\), che non compare esplicitamente in \(R\).  
Però, per il teorema di Cayley--Hamilton, \(A^n\) si può esprimere come combinazione lineare di \(I,A,\dots,A^{n-1}\). Quindi anche \(A^nB\) è combinazione lineare di
\[
B,\,AB,\,A^2B,\dots,A^{n-1}B.
\]
Ne segue che \(Ax\) è ancora combinazione lineare delle colonne di \(R\), cioè
\[
Ax\in\operatorname{Im}(R)=\mathcal{R}.
\]
Pertanto \(\mathcal{R}\) è \(A\)-invariante.

\subsubsection*{3. \(\mathcal{R}\) è il più piccolo tra questi sottospazi}

Sia \(S\subseteq\R^n\) un sottospazio \(A\)-invariante tale che
\[
\operatorname{Im}(B)\subseteq S.
\]
Poiché \(S\) contiene \(\operatorname{Im}(B)\), contiene le colonne di \(B\).  
Essendo \(S\) \(A\)-invariante, contiene anche le colonne di \(AB\), poi di \(A^2B\), e così via. Per induzione,
\[
\operatorname{Im}(A^kB)\subseteq S
\qquad \forall k=0,1,\dots,n-1.
\]
Quindi tutte le colonne di \(R\) appartengono a \(S\), e dunque
\[
\mathcal{R}=\operatorname{Im}(R)\subseteq S.
\]

Abbiamo quindi dimostrato che \(\mathcal{R}\) è il più piccolo sottospazio \(A\)-invariante che contiene \(\operatorname{Im}(B)\).
\qed

\section{Interpretazione geometrica}

Il significato intuitivo del teorema è molto chiaro:

\begin{itemize}
    \item l’ingresso entra nel sistema attraverso \(B\), quindi inizialmente può generare solo vettori in \(\operatorname{Im}(B)\);
    \item poi la dinamica libera \(A\) prende quei vettori e li evolve nel tempo;
    \item quindi lo spazio effettivamente esplorabile dall’ingresso è il più piccolo sottospazio che contiene \(\operatorname{Im}(B)\) e che sia chiuso rispetto all’azione di \(A\).
\end{itemize}

Possiamo visualizzarlo così.

\begin{figure}[H]
\centering
\begin{tikzpicture}[scale=1]
    % assi
    \draw[->] (-0.3,0) -- (5.2,0) node[right] {$x_1$};
    \draw[->] (0,-0.3) -- (0,4.0) node[above] {$x_2$};

    % parallelogramma dello spazio raggiungibile
    \fill[blue!8] (0.4,0.5) -- (4.2,0.5) -- (3.4,3.0) -- (-0.4,3.0) -- cycle;
    \draw[blue!70!black, thick] (0.4,0.5) -- (4.2,0.5) -- (3.4,3.0) -- (-0.4,3.0) -- cycle;

    % immagine di B
    \draw[red!80!black, very thick, ->] (0,0) -- (1.2,2.2);
    \node[red!80!black] at (1.8,2.4) {$\operatorname{Im}(B)$};

    % label R
    \node[blue!70!black] at (2.2,1.5) {$\mathcal{R}$};

    % freccia dinamica A
    \draw[->, thick] (3.8,1.5) to[bend left=15] (3.0,2.5);
    \node at (4.35,2.1) {\small azione di \(A\)};
\end{tikzpicture}
\caption{Interpretazione geometrica dello spazio di raggiungibilità.}
\end{figure}

\section{Invarianza rispetto al cambio di coordinate}

La raggiungibilità è una proprietà strutturale della coppia \((A,B)\), quindi non dipende dal sistema di coordinate usato per descrivere lo stato.

Consideriamo un cambio di coordinate invertibile
\begin{equation}
x=Tz,
\qquad T\in\R^{n\times n}\ \text{invertibile}.
\end{equation}

Nel discreto:
\[
x^{+}=Ax+Bu
\quad \Longrightarrow \quad
Tz^{+}=ATz+Bu.
\]
Moltiplicando per \(T^{-1}\),
\[
z^{+}=T^{-1}AT\,z+T^{-1}Bu.
\]
Quindi il sistema trasformato è
\begin{equation}
z^{+}=\tilde A z+\tilde B u,
\qquad
\tilde A=T^{-1}AT,\qquad \tilde B=T^{-1}B.
\end{equation}

La matrice di raggiungibilità nelle nuove coordinate è
\[
R_z=
\begin{bmatrix}
\tilde B & \tilde A\tilde B & \cdots & \tilde A^{n-1}\tilde B
\end{bmatrix}.
\]
Sostituendo,
\[
R_z=
\begin{bmatrix}
T^{-1}B & T^{-1}AB & \cdots & T^{-1}A^{n-1}B
\end{bmatrix}
= T^{-1}R.
\]

Poiché \(T^{-1}\) è invertibile,
\[
\rank(R_z)=\rank(R).
\]
Quindi la dimensione dello spazio di raggiungibilità non cambia.

\begin{osservazione}
La raggiungibilità è una proprietà del sistema, non della rappresentazione scelta.  
Questo è il motivo per cui ha senso cercare un cambio di coordinate che renda la struttura del sistema più trasparente.
\end{osservazione}

\section{Separazione tra parte raggiungibile e non raggiungibile}

Sia
\[
r=\dim(\mathcal{R})=\rank(R).
\]
Se \(r<n\), il sistema non è completamente raggiungibile.

In questo caso scegliamo:
\begin{itemize}
    \item una matrice \(T_R\in\R^{n\times r}\) le cui colonne formino una base di \(\mathcal{R}\);
    \item una matrice \(T_N\in\R^{n\times(n-r)}\) che completi tale base a una base di tutto \(\R^n\).
\end{itemize}

Definiamo allora
\begin{equation}
T=\begin{bmatrix}
T_R & T_N
\end{bmatrix}.
\end{equation}
Essendo le colonne di \(T\) una base di \(\R^n\), la matrice \(T\) è invertibile.

Passando alle nuove coordinate \(x=Tz\), possiamo scrivere
\[
z=
\begin{pmatrix}
z_R\\
z_N
\end{pmatrix},
\]
dove:
\begin{itemize}
    \item \(z_R\in\R^r\) rappresenta la parte raggiungibile;
    \item \(z_N\in\R^{n-r}\) rappresenta la parte non raggiungibile.
\end{itemize}

\section{Forma a blocchi della dinamica}

Poiché \(\mathcal{R}\) è \(A\)-invariante, se uno stato appartiene a \(\mathcal{R}\), allora anche la sua evoluzione libera appartiene a \(\mathcal{R}\).

Nelle nuove coordinate, i vettori di \(\mathcal{R}\) hanno la forma
\[
z=
\begin{pmatrix}
z_R\\
0
\end{pmatrix}.
\]

Scriviamo allora la matrice trasformata \(\tilde A=T^{-1}AT\) in forma a blocchi:
\begin{equation}
\tilde A=
\begin{bmatrix}
A_R & A_{RN}\\
A_{NR} & A_N
\end{bmatrix}.
\end{equation}

Applicandola a un vettore della forma \(\begin{pmatrix} z_R \\ 0 \end{pmatrix}\), otteniamo
\[
\tilde A
\begin{pmatrix}
z_R\\
0
\end{pmatrix}
=
\begin{pmatrix}
A_R z_R\\
A_{NR} z_R
\end{pmatrix}.
\]
Ma questo vettore deve appartenere ancora a \(\mathcal{R}\), quindi la sua parte inferiore deve essere nulla per ogni \(z_R\). Ne segue
\[
A_{NR}=0.
\]

Dunque la matrice assume necessariamente la forma
\begin{equation}
\tilde A=
\begin{bmatrix}
A_R & A_{RN}\\
0 & A_N
\end{bmatrix}.
\label{eq:L12_Atriangolare}
\end{equation}

\section{Forma di \(\tilde B\)}

Poiché \(\operatorname{Im}(B)\subseteq\mathcal{R}\), nelle coordinate adattate alla decomposizione dello spazio di stato il vettore \(Bu\) deve avere componenti solo nella parte raggiungibile.

Di conseguenza
\begin{equation}
\tilde B=T^{-1}B=
\begin{pmatrix}
B_R\\
0
\end{pmatrix}.
\label{eq:L12_Btriangolare}
\end{equation}

\section{Forma standard di raggiungibilità}

Abbiamo così ottenuto la \textbf{forma standard di raggiungibilità}.

Nel discreto:
\begin{equation}
z^{+}=
\begin{bmatrix}
A_R & A_{RN}\\
0 & A_N
\end{bmatrix}
z+
\begin{pmatrix}
B_R\\
0
\end{pmatrix}u.
\label{eq:L12_formastandard_td}
\end{equation}

Nel continuo:
\begin{equation}
\dot z=
\begin{bmatrix}
A_R & A_{RN}\\
0 & A_N
\end{bmatrix}
z+
\begin{pmatrix}
B_R\\
0
\end{pmatrix}u.
\label{eq:L12_formastandard_tc}
\end{equation}

\subsection*{Interpretazione}

Questa forma dice esattamente che:
\begin{itemize}
    \item il sottosistema \(z_R\) è la parte raggiungibile;
    \item il sottosistema \(z_N\) è la parte non raggiungibile;
    \item l’ingresso agisce solo sul primo blocco.
\end{itemize}

Più precisamente, la dinamica è
\[
\begin{cases}
z_R^{+}=A_R z_R + A_{RN} z_N + B_R u,\\
z_N^{+}=A_N z_N.
\end{cases}
\]
Quindi la parte non raggiungibile evolve in modo autonomo.

\begin{figure}[H]
\centering
\begin{tikzpicture}[>=Latex, node distance=2.8cm]
    \node[draw, rounded corners, minimum width=3.2cm, minimum height=1.2cm, fill=blue!8] (zr) {$z_R$};
    \node[draw, rounded corners, minimum width=3.2cm, minimum height=1.2cm, fill=gray!10, right of=zr] (zn) {$z_N$};

    \draw[->, thick] (-2.2,0) -- (zr.west) node[midway, above] {$u$};
    \draw[loop above, thick] (zr) to node[above] {$A_R$} (zr);
    \draw[loop above, thick] (zn) to node[above] {$A_N$} (zn);
    \draw[->, thick] (zr.east) -- (zn.west) node[midway, above] {$A_{RN}$};

    \node[below=0.7cm of zr] {\small sottosistema raggiungibile};
    \node[below=0.7cm of zn] {\small sottosistema non raggiungibile};
\end{tikzpicture}
\caption{Interpretazione della forma standard di raggiungibilità.}
\end{figure}

\section{Sottosistema raggiungibile e sottosistema non raggiungibile}

La decomposizione precedente ha un significato molto importante.

\subsection*{Parte raggiungibile}

Il blocco \((A_R,B_R)\) rappresenta il sottosistema su cui l’ingresso può effettivamente agire.  
Infatti in quella dinamica compare il termine \(B_Ru\).

\subsection*{Parte non raggiungibile}

Il blocco \(A_N\) rappresenta invece il sottosistema che evolve in modo completamente autonomo:
\[
z_N^{+}=A_N z_N.
\]
Qui il controllo non compare né direttamente né indirettamente.

Questo è il contenuto matematico preciso della frase:
\begin{quote}
\emph{La parte non raggiungibile evolve in anello aperto.}
\end{quote}

\subsection*{Conclusione importante}

Sul sottosistema raggiungibile posso provare a fare quello che voglio.  
Sul sottosistema non raggiungibile, invece, con l’ingresso non posso fare assolutamente nulla.

\section{Il sottosistema raggiungibile è completamente raggiungibile}

Una volta portato il sistema nella forma standard di raggiungibilità, la coppia \((A_R,B_R)\) risulta completamente raggiungibile.

Il motivo è concettualmente semplice: se il blocco raggiungibile non fosse completamente raggiungibile, allora al suo interno esisterebbe ancora una parte non raggiungibile. Ma questo contraddirebbe il fatto di aver già isolato tutta la parte raggiungibile nel primo blocco.

In altre parole, la decomposizione è costruita proprio in modo da raccogliere tutta la dinamica raggiungibile nel blocco \(A_R\).

\section{Autovalori raggiungibili e autovalori non raggiungibili}

Poiché \(\tilde A\) è triangolare a blocchi,
\[
\tilde A=
\begin{bmatrix}
A_R & A_{RN}\\
0 & A_N
\end{bmatrix},
\]
i suoi autovalori sono l’unione degli autovalori dei blocchi diagonali:
\[
\sigma(\tilde A)=\sigma(A_R)\cup\sigma(A_N).
\]

Da qui nasce la terminologia:
\begin{itemize}
    \item gli autovalori di \(A_R\) si chiamano \textbf{autovalori raggiungibili};
    \item gli autovalori di \(A_N\) si chiamano \textbf{autovalori non raggiungibili}.
\end{itemize}

Questa distinzione è fondamentale in controllo, perché:
\begin{quote}
\emph{gli autovalori non raggiungibili non possono essere modificati tramite l’ingresso.}
\end{quote}

\section{Effetto sulla funzione di trasferimento}

Consideriamo ora il sistema trasformato, nel caso continuo:
\[
\dot z=
\begin{bmatrix}
A_R & A_{RN}\\
0 & A_N
\end{bmatrix}z+
\begin{pmatrix}
B_R\\
0
\end{pmatrix}u,
\qquad
y=
\begin{bmatrix}
C_R & C_N
\end{bmatrix}z+Du.
\]

La funzione di trasferimento vale
\[
G(s)=C(sI-A)^{-1}B+D.
\]
Passando alle coordinate trasformate,
\[
G(s)=\tilde C(sI-\tilde A)^{-1}\tilde B + D.
\]

Ora
\[
sI-\tilde A=
\begin{bmatrix}
sI_r-A_R & -A_{RN}\\
0 & sI_{n-r}-A_N
\end{bmatrix},
\]
che è una matrice triangolare a blocchi. La sua inversa ha quindi la forma
\[
(sI-\tilde A)^{-1}=
\begin{bmatrix}
(sI_r-A_R)^{-1} & *\\
0 & (sI_{n-r}-A_N)^{-1}
\end{bmatrix}.
\]

Moltiplicando per \(\tilde B\),
\[
(sI-\tilde A)^{-1}\tilde B=
\begin{pmatrix}
(sI_r-A_R)^{-1}B_R\\
0
\end{pmatrix}.
\]
Infine
\begin{align}
G(s)
&=
\begin{bmatrix}
C_R & C_N
\end{bmatrix}
\begin{pmatrix}
(sI_r-A_R)^{-1}B_R\\
0
\end{pmatrix}+D\\
&=
C_R(sI_r-A_R)^{-1}B_R+D.
\end{align}

\begin{osservazione}
La funzione di trasferimento dipende solo dal sottosistema raggiungibile.  
La parte non raggiungibile non compare nella dinamica ingresso--uscita.
\end{osservazione}

Questa osservazione è molto importante: se una dinamica interna non è raggiungibile, non la si vede pilotando il sistema tramite l’ingresso.

\section{Esempio}

Consideriamo il sistema
\[
A=
\begin{bmatrix}
1 & 0 & 0\\
1 & 2 & 1\\
0 & 1 & -2
\end{bmatrix},
\qquad
B=
\begin{bmatrix}
0\\
1\\
0
\end{bmatrix},
\qquad
C=
\begin{bmatrix}
0 & 0 & 1
\end{bmatrix}.
\]

\subsection*{1. Calcolo della matrice di raggiungibilità}

Calcoliamo:
\[
AB=
\begin{bmatrix}
0\\
2\\
1
\end{bmatrix},
\qquad
A^2B=A(AB)=
\begin{bmatrix}
0\\
5\\
0
\end{bmatrix}.
\]
Quindi
\[
R=
\begin{bmatrix}
B & AB & A^2B
\end{bmatrix}
=
\begin{bmatrix}
0 & 0 & 0\\
1 & 2 & 5\\
0 & 1 & 0
\end{bmatrix}.
\]

Le prime due colonne sono linearmente indipendenti, mentre la terza è combinazione lineare delle prime due. Quindi
\[
\rank(R)=2.
\]
Il sistema non è completamente raggiungibile.

\subsection*{2. Base dello spazio raggiungibile}

Una base comoda di \(\mathcal{R}\) è
\[
v_1=
\begin{bmatrix}
0\\
1\\
0
\end{bmatrix},
\qquad
v_2=
\begin{bmatrix}
0\\
0\\
1
\end{bmatrix}.
\]
Completiamo questa base con
\[
v_3=
\begin{bmatrix}
1\\
0\\
0
\end{bmatrix}.
\]

Poniamo quindi
\[
T=
\begin{bmatrix}
0 & 0 & 1\\
1 & 0 & 0\\
0 & 1 & 0
\end{bmatrix},
\qquad
T^{-1}=
\begin{bmatrix}
0 & 1 & 0\\
0 & 0 & 1\\
1 & 0 & 0
\end{bmatrix}.
\]

\subsection*{3. Sistema trasformato}

Calcoliamo
\[
\tilde A=T^{-1}AT=
\begin{bmatrix}
2 & 1 & 1\\
1 & -2 & 0\\
0 & 0 & 1
\end{bmatrix},
\qquad
\tilde B=T^{-1}B=
\begin{bmatrix}
1\\
0\\
0
\end{bmatrix},
\qquad
\tilde C=CT=
\begin{bmatrix}
0 & 1 & 0
\end{bmatrix}.
\]

Si vede subito che il sistema è nella forma standard di raggiungibilità:
\[
\tilde A=
\begin{bmatrix}
A_R & A_{RN}\\
0 & A_N
\end{bmatrix},
\qquad
\tilde B=
\begin{pmatrix}
B_R\\
0
\end{pmatrix},
\]
con
\[
A_R=
\begin{bmatrix}
2 & 1\\
1 & -2
\end{bmatrix},
\qquad
A_N=
\begin{bmatrix}
1
\end{bmatrix},
\qquad
B_R=
\begin{bmatrix}
1\\
0
\end{bmatrix}.
\]

\subsection*{4. Commento}

Il sistema si divide quindi in:
\begin{itemize}
    \item un sottosistema raggiungibile di dimensione \(2\);
    \item un sottosistema non raggiungibile di dimensione \(1\).
\end{itemize}

L’autovalore non raggiungibile è
\[
\lambda_N=1.
\]
Questo significa che esiste un modo dinamico del sistema che non può essere influenzato dall’ingresso.

\section{Messaggio conclusivo della lezione}

La lezione ci consegna un risultato centrale:

\begin{center}
\fbox{\parbox{0.9\textwidth}{
\centering
Lo spazio di raggiungibilità è il più piccolo sottospazio \(A\)-invariante che contiene \(\operatorname{Im}(B)\).  
Per mezzo di un opportuno cambio di coordinate, il sistema può essere riscritto in una forma che separa nettamente:
\begin{itemize}
    \item la parte raggiungibile;
    \item la parte non raggiungibile.
\end{itemize}
}}
\end{center}

Da questa forma si capisce immediatamente che:
\begin{itemize}
    \item l’ingresso agisce solo sulla parte raggiungibile;
    \item la parte non raggiungibile evolve autonomamente;
    \item la funzione di trasferimento vede soltanto il sottosistema raggiungibile;
    \item gli autovalori del blocco non raggiungibile non possono essere spostati tramite il controllo.
\end{itemize}

Questa è una delle basi teoriche che porteranno poi alla \textbf{decomposizione di Kalman}.
\chapter{L13 -- Criterio di Popov--Belevitch--Hautus per la raggiungibilità}

\section{Obiettivo della lezione}

Nella lezione precedente abbiamo visto che, dopo aver separato il sistema nella sua parte raggiungibile e nella sua parte non raggiungibile, gli autovalori del blocco non raggiungibile non possono essere influenzati dall’ingresso.  
Nasce allora una domanda naturale:

\begin{quote}
\emph{Esiste un criterio puramente algebrico che permetta di capire, guardando \(A\) e \(B\), se un sistema è completamente raggiungibile?}
\end{quote}

La risposta è sì, ed è data dal \textbf{criterio PBH} (Popov--Belevitch--Hautus).

L’idea profonda del criterio è la seguente:
\begin{itemize}
    \item non basta sapere quali sono gli autovalori di \(A\);
    \item bisogna capire se l’ingresso \(B\) riesce davvero ad agire su tutti i modi associati a quegli autovalori;
    \item questo controllo si fa studiando il rango della matrice aumentata
    \[
    \begin{bmatrix}
    A-\lambda I & B
    \end{bmatrix}.
    \]
\end{itemize}

In questa lezione vedremo:
\begin{enumerate}
    \item l’enunciato del criterio PBH;
    \item il suo significato geometrico e modale;
    \item il ruolo della forma di Jordan;
    \item esempi in cui il test si applica in modo molto rapido.
\end{enumerate}

\section{Richiamo: forma standard di raggiungibilità}

Ricordiamo che, se il sistema non è completamente raggiungibile, esiste un cambio di coordinate che porta la coppia \((A,B)\) nella forma standard di raggiungibilità:
\begin{equation}
\tilde A=
\begin{bmatrix}
A_R & A_{RN}\\
0 & A_N
\end{bmatrix},
\qquad
\tilde B=
\begin{bmatrix}
B_R\\
0
\end{bmatrix},
\label{eq:L13_standard_reach}
\end{equation}
dove:
\begin{itemize}
    \item \((A_R,B_R)\) descrive il sottosistema raggiungibile;
    \item \(A_N\) descrive il sottosistema non raggiungibile.
\end{itemize}

Questa scrittura è importantissima, perché ci dice che la parte non raggiungibile evolve da sola e non può essere modificata dall’ingresso.

\begin{figure}[H]
\centering
\begin{tikzpicture}[>=Latex, node distance=3.4cm]
    \node[draw, rounded corners, minimum width=3.4cm, minimum height=1.2cm, fill=blue!8] (zr) {$z_R$};
    \node[draw, rounded corners, minimum width=3.4cm, minimum height=1.2cm, fill=gray!10, right of=zr] (zn) {$z_N$};

    \draw[->, thick] (-2.2,0) -- (zr.west) node[midway, above] {$u$};
    \draw[loop above, thick] (zr) to node[above] {$A_R$} (zr);
    \draw[loop above, thick] (zn) to node[above] {$A_N$} (zn);
    \draw[->, thick] (zr.east) -- (zn.west) node[midway, above] {$A_{RN}$};

    \node[below=0.7cm of zr] {\small parte raggiungibile};
    \node[below=0.7cm of zn] {\small parte non raggiungibile};
\end{tikzpicture}
\caption{Forma standard di raggiungibilità: l’ingresso agisce solo sul sottosistema raggiungibile.}
\end{figure}

\section{Enunciato del criterio PBH}

\begin{lemma}[Criterio PBH per la raggiungibilità]
Sia data la coppia \((A,B)\), con
\[
A\in\R^{n\times n}, \qquad B\in\R^{n\times m}.
\]
La coppia \((A,B)\) è \textbf{completamente raggiungibile} se e solo se
\begin{equation}
\operatorname{rank}
\begin{bmatrix}
A-\lambda I & B
\end{bmatrix}
=n
\qquad \forall \lambda\in\C.
\label{eq:L13_PBH}
\end{equation}
\end{lemma}

\begin{osservazione}
Anche se \(A\) e \(B\) sono reali, il test va formulato per \(\lambda\in\C\), perché gli autovalori di \(A\) possono essere complessi.
\end{osservazione}

\begin{osservazione}
Questo criterio è spesso chiamato anche \emph{criterio PBH di controllabilità}.  
Nel contesto di queste lezioni lo leggiamo come criterio di \textbf{raggiungibilità} della coppia \((A,B)\).
\end{osservazione}

\section{Perché il criterio ha senso}

La matrice
\[
\begin{bmatrix}
A-\lambda I & B
\end{bmatrix}
\]
controlla simultaneamente due cose:
\begin{itemize}
    \item la struttura modale associata a \(\lambda\), tramite \(A-\lambda I\);
    \item la capacità dell’ingresso di agire su quella struttura, tramite \(B\).
\end{itemize}

Se questa matrice perde rango, significa che esiste almeno una direzione modale associata a \(\lambda\) che l’ingresso non riesce a eccitare.

In altre parole:
\begin{quote}
\emph{un autovalore è problematico non solo perché esiste, ma perché può essere associato a un modo che l’ingresso non riesce a raggiungere.}
\end{quote}

\section{Solo gli autovalori sono critici}

Il primo fatto utile è che non serve davvero controllare \eqref{eq:L13_PBH} per tutti i complessi in modo cieco.

\begin{proposizione}
Se \(\lambda\notin\sigma(A)\), allora
\[
\operatorname{rank}
\begin{bmatrix}
A-\lambda I & B
\end{bmatrix}=n.
\]
\end{proposizione}

\begin{proof}
Se \(\lambda\notin\sigma(A)\), allora \(A-\lambda I\) è invertibile, quindi ha già rango \(n\).  
Aggiungere le colonne di \(B\) non può diminuire il rango. Dunque la matrice aumentata ha rango \(n\).
\end{proof}

\begin{osservazione}
Quindi, in pratica, il criterio PBH si controlla solo sugli \textbf{autovalori di \(A\)}.
\end{osservazione}

\section{Dimostrazione del criterio PBH}

\subsection{Necessità}

Supponiamo che il sistema sia completamente raggiungibile.  
Vogliamo dimostrare che allora \eqref{eq:L13_PBH} vale per ogni \(\lambda\in\C\).

Procediamo per assurdo. Supponiamo che esista \(\hat\lambda\in\C\) tale che
\[
\operatorname{rank}
\begin{bmatrix}
A-\hat\lambda I & B
\end{bmatrix}
<n.
\]
Allora esiste un vettore riga non nullo \(w^\top\in\C^{1\times n}\) tale che
\begin{equation}
w^\top
\begin{bmatrix}
A-\hat\lambda I & B
\end{bmatrix}
=0.
\label{eq:L13_left_annihilator}
\end{equation}
Questo equivale a scrivere il sistema
\begin{equation}
\begin{cases}
w^\top(A-\hat\lambda I)=0,\\
w^\top B=0.
\end{cases}
\label{eq:L13_left_eig_system}
\end{equation}
La prima relazione dice che
\[
w^\top A=\hat\lambda\, w^\top,
\]
cioè \(w^\top\) è un \textbf{autovettore sinistro} di \(A\) associato a \(\hat\lambda\).

A questo punto, moltiplicando ripetutamente per \(A\), otteniamo:
\[
w^\top AB=\hat\lambda\, w^\top B=0,
\]
\[
w^\top A^2B=\hat\lambda\, w^\top AB=0,
\]
e in generale
\[
w^\top A^kB=0
\qquad \forall k\ge 0.
\]
Dunque \(w^\top\) annulla tutte le colonne della matrice di raggiungibilità
\[
R=\begin{bmatrix}
B & AB & A^2B & \cdots & A^{n-1}B
\end{bmatrix},
\]
cioè
\[
w^\top R=0.
\]
Ma se il sistema è completamente raggiungibile, allora \(\operatorname{rank}(R)=n\), quindi le colonne di \(R\) generano tutto \(\R^n\), e non può esistere un vettore riga non nullo che le annulli tutte.

Assurdo.  
Quindi necessariamente
\[
\operatorname{rank}
\begin{bmatrix}
A-\lambda I & B
\end{bmatrix}=n
\qquad \forall \lambda\in\C.
\]

\subsection{Sufficienza}

Supponiamo ora che
\[
\operatorname{rank}
\begin{bmatrix}
A-\lambda I & B
\end{bmatrix}=n
\qquad \forall \lambda\in\C.
\]
Vogliamo mostrare che il sistema è completamente raggiungibile.

Procediamo ancora per assurdo. Supponiamo che il sistema \emph{non} sia completamente raggiungibile.  
Allora, per la lezione precedente, esiste un cambio di coordinate che porta il sistema nella forma standard di raggiungibilità \eqref{eq:L13_standard_reach}:
\[
\tilde A=
\begin{bmatrix}
A_R & A_{RN}\\
0 & A_N
\end{bmatrix},
\qquad
\tilde B=
\begin{bmatrix}
B_R\\
0
\end{bmatrix},
\]
con \(\dim(A_N)>0\).

Sia \(\hat\lambda\) un autovalore di \(A_N\).  
Allora esiste un vettore riga non nullo \(\hat x^\top\) tale che
\[
\hat x^\top A_N=\hat\lambda\,\hat x^\top.
\]
Consideriamo allora il vettore riga
\[
\begin{bmatrix}
0^\top & \hat x^\top
\end{bmatrix}.
\]
Moltiplicandolo per la matrice PBH del sistema trasformato otteniamo
\[
\begin{bmatrix}
0^\top & \hat x^\top
\end{bmatrix}
\begin{bmatrix}
\tilde A-\hat\lambda I & \tilde B
\end{bmatrix}
=
\begin{bmatrix}
0^\top & \hat x^\top(A_N-\hat\lambda I) & 0
\end{bmatrix}
=0.
\]
Quindi
\[
\operatorname{rank}
\begin{bmatrix}
\tilde A-\hat\lambda I & \tilde B
\end{bmatrix}<n.
\]
Ma il rango è invariante per similitudine, dunque anche per il sistema originario si avrebbe
\[
\operatorname{rank}
\begin{bmatrix}
A-\hat\lambda I & B
\end{bmatrix}<n,
\]
in contraddizione con l’ipotesi.

Quindi il sistema deve essere completamente raggiungibile.
\qed

\section{Lettura modale del criterio}

Il criterio PBH ci dice che un modo associato all’autovalore \(\lambda\) è \emph{raggiungibile} se e solo se le colonne di \(B\) riescono a rompere la dipendenza lineare eventualmente presente in \(A-\lambda I\).

Dal punto di vista intuitivo:
\begin{itemize}
    \item \(A-\lambda I\) perde rango esattamente quando \(\lambda\) è un autovalore;
    \item allora compaiono direzioni modali “speciali”;
    \item se \(B\) non riesce ad agire in quelle direzioni, l’autovalore corrispondente genera un modo non raggiungibile.
\end{itemize}

Questo collega direttamente il criterio PBH al linguaggio degli autovalori raggiungibili e non raggiungibili visto nella lezione precedente.

\section{Legame con la forma di Jordan}

Per capire ancora meglio il criterio, è utile ragionare in forma di Jordan.

Supponiamo di avere una matrice in forma di Jordan
\[
J=\operatorname{diag}\!\bigl(
J_{\lambda}^{(q_1)},
J_{\lambda}^{(q_2)},
\dots,
J_{\lambda}^{(q_k)},
J_{\mu_1}^{(p_1)},
\dots
\bigr),
\]
dove i primi \(k\) blocchi sono associati allo stesso autovalore \(\lambda\).

\begin{figure}[H]
\centering
\begin{tikzpicture}[node distance=1.7cm]
    \node[draw, minimum width=1.8cm, minimum height=1.8cm] (j1) {$J_\lambda^{(q_1)}$};
    \node[draw, minimum width=1.6cm, minimum height=1.6cm, right=of j1] (j2) {$J_\lambda^{(q_2)}$};
    \node[right=0.9cm of j2] (dots) {$\cdots$};
    \node[draw, minimum width=1.8cm, minimum height=1.8cm, right=0.9cm of dots] (jk) {$J_\lambda^{(q_k)}$};

    \draw[rounded corners, thick]
    ($(j1.north west)+(-0.5,0.5)$) rectangle ($(jk.south east)+(0.5,-0.5)$);

    \node[above=0.4cm of dots] {\(\lambda\)};
\end{tikzpicture}
\caption{Più blocchi di Jordan associati allo stesso autovalore \(\lambda\). Il loro numero è \(m_g(\lambda)\).}
\end{figure}

Allora
\[
J-\lambda I=
\operatorname{diag}\!\bigl(
J_\lambda^{(q_1)}-\lambda I,
J_\lambda^{(q_2)}-\lambda I,
\dots,
J_\lambda^{(q_k)}-\lambda I,
J_{\mu_1}^{(p_1)}-\lambda I,
\dots
\bigr).
\]
Ora:
\begin{itemize}
    \item se \(\mu_i\neq \lambda\), il blocco \(J_{\mu_i}^{(p_i)}-\lambda I\) è invertibile;
    \item se il blocco è \(J_\lambda^{(q)}\), allora
    \[
    J_\lambda^{(q)}-\lambda I
    =
    \begin{bmatrix}
    0 & 1 & 0 & \cdots & 0\\
    0 & 0 & 1 & \cdots & 0\\
    \vdots & \vdots & \ddots & \ddots & \vdots\\
    0 & 0 & \cdots & 0 & 1\\
    0 & 0 & \cdots & 0 & 0
    \end{bmatrix},
    \]
    che ha rango \(q-1\).
\end{itemize}

Quindi ciascun blocco di Jordan associato a \(\lambda\) fa perdere esattamente \emph{un’unità} di rango.

\begin{proposizione}
Se \(\lambda\) è autovalore di \(A\) con molteplicità geometrica \(m_g(\lambda)=k\), allora
\begin{equation}
\operatorname{rank}(A-\lambda I)=n-k.
\label{eq:L13_rank_deficiency}
\end{equation}
\end{proposizione}

\begin{proof}
La quantità \(k\) coincide col numero di blocchi di Jordan associati a \(\lambda\).  
Ogni blocco \(J_\lambda^{(q_i)}-\lambda I\) ha rango \(q_i-1\), mentre tutti gli altri blocchi hanno rango pieno. Sommando i contributi si ottiene una perdita totale di rango pari a \(k\), cioè
\[
\operatorname{rank}(A-\lambda I)=n-k.
\]
\end{proof}

\begin{osservazione}
La perdita di rango di \(A-\lambda I\) non dipende dalla molteplicità algebrica di \(\lambda\), ma dalla sua \textbf{molteplicità geometrica}, cioè dal numero di blocchi di Jordan associati a \(\lambda\).
\end{osservazione}

\section{Conseguenza importante: numero minimo di ingressi}

Supponiamo che \(\lambda\) sia un autovalore con
\[
m_g(\lambda)=k.
\]
Allora \(A-\lambda I\) ha rango \(n-k\). Per fare in modo che la matrice aumentata
\[
\begin{bmatrix}
A-\lambda I & B
\end{bmatrix}
\]
raggiunga rango \(n\), bisogna recuperare almeno \(k\) direzioni indipendenti.

Questo implica il seguente fatto molto utile.

\begin{proposizione}
Se esiste un autovalore \(\lambda\) di \(A\) con molteplicità geometrica \(k\), allora per poter sperare che il sistema sia completamente raggiungibile è necessario avere almeno \(k\) ingressi indipendenti.
\end{proposizione}

\begin{proof}
Ogni colonna di \(B\) può aumentare il rango della matrice aumentata di al più una unità.  
Se \(A-\lambda I\) parte da rango \(n-k\), per arrivare a rango \(n\) occorre recuperare almeno \(k\) unità di rango, quindi servono almeno \(k\) colonne efficaci di \(B\).
\end{proof}

\begin{osservazione}
In particolare, con \textbf{un solo ingresso} è impossibile avere completa raggiungibilità se esiste un autovalore con molteplicità geometrica maggiore di \(1\).
\end{osservazione}

\section{Esempio 1: sistema già in forma standard di raggiungibilità}

Consideriamo
\begin{equation}
A=
\begin{bmatrix}
-1 & 0 & 2\\
0 & 3 & 1\\
0 & 0 & 2
\end{bmatrix},
\qquad
B=
\begin{bmatrix}
1\\
0\\
0
\end{bmatrix}.
\label{eq:L13_es1}
\end{equation}

La matrice \(A\) è già triangolare superiore, e anzi è già nella forma standard di raggiungibilità: l’ingresso agisce soltanto sul primo stato, mentre gli autovalori \(3\) e \(2\) stanno nella parte non raggiungibile.

\subsection*{Test PBH sugli autovalori}

Gli autovalori sono
\[
\lambda_1=-1,\qquad \lambda_2=3,\qquad \lambda_3=2.
\]

\paragraph{Per \(\lambda=-1\):}
\[
\begin{bmatrix}
A+I & B
\end{bmatrix}
=
\begin{bmatrix}
0 & 0 & 2 & 1\\
0 & 4 & 1 & 0\\
0 & 0 & 3 & 0
\end{bmatrix}.
\]
Questa matrice ha rango \(3\).

\paragraph{Per \(\lambda=3\):}
\[
\begin{bmatrix}
A-3I & B
\end{bmatrix}
=
\begin{bmatrix}
-4 & 0 & 2 & 1\\
0 & 0 & 1 & 0\\
0 & 0 & -1 & 0
\end{bmatrix}.
\]
Questa matrice ha rango \(2\).

\paragraph{Per \(\lambda=2\):}
\[
\begin{bmatrix}
A-2I & B
\end{bmatrix}
=
\begin{bmatrix}
-3 & 0 & 2 & 1\\
0 & 1 & 1 & 0\\
0 & 0 & 0 & 0
\end{bmatrix}.
\]
Questa matrice ha rango \(2\).

\subsection*{Conclusione}

Il criterio PBH fallisce per \(\lambda=3\) e per \(\lambda=2\).  
Dunque il sistema \eqref{eq:L13_es1} \textbf{non è completamente raggiungibile}.

Più precisamente:
\begin{itemize}
    \item il modo associato a \(-1\) è raggiungibile;
    \item i modi associati a \(3\) e \(2\) non sono raggiungibili.
\end{itemize}

\section{Esempio 2: due blocchi di Jordan associati allo stesso autovalore}

Consideriamo ora
\begin{equation}
A=
\begin{bmatrix}
-1 & 1 & 0\\
0 & -1 & 0\\
0 & 0 & -1
\end{bmatrix},
\qquad
B_1=
\begin{bmatrix}
0\\
0\\
1
\end{bmatrix}.
\label{eq:L13_es2}
\end{equation}

La matrice \(A\) ha un autovalore unico \(-1\), con:
\[
m_a(-1)=3,
\qquad
m_g(-1)=2.
\]
Infatti i blocchi di Jordan associati a \(-1\) sono due: uno di ordine \(2\) e uno di ordine \(1\).

\subsection*{Test PBH}

Poiché l’unico autovalore è \(-1\), basta controllare
\[
\begin{bmatrix}
A+I & B_1
\end{bmatrix}
=
\begin{bmatrix}
0 & 1 & 0 & 0\\
0 & 0 & 0 & 0\\
0 & 0 & 0 & 1
\end{bmatrix}.
\]
Il rango è
\[
\operatorname{rank}
\begin{bmatrix}
A+I & B_1
\end{bmatrix}
=2<3.
\]
Quindi il sistema non è completamente raggiungibile.

\subsection*{Interpretazione}

Qui il problema non è l’autovalore in sé, ma il fatto che \(-1\) ha due blocchi di Jordan distinti, cioè due catene indipendenti. Con un solo ingresso non si riesce a pilotare entrambe.

\begin{osservazione}
Questo esempio mostra bene la differenza tra:
\begin{itemize}
    \item \emph{molteplicità algebrica} dell’autovalore;
    \item \emph{numero di catene modali indipendenti} associate a quell’autovalore.
\end{itemize}
Nel criterio PBH conta la seconda, cioè la molteplicità geometrica.
\end{osservazione}

\section{Esempio 3: scegliere un altro ingresso non basta necessariamente}

Riprendiamo la stessa matrice \(A\), ma scegliamo ora
\[
B_2=
\begin{bmatrix}
0\\
1\\
0
\end{bmatrix}.
\]
Allora
\[
\begin{bmatrix}
A+I & B_2
\end{bmatrix}
=
\begin{bmatrix}
0 & 1 & 0 & 0\\
0 & 0 & 0 & 1\\
0 & 0 & 0 & 0
\end{bmatrix}.
\]
Anche in questo caso
\[
\operatorname{rank}
\begin{bmatrix}
A+I & B_2
\end{bmatrix}
=2<3.
\]
Quindi il sistema continua a non essere completamente raggiungibile.

\subsection*{Commento}

Con \(B_2\) si riesce a influenzare una catena diversa, ma rimane comunque una direzione modale fuori portata.  
Il punto cruciale è che con un solo ingresso non si possono recuperare le due unità di rango perse da \(A+I\).

\section{Esempio 4: un caso completamente raggiungibile}

Consideriamo ora
\begin{equation}
A=
\begin{bmatrix}
-1 & 0 & 2\\
0 & 3 & 1\\
0 & 0 & 2
\end{bmatrix},
\qquad
B=
\begin{bmatrix}
1\\
0\\
1
\end{bmatrix}.
\label{eq:L13_es4}
\end{equation}

Gli autovalori sono ancora
\[
-1,\quad 3,\quad 2.
\]
Ma il diverso ingresso cambia radicalmente la situazione.

\paragraph{Per \(\lambda=-1\):}
\[
\begin{bmatrix}
A+I & B
\end{bmatrix}
=
\begin{bmatrix}
0 & 0 & 2 & 1\\
0 & 4 & 1 & 0\\
0 & 0 & 3 & 1
\end{bmatrix},
\qquad
\operatorname{rank}=3.
\]

\paragraph{Per \(\lambda=3\):}
\[
\begin{bmatrix}
A-3I & B
\end{bmatrix}
=
\begin{bmatrix}
-4 & 0 & 2 & 1\\
0 & 0 & 1 & 0\\
0 & 0 & -1 & 1
\end{bmatrix},
\qquad
\operatorname{rank}=3.
\]

\paragraph{Per \(\lambda=2\):}
\[
\begin{bmatrix}
A-2I & B
\end{bmatrix}
=
\begin{bmatrix}
-3 & 0 & 2 & 1\\
0 & 1 & 1 & 0\\
0 & 0 & 0 & 1
\end{bmatrix},
\qquad
\operatorname{rank}=3.
\]

\subsection*{Conclusione}

Il criterio PBH è soddisfatto per tutti gli autovalori, quindi il sistema \eqref{eq:L13_es4} è \textbf{completamente raggiungibile}.

\begin{osservazione}
Questo esempio fa capire che non conta solo la presenza di certi autovalori, ma \emph{come} l’ingresso è collegato al sistema.
\end{osservazione}

\section{Interpretazione tramite autovettori sinistri}

Il criterio PBH può essere letto anche così:

\begin{quote}
\emph{il sistema è completamente raggiungibile se non esiste alcun autovettore sinistro di \(A\) ortogonale a tutte le colonne di \(B\).}
\end{quote}

Infatti, come visto nella dimostrazione:
\[
w^\top(A-\lambda I)=0,
\qquad
w^\top B=0
\]
significa contemporaneamente che:
\begin{itemize}
    \item \(w^\top\) è un autovettore sinistro di \(A\) associato a \(\lambda\);
    \item la direzione modale corrispondente non viene eccitata dall’ingresso.
\end{itemize}

Questa è forse la lettura più concettuale dell’intero lemma PBH.

\section{Messaggio pratico della lezione}

Il criterio PBH è potentissimo perché evita, in molti casi, di calcolare esplicitamente la matrice di raggiungibilità
\[
R=\begin{bmatrix}
B & AB & \cdots & A^{n-1}B
\end{bmatrix}.
\]

In pratica:
\begin{itemize}
    \item si trovano gli autovalori di \(A\);
    \item per ciascun autovalore \(\lambda\), si calcola il rango di
    \[
    \begin{bmatrix}
    A-\lambda I & B
    \end{bmatrix};
    \]
    \item se il rango è sempre \(n\), il sistema è completamente raggiungibile;
    \item se per almeno un autovalore il rango scende, allora esiste almeno un modo non raggiungibile.
\end{itemize}

\section{Schema riassuntivo}

\begin{center}
\fbox{
\parbox{0.93\textwidth}{
\textbf{Criterio PBH per la raggiungibilità}

\vspace{0.3cm}

La coppia \((A,B)\) è completamente raggiungibile se e solo se
\[
\operatorname{rank}
\begin{bmatrix}
A-\lambda I & B
\end{bmatrix}=n
\qquad \forall \lambda\in\C.
\]

\begin{itemize}
    \item Se \(\lambda\notin\sigma(A)\), il rango è automaticamente massimo.
    \item I valori davvero critici sono quindi solo gli autovalori di \(A\).
    \item Se \(m_g(\lambda)=k\), allora \(A-\lambda I\) perde \(k\) unità di rango.
    \item Per recuperarle, \(B\) deve fornire almeno \(k\) direzioni indipendenti.
    \item Se esiste un autovettore sinistro \(w^\top\neq 0\) tale che
    \[
    w^\top A=\lambda w^\top,
    \qquad
    w^\top B=0,
    \]
    allora il modo associato a \(\lambda\) non è raggiungibile.
\end{itemize}
}}
\end{center}

\section{Conclusione}

Questa lezione chiude idealmente il discorso iniziato con la forma standard di raggiungibilità:
\begin{itemize}
    \item nella lezione precedente abbiamo separato \emph{geometricamente} la parte raggiungibile da quella non raggiungibile;
    \item in questa lezione abbiamo ottenuto un test \emph{algebrico} molto rapido per riconoscere la stessa proprietà.
\end{itemize}

Il criterio PBH è uno degli strumenti più importanti dell’intero corso, perché mette in relazione:
\begin{itemize}
    \item autovalori;
    \item blocchi di Jordan;
    \item ingressi;
    \item modi effettivamente pilotabili.
\end{itemize}

Nelle lezioni successive questo criterio tornerà continuamente, sia nella progettazione dei controllori sia nell’analisi strutturale dei sistemi.
\chapter{L14 -- Forma standard di raggiungibilità, forma canonica di controllo e introduzione all'osservabilità}

\section{Obiettivo della lezione}

In questa lezione chiudiamo il blocco sulla raggiungibilità e iniziamo quello sull’osservabilità.

Più precisamente:
\begin{itemize}
    \item usiamo un esempio per costruire esplicitamente la \textbf{forma standard di raggiungibilità};
    \item vediamo quando conviene passare alla \textbf{forma canonica di controllo};
    \item capiamo che gli \textbf{autovalori non raggiungibili} non possono essere spostati con una retroazione di stato;
    \item introduciamo il problema duale dell’\textbf{osservabilità}, cioè la ricostruzione dello stato a partire da ingressi e uscite misurate.
\end{itemize}

\bigskip

Il filo logico è questo:
\[
\text{raggiungibilità} \quad \Longleftrightarrow \quad \text{quali componenti dello stato posso muovere con } u,
\]
\[
\text{osservabilità} \quad \Longleftrightarrow \quad \text{quali componenti dello stato posso ricostruire da } (u,y).
\]

%=========================================================
\section{Esempio guida: forma standard di raggiungibilità}

Consideriamo il sistema discreto
\[
x(k+1)=Ax(k)+Bu(k),\qquad y(k)=Cx(k),
\]
con
\[
A=
\begin{bmatrix}
0 & 0 & 2\\
0 & \dfrac{\alpha}{10} & 0\\
0 & 0 & 0
\end{bmatrix},
\qquad
B=
\begin{bmatrix}
0\\
0\\
1
\end{bmatrix},
\qquad
C=
\begin{bmatrix}
0 & 1 & 1
\end{bmatrix},
\]
dove \(\alpha\in\{0,1,\dots,9\}\).

\subsection{Autovalori e struttura di Jordan}

Gli autovalori di \(A\) sono
\[
\lambda_1=0,\qquad \lambda_2=0,\qquad \lambda_3=\frac{\alpha}{10}.
\]

Quindi ci sono due casi distinti.

\subsubsection*{Caso 1: \(\alpha\neq 0\)}

In questo caso:
\[
m_a(0)=2,\qquad m_a\!\left(\frac{\alpha}{10}\right)=1.
\]

Poiché il blocco in alto a destra è non nullo, i due autovalori nulli non corrispondono a due autovettori indipendenti: infatti
\[
m_g(0)=\dim\ker(A)=1.
\]
Dunque l’autovalore \(0\) è associato a \textbf{un unico blocco di Jordan di ordine 2}.

Per l’autovalore \(\alpha/10\neq 0\) si ha invece
\[
m_g\!\left(\frac{\alpha}{10}\right)=1,
\]
quindi esso corrisponde a un blocco \(1\times 1\).

In sintesi, per \(\alpha\neq 0\), la forma di Jordan è simile a
\[
J=
\begin{bmatrix}
0 & 1 & 0\\
0 & 0 & 0\\
0 & 0 & \dfrac{\alpha}{10}
\end{bmatrix}.
\]

\subsubsection*{Caso 2: \(\alpha=0\)}

Se \(\alpha=0\), allora l’unico autovalore è
\[
\lambda=0,\qquad m_a(0)=3.
\]
Inoltre
\[
m_g(0)=\dim\ker(A)=2,
\]
perciò ci sono \textbf{due blocchi di Jordan} associati all’autovalore \(0\): uno di ordine \(2\) e uno di ordine \(1\).

Una possibile forma di Jordan è
\[
J=
\begin{bmatrix}
0 & 1 & 0\\
0 & 0 & 0\\
0 & 0 & 0
\end{bmatrix}.
\]

\subsection{Matrice di raggiungibilità}

Calcoliamo
\[
R_3=\begin{bmatrix}B & AB & A^2B\end{bmatrix}.
\]

Poiché
\[
B=
\begin{bmatrix}
0\\0\\1
\end{bmatrix},
\qquad
AB=
\begin{bmatrix}
2\\0\\0
\end{bmatrix},
\qquad
A^2B=A(AB)=
\begin{bmatrix}
0\\0\\0
\end{bmatrix},
\]
si ottiene
\[
R_3=
\begin{bmatrix}
0 & 2 & 0\\
0 & 0 & 0\\
1 & 0 & 0
\end{bmatrix}.
\]

Dunque
\[
\operatorname{Im}(R_3)=\operatorname{span}\!\left\{
\begin{bmatrix}0\\0\\1\end{bmatrix},
\begin{bmatrix}1\\0\\0\end{bmatrix}
\right\},
\qquad
\operatorname{rank}(R_3)=2.
\]

Ne segue che il sistema \textbf{non è completamente raggiungibile}, perché lo spazio di stato ha dimensione \(3\) ma lo spazio raggiungibile ha dimensione \(2\).

\subsection{Catena degli spazi raggiungibili}

Si ha
\[
\mathcal R_1=\operatorname{Im}(B)=
\operatorname{span}\!\left\{
\begin{bmatrix}0\\0\\1\end{bmatrix}
\right\},
\]
e
\[
\mathcal R_2=\operatorname{Im}\!\begin{bmatrix}B & AB\end{bmatrix}
=
\operatorname{span}\!\left\{
\begin{bmatrix}0\\0\\1\end{bmatrix},
\begin{bmatrix}1\\0\\0\end{bmatrix}
\right\}.
\]
Poiché \(A^2B=0\), si ha già
\[
\mathcal R_2=\mathcal R_3.
\]

Quindi la parte raggiungibile è il piano
\[
\mathcal R=\operatorname{span}\{e_1,e_3\}.
\]

\subsection{Scelta di una base adattata}

Scegliamo come base di \(\mathcal R\)
\[
T_R=\begin{bmatrix}e_1 & e_3\end{bmatrix},
\]
e completiamola con un vettore della parte non raggiungibile, ad esempio
\[
T_N=\begin{bmatrix}e_2\end{bmatrix}.
\]
Allora una matrice di cambio di base è
\[
T=\begin{bmatrix}T_R & T_N\end{bmatrix}
=
\begin{bmatrix}
1 & 0 & 0\\
0 & 0 & 1\\
0 & 1 & 0
\end{bmatrix}.
\]
Essendo una matrice di permutazione, vale
\[
T^{-1}=T^\top.
\]

Passando alle nuove coordinate \(x=Tz\), otteniamo
\[
\widetilde A=T^{-1}AT,
\qquad
\widetilde B=T^{-1}B,
\qquad
\widetilde C=CT.
\]

Un calcolo diretto fornisce
\[
\widetilde A=
\begin{bmatrix}
0 & 2 & 0\\
0 & 0 & 0\\
0 & 0 & \dfrac{\alpha}{10}
\end{bmatrix},
\qquad
\widetilde B=
\begin{bmatrix}
0\\1\\0
\end{bmatrix},
\qquad
\widetilde C=
\begin{bmatrix}
0 & 1 & 1
\end{bmatrix}.
\]

Questa è esattamente la \textbf{forma standard di raggiungibilità}:
\[
\widetilde A=
\begin{bmatrix}
A_R & A_{RN}\\
0 & A_N
\end{bmatrix},
\qquad
\widetilde B=
\begin{bmatrix}
B_R\\
0
\end{bmatrix}.
\]

Nel nostro caso:
\[
A_R=
\begin{bmatrix}
0 & 2\\
0 & 0
\end{bmatrix},
\qquad
A_N=
\begin{bmatrix}
\dfrac{\alpha}{10}
\end{bmatrix}.
\]

Quindi:
\begin{itemize}
    \item la dinamica su \(z_R\) è la parte raggiungibile;
    \item la dinamica su \(z_N\) è la parte non raggiungibile.
\end{itemize}

\subsection{Interpretazione degli autovalori}

In questa decomposizione:
\begin{itemize}
    \item gli autovalori di \(A_R\) sono gli \textbf{autovalori raggiungibili};
    \item gli autovalori di \(A_N\) sono gli \textbf{autovalori non raggiungibili}.
\end{itemize}

Nel nostro esempio:
\begin{itemize}
    \item l’autovalore \(0\) associato al blocco \(A_R\) è raggiungibile;
    \item l’autovalore \(\alpha/10\) è non raggiungibile.
\end{itemize}

Se \(\alpha=0\), allora anche uno dei modi nulli è non raggiungibile: infatti un autovalore può comparire sia nella parte raggiungibile sia nella parte non raggiungibile.

%=========================================================
\section{Verifica rapida con PBH}

Il lemma di Popov--Belevitch--Hautus (PBH) dice che la coppia \((A,B)\) è completamente raggiungibile se e solo se
\[
\operatorname{rank}\!\begin{bmatrix}A-\lambda I & B\end{bmatrix}=n
\qquad
\forall \lambda\in\mathbb C.
\]

Nel nostro esempio, l’autovalore sospetto è
\[
\lambda=\frac{\alpha}{10}.
\]
Infatti
\[
A-\frac{\alpha}{10}I=
\begin{bmatrix}
-\frac{\alpha}{10} & 0 & 2\\
0 & 0 & 0\\
0 & 0 & -\frac{\alpha}{10}
\end{bmatrix},
\]
e quindi
\[
\begin{bmatrix}
A-\frac{\alpha}{10}I & B
\end{bmatrix}
=
\begin{bmatrix}
-\frac{\alpha}{10} & 0 & 2 & 0\\
0 & 0 & 0 & 0\\
0 & 0 & -\frac{\alpha}{10} & 1
\end{bmatrix}.
\]
Il rango è \(2\), non \(3\), dunque \(\alpha/10\) è effettivamente \textbf{non raggiungibile}.

\bigskip

Questa verifica conferma quanto avevamo già visto dalla forma standard di raggiungibilità.

%=========================================================
\section{Ricostruzione di \(x(0)\) dalle uscite: un esempio in tempo discreto}

Restiamo nello stesso esempio, e fissiamo
\[
\alpha=5,
\qquad
u(0)=u(1)=-1.
\]
Supponiamo di volere trovare gli stati iniziali \(x(0)\) tali che
\[
y(0)=3,\qquad y(1)=1,\qquad y(2)=0.
\]

\subsection{Espansione delle uscite}

Poiché il sistema è discreto e \(D=0\), si ha
\[
y(k)=Cx(k).
\]
Inoltre
\[
x(1)=Ax(0)+Bu(0),
\]
\[
x(2)=A^2x(0)+ABu(0)+Bu(1).
\]
Di conseguenza
\[
y(0)=Cx(0),
\]
\[
y(1)=CAx(0)+CB\,u(0),
\]
\[
y(2)=CA^2x(0)+CAB\,u(0)+CB\,u(1).
\]

Portando a sinistra i termini noti legati agli ingressi, otteniamo
\[
\begin{bmatrix}
y(0)\\
y(1)-CBu(0)\\
y(2)-CABu(0)-CBu(1)
\end{bmatrix}
=
\begin{bmatrix}
C\\
CA\\
CA^2
\end{bmatrix}
x(0).
\]

\subsection{Calcolo esplicito}

Per \(\alpha=5\), cioè \(\alpha/10=1/2\), si ha
\[
A=
\begin{bmatrix}
0 & 0 & 2\\
0 & \frac12 & 0\\
0 & 0 & 0
\end{bmatrix},
\qquad
C=
\begin{bmatrix}
0 & 1 & 1
\end{bmatrix},
\qquad
B=
\begin{bmatrix}
0\\0\\1
\end{bmatrix}.
\]
Calcoliamo:
\[
CB=
\begin{bmatrix}
0 & 1 & 1
\end{bmatrix}
\begin{bmatrix}
0\\0\\1
\end{bmatrix}=1,
\]
\[
CAB=CAB=0.
\]
Quindi il vettore corretto delle uscite è
\[
\Delta y=
\begin{bmatrix}
3\\
1-1\cdot(-1)\\
0-0\cdot(-1)-1\cdot(-1)
\end{bmatrix}
=
\begin{bmatrix}
3\\2\\1
\end{bmatrix}.
\]

La matrice di osservabilità ai primi tre istanti è
\[
\mathcal O=
\begin{bmatrix}
C\\CA\\CA^2
\end{bmatrix}
=
\begin{bmatrix}
0 & 1 & 1\\
0 & \frac12 & 0\\
0 & \frac14 & 0
\end{bmatrix}.
\]
Dobbiamo quindi risolvere
\[
\begin{bmatrix}
0 & 1 & 1\\
0 & \frac12 & 0\\
0 & \frac14 & 0
\end{bmatrix}
x(0)=
\begin{bmatrix}
3\\2\\1
\end{bmatrix}.
\]

Se scriviamo
\[
x(0)=
\begin{bmatrix}
x_1(0)\\x_2(0)\\x_3(0)
\end{bmatrix},
\]
le equazioni diventano
\[
x_2(0)+x_3(0)=3,
\]
\[
\frac12 x_2(0)=2,
\]
\[
\frac14 x_2(0)=1.
\]
Le ultime due sono coerenti e forniscono
\[
x_2(0)=4.
\]
Sostituendo nella prima:
\[
x_3(0)=3-4=-1.
\]
La variabile \(x_1(0)\) non compare mai, quindi resta libera.

Conclusione:
\[
x(0)=
\begin{bmatrix}
\xi\\
4\\
-1
\end{bmatrix},
\qquad \xi\in\mathbb R.
\]

\subsection{Interpretazione}

Questo esempio è molto istruttivo:
\begin{itemize}
    \item le misure di uscita permettono di ricostruire \(x_2(0)\) e \(x_3(0)\);
    \item non permettono invece di ricostruire \(x_1(0)\).
\end{itemize}

Quindi il sistema \emph{non è completamente osservabile}: esiste almeno una direzione dello stato che non lascia traccia sull’uscita.

%=========================================================
\section{Osservazione strutturale: autovalori multipli e numero minimo di ingressi}

Un punto importante, che negli appunti compare come osservazione sui blocchi di Jordan, è il seguente.

\subsection*{Osservazione}
Se un autovalore \(\lambda\) è associato a \(\nu\) blocchi di Jordan distinti, allora la sua molteplicità geometrica è
\[
m_g(\lambda)=\nu.
\]

Perché la coppia \((A,B)\) possa essere completamente raggiungibile, la matrice
\[
\begin{bmatrix}
A-\lambda I & B
\end{bmatrix}
\]
deve avere rango massimo.

Ma in forma di Jordan, il blocco \(J-\lambda I\) perde rango esattamente di \(\nu\), cioè di tanti quanti sono i blocchi di Jordan associati a \(\lambda\). Per recuperare quel rango servono abbastanza colonne in \(B\) che agiscano in modo indipendente sulle direzioni giuste.

Ne segue la regola pratica:

\begin{quote}
\textbf{se per un autovalore \(\lambda\) ci sono \(\nu\) blocchi di Jordan, allora servono almeno \(\nu\) ingressi indipendenti per avere la possibilità di rendere raggiungibile quel modo.}
\end{quote}

In particolare, se il numero di ingressi è troppo piccolo rispetto al numero di blocchi associati a uno stesso autovalore, la completa raggiungibilità è impossibile.

Questa osservazione non sostituisce PBH, ma lo rende intuitivo.

%=========================================================
\section{Se \((A,B)\) non è completamente raggiungibile e se \((A,B)\) è completamente raggiungibile}

A questo punto si chiarisce bene il significato pratico delle due forme.

\begin{itemize}
    \item Se \((A,B)\) \textbf{non è completamente raggiungibile}, la forma utile è la
    \[
    \boxed{\text{forma standard di raggiungibilità}}
    \]
    perché separa esplicitamente la parte raggiungibile dalla parte non raggiungibile.

    \item Se \((A,B)\) \textbf{è completamente raggiungibile}, la forma più utile è la
    \[
    \boxed{\text{forma canonica di controllo}}
    \]
    perché permette di leggere direttamente il polinomio caratteristico e di progettare la retroazione in modo naturale.
\end{itemize}

%=========================================================
\section{Forma canonica di controllo}

\subsection{Punto di partenza: un’equazione scalare di ordine \(n\)}

Consideriamo una sequenza scalare \(z(k)\) che soddisfa
\begin{equation}
D^n z(k)+\sum_{i=0}^{n-1} a_i D^i z(k)=u(k),
\label{eq:L14_eq_scalare}
\end{equation}
dove \(D\) è l’operatore di avanzamento:
\[
Dz(k)=z(k+1),\qquad D^2z(k)=z(k+2),\quad \dots
\]

Introduciamo lo stato
\[
x_1=z,\qquad
x_2=Dz,\qquad
x_3=D^2z,\qquad \dots,\qquad
x_n=D^{n-1}z.
\]
Allora
\[
x_1^+=x_2,\qquad
x_2^+=x_3,\qquad
\dots,\qquad
x_{n-1}^+=x_n.
\]
Inoltre, dalla \eqref{eq:L14_eq_scalare},
\[
x_n^+=D^n z
=
-\sum_{i=0}^{n-1} a_i D^i z + u
=
-a_0x_1-a_1x_2-\cdots-a_{n-1}x_n+u.
\]

Quindi il sistema in forma di stato è
\[
x^+=A_cx+B_cu,
\]
con
\[
A_c=
\begin{bmatrix}
0 & 1 & 0 & \cdots & 0\\
0 & 0 & 1 & \cdots & 0\\
\vdots & \vdots & \vdots & \ddots & \vdots\\
0 & 0 & 0 & \cdots & 1\\
-a_0 & -a_1 & -a_2 & \cdots & -a_{n-1}
\end{bmatrix},
\qquad
B_c=
\begin{bmatrix}
0\\0\\\vdots\\0\\1
\end{bmatrix}.
\]

Questa è la \textbf{forma canonica di controllo}.

\subsection{Uscita associata}

Se l’uscita è della forma
\[
y(k)=\sum_{j=0}^{p} b_j D^j z(k),
\qquad p\le n-1,
\]
allora, per come sono state definite le variabili di stato,
\[
y(k)=b_0x_1(k)+b_1x_2(k)+\cdots+b_p x_{p+1}(k).
\]
Quindi
\[
C_c=
\begin{bmatrix}
b_0 & b_1 & \dots & b_p & 0 & \dots & 0
\end{bmatrix},
\qquad
D_c=0.
\]

\subsection{Polinomio caratteristico}

La scelta dei coefficienti nell’ultima riga di \(A_c\) non è casuale.  
Infatti
\[
p(\lambda)=\det(\lambda I-A_c)
=
\lambda^n+a_{n-1}\lambda^{n-1}+\cdots+a_1\lambda+a_0.
\]
Quindi i coefficienti del polinomio caratteristico compaiono direttamente in \(A_c\).

%=========================================================
\section{La forma canonica di controllo è sempre completamente raggiungibile}

Calcoliamo la matrice di raggiungibilità della forma canonica di controllo:
\[
R_c=
\begin{bmatrix}
B_c & A_cB_c & A_c^2B_c & \cdots & A_c^{n-1}B_c
\end{bmatrix}.
\]

Questa matrice ha sempre rango massimo.  
Infatti, osservando le colonne:
\begin{itemize}
    \item la prima è \(e_n\);
    \item la seconda contiene \(e_{n-1}\) più combinazioni delle precedenti;
    \item la terza contiene \(e_{n-2}\) più combinazioni delle precedenti;
    \item e così via.
\end{itemize}

Con una semplice eliminazione si vede che le colonne sono linearmente indipendenti. Quindi
\[
\det(R_c)\neq 0,
\qquad
\operatorname{rank}(R_c)=n.
\]

Conclusione:
\[
\boxed{\text{la forma canonica di controllo è sempre completamente raggiungibile.}}
\]

%=========================================================
\section{Come costruire il cambio di base verso la forma canonica di controllo}

Se il sistema \((A,B)\) è completamente raggiungibile, allora può essere portato in forma canonica di controllo.

Sia
\[
R=\begin{bmatrix}B & AB & \cdots & A^{n-1}B\end{bmatrix}
\]
la matrice di raggiungibilità del sistema dato, e sia
\[
R_c
\]
la matrice di raggiungibilità della forma canonica di controllo avente lo stesso polinomio caratteristico.

Poiché entrambe sono invertibili, il cambio di base si può scegliere come
\begin{equation}
T=RR_c^{-1}.
\label{eq:L14_T_RRc}
\end{equation}

Allora, ponendo \(x=Tz\), si ottiene
\[
A_c=T^{-1}AT,\qquad
B_c=T^{-1}B.
\]

\subsection{Perché questa formula funziona}

L’idea è semplice:
\begin{itemize}
    \item \(R\) descrive la catena dei vettori di raggiungibilità del sistema originale;
    \item \(R_c\) descrive la stessa catena per la forma canonica.
\end{itemize}

Moltiplicare per \(R_c^{-1}\) e poi per \(R\) significa costruire il cambio di coordinate che manda la base canonica della forma di controllo nella base di raggiungibilità del sistema originario.

\subsection{Esempio}

Consideriamo
\[
A=
\begin{bmatrix}
-1 & 1 & 0\\
0 & -1 & 1\\
0 & 0 & 2
\end{bmatrix},
\qquad
B=
\begin{bmatrix}
0\\1\\1
\end{bmatrix}.
\]
Il polinomio caratteristico è
\[
p(\lambda)=(\lambda+1)^2(\lambda-2)=\lambda^3-3\lambda-2.
\]
Quindi la corrispondente forma canonica di controllo è
\[
A_c=
\begin{bmatrix}
0 & 1 & 0\\
0 & 0 & 1\\
2 & 3 & 0
\end{bmatrix},
\qquad
B_c=
\begin{bmatrix}
0\\0\\1
\end{bmatrix}.
\]

La matrice di raggiungibilità del sistema originale è
\[
R=
\begin{bmatrix}
0 & 1 & -1\\
1 & 0 & 2\\
1 & 2 & 4
\end{bmatrix},
\]
mentre
\[
R_c=
\begin{bmatrix}
0 & 0 & 1\\
0 & 1 & 0\\
1 & 0 & 3
\end{bmatrix}.
\]
Da cui
\[
R_c^{-1}=
\begin{bmatrix}
-3 & 0 & 1\\
0 & 1 & 0\\
1 & 0 & 0
\end{bmatrix}.
\]
Quindi
\[
T=RR_c^{-1}
=
\begin{bmatrix}
-1 & 1 & 0\\
-1 & 0 & 1\\
-1 & 2 & 1
\end{bmatrix}.
\]

Questo è il cambio di base che porta il sistema in forma canonica di controllo.

%=========================================================
\section{Retroazione di stato e autovalori non raggiungibili}

Supponiamo ora di avere il sistema già in forma standard di raggiungibilità:
\[
\begin{cases}
z_R^+=A_R z_R + A_{RN} z_N + B_R u,\\
z_N^+=A_N z_N.
\end{cases}
\]

Se imponiamo una retroazione di stato
\[
u=Kx,
\qquad
K=\begin{bmatrix}K_R & K_N\end{bmatrix},
\]
allora la matrice di stato a ciclo chiuso diventa
\[
A+BK
=
\begin{bmatrix}
A_R & A_{RN}\\
0 & A_N
\end{bmatrix}
+
\begin{bmatrix}
B_R\\0
\end{bmatrix}
\begin{bmatrix}
K_R & K_N
\end{bmatrix}
=
\begin{bmatrix}
A_R+B_RK_R & A_{RN}+B_RK_N\\
0 & A_N
\end{bmatrix}.
\]

La cosa fondamentale è che il blocco \(A_N\) \textbf{non cambia}.  
Quindi gli autovalori della parte non raggiungibile restano invariati sotto retroazione di stato.

\begin{quote}
In altre parole: con una retroazione di stato puoi spostare solo gli autovalori raggiungibili.
\end{quote}

Questo è uno dei motivi più profondi per cui ha senso distinguere la parte raggiungibile da quella non raggiungibile.

%=========================================================
\section{Introduzione all’osservabilità}

Passiamo ora al problema duale.

\subsection{Domanda fondamentale}

Dato il modello del sistema
\[
(A,B,C,D),
\]
e supponendo noti:
\begin{itemize}
    \item l’ingresso \(u(\cdot)\),
    \item l’uscita misurata \(y(\cdot)\),
\end{itemize}
si possono ricostruire informazioni sullo stato iniziale \(x(0)\)?  
E più in generale sullo stato \(x(t)\)?

Questo è il \textbf{problema di osservazione dello stato}.

\subsection{Stati indistinguibili}

Due stati iniziali \(x_0\) e \(\bar x_0\) si dicono \textbf{indistinguibili} in un intervallo \([0,t_f]\) se, a parità di ingresso, producono la stessa uscita per tutto l’intervallo:
\[
y(\tau; x_0,u)=y(\tau; \bar x_0,u)
\qquad
\forall \tau\in[0,t_f],\ \forall u.
\]

Indichiamo con
\[
\mathcal N_{t_f}(x_0)
=
\left\{
\bar x_0\in\mathbb R^n:
y(\tau; x_0,u)=y(\tau; \bar x_0,u)\ \ \forall \tau,\forall u
\right\}
\]
l’insieme degli stati indistinguibili da \(x_0\).

Il sistema è \textbf{completamente osservabile} se l’unico stato indistinguibile da \(x_0\) è \(x_0\) stesso, cioè se l’unico stato indistinguibile dall’origine è l’origine:
\[
\mathcal N_{t_f}(0)=\{0\}.
\]

%=========================================================
\section{Caso continuo: formula esplicita del sottospazio di inosservabilità}

Per il sistema continuo
\[
\dot x=Ax+Bu,\qquad y=Cx+Du,
\]
con stato iniziale \(x(0)=x_0\), si ha
\[
y(\tau)=Ce^{A\tau}x_0+\int_0^\tau Ce^{A(\tau-\sigma)}Bu(\sigma)\,d\sigma + Du(\tau).
\]

Se \(x_0\) e \(\bar x_0\) sono indistinguibili, allora, sottraendo le due uscite a parità di ingresso, i termini forzati si cancellano e rimane
\[
Ce^{A\tau}(x_0-\bar x_0)=0
\qquad
\forall \tau\in[0,t_f].
\]

Dunque
\[
x_0-\bar x_0 \in \ker\!\big(Ce^{A\tau}\big)\qquad \forall \tau\in[0,t_f].
\]
Ne segue che
\begin{equation}
\mathcal N_{t_f}(0)=\bigcap_{\tau\in[0,t_f]} \ker\!\big(Ce^{A\tau}\big).
\label{eq:L14_Ntf}
\end{equation}

Questo è il \textbf{sottospazio di inosservabilità}.

%=========================================================
\section{Operatore di osservazione e gramiano di osservabilità}

Definiamo l’operatore
\[
\psi : \mathbb R^n \to L^2([0,t_f],\mathbb R^p),
\qquad
\psi(x)=Ce^{A(\cdot)}x.
\]
Questo operatore associa a uno stato iniziale la corrispondente \emph{uscita libera}.

L’operatore aggiunto soddisfa
\[
\langle \psi(x), y\rangle=\langle x,\psi^\ast(y)\rangle,
\]
e si calcola come
\[
\psi^\ast(y)=\int_0^{t_f} e^{A^\top \tau} C^\top y(\tau)\,d\tau.
\]

La composizione \(\psi^\ast\psi\) è allora
\[
\psi^\ast\psi(x)
=
\int_0^{t_f} e^{A^\top \tau}C^\top C e^{A\tau}\,d\tau \, x.
\]
Introduciamo quindi il \textbf{gramiano di osservabilità}
\begin{equation}
G_O(0,t_f)=\int_0^{t_f} e^{A^\top \tau}C^\top C e^{A\tau}\,d\tau.
\label{eq:L14_gramiano_oss}
\end{equation}

Poiché
\[
\psi^\ast\psi=G_O(0,t_f),
\]
vale
\[
\ker(\psi)=\ker(G_O(0,t_f)).
\]
Ma \(\ker(\psi)\) coincide proprio con il sottospazio degli stati che non lasciano traccia in uscita. Quindi
\[
\boxed{
\mathcal N_{t_f}(0)=\ker(G_O(0,t_f)).
}
\]

Ne segue che il sistema è completamente osservabile se e solo se
\[
G_O(0,t_f)\ \text{è invertibile}.
\]

%=========================================================
\section{Matrice di osservabilità e parallelismo con la raggiungibilità}

Come per la raggiungibilità, esiste anche una caratterizzazione puramente algebrica.

Definiamo la \textbf{matrice di osservabilità}
\[
\mathcal O=
\begin{bmatrix}
C\\
CA\\
CA^2\\
\vdots\\
CA^{n-1}
\end{bmatrix}.
\]

Allora vale il parallelo perfetto con la parte precedente:
\[
\ker(G_O(0,t_f))=\ker(\mathcal O).
\]
Quindi il sistema è completamente osservabile se e solo se
\[
\operatorname{rank}(\mathcal O)=n.
\]

\bigskip

Confronto diretto:

\[
\text{raggiungibilità}
\quad\longleftrightarrow\quad
R=\begin{bmatrix}B & AB & \cdots & A^{n-1}B\end{bmatrix},
\]
\[
\text{osservabilità}
\quad\longleftrightarrow\quad
\mathcal O=
\begin{bmatrix}
C\\CA\\ \vdots\\ CA^{n-1}
\end{bmatrix}.
\]

Questo forte parallelismo è il preludio alla \textbf{dualità} tra raggiungibilità e osservabilità.

%=========================================================
\section{Conclusione della lezione}

In questa lezione abbiamo visto quattro idee fondamentali.

\begin{enumerate}
    \item La \textbf{forma standard di raggiungibilità} separa lo stato in una parte raggiungibile e una non raggiungibile.

    \item La \textbf{forma canonica di controllo} è il modello di riferimento dei sistemi completamente raggiungibili:
    \[
    A_c=
    \begin{bmatrix}
    0 & 1 & 0 & \cdots & 0\\
    0 & 0 & 1 & \cdots & 0\\
    \vdots & \vdots & \vdots & \ddots & \vdots\\
    0 & 0 & 0 & \cdots & 1\\
    -a_0 & -a_1 & -a_2 & \cdots & -a_{n-1}
    \end{bmatrix},
    \qquad
    B_c=e_n.
    \]

    \item Gli \textbf{autovalori non raggiungibili} non possono essere spostati tramite retroazione di stato.

    \item L’\textbf{osservabilità} è il problema duale: capire se dallo storico di ingressi e uscite è possibile ricostruire lo stato.
\end{enumerate}

\bigskip

La prossima lezione potrà riprendere naturalmente da qui, sviluppando:
\begin{itemize}
    \item la forma standard di osservabilità;
    \item la dualità con la raggiungibilità;
    \item gli osservatori di stato.
\end{itemize}
\chapter{L15 -- Osservabilità, ricostruibilità e forma standard di osservabilità}

\section{Idea generale della lezione}

Dopo aver studiato la raggiungibilità, passiamo al problema duale: capire quante informazioni sullo stato si possono ricavare osservando l'uscita del sistema.

Le due domande fondamentali sono:

\begin{itemize}
    \item \textbf{Osservabilità}: conoscendo \(u(\cdot)\) e \(y(\cdot)\), posso ricostruire lo \emph{stato iniziale} \(x(0)\)?
    \item \textbf{Ricostruibilità}: conoscendo \(u(\cdot)\) e \(y(\cdot)\), posso almeno ricostruire lo \emph{stato attuale} \(x(t)\), anche se lo stato iniziale non è identificabile in modo univoco?
\end{itemize}

Questa distinzione è importante soprattutto in tempo discreto.  
In tempo continuo, come vedremo, osservabilità e ricostruibilità finiscono per coincidere.

\bigskip

La logica è perfettamente duale rispetto alla raggiungibilità:

\[
\text{raggiungibilità} \quad \Longleftrightarrow \quad \text{posso muovere lo stato con l'ingresso},
\]
\[
\text{osservabilità} \quad \Longleftrightarrow \quad \text{posso ricostruire lo stato guardando l'uscita}.
\]

%=========================================================
\section{Osservabilità in tempo continuo: richiamo}

Consideriamo un sistema lineare tempo invariante in tempo continuo:
\[
\dot x(t)=Ax(t)+Bu(t),\qquad y(t)=Cx(t)+Du(t).
\]

Fissiamo due stati iniziali \(x_0\) e \(\bar x_0\).  
Diremo che \(x_0\) e \(\bar x_0\) sono \textbf{indistinguibili} in un intervallo \([0,t_f]\) se, a parità di ingresso, producono la stessa uscita per tutto l'intervallo:
\[
y(\tau; x_0,u)=y(\tau; \bar x_0,u)
\qquad
\forall \tau\in[0,t_f],\ \forall u.
\]

Poiché la soluzione dell'uscita è
\[
y(\tau)=Ce^{A\tau}x_0+\int_0^\tau Ce^{A(\tau-\sigma)}Bu(\sigma)\,d\sigma + Du(\tau),
\]
sottraendo le due uscite si ottiene
\[
Ce^{A\tau}(x_0-\bar x_0)=0
\qquad
\forall \tau\in[0,t_f].
\]

Quindi il sottospazio degli stati indistinguibili dall'origine è
\[
\mathcal I_{t_f}(0)
=
\bigcap_{\tau\in[0,t_f]} \ker\!\big(Ce^{A\tau}\big).
\]

Dalla teoria già vista si sa che questo insieme coincide con il nucleo della matrice di osservabilità classica:
\[
\mathcal O=
\begin{bmatrix}
C\\
CA\\
CA^2\\
\vdots\\
CA^{n-1}
\end{bmatrix}.
\]

In particolare:
\[
\mathcal I_{t_f}(0)=\ker(\mathcal O).
\]

Quindi il sistema è \textbf{completamente osservabile} se e solo se
\[
\ker(\mathcal O)=\{0\}
\qquad
\Longleftrightarrow
\qquad
\rank(\mathcal O)=n.
\]

\medskip

Osservazione importante: nel caso continuo, il concetto di osservabilità è una proprietà strutturale del sistema e non dipende dalla scelta dell'istante finale \(t_f>0\).

%=========================================================
\section{Osservabilità in tempo discreto}

Consideriamo ora un sistema lineare tempo invariante in tempo discreto:
\[
\begin{cases}
x(k+1)=Ax(k)+Bu(k),\\[2mm]
y(k)=Cx(k)+Du(k).
\end{cases}
\]

La risposta del sistema al tempo \(\tau\), partendo da \(x_0\) all'istante \(0\), è:
\[
y(\tau;0,x_0,u)
=
CA^\tau x_0
+
\sum_{k=0}^{\tau-1} CA^{\tau-k-1}Bu(k)
+
Du(\tau),
\qquad \tau\in\{0,1,2,\dots\}.
\]

Se due stati iniziali \(x_0\) e \(\bar x_0\) sono indistinguibili in \(t\) passi, allora per ogni \(\tau=0,\dots,t-1\) deve valere
\[
y(\tau;0,x_0,u)=y(\tau;0,\bar x_0,u).
\]

Sottraendo membro a membro si cancellano tutti i termini forzati e rimane:
\[
\Delta y(\tau)
=
y(\tau;0,x_0,u)-y(\tau;0,\bar x_0,u)
=
CA^\tau (x_0-\bar x_0).
\]

Se ora raccogliamo tutte le differenze di uscita in un unico vettore, otteniamo
\[
\begin{bmatrix}
\Delta y(0)\\
\Delta y(1)\\
\Delta y(2)\\
\vdots\\
\Delta y(t-1)
\end{bmatrix}
=
\underbrace{
\begin{bmatrix}
C\\
CA\\
CA^2\\
\vdots\\
CA^{t-1}
\end{bmatrix}}_{\displaystyle \mathcal O_t}
(x_0-\bar x_0).
\]

La matrice
\[
\mathcal O_t=
\begin{bmatrix}
C\\
CA\\
CA^2\\
\vdots\\
CA^{t-1}
\end{bmatrix}
\]
si chiama \textbf{matrice di osservabilità in \(t\) osservazioni}.

\subsection{Sottospazio di inosservabilità in \(t\) passi}

Definiamo
\[
\bar{\mathcal O}_t:=\ker(\mathcal O_t).
\]
Questo è il \textbf{sottospazio di inosservabilità in \(t\) passi}.

Esso contiene tutte le differenze iniziali \(x_0-\bar x_0\) che, per \(t\) passi, non lasciano alcuna traccia in uscita.

In altri termini:
\[
x_0 \sim_t \bar x_0
\qquad \Longleftrightarrow \qquad
x_0-\bar x_0\in \bar{\mathcal O}_t.
\]

Quindi:
\[
\bar{\mathcal O}_t=\{0\}
\qquad \Longleftrightarrow \qquad
\text{il sistema è completamente osservabile in \(t\) passi.}
\]

%=========================================================
\section{Catena dei sottospazi di inosservabilità}

Aumentando il numero di osservazioni, si aggiungono righe alla matrice \(\mathcal O_t\).  
Di conseguenza il nucleo può solo restringersi:
\[
\bar{\mathcal O}_1 \supseteq \bar{\mathcal O}_2 \supseteq \bar{\mathcal O}_3 \supseteq \cdots
\]

Questo è intuitivo: più a lungo osservo il sistema, più informazione raccolgo, quindi meno stati iniziali restano indistinguibili.

\subsection{Quando si ferma questa catena?}

La risposta è: \emph{al più dopo \(n\) osservazioni}.

Infatti, per il teorema di Cayley--Hamilton, ogni potenza \(A^k\) con \(k\ge n\) si può scrivere come combinazione lineare di
\[
I,\ A,\ A^2,\ \dots,\ A^{n-1}.
\]
Quindi anche ogni riga \(CA^k\) con \(k\ge n\) è combinazione lineare delle righe precedenti.  
Ne segue che oltre \(n\) osservazioni non compare più nessuna informazione nuova.

Perciò:
\[
\bar{\mathcal O}_n=\bar{\mathcal O}_{n+1}=\bar{\mathcal O}_{n+2}=\cdots
\]

e dunque, per verificare la completa osservabilità, basta controllare
\[
\rank(\mathcal O_n)=n.
\]

%=========================================================
\section{Primo esempio: nessun miglioramento aumentando le osservazioni}

Consideriamo
\[
A=
\begin{bmatrix}
-1 & 2 & 0\\
0 & -1 & 0\\
1 & -1 & 2
\end{bmatrix},
\qquad
C=
\begin{bmatrix}
0 & 1 & 0
\end{bmatrix}.
\]

\subsection{Prima osservazione}

Si ha
\[
\mathcal O_1=C=
\begin{bmatrix}
0 & 1 & 0
\end{bmatrix}.
\]
Quindi
\[
\bar{\mathcal O}_1=\ker(\mathcal O_1)
=
\left\{
\begin{bmatrix}
\alpha\\
0\\
\beta
\end{bmatrix}
:\ \alpha,\beta\in\mathbb R
\right\}
=
\operatorname{span}\left\{
\begin{bmatrix}
1\\0\\0
\end{bmatrix},
\begin{bmatrix}
0\\0\\1
\end{bmatrix}
\right\}.
\]
Dopo un solo passo, quindi, conosco solo la seconda componente dello stato.

\subsection{Seconda osservazione}

Calcoliamo
\[
CA=
\begin{bmatrix}
0 & 1 & 0
\end{bmatrix}
\begin{bmatrix}
-1 & 2 & 0\\
0 & -1 & 0\\
1 & -1 & 2
\end{bmatrix}
=
\begin{bmatrix}
0 & -1 & 0
\end{bmatrix}.
\]
Quindi
\[
\mathcal O_2=
\begin{bmatrix}
C\\
CA
\end{bmatrix}
=
\begin{bmatrix}
0 & 1 & 0\\
0 & -1 & 0
\end{bmatrix}.
\]
Ma le due righe sono linearmente dipendenti, quindi
\[
\bar{\mathcal O}_2=\bar{\mathcal O}_1.
\]

\subsection{Terza osservazione}

Analogamente:
\[
CA^2=
\begin{bmatrix}
0 & 1 & 0
\end{bmatrix},
\]
quindi anche
\[
\bar{\mathcal O}_3=\bar{\mathcal O}_2.
\]

\medskip

Conclusione: in questo esempio
\[
\bar{\mathcal O}_1=\bar{\mathcal O}_2=\bar{\mathcal O}_3,
\]
cioè già dopo la prima osservazione ho raccolto tutta l'informazione che l'uscita può dare.

%=========================================================
\section{Secondo esempio: le osservazioni riducono davvero il sottospazio di inosservabilità}

Consideriamo lo stesso \(A\), ma adesso
\[
C=
\begin{bmatrix}
1 & 1 & 0
\end{bmatrix}.
\]

\subsection{Prima osservazione}

La matrice di osservabilità a un passo è
\[
\mathcal O_1=
\begin{bmatrix}
1 & 1 & 0
\end{bmatrix}.
\]
Quindi
\[
\bar{\mathcal O}_1=\ker(\mathcal O_1)
=
\left\{
\begin{bmatrix}
\alpha\\
-\alpha\\
\beta
\end{bmatrix}
:\ \alpha,\beta\in\mathbb R
\right\}
=
\operatorname{span}\left\{
\begin{bmatrix}
1\\-1\\0
\end{bmatrix},
\begin{bmatrix}
0\\0\\1
\end{bmatrix}
\right\}.
\]

\subsection{Seconda osservazione}

Calcoliamo
\[
CA=
\begin{bmatrix}
1 & 1 & 0
\end{bmatrix}
\begin{bmatrix}
-1 & 2 & 0\\
0 & -1 & 0\\
1 & -1 & 2
\end{bmatrix}
=
\begin{bmatrix}
-1 & 1 & 0
\end{bmatrix}.
\]
Quindi
\[
\mathcal O_2=
\begin{bmatrix}
1 & 1 & 0\\
-1 & 1 & 0
\end{bmatrix}.
\]
La condizione
\[
\mathcal O_2 x=0
\]
equivale al sistema
\[
\begin{cases}
x_1+x_2=0,\\
-x_1+x_2=0,
\end{cases}
\]
da cui segue
\[
x_1=x_2=0.
\]
Quindi
\[
\bar{\mathcal O}_2=
\operatorname{span}\left\{
\begin{bmatrix}
0\\0\\1
\end{bmatrix}
\right\}.
\]

\subsection{Terza osservazione}

Calcoliamo ancora
\[
CA^2=
\begin{bmatrix}
1 & -3 & 0
\end{bmatrix}.
\]
Quindi
\[
\mathcal O_3=
\begin{bmatrix}
1 & 1 & 0\\
-1 & 1 & 0\\
1 & -3 & 0
\end{bmatrix}.
\]
Ma la terza riga non introduce una nuova direzione indipendente, quindi
\[
\bar{\mathcal O}_3=\bar{\mathcal O}_2.
\]

\medskip

Conclusione: in questo caso
\[
\dim \bar{\mathcal O}_1=2,\qquad
\dim \bar{\mathcal O}_2=1,\qquad
\dim \bar{\mathcal O}_3=1.
\]
Dunque la seconda osservazione riduce il sottospazio di inosservabilità, mentre la terza non aggiunge più informazione.

%=========================================================
\section{Ricostruibilità in tempo discreto}

L'osservabilità riguarda la ricostruzione dello \emph{passato}, cioè dello stato iniziale \(x(0)\).  
La ricostruibilità riguarda invece lo \emph{stato attuale} \(x(t)\).

\medskip

La domanda è:

\begin{quote}
Se due stati iniziali \(x_0\) e \(\bar x_0\) sono indistinguibili in \(t\) passi, allora dopo \(t\) passi conducono necessariamente allo stesso stato?
\end{quote}

Se la risposta è sì, allora anche se non so distinguere esattamente da quale stato iniziale sono partito, posso comunque sapere in modo univoco dove si trova il sistema al tempo \(t\).

\subsection{Definizione}

Il sistema è detto \textbf{ricostruibile in \(t\) passi} se per ogni coppia di stati indistinguibili \(x_0,\bar x_0\) in \(t\) passi si ha
\[
x(t;x_0,u)=x(t;\bar x_0,u)
\qquad
\forall u.
\]

\subsection{Condizione equivalente}

Le due traiettorie sono
\[
x_1(t)=A^t x_0+\sum_{k=0}^{t-1} A^{t-k-1}Bu(k),
\]
\[
x_2(t)=A^t \bar x_0+\sum_{k=0}^{t-1} A^{t-k-1}Bu(k).
\]
Sottraendo:
\[
x_2(t)-x_1(t)=A^t(\bar x_0-x_0).
\]

Se \(x_0\) e \(\bar x_0\) sono indistinguibili in \(t\) passi, allora
\[
\bar x_0-x_0\in \ker(\mathcal O_t)=\bar{\mathcal O}_t.
\]
Perché gli stati al tempo \(t\) coincidano, serve allora che
\[
A^t(\bar x_0-x_0)=0.
\]
Quindi il sistema è ricostruibile in \(t\) passi se e solo se
\[
\boxed{
\ker(\mathcal O_t)\subseteq \ker(A^t).
}
\]

Questa è la caratterizzazione fondamentale della ricostruibilità nel discreto.

\subsection{Relazione con l'osservabilità}

Se il sistema è osservabile in \(t\) passi, allora
\[
\ker(\mathcal O_t)=\{0\},
\]
e dunque automaticamente
\[
\ker(\mathcal O_t)\subseteq \ker(A^t).
\]
Quindi:
\[
\boxed{
\text{osservabilità} \ \Longrightarrow\ \text{ricostruibilità}.
}
\]

Il viceversa, in generale, \textbf{non è vero} nel tempo discreto.

\subsection{Quando coincidono?}

Se \(A\) è invertibile, allora anche \(A^t\) è invertibile e quindi
\[
\ker(A^t)=\{0\}.
\]
In questo caso la condizione
\[
\ker(\mathcal O_t)\subseteq \ker(A^t)
\]
diventa
\[
\ker(\mathcal O_t)\subseteq \{0\},
\]
cioè
\[
\ker(\mathcal O_t)=\{0\}.
\]
Quindi, se \(A\) è invertibile,
\[
\boxed{
\text{ricostruibilità} \ \Longleftrightarrow\ \text{osservabilità}.
}
\]

\subsection{Caso continuo}

Nel caso continuo, al posto di \(A^t\) compare \(e^{At}\), ma \(e^{At}\) è sempre invertibile.  
Per questo motivo, in tempo continuo, osservabilità e ricostruibilità coincidono sempre.

\[
\boxed{
\text{nel tempo continuo: ricostruibilità } \Longleftrightarrow \text{ osservabilità.}
}
\]

%=========================================================
\section{Invarianza per cambi di coordinate}

Come per la raggiungibilità, anche l'osservabilità è una proprietà strutturale del sistema e non dipende dalla scelta delle coordinate di stato.

Consideriamo il cambio di coordinate
\[
x=Tz,
\qquad T \ \text{invertibile}.
\]
Allora il sistema diventa
\[
z(k+1)=T^{-1}AT\,z(k)+T^{-1}Bu(k),
\qquad
y(k)=CT\,z(k)+Du(k).
\]

La nuova matrice di osservabilità è
\[
\mathcal O_z=
\begin{bmatrix}
CT\\
CTT^{-1}AT\\
CT(T^{-1}AT)^2\\
\vdots\\
CT(T^{-1}AT)^{n-1}
\end{bmatrix}.
\]
Semplificando:
\[
\mathcal O_z=
\begin{bmatrix}
C\\
CA\\
CA^2\\
\vdots\\
CA^{n-1}
\end{bmatrix}T
=
\mathcal O_x\,T.
\]

Quindi
\[
\rank(\mathcal O_z)=\rank(\mathcal O_x),
\]
e inoltre
\[
\ker(\mathcal O_z)=T^{-1}\ker(\mathcal O_x).
\]

Conclusione:
\[
\boxed{
\text{l'osservabilità è una proprietà strutturale del sistema.}
}
\]

%=========================================================
\section{Dualità tra raggiungibilità e osservabilità}

Questa è una delle idee più importanti del corso.

Consideriamo il sistema
\[
\begin{cases}
x^+=Ax+Bu,\\
y=Cx+Du.
\end{cases}
\]
Il \textbf{sistema duale} è
\[
\begin{cases}
\xi^+=A^\top \xi + C^\top v,\\
\eta=B^\top \xi + D^\top v.
\end{cases}
\]

La matrice di raggiungibilità del sistema duale è
\[
R^\star=
\begin{bmatrix}
C^\top & A^\top C^\top & (A^\top)^2 C^\top & \cdots & (A^\top)^{n-1}C^\top
\end{bmatrix}.
\]
Ma questa è esattamente la trasposta della matrice di osservabilità del sistema originale:
\[
(R^\star)^\top=
\begin{bmatrix}
C\\
CA\\
CA^2\\
\vdots\\
CA^{n-1}
\end{bmatrix}
=
\mathcal O.
\]

Di conseguenza:
\[
\boxed{
(A,C)\ \text{osservabile}
\quad \Longleftrightarrow \quad
(A^\top,C^\top)\ \text{raggiungibile}.
}
\]

Questo parallelismo permette di trasferire quasi automaticamente molti risultati dalla teoria della raggiungibilità a quella dell'osservabilità.

%=========================================================
\section{Teorema fondamentale: caratterizzazione del sottospazio di inosservabilità}

Sia
\[
\bar{\mathcal O}:=\ker(\mathcal O)
\]
il sottospazio di inosservabilità del sistema.

Vale il seguente risultato fondamentale:

\medskip

\noindent
\textbf{Teorema.}
\emph{
\(\bar{\mathcal O}\) è il più grande sottospazio \(A\)-invariante contenuto in \(\ker(C)\).
}

\subsection*{Dimostrazione}

La dimostrazione ha due parti.

\paragraph{1) \(\bar{\mathcal O}\) è contenuto in \(\ker(C)\) ed è \(A\)-invariante.}

Se \(x\in\bar{\mathcal O}\), allora
\[
\mathcal O x=0,
\]
cioè
\[
Cx=0,\qquad CAx=0,\qquad CA^2x=0,\qquad \dots,\qquad CA^{n-1}x=0.
\]
In particolare \(Cx=0\), quindi
\[
x\in\ker(C).
\]

Mostriamo ora l'invarianza.  
Se \(x\in\bar{\mathcal O}\), allora
\[
CA^k x=0\qquad \forall k=0,\dots,n-1.
\]
Vogliamo verificare che anche \(Ax\in\bar{\mathcal O}\), cioè che
\[
CA^k(Ax)=0
\qquad \forall k=0,\dots,n-1.
\]
Ma
\[
CA^k(Ax)=CA^{k+1}x.
\]
Per \(k=0,\dots,n-2\) questo è nullo per definizione di \(\bar{\mathcal O}\).  
Per \(k=n-1\) compare \(CA^n x\), ma per Cayley--Hamilton \(A^n\) è combinazione lineare di \(I,A,\dots,A^{n-1}\), quindi anche \(CA^n x\) è combinazione lineare di
\[
Cx,\ CAx,\ \dots,\ CA^{n-1}x,
\]
tutti nulli. Dunque
\[
CA^n x=0.
\]
Quindi \(Ax\in\bar{\mathcal O}\), ossia \(\bar{\mathcal O}\) è \(A\)-invariante.

\paragraph{2) È il più grande con questa proprietà.}

Sia \(S\subseteq \mathbb R^n\) un sottospazio tale che:
\[
S\subseteq \ker(C),
\qquad
AS\subseteq S.
\]
Prendiamo \(x\in S\). Siccome \(S\subseteq\ker(C)\), si ha \(Cx=0\).  
Poiché \(AS\subseteq S\), allora \(Ax\in S\), e quindi
\[
CAx=0.
\]
Ripetendo:
\[
A^2x\in S \Rightarrow CA^2x=0,
\qquad \dots,
\qquad A^{n-1}x\in S \Rightarrow CA^{n-1}x=0.
\]
Dunque
\[
\mathcal O x=0,
\]
cioè \(x\in \ker(\mathcal O)=\bar{\mathcal O}\).

Quindi
\[
S\subseteq \bar{\mathcal O}.
\]
Abbiamo dunque provato che \(\bar{\mathcal O}\) è il più grande sottospazio \(A\)-invariante contenuto in \(\ker(C)\). \qed

%=========================================================
\section{Forma standard di osservabilità}

Il teorema precedente suggerisce di scegliere nuove coordinate che separino esplicitamente la parte osservabile da quella non osservabile.

\subsection{Scelta della base}

Sia
\[
\bar o := \dim(\bar{\mathcal O}).
\]
Scegliamo:
\begin{itemize}
    \item \(T_{\bar o}\in\mathbb R^{n\times \bar o}\): una matrice le cui colonne formano una base di \(\bar{\mathcal O}\);
    \item \(T_o\in\mathbb R^{n\times (n-\bar o)}\): un completamento a base di \(\mathbb R^n\).
\end{itemize}

Poniamo quindi
\[
T=\begin{bmatrix}T_o & T_{\bar o}\end{bmatrix},
\qquad
x=Tz.
\]

\subsection{Forma del sistema trasformato}

Nelle nuove coordinate:
\[
z^+=T^{-1}AT\,z+T^{-1}Bu,
\qquad
y=CTz+Du.
\]

Poiché \(\bar{\mathcal O}\) è \(A\)-invariante, la matrice trasformata assume la struttura a blocchi
\[
\widetilde A=
\begin{bmatrix}
A_o & 0\\
A_{\bar o o} & A_{\bar o}
\end{bmatrix}.
\]
Inoltre, poiché tutte le colonne di \(T_{\bar o}\) appartengono a \(\ker(C)\), si ha
\[
CT=
\begin{bmatrix}
C_o & 0
\end{bmatrix}.
\]

Quindi il sistema assume la \textbf{forma standard di osservabilità}:
\[
\boxed{
\begin{cases}
z_o^+ = A_o z_o + B_o u,\\
z_{\bar o}^+ = A_{\bar o o} z_o + A_{\bar o} z_{\bar o} + B_{\bar o} u,\\
y = C_o z_o + Du.
\end{cases}
}
\]

In forma matriciale:
\[
\boxed{
z^+=
\begin{bmatrix}
A_o & 0\\
A_{\bar o o} & A_{\bar o}
\end{bmatrix} z
+
\begin{bmatrix}
B_o\\
B_{\bar o}
\end{bmatrix}u,
\qquad
y=
\begin{bmatrix}
C_o & 0
\end{bmatrix}z+Du.
}
\]

\subsection{Interpretazione}

La parte \(z_o\) è la parte \textbf{osservabile} del sistema.  
La parte \(z_{\bar o}\) è la parte \textbf{non osservabile}: non compare direttamente in uscita.

Anzi, poiché nella legge di uscita compare solo \(z_o\),
\[
y=C_oz_o+Du,
\]
il blocco \(z_{\bar o}\) non può essere ricostruito dalle sole misure di uscita.

Inoltre, il sottosistema \((A_o,C_o)\) è completamente osservabile.

\medskip

Gli autovalori di \(A_o\) sono gli \textbf{autovalori osservabili}, mentre quelli di \(A_{\bar o}\) sono gli \textbf{autovalori non osservabili}.

%=========================================================
\section{Parallelismo completo con la forma standard di raggiungibilità}

A questo punto il parallelismo è perfetto.

\subsection*{Forma standard di raggiungibilità}
\[
\widetilde A=
\begin{bmatrix}
A_R & A_{RN}\\
0 & A_N
\end{bmatrix},
\qquad
\widetilde B=
\begin{bmatrix}
B_R\\
0
\end{bmatrix}.
\]
Qui la parte \emph{non raggiungibile} non è direttamente influenzabile dall'ingresso.

\subsection*{Forma standard di osservabilità}
\[
\widetilde A=
\begin{bmatrix}
A_o & 0\\
A_{\bar o o} & A_{\bar o}
\end{bmatrix},
\qquad
\widetilde C=
\begin{bmatrix}
C_o & 0
\end{bmatrix}.
\]
Qui la parte \emph{non osservabile} non è direttamente visibile in uscita.

\medskip

Questa simmetria non è casuale: deriva dalla dualità tra osservabilità e raggiungibilità.

%=========================================================
\section{Conclusioni della lezione}

In questa lezione abbiamo fissato le idee fondamentali su osservabilità e ricostruibilità.

\begin{enumerate}
    \item In tempo discreto, la matrice di osservabilità in \(t\) osservazioni è
    \[
    \mathcal O_t=
    \begin{bmatrix}
    C\\
    CA\\
    \vdots\\
    CA^{t-1}
    \end{bmatrix},
    \]
    e il sottospazio di inosservabilità in \(t\) passi è
    \[
    \bar{\mathcal O}_t=\ker(\mathcal O_t).
    \]

    \item La catena
    \[
    \bar{\mathcal O}_1 \supseteq \bar{\mathcal O}_2 \supseteq \cdots
    \]
    è decrescente e si stabilizza al più dopo \(n\) passi.

    \item La ricostruibilità in \(t\) passi è equivalente alla condizione
    \[
    \ker(\mathcal O_t)\subseteq \ker(A^t).
    \]

    \item In tempo continuo, ricostruibilità e osservabilità coincidono; in tempo discreto no, salvo casi particolari, ad esempio quando \(A\) è invertibile.

    \item L'osservabilità è una proprietà strutturale, invariante per cambi di coordinate.

    \item Il sottospazio di inosservabilità
    \[
    \bar{\mathcal O}=\ker(\mathcal O)
    \]
    è il più grande sottospazio \(A\)-invariante contenuto in \(\ker(C)\).

    \item In opportune coordinate, il sistema si separa in una parte osservabile e una non osservabile:
    questa è la \textbf{forma standard di osservabilità}.
\end{enumerate}

\bigskip

La lezione successiva potrà proseguire in modo naturale verso:
\begin{itemize}
    \item gli autovalori osservabili e non osservabili;
    \item il lemma PBH per l'osservabilità;
    \item gli osservatori di stato e, in particolare, l'osservatore di Luenberger.
\end{itemize}
\chapter{L16 -- Osservabilità, forma standard di osservabilità e decomposizione di Kalman}
\label{chap:L16}

\section{Idea generale della lezione}

In questa lezione completiamo il quadro iniziato nelle lezioni precedenti su:
\begin{itemize}
    \item sottospazio di inosservabilità;
    \item forma standard di osservabilità;
    \item criterio PBH per l'osservabilità;
    \item relazione tra autovalori, osservabilità e poli della funzione di trasferimento;
    \item decomposizione canonica di Kalman;
    \item interpretazione dei modi del sistema in termini di raggiungibilità e osservabilità.
\end{itemize}

L'idea di fondo è la seguente:

\begin{quote}
\emph{non tutti gli autovalori della matrice di stato \(A\) compaiono necessariamente nella funzione di trasferimento \(G(s)\): compaiono solo i modi che sono sia raggiungibili sia osservabili.}
\end{quote}

Questo è il cuore della lezione.

\bigskip

\section{Parte osservabile e parte inosservabile}

Ricordiamo che, per un sistema lineare stazionario

\[
\Sigma:
\begin{cases}
\dot x = Ax + Bu \\
y = Cx + Du
\end{cases}
\qquad\text{(tempo continuo)}
\]

oppure

\[
\Sigma:
\begin{cases}
x(k+1) = Ax(k) + Bu(k) \\
y(k) = Cx(k) + Du(k)
\end{cases}
\qquad\text{(tempo discreto)},
\]

il \emph{sottospazio di inosservabilità} è

\[
\bar{\mathcal O}
=
\ker(\mathcal O)
=
\bigcap_{i=0}^{n-1}\ker(CA^i),
\]

dove

\[
\mathcal O=
\begin{bmatrix}
C \\ CA \\ CA^2 \\ \vdots \\ CA^{n-1}
\end{bmatrix}
\]

è la matrice di osservabilità.

\medskip

Il significato geometrico è il seguente:

\begin{itemize}
    \item se \(x_0\in\bar{\mathcal O}\), allora l'evoluzione libera a partire da \(x_0\) non produce alcun effetto misurabile in uscita;
    \item quindi \(\bar{\mathcal O}\) raccoglie tutti gli stati che il sensore non riesce a distinguere dall'origine;
    \item inoltre \(\bar{\mathcal O}\) è il più grande sottospazio \(A\)-invariante contenuto in \(\ker(C)\).
\end{itemize}

\medskip

Scegliendo una base adattata a \(\bar{\mathcal O}\), il sistema può essere scritto in \emph{forma standard di osservabilità}:

\[
T^{-1}AT=
\begin{bmatrix}
A_o & 0 \\
A_{\bar o,o} & A_{\bar o}
\end{bmatrix},
\qquad
T^{-1}B=
\begin{bmatrix}
B_o \\ B_{\bar o}
\end{bmatrix},
\qquad
CT=
\begin{bmatrix}
C_o & 0
\end{bmatrix},
\qquad
D=\tilde D.
\]

Qui:
\begin{itemize}
    \item \(A_o\) descrive la \emph{parte osservabile};
    \item \(A_{\bar o}\) descrive la \emph{parte inosservabile};
    \item l'uscita dipende solo dalle coordinate associate alla parte osservabile.
\end{itemize}

\section{Autovalori osservabili e funzione di trasferimento}

Nella forma standard di osservabilità si ha

\[
G(s)=C(sI-A)^{-1}B + D.
\]

Sostituendo la struttura a blocchi precedente:

\[
sI-\tilde A=
\begin{bmatrix}
sI-A_o & 0 \\
-A_{\bar o,o} & sI-A_{\bar o}
\end{bmatrix}.
\]

Essendo questa matrice triangolare a blocchi inferiori, la sua inversa ha la forma

\[
(sI-\tilde A)^{-1}
=
\begin{bmatrix}
(sI-A_o)^{-1} & 0 \\
\ast & (sI-A_{\bar o})^{-1}
\end{bmatrix}.
\]

Di conseguenza

\[
G(s)
=
\begin{bmatrix}
C_o & 0
\end{bmatrix}
\begin{bmatrix}
(sI-A_o)^{-1} & 0 \\
\ast & (sI-A_{\bar o})^{-1}
\end{bmatrix}
\begin{bmatrix}
B_o \\ B_{\bar o}
\end{bmatrix}
+ D
=
C_o(sI-A_o)^{-1}B_o + D.
\]

Questa formula è fondamentale: la funzione di trasferimento \emph{non dipende} dal blocco \(A_{\bar o}\).

\medskip

Quindi:

\begin{quote}
\emph{Gli autovalori del blocco inosservabile \(A_{\bar o}\) non compaiono nella funzione di trasferimento.}
\end{quote}

Equivalentemente, i poli della funzione di trasferimento sono legati solo agli autovalori osservabili.

\medskip

Se il sistema originario ha ordine \(n\), ma la parte osservabile ha dimensione \(n-\bar o\), allora nella funzione di trasferimento il denominatore effettivo ha grado \(n-\bar o\): ciò significa che sono avvenute cancellazioni polo-zero.

\bigskip

\section{Interpretazione: perché compaiono cancellazioni?}

Supponiamo che il sistema completo abbia polinomio caratteristico

\[
\chi_A(s)=\det(sI-A)
\]

di grado \(n\). Se una parte degli autovalori è inosservabile, allora nella funzione di trasferimento non compaiono tutti i modi del sistema, ma solo quelli effettivamente ``visibili'' dall'uscita.

Perciò si può avere

\[
G(s)=\frac{m(s)}{d(s)}, \qquad \deg d(s)=n-\bar o<n.
\]

In altre parole:
\begin{itemize}
    \item il modello di stato ha un certo ordine;
    \item la funzione di trasferimento può avere ordine minore;
    \item la differenza è dovuta a modi interni nascosti, cioè non osservabili.
\end{itemize}

\bigskip

\section{Lemma PBH per l'osservabilità}

Esiste un criterio molto comodo, duale di quello di raggiungibilità.

\medskip

\noindent\textbf{Lemma di Popov-Belevitch-Hautus (PBH) per l'osservabilità.}

La coppia \((A,C)\) è osservabile se e solo se

\[
\operatorname{rank}
\begin{bmatrix}
\lambda I-A\\
C
\end{bmatrix}
= n
\qquad \forall \lambda\in\mathbb C.
\]

In forma equivalente, trasponendo,

\[
\operatorname{rank}
\begin{bmatrix}
\lambda I-A^T & C^T
\end{bmatrix}
= n
\qquad \forall \lambda\in\mathbb C.
\]

\medskip

\noindent\textbf{Interpretazione.}
Se per un certo autovalore \(\lambda\) il rango scende, allora esiste un vettore non nullo \(v\neq 0\) tale che

\[
(\lambda I-A)v=0,
\qquad
Cv=0.
\]

La prima relazione dice che \(v\) è un autovettore associato a \(\lambda\), mentre la seconda dice che quell'autodirezione è invisibile in uscita.  
Quindi il modo \(\lambda\) non è osservabile.

\medskip

\noindent\textbf{Lettura concettuale del lemma.}
\begin{quote}
\emph{Un sistema è osservabile se nessun autovettore di \(A\) è annullato da \(C\).}
\end{quote}

Quando ci sono autovalori multipli o blocchi di Jordan non banali, il criterio PBH è ancora più utile della sola matrice di osservabilità, perché consente di capire subito quali modi sono nascosti.

\section{Un primo esempio: osservabilità con più o meno sensori}

Consideriamo

\[
A=
\begin{bmatrix}
2 & 1 & 0\\
-1 & 0 & 1\\
0 & 2 & 1
\end{bmatrix},
\qquad
C=
\begin{bmatrix}
1 & 0 & 0\\
0 & 1 & 0
\end{bmatrix}.
\]

La matrice di osservabilità è

\[
\mathcal O=
\begin{bmatrix}
C\\
CA\\
CA^2
\end{bmatrix}.
\]

Poiché

\[
CA=
\begin{bmatrix}
2 & 1 & 0\\
-1 & 0 & 1
\end{bmatrix},
\]

già le righe di \(C\) e la seconda riga di \(CA\) sono linearmente indipendenti; quindi si ottiene rango \(3\), e il sistema è completamente osservabile.

\medskip

Ora supponiamo invece di usare un solo sensore:

\[
C_1=
\begin{bmatrix}
1 & 0 & 0
\end{bmatrix}.
\]

Allora

\[
\mathcal O_1=
\begin{bmatrix}
C_1\\
C_1A\\
C_1A^2
\end{bmatrix}
=
\begin{bmatrix}
1 & 0 & 0\\
2 & 1 & 0\\
3 & 2 & 1
\end{bmatrix}.
\]

Si vede immediatamente che \(\mathcal O_1\) è triangolare inferiore con diagonale non nulla, dunque

\[
\det(\mathcal O_1)=1\neq 0,
\]

e quindi anche con un solo sensore il sistema resta completamente osservabile.

\medskip

Questo esempio mostra una cosa importante:

\begin{quote}
\emph{non è vero che serva necessariamente misurare tutte le componenti di stato; basta che le uscite contengano informazione sufficiente a ricostruire lo stato.}
\end{quote}

\bigskip

\section{Esempio fondamentale: sistema raggiungibile ma non osservabile}

Consideriamo ora

\[
A=
\begin{bmatrix}
2 & 1 & 0\\
1 & 0 & 0\\
0 & 2 & 1
\end{bmatrix},
\qquad
B=
\begin{bmatrix}
0\\
1\\
0
\end{bmatrix},
\qquad
C=
\begin{bmatrix}
1 & 1 & 0
\end{bmatrix}.
\]

\subsection{Studio dell'osservabilità}

La matrice di osservabilità è

\[
\mathcal O=
\begin{bmatrix}
C\\
CA\\
CA^2
\end{bmatrix}.
\]

Calcoliamo:

\[
CA=
\begin{bmatrix}
3 & 1 & 0
\end{bmatrix},
\qquad
CA^2=
\begin{bmatrix}
7 & 3 & 0
\end{bmatrix}.
\]

Dunque

\[
\mathcal O=
\begin{bmatrix}
1 & 1 & 0\\
3 & 1 & 0\\
7 & 3 & 0
\end{bmatrix}.
\]

La terza colonna è nulla, quindi

\[
\operatorname{rank}(\mathcal O)=2.
\]

Il sistema non è completamente osservabile.  
Inoltre

\[
\bar{\mathcal O}
=
\ker(\mathcal O)
=
\operatorname{span}\left\{
\begin{bmatrix}
0\\0\\1
\end{bmatrix}
\right\}.
\]

Quindi la direzione \(e_3\) è inosservabile.

\subsection{Studio della raggiungibilità}

La matrice di raggiungibilità è

\[
\mathcal R=
\begin{bmatrix}
B & AB & A^2B
\end{bmatrix}.
\]

Si calcola:

\[
AB=
\begin{bmatrix}
1\\0\\2
\end{bmatrix},
\qquad
A^2B=
\begin{bmatrix}
2\\1\\2
\end{bmatrix},
\]

per cui

\[
\mathcal R=
\begin{bmatrix}
0 & 1 & 2\\
1 & 0 & 1\\
0 & 2 & 2
\end{bmatrix}.
\]

Si verifica che

\[
\operatorname{rank}(\mathcal R)=3,
\]

quindi il sistema è completamente raggiungibile.

\subsection{Conclusione}

Abbiamo dunque una situazione molto istruttiva:

\[
\text{sistema completamente raggiungibile, ma non completamente osservabile.}
\]

Questo significa che:
\begin{itemize}
    \item con l'ingresso posso eccitare tutti i modi;
    \item ma uno di questi modi non viene visto in uscita.
\end{itemize}

\subsection{Verifica immediata con PBH}

L'autovalore \(\lambda=1\) è chiaramente un autovalore di \(A\), perché

\[
Ae_3=e_3,
\qquad
Ce_3=0.
\]

Quindi il modo \(\lambda=1\) è inosservabile.  
Infatti

\[
\operatorname{rank}
\begin{bmatrix}
I-A\\
C
\end{bmatrix}
<3.
\]

\subsection{Il sistema è già in forma standard di osservabilità}

Osserviamo che \(e_3\) è proprio l'ultima coordinata, e la matrice \(A\) ha la forma

\[
A=
\begin{bmatrix}
2 & 1 & 0\\
1 & 0 & 0\\
0 & 2 & 1
\end{bmatrix}
=
\left[
\begin{array}{c|c}
A_o & 0\\
\hline
A_{\bar o,o} & A_{\bar o}
\end{array}
\right]
\]

con

\[
A_o=
\begin{bmatrix}
2 & 1\\
1 & 0
\end{bmatrix},
\qquad
A_{\bar o,o}=
\begin{bmatrix}
0 & 2
\end{bmatrix},
\qquad
A_{\bar o}=
\begin{bmatrix}
1
\end{bmatrix}.
\]

Inoltre

\[
B=
\begin{bmatrix}
0\\1\\0
\end{bmatrix}
=
\begin{bmatrix}
B_o\\
B_{\bar o}
\end{bmatrix},
\qquad
B_o=
\begin{bmatrix}
0\\1
\end{bmatrix},
\quad
B_{\bar o}=0,
\]

e

\[
C=
\begin{bmatrix}
1 & 1 & 0
\end{bmatrix}
=
\begin{bmatrix}
C_o & 0
\end{bmatrix},
\qquad
C_o=
\begin{bmatrix}
1 & 1
\end{bmatrix}.
\]

Perciò la funzione di trasferimento dipende solo dal blocco osservabile \(A_o\):

\[
G(s)=C_o(sI-A_o)^{-1}B_o.
\]

Calcolando si ottiene

\[
G(s)=\frac{s-1}{s^2-2s-1}.
\]

D'altra parte il polinomio caratteristico della matrice completa \(A\) è

\[
\chi_A(s)=\det(sI-A)=(s-1)(s^2-2s-1).
\]

Quindi il polo \(s=1\), associato al modo inosservabile, si cancella.

\medskip

Questo esempio riassume perfettamente il significato della lezione:

\begin{quote}
\emph{un autovalore può essere presente nello stato ma assente nella funzione di trasferimento, se è inosservabile.}
\end{quote}

\bigskip

\section{Forma canonica di osservabilità}

Per sistemi SISO strettamente propri, la forma canonica di osservabilità è la trasposta della forma canonica di controllo.

Supponiamo di avere una funzione di trasferimento

\[
G(s)=
\frac{b_0+b_1s+\cdots+b_ps^p}
{s^n+a_{n-1}s^{n-1}+\cdots+a_1s+a_0},
\qquad p<n.
\]

Una realizzazione canonica osservabile è

\[
A_o=
\begin{bmatrix}
0 & 0 & \cdots & 0 & -a_0\\
1 & 0 & \cdots & 0 & -a_1\\
0 & 1 & \cdots & 0 & -a_2\\
\vdots & \vdots & \ddots & \vdots & \vdots\\
0 & 0 & \cdots & 1 & -a_{n-1}
\end{bmatrix},
\]

\[
B_o=
\begin{bmatrix}
b_0\\
b_1\\
\vdots\\
b_p\\
0\\
\vdots\\
0
\end{bmatrix},
\qquad
C_o=
\begin{bmatrix}
0 & 0 & \cdots & 0 & 1
\end{bmatrix},
\qquad
D=0.
\]

\medskip

Questa forma è utile perché:
\begin{itemize}
    \item è \emph{sempre completamente osservabile};
    \item rappresenta direttamente la funzione di trasferimento;
    \item è duale della forma canonica di controllo.
\end{itemize}

\section{Dualità tra raggiungibilità e osservabilità}

La dualità è un principio chiave.

Consideriamo il sistema originario

\[
\Sigma:
\begin{cases}
x^+ = Ax + Bu \\
y = Cx + Du
\end{cases}
\]

e il sistema duale

\[
\Sigma^d:
\begin{cases}
z^+ = A^Tz + C^Tv \\
w = B^Tz + D^Tv.
\end{cases}
\]

Allora:
\begin{itemize}
    \item la raggiungibilità del sistema duale corrisponde all'osservabilità del sistema originario;
    \item la matrice di raggiungibilità del duale è la trasposta della matrice di osservabilità dell'originale.
\end{itemize}

Infatti

\[
\mathcal R_d=
\begin{bmatrix}
C^T & A^TC^T & (A^T)^2C^T & \cdots & (A^T)^{n-1}C^T
\end{bmatrix}
=
\mathcal O^T.
\]

Quindi

\[
(A,C)\ \text{osservabile}
\iff
(A^T,C^T)\ \text{raggiungibile}.
\]

Questa osservazione spiega perché:
\begin{itemize}
    \item la forma standard di osservabilità è la versione duale della forma standard di raggiungibilità;
    \item il lemma PBH per l'osservabilità si ottiene trasponendo quello per la raggiungibilità;
    \item la forma canonica di osservabilità è la trasposta della forma canonica di controllo.
\end{itemize}

\bigskip

\section{Decomposizione canonica di Kalman}

La forma standard di raggiungibilità separa:
\[
\mathcal R \quad\text{e}\quad \bar{\mathcal R},
\]
mentre la forma standard di osservabilità separa:
\[
\mathcal O \quad\text{e}\quad \bar{\mathcal O}.
\]

Se vogliamo tenere conto \emph{simultaneamente} di entrambe le proprietà, dobbiamo combinare le due decomposizioni.

\medskip

L'idea è considerare i quattro tipi di modi:
\begin{enumerate}
    \item raggiungibili e osservabili;
    \item raggiungibili ma non osservabili;
    \item non raggiungibili ma osservabili;
    \item non raggiungibili e non osservabili.
\end{enumerate}

\subsection{Schema geometrico}

\begin{center}
\begin{tikzpicture}[scale=1.0]
    \draw[thick] (0,0) ellipse (4.2 and 2.5);
    \node at (4.8,1.9) {$\mathbb R^n$};

    \fill[red!25] (-1.1,0) ellipse (1.8 and 1.2);
    \draw[thick,red!70!black] (-1.1,0) ellipse (1.8 and 1.2);
    \node[red!70!black] at (-1.8,1.45) {$\mathcal R$};

    \fill[blue!20] (1.1,0) ellipse (1.8 and 1.2);
    \draw[thick,blue!70!black] (1.1,0) ellipse (1.8 and 1.2);
    \node[blue!70!black] at (2.1,1.45) {$\bar{\mathcal O}$};

    \fill[orange!60,opacity=0.9] (0,0) ellipse (0.45 and 1.0);
    \node[orange!80!black] at (0,-1.55) {$\mathcal R\cap\bar{\mathcal O}$};
\end{tikzpicture}
\end{center}

\medskip

Il primo passo pratico è costruire una base di \(\mathcal R\cap\bar{\mathcal O}\).  
Poi si completa:
\begin{itemize}
    \item dentro \(\mathcal R\);
    \item dentro \(\bar{\mathcal O}\);
    \item infine in tutto lo spazio di stato.
\end{itemize}

\subsection{Forma strutturale}

Esiste allora un cambio di coordinate \(x=Tz\) tale che il sistema assuma una forma a quattro blocchi.  
A seconda dell'ordine scelto nella base, la disposizione dei blocchi può cambiare; una forma tipica è

\[
\tilde A=
\begin{bmatrix}
A_{RO} & * & * & *\\
0 & A_{R\bar O} & * & *\\
0 & 0 & A_{\bar R O} & *\\
0 & 0 & 0 & A_{\bar R\bar O}
\end{bmatrix},
\qquad
\tilde B=
\begin{bmatrix}
B_{RO}\\
B_{R\bar O}\\
0\\
0
\end{bmatrix},
\qquad
\tilde C=
\begin{bmatrix}
C_{RO} & 0 & C_{\bar R O} & 0
\end{bmatrix}.
\]

La lettura è immediata:
\begin{itemize}
    \item l'ingresso agisce solo sui blocchi raggiungibili;
    \item l'uscita vede solo i blocchi osservabili;
    \item il blocco veramente rilevante per la dinamica ingresso--uscita è quello \emph{raggiungibile e osservabile}, cioè \(A_{RO}\).
\end{itemize}

\subsection{Riduzione della funzione di trasferimento}

La conseguenza più importante è:

\[
G(s)=C_{RO}(sI-A_{RO})^{-1}B_{RO}+D.
\]

Quindi i poli di \(G(s)\) sono esattamente gli autovalori del blocco \(A_{RO}\), cioè:

\begin{quote}
\emph{i modi che sono contemporaneamente raggiungibili e osservabili.}
\end{quote}

\section{Riassunto dei modi che compaiono}

Per fissare bene le idee, è utile distinguere quattro livelli.

\begin{center}
\renewcommand{\arraystretch}{1.4}
\begin{tabular}{|c|c|c|}
\hline
\textbf{Oggetto studiato} & \textbf{Modi che compaiono} & \textbf{Commento} \\
\hline
Stato libero \(x(t)\) & autovalori di \(A\) & tutti i modi interni \\
\hline
Stato forzato & autovalori raggiungibili di \(A\) & solo quelli eccitabili dall'ingresso \\
\hline
Uscita libera \(y(t)\) & autovalori osservabili di \(A\) & solo quelli visibili dal sensore \\
\hline
Funzione di trasferimento \(G(s)\) & autovalori sia raggiungibili sia osservabili & modi ingresso--uscita \\
\hline
\end{tabular}
\end{center}

\medskip

Questo schema è una delle sintesi più importanti di tutto il corso.

\bigskip

\section{Esempio finale: raggiungibile e non osservabile}

Riprendiamo l'esempio

\[
A=
\begin{bmatrix}
2 & 1 & 0\\
1 & 0 & 0\\
0 & 2 & 1
\end{bmatrix},
\qquad
B=
\begin{bmatrix}
0\\1\\0
\end{bmatrix},
\qquad
C=
\begin{bmatrix}
1 & 1 & 0
\end{bmatrix}.
\]

Abbiamo trovato:

\[
\mathcal R=\mathbb R^3,
\qquad
\bar{\mathcal O}
=
\operatorname{span}\left\{
\begin{bmatrix}
0\\0\\1
\end{bmatrix}
\right\}.
\]

Poiché \(\mathcal R=\mathbb R^3\), si ha

\[
\mathcal R\cap\bar{\mathcal O}
=
\bar{\mathcal O}
=
\operatorname{span}\left\{
\begin{bmatrix}
0\\0\\1
\end{bmatrix}
\right\}.
\]

Quindi:
\begin{itemize}
    \item esiste una parte raggiungibile ma non osservabile;
    \item non esiste una parte non raggiungibile;
    \item la decomposizione di Kalman qui si riduce di fatto alla forma standard di osservabilità.
\end{itemize}

In particolare il sistema è già scritto come

\[
A=
\left[
\begin{array}{c|c}
A_o & 0\\
\hline
A_{\bar o,o} & A_{\bar o}
\end{array}
\right],
\qquad
C=
\begin{bmatrix}
C_o & 0
\end{bmatrix},
\]

con

\[
A_o=
\begin{bmatrix}
2 & 1\\
1 & 0
\end{bmatrix},
\qquad
A_{\bar o}=
\begin{bmatrix}
1
\end{bmatrix}.
\]

Perciò il modo \(1\) rimane nello stato ma scompare dalla funzione di trasferimento.

\bigskip

\section{Osservazione finale}

Le due forme standard che abbiamo costruito nel corso hanno significati complementari:

\begin{itemize}
    \item \textbf{Forma standard di raggiungibilità (FSR):} separa ciò che l'ingresso può muovere da ciò che non può muovere.
    \item \textbf{Forma standard di osservabilità (FSO):} separa ciò che l'uscita riesce a vedere da ciò che non riesce a vedere.
\end{itemize}

La \textbf{decomposizione di Kalman} sovrappone queste due informazioni e separa completamente:
\begin{itemize}
    \item parte raggiungibile e osservabile;
    \item parte raggiungibile e non osservabile;
    \item parte non raggiungibile e osservabile;
    \item parte non raggiungibile e non osservabile.
\end{itemize}

La morale è:

\begin{quote}
\emph{dal punto di vista ingresso--uscita conta solo la parte minima del sistema, cioè quella contemporaneamente raggiungibile e osservabile.}
\end{quote}

Tutto il resto è dinamica interna nascosta: può esistere nello stato, ma non compare nella funzione di trasferimento.

\bigskip

\section*{Promemoria operativo per gli esercizi}

Quando in un compito devi studiare osservabilità e poli della funzione di trasferimento, conviene seguire questa scaletta:

\begin{enumerate}
    \item calcolare \(\mathcal O\) e il suo rango;
    \item ricavare \(\bar{\mathcal O}=\ker(\mathcal O)\);
    \item se utile, verificare i modi inosservabili con PBH;
    \item cercare una base adattata a \(\bar{\mathcal O}\) per scrivere la FSO;
    \item leggere il blocco osservabile \(A_o\);
    \item concludere che i poli della FDT sono gli autovalori del blocco osservabile;
    \item se serve combinare tutto con la raggiungibilità, passare alla decomposizione di Kalman.
\end{enumerate}
\chapter{L17 -- Recap su decomposizione di Kalman, esercizi su raggiungibilità/osservabilità e forma minima}
\label{chap:L17}

\section{Obiettivo della lezione}

Questa lezione fa da raccordo tra i risultati già visti su:
\begin{itemize}
    \item sottospazio di raggiungibilità \(\mathcal R\),
    \item sottospazio di inosservabilità \(\bar{\mathcal O}\),
    \item forma standard di raggiungibilità,
    \item forma standard di osservabilità,
    \item decomposizione di Kalman,
\end{itemize}
e una serie di esercizi nei quali questi concetti vengono usati in modo operativo.

L'idea di fondo è sempre la stessa:
\begin{quote}
\emph{la dinamica interna del sistema può essere scomposta in parti che l'ingresso riesce o non riesce ad eccitare, e in parti che l'uscita riesce o non riesce a vedere.}
\end{quote}

In particolare, questa lezione insiste su tre aspetti:
\begin{enumerate}
    \item come si costruisce concretamente l'intersezione \(\mathcal R\cap\bar{\mathcal O}\);
    \item come si riconoscono i modi raggiungibili e osservabili in esercizi concreti;
    \item perché la funzione di trasferimento ``vede'' solo la parte minima, cioè quella sia raggiungibile sia osservabile.
\end{enumerate}

\bigskip

\section{Richiamo: costruzione dell'intersezione \(\mathcal R\cap\bar{\mathcal O}\)}

Riprendiamo un esempio già discusso nella lezione precedente, nel quale erano note due basi:

\[
T_R=
\begin{bmatrix}
2 & -1\\
3 & -3\\
1 & -1
\end{bmatrix},
\qquad
T_{\bar O}=
\begin{bmatrix}
2\\
0\\
0
\end{bmatrix}.
\]

Le colonne di \(T_R\) generano il sottospazio di raggiungibilità \(\mathcal R\), mentre le colonne di \(T_{\bar O}\) generano il sottospazio di inosservabilità \(\bar{\mathcal O}\).

Vogliamo trovare
\[
\mathcal R\cap\bar{\mathcal O}.
\]

\subsection{Metodo generale}

Un vettore \(v\) appartiene all'intersezione se e solo se può essere scritto sia come combinazione lineare delle colonne di \(T_R\), sia come combinazione lineare delle colonne di \(T_{\bar O}\). Dunque dobbiamo imporre

\[
T_R y = T_{\bar O} z,
\]
dove in questo caso
\[
y=
\begin{bmatrix}
y_1\\y_2
\end{bmatrix},
\qquad
z=
\begin{bmatrix}
z_1
\end{bmatrix}.
\]

Equivalentemente:

\[
\begin{bmatrix}
T_R & -T_{\bar O}
\end{bmatrix}
\begin{bmatrix}
y\\z
\end{bmatrix}
=0.
\]

Sostituendo i dati:

\[
\begin{bmatrix}
2 & -1 & -2\\
3 & -3 & 0\\
1 & -1 & 0
\end{bmatrix}
\begin{bmatrix}
y_1\\y_2\\z_1
\end{bmatrix}
=0.
\]

Il sistema lineare associato è

\[
\begin{cases}
2y_1-y_2-2z_1=0,\\
3y_1-3y_2=0,\\
y_1-y_2=0.
\end{cases}
\]

Le ultime due equazioni danno subito
\[
y_1=y_2.
\]
Sostituendo nella prima si ottiene
\[
2y_1-y_1-2z_1=0
\quad\Longrightarrow\quad
z_1=\frac{y_1}{2}.
\]

Scegliendo, ad esempio, \(y_1=y_2=1\), si ha
\[
z_1=\frac12.
\]

Quindi un vettore dell'intersezione è

\[
v=T_R
\begin{bmatrix}
1\\1
\end{bmatrix}
=
\begin{bmatrix}
2 & -1\\
3 & -3\\
1 & -1
\end{bmatrix}
\begin{bmatrix}
1\\1
\end{bmatrix}
=
\begin{bmatrix}
1\\0\\0
\end{bmatrix}.
\]

D'altra parte

\[
T_{\bar O}\frac12
=
\begin{bmatrix}
2\\0\\0
\end{bmatrix}\frac12
=
\begin{bmatrix}
1\\0\\0
\end{bmatrix},
\]
quindi è tutto coerente.

\subsection{Conclusione}

Abbiamo trovato

\[
\mathcal R\cap\bar{\mathcal O}
=
\operatorname{span}\left\{
\begin{bmatrix}
1\\0\\0
\end{bmatrix}
\right\}.
\]

Questa procedura è importante perché nella decomposizione di Kalman bisogna prima individuare proprio questa intersezione, cioè la parte del sistema che è:
\begin{itemize}
    \item raggiungibile,
    \item ma non osservabile.
\end{itemize}

\bigskip

\section{Esercizio 1: prova del 18/12/2023, sistema SISO a tempo discreto}

Consideriamo il sistema a tempo discreto

\[
\Sigma:
\begin{cases}
x(k+1)=Ax(k)+Bu(k),\\
y(k)=Cx(k)+Du(k),
\end{cases}
\]
con

\[
A=
\begin{bmatrix}
0 & 1 & 1 & 0\\
0 & 1 & 0 & 1\\
1 & 0 & 0 & 1\\
1 & 0 & 1 & 0
\end{bmatrix},
\qquad
B=
\begin{bmatrix}
1\\
1\\
0\\
0
\end{bmatrix},
\qquad
C=
\begin{bmatrix}
0 & 1 & 0 & 0
\end{bmatrix},
\qquad
D=0.
\]

Il compito consiste nel determinare:
\begin{enumerate}
    \item gli stati raggiungibili dall'origine in \(1,2,3,4\) passi;
    \item gli stati non distinguibili dall'origine in \(1,2,3,4\) passi, cioè il sottospazio di inosservabilità progressiva;
    \item una forma di Kalman del sistema;
    \item gli spazi nei quali evolve l'evoluzione libera per certe condizioni iniziali assegnate.
\end{enumerate}

\subsection{Stati raggiungibili dall'origine in \(1,2,3,4\) passi}

Nel tempo discreto, gli stati raggiungibili in \(t\) passi dall'origine sono generati dalle colonne della matrice

\[
R_t=
\begin{bmatrix}
B & AB & A^2B & \cdots & A^{t-1}B
\end{bmatrix}.
\]

\paragraph{Passo 1.}
Si ha

\[
R_1 = \operatorname{Im}(B)
=
\operatorname{span}\left\{
\begin{bmatrix}
1\\1\\0\\0
\end{bmatrix}
\right\}.
\]

\paragraph{Passo 2.}
Calcoliamo

\[
AB=
\begin{bmatrix}
1\\1\\1\\1
\end{bmatrix}.
\]

Quindi

\[
R_2=
\operatorname{Im}
\begin{bmatrix}
1 & 1\\
1 & 1\\
0 & 1\\
0 & 1
\end{bmatrix}
=
\operatorname{span}\left\{
\begin{bmatrix}
1\\1\\0\\0
\end{bmatrix},
\begin{bmatrix}
0\\0\\1\\1
\end{bmatrix}
\right\}.
\]

Infatti
\[
\begin{bmatrix}
1\\1\\1\\1
\end{bmatrix}
=
\begin{bmatrix}
1\\1\\0\\0
\end{bmatrix}
+
\begin{bmatrix}
0\\0\\1\\1
\end{bmatrix}.
\]

\paragraph{Passo 3.}
Calcoliamo

\[
A^2B=
A(AB)=
A
\begin{bmatrix}
1\\1\\1\\1
\end{bmatrix}
=
\begin{bmatrix}
2\\2\\2\\2
\end{bmatrix}.
\]

Questa colonna è combinazione lineare delle precedenti, quindi

\[
R_3 = R_2.
\]

\paragraph{Passo 4.}
Di conseguenza anche

\[
R_4 = R_3 = R_2.
\]

\paragraph{Conclusione sul sottospazio di raggiungibilità.}
Lo spazio di raggiungibilità è

\[
\mathcal R = R_2 = R_3 = R_4
=
\operatorname{span}\left\{
\begin{bmatrix}
1\\1\\0\\0
\end{bmatrix},
\begin{bmatrix}
0\\0\\1\\1
\end{bmatrix}
\right\},
\]
e quindi
\[
\dim \mathcal R = 2.
\]

Il sistema \emph{non} è completamente raggiungibile, perché lo spazio di stato è \(\mathbb R^4\).

\bigskip

\subsection{Stati non osservabili in \(1,2,3,4\) passi}

Nel tempo discreto, la matrice di osservabilità in \(t\) passi è

\[
O_t=
\begin{bmatrix}
C\\
CA\\
CA^2\\
\vdots\\
CA^{t-1}
\end{bmatrix},
\]
e il sottospazio di inosservabilità in \(t\) passi è

\[
\bar O_t = \ker(O_t).
\]

\paragraph{Passo 1.}
Poiché

\[
C=
\begin{bmatrix}
0 & 1 & 0 & 0
\end{bmatrix},
\]
si ha

\[
\bar O_1
=
\ker(C)
=
\left\{
x\in\mathbb R^4 : x_2=0
\right\}
=
\operatorname{span}\left\{
\begin{bmatrix}
1\\0\\0\\0
\end{bmatrix},
\begin{bmatrix}
0\\0\\1\\0
\end{bmatrix},
\begin{bmatrix}
0\\0\\0\\1
\end{bmatrix}
\right\}.
\]

\paragraph{Passo 2.}
Calcoliamo

\[
CA=
\begin{bmatrix}
0 & 1 & 0 & 1
\end{bmatrix}.
\]

Quindi

\[
O_2=
\begin{bmatrix}
0 & 1 & 0 & 0\\
0 & 1 & 0 & 1
\end{bmatrix}.
\]

Le condizioni \(O_2x=0\) diventano

\[
\begin{cases}
x_2=0,\\
x_2+x_4=0.
\end{cases}
\]

Dunque
\[
x_2=0,\qquad x_4=0,
\]
e quindi

\[
\bar O_2
=
\operatorname{span}\left\{
\begin{bmatrix}
1\\0\\0\\0
\end{bmatrix},
\begin{bmatrix}
0\\0\\1\\0
\end{bmatrix}
\right\}.
\]

\paragraph{Passo 3.}
Calcoliamo

\[
CA^2=
\begin{bmatrix}
1 & 1 & 1 & 1
\end{bmatrix}.
\]

Allora

\[
O_3=
\begin{bmatrix}
0 & 1 & 0 & 0\\
0 & 1 & 0 & 1\\
1 & 1 & 1 & 1
\end{bmatrix}.
\]

Le equazioni \(O_3x=0\) sono

\[
\begin{cases}
x_2=0,\\
x_2+x_4=0,\\
x_1+x_2+x_3+x_4=0.
\end{cases}
\]

Dalle prime due:
\[
x_2=0,\qquad x_4=0.
\]
Sostituendo nella terza:
\[
x_1+x_3=0.
\]

Quindi

\[
\bar O_3
=
\operatorname{span}\left\{
\begin{bmatrix}
1\\0\\-1\\0
\end{bmatrix}
\right\}.
\]

\paragraph{Passo 4.}
Poiché
\[
CA^3=
\begin{bmatrix}
2 & 2 & 2 & 2
\end{bmatrix}
\]
è multiplo di \(CA^2\), non si ottiene alcuna informazione nuova. Dunque

\[
\bar O_4 = \bar O_3.
\]

\paragraph{Conclusione sul sottospazio di inosservabilità.}
Si ha

\[
\bar{\mathcal O}
=
\bar O_3
=
\operatorname{span}\left\{
\begin{bmatrix}
1\\0\\-1\\0
\end{bmatrix}
\right\},
\qquad
\dim \bar{\mathcal O}=1.
\]

\bigskip

\subsection{Forma di Kalman del sistema}

Abbiamo trovato:

\[
\mathcal R
=
\operatorname{span}\left\{
\begin{bmatrix}
1\\1\\0\\0
\end{bmatrix},
\begin{bmatrix}
0\\0\\1\\1
\end{bmatrix}
\right\},
\qquad
\bar{\mathcal O}
=
\operatorname{span}\left\{
\begin{bmatrix}
1\\0\\-1\\0
\end{bmatrix}
\right\}.
\]

Per costruire la forma di Kalman bisogna prima controllare l'intersezione

\[
\mathcal R\cap\bar{\mathcal O}.
\]

Cerchiamo dunque se
\[
\alpha
\begin{bmatrix}
1\\1\\0\\0
\end{bmatrix}
+
\beta
\begin{bmatrix}
0\\0\\1\\1
\end{bmatrix}
=
\gamma
\begin{bmatrix}
1\\0\\-1\\0
\end{bmatrix}.
\]

Confrontando le componenti:

\[
\begin{cases}
\alpha=\gamma,\\
\alpha=0,\\
\beta=-\gamma,\\
\beta=0.
\end{cases}
\]

Ne segue
\[
\alpha=\beta=\gamma=0,
\]
quindi

\[
\mathcal R\cap\bar{\mathcal O}=\{0\}.
\]

Questo vuol dire che \emph{non esiste una parte del sistema che sia contemporaneamente raggiungibile e non osservabile}.

\medskip

Possiamo allora scegliere una base adattata nel modo seguente:
\[
T_{RO}=T_R
=
\begin{bmatrix}
1 & 0\\
1 & 0\\
0 & 1\\
0 & 1
\end{bmatrix},
\qquad
T_{\bar R\bar O}=T_{\bar O}
=
\begin{bmatrix}
1\\
0\\
-1\\
0
\end{bmatrix},
\qquad
T_{\bar R O}=
\begin{bmatrix}
1\\0\\0\\0
\end{bmatrix}.
\]

Una possibile matrice di cambiamento di base è dunque

\[
T=
\begin{bmatrix}
1 & 0 & 1 & 1\\
1 & 0 & 0 & 0\\
0 & 1 & -1 & 0\\
0 & 1 & 0 & 0
\end{bmatrix}.
\]

In queste coordinate il sistema assume una forma di Kalman nella quale compaiono:
\begin{itemize}
    \item un blocco raggiungibile e osservabile di dimensione \(2\);
    \item un blocco non raggiungibile ma osservabile di dimensione \(1\);
    \item un blocco non raggiungibile e non osservabile di dimensione \(1\).
\end{itemize}

Poiché l'intersezione \(\mathcal R\cap\bar{\mathcal O}\) è nulla, il blocco ``raggiungibile e non osservabile'' non compare.

\bigskip

\subsection{Sottospazi di evoluzione libera per alcune condizioni iniziali}

Per un sistema a tempo discreto l'evoluzione libera è

\[
x(k)=A^k x_0.
\]

Il sottospazio generato dall'evoluzione libera a partire da \(x_0\) è il sottospazio ciclico

\[
S(x_0)=\operatorname{span}\{x_0,Ax_0,A^2x_0,\dots\}.
\]

Consideriamo tre condizioni iniziali.

\paragraph{Caso 1:}
\[
x_{0a}=
\begin{bmatrix}
-1\\0\\1\\0
\end{bmatrix}.
\]

Si calcola

\[
Ax_{0a}=
\begin{bmatrix}
1\\0\\-1\\0
\end{bmatrix}
=
-x_{0a}.
\]

Quindi \(x_{0a}\) è autovettore associato all'autovalore \(-1\), e pertanto

\[
S(x_{0a})=\operatorname{span}\{x_{0a}\}.
\]

\paragraph{Caso 2:}
\[
x_{0b}=
\begin{bmatrix}
1\\1\\-1\\-1
\end{bmatrix}.
\]

Si verifica che

\[
Ax_{0b}=0.
\]

Quindi per \(k\ge 1\) l'evoluzione si annulla, e il sottospazio generato è

\[
S(x_{0b})=\operatorname{span}\{x_{0b}\}.
\]

\paragraph{Caso 3:}
\[
x_{0c}=
\begin{bmatrix}
1\\-1\\0\\0
\end{bmatrix}.
\]

Si ottiene

\[
Ax_{0c}=
\begin{bmatrix}
-1\\-1\\1\\1
\end{bmatrix},
\qquad
A^2x_{0c}=0.
\]

Quindi

\[
S(x_{0c})=\operatorname{span}\{x_{0c},Ax_{0c}\}.
\]

Questo terzo caso è importante perché mostra una dinamica \emph{nilpotente}: dopo due passi l'evoluzione libera si annulla.

\bigskip

\section{Esercizio 2: modi del sistema, forma di Jordan e criteri PBH}

Consideriamo ora il sistema

\[
A=
\begin{bmatrix}
-1 & 0 & 1\\
-1 & 0 & 0\\
0 & 0 & -1
\end{bmatrix},
\qquad
B=
\begin{bmatrix}
1\\
0\\
-1
\end{bmatrix},
\qquad
C=
\begin{bmatrix}
-2 & 2 & 0
\end{bmatrix}.
\]

L'esercizio chiede di:
\begin{enumerate}
    \item determinare i modi del sistema;
    \item costruire una base di Jordan;
    \item stabilire quali autovalori sono interni al sottospazio di raggiungibilità e quali sono interni al sottospazio di inosservabilità.
\end{enumerate}

\subsection{Autovalori e struttura di Jordan}

Poiché \(A\) è triangolare a blocchi, i suoi autovalori si leggono facilmente:

\[
\sigma(A)=\{0,-1,-1\}.
\]

Quindi:
\begin{itemize}
    \item l'autovalore \(0\) ha molteplicità algebrica \(1\);
    \item l'autovalore \(-1\) ha molteplicità algebrica \(2\).
\end{itemize}

Studiamo la geometricità di \(-1\). Si ha

\[
A+I=
\begin{bmatrix}
0 & 0 & 1\\
-1 & 1 & 0\\
0 & 0 & 0
\end{bmatrix}.
\]

Il rango è \(2\), quindi

\[
\dim\ker(A+I)=1.
\]

Dunque l'autovalore \(-1\) ha \emph{un solo autovettore} e quindi è associato a un unico blocco di Jordan di ordine \(2\).

Un autovettore è

\[
q_1=
\begin{bmatrix}
1\\1\\0
\end{bmatrix},
\qquad
(A+I)q_1=0.
\]

Serve poi un autovettore generalizzato \(q_2\) tale che

\[
(A+I)q_2=q_1.
\]

Una possibile scelta è

\[
q_2=
\begin{bmatrix}
0\\1\\1
\end{bmatrix},
\]
infatti

\[
(A+I)
\begin{bmatrix}
0\\1\\1
\end{bmatrix}
=
\begin{bmatrix}
1\\1\\0
\end{bmatrix}
=q_1.
\]

Per l'autovalore \(0\) possiamo prendere un autovettore \(q_3\) soluzione di
\[
Aq_3=0.
\]
Ad esempio:

\[
q_3=
\begin{bmatrix}
1\\
1\\
1
\end{bmatrix}.
\]

Quindi una possibile base di Jordan è

\[
T=
\begin{bmatrix}
1 & 0 & 1\\
1 & 1 & 1\\
0 & 1 & 1
\end{bmatrix}.
\]

\bigskip

\subsection{Autovalori interni al sottospazio di raggiungibilità}

Usiamo il criterio PBH per la raggiungibilità.

La coppia \((A,B)\) è raggiungibile su un autovalore \(\lambda\) se

\[
\operatorname{rank}\begin{bmatrix}
A-\lambda I & B
\end{bmatrix}=n.
\]

\paragraph{Caso \(\lambda=0\).}
Si ha

\[
\begin{bmatrix}
A & B
\end{bmatrix}
=
\begin{bmatrix}
-1 & 0 & 1 & 1\\
-1 & 0 & 0 & 0\\
0 & 0 & -1 & -1
\end{bmatrix}.
\]

Il rango è \(2<3\), quindi l'autovalore \(0\) \emph{non è raggiungibile}.

\paragraph{Caso \(\lambda=-1\).}
Si ha

\[
\begin{bmatrix}
A+I & B
\end{bmatrix}
=
\begin{bmatrix}
0 & 0 & 1 & 1\\
-1 & 1 & 0 & 0\\
0 & 0 & 0 & -1
\end{bmatrix},
\]
che ha rango \(3\). Quindi l'autovalore \(-1\) è raggiungibile.

\paragraph{Conclusione.}
L'unico autovalore \emph{interno al sottospazio di raggiungibilità} è

\[
\boxed{\lambda=-1.}
\]

L'autovalore \(0\) è esterno a \(\mathcal R\).

\bigskip

\subsection{Autovalori interni al sottospazio di inosservabilità}

Usiamo ora il criterio PBH per l'osservabilità.

La coppia \((A,C)\) è osservabile su \(\lambda\) se

\[
\operatorname{rank}
\begin{bmatrix}
A-\lambda I\\
C
\end{bmatrix}=n.
\]

\paragraph{Caso \(\lambda=0\).}
Si ottiene

\[
\begin{bmatrix}
A\\
C
\end{bmatrix}
=
\begin{bmatrix}
-1 & 0 & 1\\
-1 & 0 & 0\\
0 & 0 & -1\\
-2 & 2 & 0
\end{bmatrix},
\]
che ha rango \(3\). Quindi \(0\) è osservabile.

\paragraph{Caso \(\lambda=-1\).}
Si ottiene

\[
\begin{bmatrix}
A+I\\
C
\end{bmatrix}
=
\begin{bmatrix}
0 & 0 & 1\\
-1 & 1 & 0\\
0 & 0 & 0\\
-2 & 2 & 0
\end{bmatrix}.
\]

La quarta riga è multiplo della seconda, quindi il rango è \(2<3\). Dunque \(-1\) \emph{non è osservabile}.

\paragraph{Conclusione.}
L'autovalore \(-1\) è interno al sottospazio di inosservabilità:

\[
\boxed{\lambda=-1 \text{ è inosservabile}.}
\]

In sintesi:
\begin{itemize}
    \item \(0\) è osservabile ma non raggiungibile;
    \item \(-1\) è raggiungibile ma non osservabile.
\end{itemize}

Questo è un esempio molto istruttivo, perché mostra che un autovalore può essere presente nel sistema ma non comparire nella parte ingresso--uscita a seconda di come si combinano raggiungibilità e osservabilità.

\bigskip

\section{Esercizio 3: equazione differenziale, realizzazione e modi visibili in uscita}

Consideriamo l'equazione differenziale ingresso--uscita

\[
y^{(4)}(t)+y^{(3)}(t)-\ddot y(t)-\dot y(t)=\dot u(t)-u(t).
\]

\subsection{Funzione di trasferimento}

Applicando la trasformata di Laplace in condizioni iniziali nulle otteniamo

\[
\left(s^4+s^3-s^2-s\right)Y(s)=(s-1)U(s),
\]
da cui

\[
G(s)=\frac{Y(s)}{U(s)}=
\frac{s-1}{s^4+s^3-s^2-s}.
\]

Fattorizzando il denominatore:

\[
s^4+s^3-s^2-s
=
s(s^3+s^2-s-1)
=
s(s+1)^2(s-1).
\]

Quindi

\[
G(s)=\frac{s-1}{s(s+1)^2(s-1)}.
\]

Dal punto di vista della funzione di trasferimento minima si può cancellare il fattore \((s-1)\), ottenendo

\[
G_{\min}(s)=\frac{1}{s(s+1)^2}.
\]

\paragraph{Interpretazione.}
Questo significa che il modo \(\lambda=1\):
\begin{itemize}
    \item è presente in una realizzazione non minima,
    \item ma non compare nella dinamica ingresso--uscita,
    \item quindi è o non raggiungibile, o non osservabile, o entrambe le cose.
\end{itemize}

Nel caso della forma canonica di controllo che costruiremo tra poco, la realizzazione è completamente raggiungibile per costruzione, quindi il modo \(\lambda=1\) sarà necessariamente \emph{non osservabile}.

\bigskip

\subsection{Una realizzazione in forma canonica di controllo}

Poiché il denominatore \emph{non ridotto} è

\[
d(s)=s^4+s^3-s^2-s,
\]
i coefficienti sono

\[
a_0=0,\qquad a_1=-1,\qquad a_2=-1,\qquad a_3=1.
\]

Una forma canonica di controllo è

\[
A_c=
\begin{bmatrix}
0 & 1 & 0 & 0\\
0 & 0 & 1 & 0\\
0 & 0 & 0 & 1\\
0 & 1 & 1 & -1
\end{bmatrix},
\qquad
B_c=
\begin{bmatrix}
0\\0\\0\\1
\end{bmatrix}.
\]

Poiché il numeratore è
\[
n(s)=s-1=-1+s+0s^2+0s^3,
\]
si può scegliere

\[
C_c=
\begin{bmatrix}
-1 & 1 & 0 & 0
\end{bmatrix},
\qquad
D=0.
\]

\paragraph{Osservazione.}
Questa realizzazione è:
\begin{itemize}
    \item completamente raggiungibile, per costruzione;
    \item non completamente osservabile, perché nella funzione di trasferimento è presente una cancellazione polo-zero.
\end{itemize}

Gli autovalori della matrice \(A_c\) sono

\[
0,\,-1,\,-1,\,1.
\]

Dopo la cancellazione nella FDT, resta solo

\[
0,\,-1,\,-1.
\]

Quindi il modo \(\lambda=1\) è un modo interno non osservabile.

\bigskip

\subsection{Quali andamenti possono comparire nell'uscita forzata?}

Dalla funzione di trasferimento minima

\[
G_{\min}(s)=\frac{1}{s(s+1)^2}
\]
deduciamo che i modi ingresso--uscita sono associati ai poli:
\[
0,\quad -1 \text{ con molteplicità } 2.
\]

Quindi gli andamenti fondamentali che possono comparire nell'uscita forzata sono

\[
1,\qquad e^{-t},\qquad t e^{-t}.
\]

Questo perché:
\begin{itemize}
    \item il polo in \(0\) genera il modo costante \(1\);
    \item il polo doppio in \(-1\) genera i modi \(e^{-t}\) e \(t e^{-t}\).
\end{itemize}

\paragraph{Conclusioni sui tre casi tipici.}
\begin{enumerate}
    \item \(\sin t\) e \(1\): \emph{non} può comparire \(\sin t\), perché richiederebbe autovalori complessi puri \(\pm j\), che qui non ci sono. Il modo costante \(1\), invece, è possibile.
    \item \(t e^{-t}\) e \(1\): entrambi sono possibili, perché corrispondono ai poli \(0\) e \(-1\) doppio.
    \item \(e^{t}\) ed \(e^{-t}\): il modo \(e^t\) corrisponde all'autovalore \(1\), ma quel polo è cancellato nella FDT, quindi non è visibile in uscita. Il modo \(e^{-t}\) è invece possibile.
\end{enumerate}

\medskip

La regola generale da ricordare è questa:

\begin{quote}
\emph{nell'uscita forzata compaiono solo i modi associati agli autovalori che sono contemporaneamente raggiungibili e osservabili.}
\end{quote}

\bigskip

\section{Realizzazione dei sistemi e forma minima}

Questa lezione si chiude con una domanda naturale:

\begin{quote}
\emph{dati numeratore e denominatore di una funzione di trasferimento, quale realizzazione in forma di stato conviene scegliere?}
\end{quote}

Sia

\[
G(s)=\frac{n(s)}{d(s)},
\qquad
\deg d(s)=n.
\]

Una \emph{realizzazione minima} di \(G(s)\) è una realizzazione in forma di stato:
\begin{itemize}
    \item con il minimo numero di stati possibile;
    \item completamente raggiungibile;
    \item completamente osservabile.
\end{itemize}

La decomposizione di Kalman spiega esattamente come si ottiene la forma minima: basta eliminare tutti i modi che non sono contemporaneamente raggiungibili e osservabili.

Se il sistema in forma di Kalman è scritto come

\[
\tilde A=
\begin{bmatrix}
A_{RO} & * & * & *\\
0 & A_{R\bar O} & * & *\\
0 & 0 & A_{\bar R O} & *\\
0 & 0 & 0 & A_{\bar R\bar O}
\end{bmatrix},
\qquad
\tilde B=
\begin{bmatrix}
B_{RO}\\
B_{R\bar O}\\
0\\
0
\end{bmatrix},
\qquad
\tilde C=
\begin{bmatrix}
C_{RO} & 0 & C_{\bar R O} & 0
\end{bmatrix},
\]
allora la funzione di trasferimento dipende solo dal blocco
\[
(A_{RO},B_{RO},C_{RO}).
\]

Più precisamente,

\[
G(s)=C_{RO}(sI-A_{RO})^{-1}B_{RO}+D.
\]

Questa è la \emph{forma minima} ottenuta eliminando i modi nascosti dall'ingresso o dall'uscita.

\bigskip

\section{Messaggio finale della lezione}

Questa lezione, anche se fatta in gran parte di esercizi, contiene un messaggio teorico molto importante:

\begin{enumerate}
    \item lo stato completo di un modello può contenere più modi di quelli che la funzione di trasferimento mostra;
    \item la funzione di trasferimento vede solo i modi contemporaneamente raggiungibili e osservabili;
    \item la decomposizione di Kalman è lo strumento che separa sistematicamente:
    \begin{itemize}
        \item parte raggiungibile e osservabile,
        \item parte raggiungibile ma non osservabile,
        \item parte non raggiungibile ma osservabile,
        \item parte non raggiungibile e non osservabile;
    \end{itemize}
    \item la forma minima si ottiene mantenendo solo il sottosistema raggiungibile e osservabile.
\end{enumerate}

\bigskip

\section*{Promemoria operativo per gli esercizi}

Quando in un esercizio ti chiedono di studiare il sistema, la scaletta più efficiente è spesso questa:

\begin{enumerate}
    \item calcolare \(\mathcal R\) oppure \(R_t\) se il sistema è a tempo discreto;
    \item calcolare \(\bar{\mathcal O}\) oppure \(\bar O_t\);
    \item verificare se ci sono intersezioni \(\mathcal R\cap \bar{\mathcal O}\);
    \item usare PBH per associare i singoli autovalori ai sottospazi;
    \item ricostruire la forma di Kalman;
    \item leggere da lì quali modi compaiono nella funzione di trasferimento;
    \item ridurre il sistema alla sua forma minima.
\end{enumerate}
\chapter{L18 -- Realizzazione dei sistemi, interconnessioni SISO, sistemi MIMO e forma minima}
\label{chap:L18}

\section{Richiamo iniziale: quando compare il termine \(t e^{-t}\)?}

Prima di entrare negli argomenti nuovi, conviene fissare bene un fatto che nei tuoi appunti compare come richiamo.

Se in una risposta temporale compare un termine del tipo
\[
t e^{-t},
\]
allora l'autovalore \(\lambda=-1\) \emph{non} è semplice nel senso dinamico: deve essere associato ad almeno un blocco di Jordan di ordine maggiore o uguale a \(2\).

Più precisamente:
\begin{itemize}
    \item se \(-1\) è autovalore semplice e \(A\) è diagonalizzabile in quel modo, allora i modi associati a \(-1\) sono solo multipli di \(e^{-t}\);
    \item il termine \(t e^{-t}\) compare solo se esiste una catena di Jordan di lunghezza almeno \(2\) associata a \(-1\);
    \item in generale, per un autovalore \(\lambda\), la dimensione del nucleo di \(A-\lambda I\) dà il numero di blocchi di Jordan associati a \(\lambda\), mentre la molteplicità algebrica dà la somma delle loro dimensioni.
\end{itemize}

Questa osservazione è importante perché, quando si studiano i modi di un sistema, non basta conoscere gli autovalori: bisogna anche capire la loro \emph{struttura di Jordan}.

\bigskip

\section{Forma minima di una realizzazione}

Sia
\[
G(s)=\frac{n(s)}{d(s)},
\qquad \deg d(s)=n.
\]

Una \emph{realizzazione minima} di \(G(s)\) è una realizzazione in forma di stato con:
\begin{itemize}
    \item numero minimo di stati;
    \item coppia \((A,B)\) completamente raggiungibile;
    \item coppia \((A,C)\) completamente osservabile.
\end{itemize}

Dal punto di vista strutturale, la forma minima si ottiene prendendo la decomposizione di Kalman e conservando solo il blocco contemporaneamente:
\begin{itemize}
    \item raggiungibile;
    \item osservabile.
\end{itemize}

Se indichiamo tale blocco con \((A_{RO},B_{RO},C_{RO})\), allora la funzione di trasferimento è

\[
G(s)=C_{RO}(sI-A_{RO})^{-1}B_{RO}+D.
\]

Questa è l'idea fondamentale:
\begin{quote}
\emph{la funzione di trasferimento ``vede'' solo la parte del sistema che è sia raggiungibile sia osservabile.}
\end{quote}

\bigskip

\section{Esercizio: equazione alle differenze e cancellazione polo-zero}

Consideriamo il sistema discreto descritto dall'equazione alle differenze
\[
y(t+2)+3y(t+1)+2y(t)=u(t+1)+2u(t),
\]
con condizioni iniziali nulle.

\subsection{Calcolo della funzione di trasferimento}

Applicando la trasformata \(Z\) con condizioni iniziali nulle si ottiene

\[
(z^2+3z+2)Y(z)=(z+2)U(z),
\]
quindi
\[
G(z)=\frac{Y(z)}{U(z)}=\frac{z+2}{z^2+3z+2}.
\]

Poiché
\[
z^2+3z+2=(z+1)(z+2),
\]
segue che
\[
G(z)=\frac{z+2}{(z+1)(z+2)}=\frac{1}{z+1}.
\]

Dunque c'è una cancellazione polo-zero nel punto \(z=-2\).

\subsection{Interpretazione della cancellazione}

La cancellazione dice che una realizzazione costruita \emph{prima} della riduzione può contenere uno stato in più rispetto al minimo necessario.

In particolare:
\begin{itemize}
    \item una realizzazione in forma canonica di controllo costruita dal denominatore non ridotto sarà completamente raggiungibile, ma non completamente osservabile;
    \item una realizzazione in forma canonica di osservazione costruita dal denominatore non ridotto sarà completamente osservabile, ma non completamente raggiungibile;
    \item se invece si realizza il sistema \emph{dopo} la cancellazione, si ottiene una realizzazione minima di ordine \(1\), che è sia raggiungibile sia osservabile.
\end{itemize}

\subsection{Una realizzazione minima}

Dalla forma ridotta
\[
G(z)=\frac{1}{z+1}
\]
si ottiene una realizzazione minima scalare, ad esempio

\[
x(t+1)=-x(t)+u(t),
\qquad
y(t)=x(t),
\]
cioè
\[
A=(-1),\qquad B=(1),\qquad C=(1),\qquad D=0.
\]

Essendo di ordine \(1\) e non avendo più cancellazioni, questa realizzazione è minima.

\bigskip

\section{Interconnessione di sistemi SISO}

Passiamo ora a un tema nuovo: come si realizza in forma di stato l'interconnessione di due sistemi.

\subsection{Connessione in serie}

Consideriamo due sistemi SISO \(\Sigma_1\) e \(\Sigma_2\), con funzioni di trasferimento
\[
G_1(s)=\frac{n_1(s)}{d_1(s)},
\qquad
G_2(s)=\frac{n_2(s)}{d_2(s)}.
\]

La connessione in serie è schematicamente

\[
u \longrightarrow \boxed{\Sigma_1} \longrightarrow \boxed{\Sigma_2} \longrightarrow y.
\]

Supponiamo che le realizzazioni siano

\[
\Sigma_1:
\begin{cases}
\dot x_1=A_1x_1+B_1u_1,\\
y_1=C_1x_1+D_1u_1,
\end{cases}
\qquad
\Sigma_2:
\begin{cases}
\dot x_2=A_2x_2+B_2u_2,\\
y_2=C_2x_2+D_2u_2.
\end{cases}
\]

Nella connessione in serie vale
\[
u_1=u,\qquad u_2=y_1,\qquad y=y_2.
\]

Quindi:
\[
u_2=C_1x_1+D_1u.
\]

Sostituendo nel secondo sistema:

\[
\dot x_2=A_2x_2+B_2(C_1x_1+D_1u)
      =A_2x_2+B_2C_1x_1+B_2D_1u.
\]

Inoltre
\[
y=y_2=C_2x_2+D_2u_2
     =C_2x_2+D_2C_1x_1+D_2D_1u.
\]

Se definiamo lo stato complessivo
\[
x=\begin{bmatrix}x_1\\x_2\end{bmatrix},
\]
si ottiene la realizzazione complessiva

\[
\begin{cases}
\dot x=
\begin{bmatrix}
A_1 & 0\\
B_2C_1 & A_2
\end{bmatrix}x+
\begin{bmatrix}
B_1\\
B_2D_1
\end{bmatrix}u,\\[1.2em]
y=
\begin{bmatrix}
D_2C_1 & C_2
\end{bmatrix}x+D_2D_1u.
\end{cases}
\]

Naturalmente, a livello di funzione di trasferimento, vale
\[
G(s)=G_2(s)G_1(s).
\]

\subsection{Cosa succede se ci sono cancellazioni?}

Se nella cascata si verifica una cancellazione polo-zero, la realizzazione ottenuta non è più minima.

Ci sono due casi concettualmente distinti.

\paragraph{Caso 1: uno zero di \(\Sigma_1\) coincide con un polo di \(\Sigma_2\).}

Allora il modo cancellato appartiene al secondo sistema, cioè a quello ``a valle''. In questo caso il sistema complessivo può perdere \emph{raggiungibilità}: esiste una dinamica interna di \(\Sigma_2\) che non riesce più a essere eccitata dall'ingresso esterno.

\paragraph{Caso 2: un polo di \(\Sigma_1\) coincide con uno zero di \(\Sigma_2\).}

In questo caso il modo cancellato appartiene al primo sistema, cioè a quello ``a monte''. Il sistema complessivo può perdere \emph{osservabilità}: esiste una dinamica interna che l'uscita finale non riesce più a vedere.

\medskip

In sintesi:
\begin{quote}
\emph{nelle interconnessioni in serie, le cancellazioni possono nascondere modi interni, rendendo il sistema complessivo non minimo.}
\end{quote}

\bigskip

\subsection{Connessione in parallelo}

La connessione in parallelo è schematicamente

\[
\begin{array}{c}
u \\
\downarrow
\end{array}
\qquad
\begin{array}{c}
\boxed{\Sigma_1}\\[0.4em]
\boxed{\Sigma_2}
\end{array}
\qquad
\longrightarrow\ y=y_1+y_2.
\]

Più precisamente:
\[
u_1=u,\qquad u_2=u,\qquad y=y_1+y_2.
\]

Se
\[
\Sigma_i:
\begin{cases}
\dot x_i=A_ix_i+B_iu,\\
y_i=C_ix_i+D_iu,
\end{cases}
\qquad i=1,2,
\]
e si pone
\[
x=\begin{bmatrix}x_1\\x_2\end{bmatrix},
\]
allora la realizzazione complessiva è

\[
\begin{cases}
\dot x=
\begin{bmatrix}
A_1 & 0\\
0 & A_2
\end{bmatrix}x+
\begin{bmatrix}
B_1\\
B_2
\end{bmatrix}u,\\[1.2em]
y=
\begin{bmatrix}
C_1 & C_2
\end{bmatrix}x+(D_1+D_2)u.
\end{cases}
\]

A livello di funzione di trasferimento:
\[
G(s)=G_1(s)+G_2(s)
=
\frac{n_1(s)}{d_1(s)}+\frac{n_2(s)}{d_2(s)}
=
\frac{n_1(s)d_2(s)+n_2(s)d_1(s)}{d_1(s)d_2(s)}.
\]

\subsection{Quando possono avvenire cancellazioni in parallelo?}

Nel collegamento in parallelo, eventuali cancellazioni nel sistema risultante possono comparire solo se i due sottosistemi hanno fattori dinamici comuni, cioè poli coincidenti nei denominatori.

In quel caso, il modo comune può sparire nella somma delle due risposte e il sistema complessivo non risulta più minimo.

Anche qui il messaggio è lo stesso:
\[
\text{cancellazione} \quad \Longrightarrow \quad \text{realizzazione non minima}.
\]

\bigskip

\section{Sistemi MIMO e matrice di trasferimento}

Passiamo ora ai sistemi con più ingressi e più uscite.

Se il sistema ha \(m\) ingressi e \(\ell\) uscite, la sua matrice di trasferimento è
\[
G(s)\in\mathbb R(s)^{\ell\times m},
\]
e si scrive come
\[
Y(s)=G(s)U(s),
\]
dove
\[
Y(s)=
\begin{bmatrix}
Y_1(s)\\ \vdots \\ Y_\ell(s)
\end{bmatrix},
\qquad
U(s)=
\begin{bmatrix}
U_1(s)\\ \vdots \\ U_m(s)
\end{bmatrix}.
\]

Gli elementi della matrice sono le funzioni di trasferimento ingresso-uscita:
\[
G(s)=
\begin{bmatrix}
g_{11}(s) & \cdots & g_{1m}(s)\\
\vdots & \ddots & \vdots\\
g_{\ell 1}(s) & \cdots & g_{\ell m}(s)
\end{bmatrix},
\]
e per ogni uscita vale
\[
Y_i(s)=g_{i1}(s)U_1(s)+g_{i2}(s)U_2(s)+\cdots+g_{im}(s)U_m(s).
\]

\subsection{Significato dei poli della matrice di trasferimento}

Per capire i poli di una matrice di trasferimento, una prima idea è spegnere tutti gli ingressi tranne uno.

Se fissiamo \(j\) e imponiamo
\[
U_k(s)=0 \qquad \forall k\neq j,
\]
allora
\[
Y(s)=G_{\bullet j}(s)\,U_j(s),
\]
dove \(G_{\bullet j}(s)\) è la \(j\)-esima colonna di \(G(s)\).

Questo significa che una colonna della matrice di trasferimento descrive \emph{tutti gli effetti del singolo ingresso \(u_j\) su tutte le uscite}. Analogamente, una riga descrive come una certa uscita dipende da tutti gli ingressi.

\subsection{Il termine diretto \(D\)}

Nel modello in forma di stato
\[
\begin{cases}
\dot x=Ax+Bu,\\
y=Cx+Du,
\end{cases}
\]
la matrice \(D\) raccoglie i termini di \emph{accoppiamento diretto} ingresso--uscita.

In una matrice di trasferimento MIMO, se un elemento \(g_{ij}(s)\) non è strettamente proprio, lo si può decomporre come
\[
g_{ij}(s)=d_{ij}+\text{parte strettamente propria},
\]
e il termine costante \(d_{ij}\) va inserito nella matrice \(D\).

\bigskip

\section{Realizzazione di una matrice di trasferimento lavorando per colonne}

Un metodo pratico per realizzare un sistema MIMO consiste nel lavorare \emph{per colonne}.

Supponiamo, ad esempio, di avere due ingressi e tre uscite, cioè una matrice di trasferimento \(G(s)\in\mathbb R(s)^{3\times 2}\). La scriviamo per colonne:
\[
G(s)=
\begin{bmatrix}
\vert & \vert\\
G^{(1)}(s) & G^{(2)}(s)\\
\vert & \vert
\end{bmatrix},
\]
dove ciascuna colonna \(G^{(j)}(s)\) contiene l'effetto dell'ingresso \(u_j\) su tutte le uscite.

Per ogni colonna:
\begin{enumerate}
    \item si cerca un denominatore comune;
    \item si costruisce una realizzazione SISO/MISO in forma canonica di controllo;
    \item si accoppiano poi i vari blocchi in modo disaccoppiato sugli ingressi.
\end{enumerate}

\subsection{Struttura generale ottenuta per colonne}

Se dalla colonna \(j\)-esima si ricava una realizzazione
\[
(A^{(j)},B^{(j)},C^{(j)}),
\]
allora una realizzazione complessiva si può scrivere come

\[
A=
\begin{bmatrix}
A^{(1)} & 0 & \cdots & 0\\
0 & A^{(2)} & \cdots & 0\\
\vdots & \vdots & \ddots & \vdots\\
0 & 0 & \cdots & A^{(m)}
\end{bmatrix},
\qquad
B=
\begin{bmatrix}
B^{(1)} & 0 & \cdots & 0\\
0 & B^{(2)} & \cdots & 0\\
\vdots & \vdots & \ddots & \vdots\\
0 & 0 & \cdots & B^{(m)}
\end{bmatrix},
\]
e
\[
C=
\begin{bmatrix}
C^{(1)} & C^{(2)} & \cdots & C^{(m)}
\end{bmatrix}.
\]

Questa costruzione è molto utile perché permette di realizzare la matrice di trasferimento in modo operativo. Tuttavia:
\begin{itemize}
    \item la realizzazione ottenuta \emph{non è in generale minima};
    \item potrebbero esserci cancellazioni tra colonne;
    \item spesso è necessario ridurre poi la realizzazione con strumenti strutturali.
\end{itemize}

\bigskip

\section{Metodo duale: lavorare per righe}

Esiste naturalmente il metodo duale, che consiste nel lavorare per \emph{righe}.

In questo caso si guarda ogni uscita separatamente:
\[
Y_i(s)=g_{i1}(s)U_1(s)+\cdots+g_{im}(s)U_m(s).
\]

Dal punto di vista realizzativo:
\begin{itemize}
    \item lavorare per colonne enfatizza il ruolo degli ingressi;
    \item lavorare per righe enfatizza il ruolo delle uscite;
    \item i due metodi sono duali.
\end{itemize}

Non c'è una regola assoluta per scegliere tra i due:
\begin{quote}
\emph{si sceglie di solito il metodo più conveniente, cioè quello che porta a sottosistemi di ordine più basso o con strutture più semplici.}
\end{quote}

\bigskip

\section{Come si individuano poli e zeri di una matrice di trasferimento?}

Per sistemi MIMO la nozione di poli e zeri non si legge più guardando semplicemente ogni singolo elemento della matrice. Serve uno strumento più robusto: la \emph{forma di Smith} di una matrice polinomiale.

\section{Forma di Smith di una matrice polinomiale}

Sia
\[
P(s)\in \mathbb R[s]^{p\times m}
\]
una matrice polinomiale.

Allora esistono matrici unimodulari \(U_1(s)\) e \(U_2(s)\) tali che
\[
P(s)=U_1(s)\,\Lambda(s)\,U_2(s),
\]
dove
\[
\Lambda(s)=
\operatorname{diag}\bigl(d_1(s),d_2(s),\dots,d_r(s),0,\dots,0\bigr),
\]
con:
\begin{itemize}
    \item \(r=\operatorname{rk}(P(s))\);
    \item i polinomi \(d_i(s)\) sono monici;
    \item \(d_i(s)\) divide \(d_{i+1}(s)\) per ogni \(i\).
\end{itemize}

Questa è la \emph{forma di Smith}.

\subsection{Come si costruiscono i polinomi invarianti}

Si definiscono:
\[
D_0(s)=1,
\]
e, per \(i=1,\dots,r\),
\[
D_i(s)=\gcd\{\text{tutti i minori di ordine } i \text{ di } P(s)\},
\]
presi monici.

Poi si pongono
\[
d_i(s)=\frac{D_i(s)}{D_{i-1}(s)}.
\]

I polinomi \(d_i(s)\) sono proprio gli \emph{invariant factors} della matrice.

\bigskip

\section{Esempio di forma di Smith}

Consideriamo
\[
P(s)=
\begin{bmatrix}
1 & -1\\
s^2+s-4 & 2s^2-s-8\\
s^2-4 & 2s^2-8
\end{bmatrix}.
\]

\subsection{Calcolo di \(D_1(s)\)}

Il massimo comun divisore di tutti gli elementi è
\[
D_1(s)=1.
\]

\subsection{Calcolo di \(D_2(s)\)}

Occorre calcolare i minori \(2\times 2\).

\paragraph{Minore delle righe 1 e 2:}
\[
\det
\begin{bmatrix}
1 & -1\\
s^2+s-4 & 2s^2-s-8
\end{bmatrix}
=
2s^2-s-8+s^2+s-4
=
3s^2-12
=
3(s^2-4).
\]

\paragraph{Minore delle righe 1 e 3:}
\[
\det
\begin{bmatrix}
1 & -1\\
s^2-4 & 2s^2-8
\end{bmatrix}
=
2s^2-8+s^2-4
=
3s^2-12
=
3(s^2-4).
\]

\paragraph{Minore delle righe 2 e 3:}
\[
\det
\begin{bmatrix}
s^2+s-4 & 2s^2-s-8\\
s^2-4 & 2s^2-8
\end{bmatrix}.
\]

Sviluppando:
\[
(s^2+s-4)(2s^2-8)-(2s^2-s-8)(s^2-4)
=
3s(s^2-4).
\]

Il massimo comun divisore monico di questi minori è quindi
\[
D_2(s)=s^2-4.
\]

\subsection{Polinomi invarianti}

Poiché
\[
D_0(s)=1,\qquad D_1(s)=1,\qquad D_2(s)=s^2-4,
\]
si ottiene
\[
d_1(s)=\frac{D_1(s)}{D_0(s)}=1,
\qquad
d_2(s)=\frac{D_2(s)}{D_1(s)}=s^2-4.
\]

Quindi la forma di Smith è

\[
\Lambda(s)=
\begin{bmatrix}
1 & 0\\
0 & s^2-4\\
0 & 0
\end{bmatrix}.
\]

\medskip

Questo esempio è importante perché mostra che, in un sistema MIMO, i poli e gli zeri strutturali non si leggono banalmente elemento per elemento, ma emergono dai fattori invarianti.

\bigskip

\section{Messaggio concettuale della lezione}

Possiamo riassumere l'intera lezione così.

\begin{enumerate}
    \item Una funzione di trasferimento può essere realizzata in molti modi diversi, ma la \emph{forma minima} è quella con il minor numero di stati ed è sia completamente raggiungibile sia completamente osservabile.
    \item Nelle interconnessioni in serie e in parallelo possono comparire cancellazioni interne, che rendono la realizzazione complessiva non minima.
    \item Nei sistemi MIMO la matrice di trasferimento va studiata come oggetto unico, non guardando soltanto i singoli elementi.
    \item Per realizzare un sistema MIMO si può lavorare per colonne o per righe, scegliendo la via più conveniente.
    \item Per studiare poli e zeri strutturali di una matrice di trasferimento si usa la forma di Smith della matrice polinomiale associata.
\end{enumerate}

\bigskip

\section*{Promemoria operativo}

Quando ti danno una matrice di trasferimento MIMO e ti chiedono una realizzazione, la strategia pratica è:

\begin{enumerate}
    \item separare l'eventuale parte diretta \(D\), cioè i termini non strettamente propri;
    \item decidere se conviene lavorare per colonne o per righe;
    \item per ciascun blocco trovato, costruire una realizzazione canonica del denominatore comune;
    \item assemblare i blocchi in una realizzazione complessiva;
    \item verificare se la realizzazione è minima oppure se restano modi cancellati;
    \item se serve, usare la decomposizione di Kalman o la forma di Smith per ridurre il sistema.
\end{enumerate}
% =========================
% L19.tex
% =========================

\chapter{L19 -- Realizzazione in forma di stato di sistemi MIMO, zeri di trasmissione e stabilità}
\label{chap:L19}

\section{Dalla matrice di trasferimento alla forma di stato}

In questa lezione affrontiamo tre temi collegati tra loro:
\begin{enumerate}
    \item come leggere la struttura di un sistema \textbf{MIMO} a partire dalla sua matrice di trasferimento;
    \item come individuare i \textbf{poli}, gli \textbf{zeri di trasmissione} e l'\textbf{ordine minimo} della realizzazione;
    \item come distinguere \textbf{stabilità interna} e \textbf{stabilità esterna (BIBO)}.
\end{enumerate}

Partiamo da una matrice di trasferimento razionale
\[
G(s)\in \mathbb{R}(s)^{\ell\times m},
\]
dove \(m\) è il numero di ingressi e \(\ell\) il numero di uscite.

L'obiettivo finale è sempre quello di costruire una realizzazione in forma di stato
\[
\begin{cases}
\dot x = Ax + Bu,\\
y = Cx + Du,
\end{cases}
\qquad
B\in\mathbb{R}^{n\times m},\quad
C\in\mathbb{R}^{\ell\times n},
\]
che rappresenti il sistema dato.

Nel caso MIMO, però, non basta guardare una singola funzione di trasferimento: la struttura del sistema è contenuta nell'intera matrice \(G(s)\). Per questo diventa naturale introdurre la \textbf{forma canonica di Smith} per matrici polinomiali e, di conseguenza, la \textbf{forma di Smith--McMillan} per matrici razionali.

\bigskip

\section{Forma canonica di Smith per matrici polinomiali}

Sia
\[
P(s)\in \mathbb{R}[s]^{\ell\times m}
\]
una matrice a coefficienti polinomiali. Allora esistono due matrici unimodulari \(U_1(s)\) e \(U_2(s)\) tali che
\[
P(s)=U_1(s)\,\Lambda(s)\,U_2(s),
\]
dove \(\Lambda(s)\) è diagonale della forma
\[
\Lambda(s)=
\operatorname{diag}\bigl(\delta_1(s),\delta_2(s),\dots,\delta_r(s),0,\dots,0\bigr),
\]
con:
\begin{itemize}
    \item \(r=\operatorname{rk}P(s)\) (rango normale della matrice polinomiale);
    \item ogni \(\delta_i(s)\) monico;
    \item \(\delta_i(s)\mid \delta_{i+1}(s)\) per ogni \(i\).
\end{itemize}

Questa è la \textbf{forma canonica di Smith}. I polinomi \(\delta_i(s)\) sono i \textbf{fattori invarianti} della matrice \(P(s)\).

\medskip

\noindent
\textbf{Osservazione importante.}
La forma di Smith si applica a una \emph{matrice polinomiale}, non direttamente a una matrice razionale. Per questo, quando partiamo da \(G(s)\), prima la riscriviamo come rapporto tra una matrice polinomiale e un denominatore comune.

\bigskip

\section{Dalla forma di Smith alla forma di Smith--McMillan}

Supponiamo che la matrice di trasferimento possa essere scritta come
\[
G(s)=\frac{1}{d(s)}\,P(s),
\]
dove:
\begin{itemize}
    \item \(d(s)\) è un polinomio scalare comune ai denominatori;
    \item \(P(s)\in\mathbb{R}[s]^{\ell\times m}\).
\end{itemize}

Se per \(P(s)\) calcoliamo la forma di Smith
\[
P(s)=U_1(s)\,\Lambda(s)\,U_2(s),
\]
allora
\[
G(s)=U_1(s)\,\Pi(s)\,U_2(s),
\qquad
\Pi(s)=\frac{\Lambda(s)}{d(s)}.
\]

Esplicitamente,
\[
\Pi(s)=
\operatorname{diag}\!\left(
\frac{\delta_1(s)}{d(s)},
\frac{\delta_2(s)}{d(s)},
\dots,
\frac{\delta_r(s)}{d(s)},
0,\dots,0
\right).
\]

Ora ogni rapporto \(\delta_i(s)/d(s)\) va semplificato. Dopo aver eseguito \emph{tutte} le cancellazioni possibili, si ottiene la \textbf{forma di Smith--McMillan}
\[
\Pi_{SM}(s)=
\operatorname{diag}\!\left(
\frac{\alpha_1(s)}{\beta_1(s)},
\dots,
\frac{\alpha_r(s)}{\beta_r(s)},
0,\dots,0
\right),
\]
dove:
\begin{itemize}
    \item i polinomi \(\alpha_i(s)\) e \(\beta_i(s)\) sono coprimi;
    \item i \(\beta_i(s)\) raccolgono i poli strutturali del sistema;
    \item gli \(\alpha_i(s)\) raccolgono gli zeri strutturali, cioè gli \textbf{zeri di trasmissione}.
\end{itemize}

\medskip

\noindent
\textbf{Informazioni che si leggono dalla forma di Smith--McMillan.}
\begin{itemize}
    \item I \textbf{poli} della matrice di trasferimento sono le radici dei \(\beta_i(s)\).
    \item Gli \textbf{zeri di trasmissione} sono le radici degli \(\alpha_i(s)\).
    \item L'\textbf{ordine minimo di una realizzazione} è la somma dei gradi dei denominatori dopo cancellazione:
    \[
    n_{\min}=\sum_{i=1}^r \deg \beta_i(s).
    \]
    Questo numero è il \textbf{grado di McMillan}.
\end{itemize}

\bigskip

\section{Esempio sulla forma di Smith--McMillan}

Consideriamo la matrice polinomiale
\[
P(s)=
\begin{bmatrix}
1 & -1\\
s^2+s-4 & 2s^2-s-8\\
s^2-4 & 2s^2-8
\end{bmatrix}.
\]

Dai calcoli della lezione precedente si sa che la sua forma di Smith è
\[
\Lambda(s)=
\begin{bmatrix}
1 & 0\\
0 & (s+2)(s-2)\\
0 & 0
\end{bmatrix}.
\]

Supponiamo ora di avere la matrice di trasferimento
\[
G(s)=\frac{1}{s^2+3s+2}\,P(s).
\]
Poiché
\[
s^2+3s+2=(s+1)(s+2),
\]
otteniamo
\[
\Pi(s)=\frac{\Lambda(s)}{s^2+3s+2}
=
\begin{bmatrix}
\dfrac{1}{(s+1)(s+2)} & 0\\[1.1em]
0 & \dfrac{(s+2)(s-2)}{(s+1)(s+2)}\\[1.1em]
0 & 0
\end{bmatrix}.
\]

Semplificando il secondo elemento diagonale:
\[
\Pi_{SM}(s)=
\begin{bmatrix}
\dfrac{1}{(s+1)(s+2)} & 0\\[1.1em]
0 & \dfrac{s-2}{s+1}\\[1.1em]
0 & 0
\end{bmatrix}.
\]

A questo punto possiamo leggere direttamente:

\subsection*{Poli}
I denominatori non banali sono
\[
(s+1)(s+2),\qquad s+1.
\]
Quindi i poli sono:
\[
s=-1 \quad \text{(doppio nel grado complessivo)},\qquad s=-2.
\]

Infatti il prodotto dei denominatori diagonali è
\[
(s+1)(s+2)\cdot (s+1)=(s+1)^2(s+2).
\]

\subsection*{Zeri di trasmissione}
I numeratori non banali sono
\[
1,\qquad s-2.
\]
Dunque compare un solo zero di trasmissione:
\[
s=2.
\]

\subsection*{Ordine minimo}
L'ordine minimo della realizzazione è
\[
n_{\min}=\deg\big((s+1)(s+2)\big)+\deg(s+1)=2+1=3.
\]

Questa è un'informazione molto importante:
\begin{quote}
\emph{anche se una qualunque realizzazione non minima può avere più stati, la forma minima del sistema ha dimensione \(3\).}
\end{quote}

\bigskip

\section{Interpretazione degli zeri di trasmissione}

Negli appunti compare la domanda fondamentale:

\begin{quote}
Dato un ingresso non nullo, può esistere una condizione iniziale tale che l'uscita sia identicamente nulla?
\end{quote}

Questa è proprio l'idea intuitiva di \textbf{zero di trasmissione}.

\subsection{Caso SISO continuo}

Consideriamo un sistema SISO in tempo continuo:
\[
\begin{cases}
\dot x = Ax + Bu,\\
y = Cx + Du.
\end{cases}
\]

Cerchiamo un ingresso del tipo
\[
u(t)=e^{\delta t}
\]
e una traiettoria dello stato della forma
\[
x(t)=x_0 e^{\delta t},
\]
con \(\delta\in\mathbb{C}\).

Sostituendo nell'equazione di stato:
\[
\delta x_0 e^{\delta t}=Ax_0 e^{\delta t}+B e^{\delta t}.
\]
Dividendo per \(e^{\delta t}\neq 0\),
\[
(\delta I-A)x_0=B.
\]

Se \(\delta\) non è autovalore di \(A\), possiamo invertire \(\delta I-A\) e scrivere
\[
x_0=(\delta I-A)^{-1}B.
\]

Allora l'uscita vale
\[
y(t)=C x(t)+D u(t)
= C(\delta I-A)^{-1}B\,e^{\delta t}+D e^{\delta t},
\]
cioè
\[
y(t)=\bigl(C(\delta I-A)^{-1}B+D\bigr)e^{\delta t}.
\]

Ma il termine tra parentesi è proprio
\[
G(\delta)=C(\delta I-A)^{-1}B+D.
\]

Quindi
\[
y(t)=G(\delta)e^{\delta t}.
\]

\medskip

\noindent
Se \(\delta\) è uno zero di \(G(s)\), allora
\[
G(\delta)=0
\quad\Longrightarrow\quad
y(t)\equiv 0,
\]
pur con ingresso non nullo.

\medskip

\noindent
\textbf{Interpretazione.}
Uno zero di trasmissione è un valore dell'esponenziale d'ingresso per cui è possibile scegliere opportunamente lo stato iniziale in modo da annullare completamente l'uscita.

\subsection{Generalizzazione MIMO}

Nel caso MIMO l'idea è analoga, ma l'ingresso ha una \emph{direzione}:
\[
u(t)=v e^{\delta t},\qquad v\neq 0.
\]
Si cercano quindi \(x_0\neq 0\) e \(v\neq 0\) tali che
\[
\begin{cases}
(\delta I-A)x_0-Bv=0,\\
Cx_0+Dv=0.
\end{cases}
\]

In forma compatta:
\[
\begin{bmatrix}
\delta I-A & -B\\
C & D
\end{bmatrix}
\begin{bmatrix}
x_0\\
v
\end{bmatrix}
=0.
\]

La matrice
\[
\mathcal{R}(\delta)=
\begin{bmatrix}
\delta I-A & -B\\
C & D
\end{bmatrix}
\]
è la \textbf{matrice di Rosenbrock}. I suoi valori \(\delta\) per cui il rango cala sono gli zeri di trasmissione del sistema.

\bigskip

\section{Stabilità interna}

La \textbf{stabilità interna} riguarda l'evoluzione dello \emph{stato}, quindi dipende dagli autovalori di \(A\), non direttamente dalla funzione di trasferimento.

\medskip

\noindent
Per il sistema autonomo
\[
\dot x = Ax,
\]
le componenti della soluzione sono costruite a partire dai modi associati agli autovalori di \(A\).

\subsection{Modi elementari nel caso continuo}

Se \(\lambda\in\mathbb{R}\) è autovalore reale di \(A\), i modi associati sono del tipo
\[
e^{\lambda t},\quad
t e^{\lambda t},\quad
t^2 e^{\lambda t},\quad \dots
\]
a seconda della dimensione dei blocchi di Jordan.

Se invece \(\lambda=\sigma\pm j\omega\) è complesso, compaiono modi del tipo
\[
e^{\sigma t}\cos(\omega t),\qquad
e^{\sigma t}\sin(\omega t),
\]
eventualmente moltiplicati per potenze di \(t\) nel caso di blocchi di Jordan non banali:
\[
t^k e^{\sigma t}\cos(\omega t),\qquad
t^k e^{\sigma t}\sin(\omega t).
\]

\subsection{Modi elementari nel caso discreto}

Per il sistema
\[
x(k+1)=Ax(k),
\]
i modi associati a un autovalore \(\lambda\) sono del tipo
\[
\lambda^k,\qquad
k\lambda^k,\qquad
k^2\lambda^k,\qquad \dots
\]
sempre a seconda della struttura di Jordan.

\subsection{Condizioni di stabilità interna asintotica}

\paragraph{Tempo continuo.}
Il sistema è asintoticamente stabile se e solo se
\[
\Re(\lambda_i)<0
\qquad \forall\, \lambda_i\in\sigma(A).
\]

\paragraph{Tempo discreto.}
Il sistema è asintoticamente stabile se e solo se
\[
|\lambda_i|<1
\qquad \forall\, \lambda_i\in\sigma(A).
\]

\medskip

\noindent
\textbf{Osservazione.}
Se un autovalore ha parte reale nulla (TC) o modulo unitario (TD), si possono avere oscillazioni persistenti: il sistema non diverge, ma non converge neppure a zero. Questo è esattamente ciò che succede nell'oscillatore armonico non smorzato.

\bigskip

\section{Stabilità esterna o BIBO}

La \textbf{stabilità esterna} riguarda invece il legame ingresso--uscita.

\medskip

\noindent
\textbf{Definizione (BIBO).}
Un sistema è \emph{Bounded Input Bounded Output} se, per ogni ingresso limitato, l'uscita corrispondente a stato iniziale nullo rimane limitata.

In simboli:
\[
\|u(t)\|\le M_u\ \forall t\ge 0
\quad\Longrightarrow\quad
\exists M_y>0:\ \|y(t)\|\le M_y\ \forall t\ge 0.
\]

Nel caso LTI, la stabilità BIBO si legge dai poli della funzione di trasferimento, non dagli autovalori nascosti di una realizzazione non minima.

\subsection{Condizione sui poli}

\paragraph{Tempo continuo.}
Un sistema razionale proprio è BIBO stabile se tutti i poli della funzione di trasferimento hanno parte reale strettamente negativa:
\[
\Re(p_i)<0.
\]

\paragraph{Tempo discreto.}
Un sistema razionale proprio è BIBO stabile se tutti i poli stanno strettamente dentro il cerchio unitario:
\[
|p_i|<1.
\]

\medskip

\noindent
\textbf{Punto chiave.}
I poli della funzione di trasferimento corrispondono solo agli autovalori di \(A\) che sono \textbf{sia raggiungibili sia osservabili}. Per questo:
\begin{itemize}
    \item la stabilità interna guarda \emph{tutti} gli autovalori di \(A\);
    \item la stabilità esterna guarda solo i modi che compaiono effettivamente in ingresso--uscita.
\end{itemize}

\bigskip

\section{Esempio: un sistema non BIBO stabile}

Consideriamo
\[
G(s)=\frac{1}{s},
\qquad
U(s)=\frac{1}{s}.
\]
Allora
\[
Y(s)=G(s)U(s)=\frac{1}{s^2}.
\]

Nel tempo, questo corrisponde a una rampa:
\[
y(t)=t.
\]
L'ingresso \(u(t)=1\) è limitato, ma l'uscita cresce senza limite. Dunque il sistema \(\frac{1}{s}\) \textbf{non è BIBO stabile}.

\medskip

Questo esempio chiarisce bene che, per avere BIBO stabilità, non basta che i modi non esplodano immediatamente: serve che l'uscita resti limitata per ogni ingresso limitato.

\bigskip

\section{Oscillatore armonico non smorzato}

Consideriamo il sistema
\[
\ddot y + y = u.
\]

Introduciamo le variabili di stato
\[
x_1=y,\qquad x_2=\dot y.
\]
Allora
\[
\begin{cases}
\dot x_1 = x_2,\\
\dot x_2 = -x_1 + u,
\end{cases}
\]
ossia
\[
\dot x =
\begin{bmatrix}
0 & 1\\
-1 & 0
\end{bmatrix}x+
\begin{bmatrix}
0\\
1
\end{bmatrix}u,
\qquad
y=
\begin{bmatrix}
1 & 0
\end{bmatrix}x.
\]

\subsection{Stabilità interna}

La matrice di stato è
\[
A=
\begin{bmatrix}
0 & 1\\
-1 & 0
\end{bmatrix},
\]
i cui autovalori sono
\[
\lambda_{1,2}=\pm j.
\]

Quindi il sistema autonomo non è asintoticamente stabile: i modi sono sinusoidali puri. Con \(u=0\), le traiettorie sono chiuse e non convergono all'origine.

Infatti:
\[
\frac{d}{dt}(x_1^2+x_2^2)
=2x_1\dot x_1+2x_2\dot x_2
=2x_1x_2+2x_2(-x_1)=0,
\]
quindi
\[
x_1^2(t)+x_2^2(t)=\text{costante}.
\]

Dunque le traiettorie stanno su circonferenze concentriche.

\begin{figure}[h!]
\centering
\begin{tikzpicture}[scale=1.1]
    % assi
    \draw[->] (-2.6,0) -- (2.6,0) node[right] {$x_1$};
    \draw[->] (0,-2.6) -- (0,2.6) node[above] {$x_2$};

    % orbite
    \draw[dashed] (0,0) circle (1.1);
    \draw (0,0) circle (2.0);

    % frecce orbita esterna (senso orario)
    \draw[->] (2.0,0) arc[start angle=0,end angle=-30,radius=2.0];
    \draw[->] (0,-2.0) arc[start angle=-90,end angle=-120,radius=2.0];
    \draw[->] (-2.0,0) arc[start angle=180,end angle=150,radius=2.0];
    \draw[->] (0,2.0) arc[start angle=90,end angle=60,radius=2.0];

    % frecce orbita interna
    \draw[->] (1.1,0) arc[start angle=0,end angle=-35,radius=1.1];
    \draw[->] (0,-1.1) arc[start angle=-90,end angle=-125,radius=1.1];
    \draw[->] (-1.1,0) arc[start angle=180,end angle=145,radius=1.1];
    \draw[->] (0,1.1) arc[start angle=90,end angle=55,radius=1.1];

    \node at (1.55,1.6) {\small \(x_1^2+x_2^2=c\)};
\end{tikzpicture}
\caption{Ritratto di fase dell'oscillatore armonico non smorzato.}
\end{figure}

\subsection{Stabilità esterna}

Calcoliamo la funzione di trasferimento:
\[
G(s)=C(sI-A)^{-1}B
=
\begin{bmatrix}
1 & 0
\end{bmatrix}
\left(
\begin{bmatrix}
s & -1\\
1 & s
\end{bmatrix}^{-1}
\right)
\begin{bmatrix}
0\\
1
\end{bmatrix}
=
\frac{1}{s^2+1}.
\]

I poli sono in
\[
s=\pm j.
\]
Essendo poli sull'asse immaginario, il sistema \textbf{non è BIBO stabile}. Infatti un ingresso sinusoidale in risonanza produce un'uscita non limitata.

Ad esempio, con
\[
u(t)=\sin t,
\]
si ottiene una risposta contenente un termine del tipo
\[
t\cos t,
\]
che cresce in modulo e quindi non resta limitato.

\medskip

\noindent
Questo esempio è molto istruttivo:
\begin{itemize}
    \item il sistema non è asintoticamente stabile internamente;
    \item il sistema non è BIBO stabile esternamente;
    \item le due nozioni però restano concettualmente distinte.
\end{itemize}

\bigskip

\section{Ricapitolazione finale}

Raccogliamo i messaggi centrali della lezione.

\subsection*{1. Realizzazione MIMO}
Per un sistema MIMO si parte dalla matrice di trasferimento \(G(s)\). Per capirne la struttura conviene scriverla come
\[
G(s)=\frac{1}{d(s)}P(s),
\]
studiare la forma di Smith di \(P(s)\) e poi la forma di Smith--McMillan di \(G(s)\).

\subsection*{2. Poli, zeri e ordine minimo}
Dalla forma di Smith--McMillan si leggono:
\begin{itemize}
    \item i poli della matrice di trasferimento;
    \item gli zeri di trasmissione;
    \item il grado di McMillan, cioè la dimensione minima di una realizzazione.
\end{itemize}

\subsection*{3. Zeri di trasmissione}
Uno zero di trasmissione è un valore \(\delta\) tale che esiste un ingresso esponenziale \(e^{\delta t}\) e una opportuna condizione iniziale che rendono l'uscita identicamente nulla.

\subsection*{4. Stabilità interna}
Dipende dagli autovalori di \(A\), quindi dai modi dello stato.

\subsection*{5. Stabilità esterna}
Dipende dai poli della funzione di trasferimento, cioè solo dai modi raggiungibili e osservabili.

\subsection*{6. Differenza tra le due}
In una realizzazione non minima si può avere:
\begin{itemize}
    \item sistema internamente instabile;
    \item ma esternamente stabile,
\end{itemize}
se il modo instabile è nascosto perché non raggiungibile o non osservabile.

\medskip

In altre parole:
\begin{quote}
\emph{la funzione di trasferimento non racconta tutto dello stato interno: racconta solo ciò che è visibile dall'ingresso e dall'uscita.}
\end{quote}
% =========================
% L20.tex
% =========================

\chapter{L20 -- Equilibri, movimenti, stabilità e linearizzazione}
\label{chap:L20}

\section{Richiamo iniziale: stato, traiettoria e rappresentazione grafica}

Prima di introdurre i concetti di equilibrio e stabilità, è utile chiarire bene il significato geometrico di una traiettoria di stato.

Consideriamo, come esempio introduttivo, l'oscillatore armonico non forzato
\[
\ddot y(t) + y(t) = 0.
\]
Introducendo le variabili di stato
\[
x_1(t)=y(t), 
\qquad
x_2(t)=\dot y(t),
\]
si ottiene il sistema
\[
\begin{cases}
\dot x_1(t)=x_2(t),\\
\dot x_2(t)=-x_1(t).
\end{cases}
\]
In forma matriciale:
\[
\dot x(t)=
\begin{bmatrix}
0 & 1\\
-1 & 0
\end{bmatrix}x(t).
\]

Se ad esempio si sceglie la condizione iniziale
\[
x(0)=
\begin{bmatrix}
1\\
0
\end{bmatrix},
\]
la soluzione è
\[
x_1(t)=\cos t,
\qquad
x_2(t)=-\sin t.
\]
Quindi, nel piano di stato \((x_1,x_2)\), la traiettoria descrive una circonferenza unitaria.

\subsection*{Interpretazione geometrica}
Questo esempio fa vedere bene la differenza tra:
\begin{itemize}
    \item il \textbf{movimento nel tempo}, cioè l'evoluzione delle componenti \(x_1(t)\), \(x_2(t)\) come funzioni di \(t\);
    \item la \textbf{traiettoria nello spazio di stato}, cioè il luogo dei punti \(x(t)\) nel piano \((x_1,x_2)\).
\end{itemize}

\begin{figure}[h!]
\centering
\begin{tikzpicture}[scale=1.25]
    % assi
    \draw[->] (-2.2,0) -- (2.4,0) node[right] {$x_1$};
    \draw[->] (0,-2.2) -- (0,2.4) node[above] {$x_2$};

    % circonferenza
    \draw[dashed] (0,0) circle (1.5);

    % frecce senso orario
    \draw[->] (1.5,0) arc[start angle=0,end angle=-35,radius=1.5];
    \draw[->] (0,-1.5) arc[start angle=-90,end angle=-125,radius=1.5];
    \draw[->] (-1.5,0) arc[start angle=180,end angle=145,radius=1.5];
    \draw[->] (0,1.5) arc[start angle=90,end angle=55,radius=1.5];

    % punto iniziale
    \fill (1.5,0) circle (1.5pt);
    \node[below right] at (1.5,0) {$x(0)$};

    \node at (1.15,1.35) {\small \(x_1^2+x_2^2=1\)};
\end{tikzpicture}
\caption{Traiettoria dell'oscillatore armonico nel piano di stato.}
\end{figure}

\bigskip

\section{Equilibrio di un sistema dinamico}

Passiamo adesso al concetto di equilibrio.

\subsection{Caso continuo}
Per un sistema non lineare tempo continuo
\[
\dot x = f(x,u),
\]
una coppia \((\bar x,\bar u)\) si dice \textbf{coppia di equilibrio} se, applicando l'ingresso costante \(\bar u\), lo stato \(\bar x\) rimane costante nel tempo. Questo equivale a imporre
\[
f(\bar x,\bar u)=0.
\]

\subsection{Caso discreto}
Per un sistema non lineare tempo discreto
\[
x(k+1)=f(x(k),u(k)),
\]
la coppia \((\bar x,\bar u)\) è di equilibrio se
\[
\bar x=f(\bar x,\bar u).
\]

\subsection*{Interpretazione}
Dire che \((\bar x,\bar u)\) è una coppia di equilibrio significa che, se il sistema parte da \(\bar x\) e l'ingresso viene mantenuto uguale a \(\bar u\), allora lo stato non si muove:
\[
x(t)\equiv \bar x
\quad \text{oppure} \quad
x(k)\equiv \bar x.
\]

In altre parole, la ``velocità'' dello stato è nulla: il sistema resta fermo in quel punto.

\bigskip

\section{Movimento e traiettoria}

Sia assegnata una soluzione del sistema, ad esempio nel continuo
\[
x(t)=\varphi(t,t_0,x_0,u),
\]
oppure nel discreto
\[
x(k)=\varphi(k,k_0,x_0,u).
\]

\begin{itemize}
    \item Il \textbf{movimento} è l'evoluzione temporale del sistema.
    \item La \textbf{traiettoria} è la curva descritta dal vettore di stato nello spazio di stato.
\end{itemize}

Se ad esempio si guarda una sola componente \(x_i(t)\), si ottiene il grafico in funzione del tempo. Se invece si rappresenta \(x(t)\) nello spazio degli stati, si ottiene la traiettoria.

\begin{figure}[h!]
\centering
\begin{tikzpicture}[scale=1.0]
    % sinistra: x(t)
    \draw[->] (0,0) -- (4.4,0) node[right] {$t$};
    \draw[->] (0,0) -- (0,2.8) node[above] {$x(t)$};
    \draw[smooth] plot coordinates {(0.2,0.4) (0.8,1.2) (1.5,1.0) (2.2,0.6) (3.1,1.8) (4.0,2.2)};
    \draw[dashed] (2.2,0) -- (2.2,0.6);
    \node[below] at (2.2,0) {$t^\star$};

    % destra: traiettoria nello spazio di stato
    \begin{scope}[xshift=7cm]
        \draw[->] (-2.2,0) -- (2.5,0) node[right] {$x_1$};
        \draw[->] (0,-2.0) -- (0,2.5) node[above] {$x_2$};
        \draw[smooth] plot coordinates {(-1.8,-1.2) (-1.0,-0.5) (-0.3,0.8) (0.4,1.3) (1.0,0.7) (1.8,1.8)};
        \fill (0.4,1.3) circle (1.2pt);
        \node[right] at (0.45,1.3) {\small \(x(t^\star)\)};
    \end{scope}
\end{tikzpicture}
\caption{A sinistra un movimento nel tempo, a destra la corrispondente traiettoria nello spazio di stato.}
\end{figure}

\bigskip

\section{Stabilità di un movimento nominale}

La teoria della stabilità non si applica solo ai punti di equilibrio: in generale si può parlare anche di stabilità di un \textbf{movimento nominale}.

Supponiamo di avere una traiettoria nominale
\[
x(t)=\varphi(t,t_0,x_0,\bar u)
\]
generata da un certo ingresso assegnato \(\bar u\), e una traiettoria perturbata
\[
\tilde x(t)=\varphi(t,t_0,x_0+\delta x_0,\bar u)
\]
ottenuta modificando leggermente la condizione iniziale.

\subsection{Movimento stabile}
Il movimento nominale si dice \textbf{stabile} se
\[
\forall \varepsilon>0 \ \exists \delta>0 \ \text{tale che}\
\|\delta x_0\|<\delta
\ \Longrightarrow\
\|\tilde x(t)-x(t)\|<\varepsilon
\quad \forall t\ge t_0.
\]

Interpretazione: se parto sufficientemente vicino alla traiettoria nominale, allora ci resto vicino per tutti i tempi successivi.

\subsection{Movimento attrattivo}
Il movimento nominale si dice \textbf{attrattivo} se esiste un intorno iniziale tale che
\[
\|\delta x_0\|<\delta
\ \Longrightarrow\
\lim_{t\to+\infty}\|\tilde x(t)-x(t)\|=0.
\]

Interpretazione: le traiettorie perturbate, magari dopo essersi un po' allontanate, alla fine tendono alla traiettoria nominale.

\subsection{Movimento asintoticamente stabile}
Il movimento nominale si dice \textbf{asintoticamente stabile} se è sia stabile sia attrattivo.

\bigskip

\section{Esempio: attrattivo ma non stabile}

Un esempio classico, presente anche nei tuoi appunti, è il sistema discreto scalare
\[
x(k+1)=
\begin{cases}
2x(k), & |x(k)|<1,\\[0.1cm]
0, & |x(k)|\ge 1.
\end{cases}
\]

Studiamo l'equilibrio nell'origine \(x=0\).

\subsection*{Perché è attrattivo}
Se la traiettoria prima o poi esce dall'intervallo \((-1,1)\), al passo successivo va a zero e poi resta a zero. Dunque molte traiettorie convergono all'origine.

\subsection*{Perché non è stabile}
Prendiamo una condizione iniziale molto piccola, ad esempio
\[
x(0)=\frac{1}{2}.
\]
Allora
\[
x(1)=2x(0)=1,
\qquad
x(2)=0.
\]
Quindi la traiettoria, pur convergendo infine all'origine, prima si allontana da un intorno arbitrariamente piccolo dello zero. Questo viola la definizione di stabilità.

\begin{figure}[h!]
\centering
\begin{tikzpicture}[scale=1.0]
    \draw[->] (0,0) -- (5.2,0) node[right] {$k$};
    \draw[->] (0,-0.4) -- (0,2.6) node[above] {$x(k)$};

    % livelli
    \draw[dashed] (0,1) -- (4.7,1);
    \draw[dashed] (0,0.5) -- (4.7,0.5);

    % punti
    \fill (0,0.5) circle (1.6pt);
    \fill (1,1.0) circle (1.6pt);
    \fill (2,0.0) circle (1.6pt);
    \fill (3,0.0) circle (1.6pt);
    \fill (4,0.0) circle (1.6pt);

    \draw (0,0.5) -- (1,1.0) -- (2,0.0) -- (4,0.0);

    \node[below] at (0,0) {$0$};
    \node[below] at (1,0) {$1$};
    \node[below] at (2,0) {$2$};
\end{tikzpicture}
\caption{Traiettoria che converge a zero ma non resta vicina allo zero per tutti i tempi: l'origine è attrattiva ma non stabile.}
\end{figure}

\bigskip

\section{Equilibrio come caso particolare di movimento}

Un punto di equilibrio è un caso particolare di traiettoria nominale: è una traiettoria costante.

Infatti, se \((\bar x,\bar u)\) è una coppia di equilibrio, allora la traiettoria nominale associata è
\[
x(t)\equiv \bar x
\qquad \text{oppure} \qquad
x(k)\equiv \bar x.
\]

Per questo motivo le definizioni di stabilità, attrattività e stabilità asintotica di un equilibrio si ottengono come caso particolare di quelle date per un movimento qualunque.

\bigskip

\section{Traslazione dell'equilibrio all'origine}

Quando si studia la stabilità di un equilibrio \(\bar x\), è quasi sempre conveniente traslare il sistema in nuove coordinate in modo che l'equilibrio diventi l'origine.

Definiamo
\[
z=x-\bar x,
\qquad
v=u-\bar u.
\]

\subsection{Caso continuo}
Se
\[
\dot x=f(x,u),
\]
allora
\[
\dot z=\dot x=f(z+\bar x,v+\bar u).
\]
Poiché \((\bar x,\bar u)\) è equilibrio, vale
\[
f(\bar x,\bar u)=0.
\]
Quindi il sistema traslato ha equilibrio nell'origine \(z=0\).

\subsection{Caso discreto}
Se
\[
x(k+1)=f(x(k),u(k)),
\]
ponendo
\[
z(k)=x(k)-\bar x,
\qquad
v(k)=u(k)-\bar u,
\]
si ottiene di nuovo un sistema equivalente con equilibrio nell'origine.

\medskip

\noindent
\textbf{Conclusione importante.}
Le proprietà di stabilità dell'equilibrio \(\bar x\) nel sistema originale coincidono con le proprietà di stabilità dell'origine nel sistema traslato.

Per questo, da qui in poi, si studia quasi sempre la stabilità dell'origine.

\bigskip

\section{Equilibri dei sistemi lineari autonomi}

Consideriamo ora un sistema lineare autonomo.

\subsection{Tempo continuo}
\[
\dot x=Ax.
\]
Un punto \(\bar x\) è equilibrio se
\[
A\bar x=0.
\]
Quindi l'insieme degli equilibri è
\[
\ker A.
\]

\subsection{Tempo discreto}
\[
x(k+1)=Ax(k).
\]
Un punto \(\bar x\) è equilibrio se
\[
\bar x=A\bar x
\quad \Longleftrightarrow \quad
(A-I)\bar x=0.
\]
Quindi l'insieme degli equilibri è
\[
\ker(A-I).
\]

\subsection*{Osservazione}
Se \(A\) non è invertibile (nel continuo), oppure se \(1\) è autovalore di \(A\) (nel discreto), allora possono esistere infiniti equilibri.

\bigskip

\section{Perché un equilibrio non isolato non può essere asintoticamente stabile}

Se un equilibrio non è isolato, allora in ogni suo intorno esistono altri punti di equilibrio.

Ma una traiettoria che parte esattamente da un altro equilibrio resta ferma in quel punto e quindi non può convergere all'equilibrio considerato. Perciò un equilibrio non isolato non può essere attrattivo, e dunque non può essere asintoticamente stabile.

\medskip

\noindent
Può però essere \textbf{stabile} nel senso di Lyapunov.

\bigskip

\section{Regione di attrazione o RAS}

Sia l'origine un equilibrio asintoticamente stabile. Si definisce \textbf{regione di attrazione} (RAS, \emph{Region of Asymptotic Stability}) l'insieme delle condizioni iniziali che generano traiettorie convergenti all'origine:
\[
\mathcal{A}(0)=\left\{x_0 \in \mathbb{R}^n \ :\ \lim_{t\to +\infty}\varphi(t,t_0,x_0)=0 \right\}.
\]

Se l'origine è \textbf{globalmente asintoticamente stabile}, allora
\[
\mathcal{A}(0)=\mathbb{R}^n.
\]

In generale, però, la RAS può essere un sottoinsieme proprio dello spazio di stato e non è sempre facile da determinare esplicitamente.

\bigskip

\section{Linearizzazione di un sistema non lineare}

Per studiare la stabilità locale di un equilibrio di un sistema non lineare, si usa la linearizzazione.

Consideriamo il sistema
\[
\dot x = f(x,u),
\qquad
y=h(x,u),
\]
e una coppia di equilibrio \((\bar x,\bar u)\), cioè
\[
f(\bar x,\bar u)=0.
\]

Introduciamo le variabili di deviazione:
\[
z=x-\bar x,
\qquad
v=u-\bar u,
\qquad
\eta=y-\bar y,
\quad \bar y=h(\bar x,\bar u).
\]

Sviluppando \(f\) e \(h\) al primo ordine attorno a \((\bar x,\bar u)\) si ottiene
\[
f(x,u)\approx
f(\bar x,\bar u)
+\left.\frac{\partial f}{\partial x}\right|_{(\bar x,\bar u)}(x-\bar x)
+\left.\frac{\partial f}{\partial u}\right|_{(\bar x,\bar u)}(u-\bar u),
\]
e analogamente
\[
h(x,u)\approx
h(\bar x,\bar u)
+\left.\frac{\partial h}{\partial x}\right|_{(\bar x,\bar u)}(x-\bar x)
+\left.\frac{\partial h}{\partial u}\right|_{(\bar x,\bar u)}(u-\bar u).
\]

Poiché \(f(\bar x,\bar u)=0\), il sistema linearizzato diventa
\[
\begin{cases}
\dot z = Az + Bv,\\
\eta = Cz + Dv,
\end{cases}
\]
dove
\[
A=\left.\frac{\partial f}{\partial x}\right|_{(\bar x,\bar u)},
\qquad
B=\left.\frac{\partial f}{\partial u}\right|_{(\bar x,\bar u)},
\qquad
C=\left.\frac{\partial h}{\partial x}\right|_{(\bar x,\bar u)},
\qquad
D=\left.\frac{\partial h}{\partial u}\right|_{(\bar x,\bar u)}.
\]

\medskip

\noindent
\textbf{Interpretazione.}
La linearizzazione fornisce una buona approssimazione del sistema non lineare \emph{solo localmente}, cioè per stati e ingressi sufficientemente vicini al punto di equilibrio considerato.

\bigskip

\section{Esempio: pendolo con attrito}

Consideriamo il pendolo con attrito viscoso e ingresso di coppia:
\[
m\ell^2\ddot\theta + c\dot\theta + mg\ell \sin\theta = u.
\]

Introduciamo le variabili di stato
\[
x_1=\theta,
\qquad
x_2=\dot\theta.
\]
Allora il sistema si scrive come
\[
\begin{cases}
\dot x_1 = x_2,\\
\dot x_2 = -\dfrac{g}{\ell}\sin x_1 - \dfrac{c}{m\ell^2}x_2 + \dfrac{1}{m\ell^2}u.
\end{cases}
\]

\subsection{Equilibri per \(u=0\)}
Se \(u=0\), gli equilibri soddisfano
\[
\begin{cases}
0=x_2,\\
0=-\dfrac{g}{\ell}\sin x_1.
\end{cases}
\]
Quindi
\[
x_2=0,
\qquad
\sin x_1=0
\quad\Longrightarrow\quad
x_1=k\pi,\qquad k\in\mathbb{Z}.
\]
Dunque ci sono infiniti equilibri:
\[
(k\pi,0), \qquad k\in\mathbb{Z}.
\]

\subsection{Jacobiana del sistema}
La matrice Jacobiana rispetto allo stato è
\[
\frac{\partial f}{\partial x}(x_1,x_2)=
\begin{bmatrix}
0 & 1\\[0.1cm]
-\dfrac{g}{\ell}\cos x_1 & -\dfrac{c}{m\ell^2}
\end{bmatrix},
\]
mentre rispetto all'ingresso
\[
\frac{\partial f}{\partial u}(x_1,x_2)=
\begin{bmatrix}
0\\[0.1cm]
\dfrac{1}{m\ell^2}
\end{bmatrix}.
\]

\subsection{Linearizzazione attorno a \((0,0)\)}
Poiché \(\cos 0=1\), si ottiene
\[
A_1=
\begin{bmatrix}
0 & 1\\[0.1cm]
-\dfrac{g}{\ell} & -\dfrac{c}{m\ell^2}
\end{bmatrix},
\qquad
B_1=
\begin{bmatrix}
0\\[0.1cm]
\dfrac{1}{m\ell^2}
\end{bmatrix}.
\]

La traccia e il determinante valgono
\[
\operatorname{tr}(A_1)=-\frac{c}{m\ell^2}<0,
\qquad
\det(A_1)=\frac{g}{\ell}>0.
\]
Quindi gli autovalori hanno parte reale negativa: l'equilibrio \((0,0)\) risulta localmente asintoticamente stabile se \(c>0\).

\subsection{Linearizzazione attorno a \((\pi,0)\)}
Poiché \(\cos \pi=-1\), si ottiene
\[
A_2=
\begin{bmatrix}
0 & 1\\[0.1cm]
\dfrac{g}{\ell} & -\dfrac{c}{m\ell^2}
\end{bmatrix},
\qquad
B_2=
\begin{bmatrix}
0\\[0.1cm]
\dfrac{1}{m\ell^2}
\end{bmatrix}.
\]

Ora
\[
\det(A_2)=-\frac{g}{\ell}<0.
\]
Quindi gli autovalori hanno segno opposto e l'equilibrio \((\pi,0)\) è instabile.

\medskip

\noindent
Poiché il pendolo è periodico in \(\theta\), gli equilibri con \(x_1=2k\pi\) hanno lo stesso comportamento di \((0,0)\), mentre quelli con \(x_1=(2k+1)\pi\) hanno lo stesso comportamento di \((\pi,0)\).

\bigskip

\section{Teorema indiretto di Lyapunov}

Il teorema indiretto di Lyapunov consente di dedurre la stabilità locale del sistema non lineare a partire dalla stabilità del sistema linearizzato.

Consideriamo il sistema autonomo
\[
\dot x=f(x),
\qquad
f(0)=0.
\]
Sia
\[
\dot z = Az
\]
la linearizzazione di \(\dot x=f(x)\) attorno all'origine, con
\[
A=\left.\frac{\partial f}{\partial x}\right|_{x=0}.
\]

Allora valgono i seguenti risultati.

\subsection*{1. Caso asintoticamente stabile}
Se la matrice \(A\) è di Hurwitz, cioè tutti i suoi autovalori hanno parte reale strettamente negativa, allora l'origine del sistema non lineare è \textbf{localmente asintoticamente stabile}.

\subsection*{2. Caso instabile}
Se \(A\) possiede almeno un autovalore con parte reale positiva, allora l'origine del sistema non lineare è \textbf{instabile}.

\subsection*{3. Caso critico}
Se tutti gli autovalori di \(A\) hanno parte reale non positiva, ma almeno uno ha parte reale nulla, allora il teorema indiretto \textbf{non è conclusivo}. In questo caso servono altri strumenti, per esempio il metodo diretto di Lyapunov.

\bigskip

\section{Caso critico: perché il teorema indiretto può non bastare}

Se, nel pendolo, si pone \(c=0\), la linearizzazione attorno a \((0,0)\) diventa
\[
A=
\begin{bmatrix}
0 & 1\\
-\dfrac{g}{\ell} & 0
\end{bmatrix}.
\]
Allora
\[
\operatorname{tr}(A)=0,
\qquad
\det(A)=\frac{g}{\ell}>0,
\]
e gli autovalori sono
\[
\lambda_{1,2}=\pm j\sqrt{\frac{g}{\ell}}.
\]

Il sistema linearizzato ha quindi orbite periodiche chiuse, e il teorema indiretto di Lyapunov non permette di concludere né stabilità asintotica né instabilità del sistema non lineare.

Questo è il tipico \textbf{caso critico}.

\bigskip

\section{Schema riassuntivo sulle classificazioni di stabilità}

Per un equilibrio dell'origine nel continuo:
\begin{itemize}
    \item se esiste un autovalore con \(\Re(\lambda)>0\), allora l'origine è \textbf{instabile};
    \item se tutti gli autovalori hanno \(\Re(\lambda)<0\), allora l'origine è \textbf{localmente asintoticamente stabile};
    \item se tutti hanno \(\Re(\lambda)\le 0\) e qualcuno ha \(\Re(\lambda)=0\), il test lineare è \textbf{inconcludente}.
\end{itemize}

Nel discreto, le condizioni si riscrivono sostituendo:
\[
\Re(\lambda)<0
\quad \longleftrightarrow \quad
|\lambda|<1,
\]
\[
\Re(\lambda)>0
\quad \text{(nel caso continuo)}
\quad \leadsto \quad
|\lambda|>1
\quad \text{(nel caso discreto)}.
\]

\bigskip

\section{Esempio geometrico: equilibrio stabile e instabile}

Un modo utile per interpretare la stabilità è osservare il comportamento delle traiettorie vicine a un punto di equilibrio.

\begin{itemize}
    \item Se tutte le traiettorie che partono abbastanza vicine restano vicine, l'equilibrio è \textbf{stabile}.
    \item Se inoltre tendono all'equilibrio, è \textbf{asintoticamente stabile}.
    \item Se esistono traiettorie arbitrariamente vicine che se ne allontanano, l'equilibrio è \textbf{instabile}.
\end{itemize}

Ad esempio, nel pendolo:
\begin{itemize}
    \item la posizione verso il basso \(\theta=0\) è stabile (e con attrito è anche asintoticamente stabile);
    \item la posizione capovolta \(\theta=\pi\) è instabile.
\end{itemize}

\bigskip

\section{Conclusioni della lezione}

In questa lezione abbiamo chiarito alcuni concetti fondamentali per lo studio qualitativo dei sistemi dinamici.

\begin{enumerate}
    \item Il concetto di \textbf{equilibrio} è il caso particolare di una traiettoria costante.
    \item La stabilità può essere definita inizialmente per un \textbf{movimento nominale} e poi specializzata al caso dell'equilibrio.
    \item Un equilibrio si studia più comodamente traslandolo nell'origine.
    \item Nei sistemi lineari autonomi, l'insieme degli equilibri è
    \[
    \ker A \quad \text{(TC)},
    \qquad
    \ker(A-I) \quad \text{(TD)}.
    \]
    \item Un equilibrio non isolato non può essere asintoticamente stabile.
    \item La \textbf{RAS} è l'insieme delle condizioni iniziali che convergono all'equilibrio.
    \item Per i sistemi non lineari, la \textbf{linearizzazione} permette di dedurre proprietà locali di stabilità.
    \item Il \textbf{teorema indiretto di Lyapunov} consente di concludere:
    \begin{itemize}
        \item stabilità asintotica locale se il linearizzato è asintoticamente stabile;
        \item instabilità se il linearizzato ha almeno un modo divergente;
        \item nessuna conclusione nel caso critico.
    \end{itemize}
\end{enumerate}

\medskip

In sostanza, la linearizzazione è il ponte che collega lo studio locale dei sistemi non lineari alla teoria, molto più semplice e potente, dei sistemi lineari.
% =========================
% L21.tex
% =========================

\chapter{L21 -- Funzioni definite in segno e teorema diretto di Lyapunov}
\label{chap:L21}

\section{Richiamo: funzioni definite in segno}

In questa lezione introduciamo uno degli strumenti più importanti per lo studio della stabilità dei sistemi non lineari: il \emph{metodo diretto di Lyapunov}.  
Prima, però, serve chiarire il significato di \emph{funzione definita in segno}.

Sia
\[
V:\mathbb{R}^n \to \mathbb{R},
\]
continua in un intorno dell'origine. Indichiamo con
\[
B_r=\{x\in\mathbb{R}^n:\|x\|<r\}
\]
la palla aperta di raggio \(r\) centrata nell'origine, e con
\[
\overline{B}_r=\{x\in\mathbb{R}^n:\|x\|\le r\}
\]
la corrispondente palla chiusa.

\subsection{Definizioni}

Diremo che \(V\) è:

\begin{itemize}
    \item \textbf{positiva definita} in \(B_r\) se
    \[
    V(0)=0,
    \qquad
    V(x)>0 \quad \forall x\in B_r\setminus\{0\};
    \]

    \item \textbf{positiva semidefinita} in \(B_r\) se
    \[
    V(0)=0,
    \qquad
    V(x)\ge 0 \quad \forall x\in B_r;
    \]

    \item \textbf{negativa definita} in \(B_r\) se \(-V\) è positiva definita, cioè
    \[
    V(0)=0,
    \qquad
    V(x)<0 \quad \forall x\in B_r\setminus\{0\};
    \]

    \item \textbf{negativa semidefinita} in \(B_r\) se \(-V\) è positiva semidefinita, cioè
    \[
    V(0)=0,
    \qquad
    V(x)\le 0 \quad \forall x\in B_r.
    \]
\end{itemize}

Se una funzione assume sia valori positivi sia valori negativi arbitrariamente vicini all'origine, allora si dice \textbf{indefinita} o \textbf{non definita in segno}.

\medskip

\noindent
L'idea intuitiva è semplice:
\begin{itemize}
    \item una funzione positiva definita ``misura una distanza energetica'' dall'origine;
    \item una funzione negativa definita è il suo opposto;
    \item una funzione semidefinita può annullarsi anche fuori dall'origine;
    \item una funzione indefinita cambia segno vicino all'origine e quindi non può essere usata direttamente come misura di energia.
\end{itemize}

\section{Esempi in \(\mathbb{R}^2\)}

\subsection*{Esempio 1: funzione positiva definita}
Consideriamo
\[
V(x_1,x_2)=x_1^2+x_2^2.
\]
Allora
\[
V(0,0)=0,
\qquad
V(x_1,x_2)>0 \quad \forall (x_1,x_2)\neq (0,0).
\]
Quindi \(V\) è positiva definita in tutto \(\mathbb{R}^2\).

\begin{figure}[h!]
\centering
\begin{tikzpicture}[scale=0.95]
    \draw[->] (-2.2,0) -- (2.3,0) node[right] {$x_1$};
    \draw[->] (0,-2.2) -- (0,2.3) node[above] {$x_2$};
    \draw (0,0) circle (0.55);
    \draw (0,0) circle (1.05);
    \draw (0,0) circle (1.55);
    \node at (0.45,0.45) {\small \(V=1\)};
    \node at (0.8,0.9) {\small \(V=2\)};
    \node at (1.1,1.25) {\small \(V=3\)};
\end{tikzpicture}
\caption{Curve di livello di \(V(x_1,x_2)=x_1^2+x_2^2\): sono circonferenze concentriche.}
\end{figure}

\subsection*{Esempio 2: funzione negativa definita in un intorno dell'origine}
Consideriamo
\[
V(x_1,x_2)=\bigl(x_1^2+x_2^2\bigr)\bigl(x_1^2+x_2^2-2\bigr).
\]
Se ci restringiamo a una palla sufficientemente piccola, ad esempio \(\overline{B}_1\), allora:
\[
x_1^2+x_2^2 \ge 0,
\qquad
x_1^2+x_2^2-2 \le -1.
\]
Di conseguenza, per ogni \(x\neq 0\) con \(\|x\|\le 1\),
\[
V(x)<0.
\]
Inoltre \(V(0)=0\). Quindi \(V\) è negativa definita in \(\overline{B}_1\).

\begin{figure}[h!]
\centering
\begin{tikzpicture}[scale=1.0]
    \draw[->] (-2.2,0) -- (2.3,0) node[right] {$x_1$};
    \draw[->] (0,-2.2) -- (0,2.3) node[above] {$x_2$};
    \fill[gray!20] (0,0) circle (1.0);
    \draw[thick] (0,0) circle (1.0);
    \draw[red,thick] (0,0) circle (1.4142);
    \node at (0.65,0.35) {\small \(B_1\)};
    \node[red] at (1.55,0.25) {\small \(\sqrt{2}\)};
\end{tikzpicture}
\caption{In \(B_1\) il fattore \(x_1^2+x_2^2-2\) è sempre negativo, quindi \(V<0\) per ogni \(x\neq 0\).}
\end{figure}

\subsection*{Esempio 3: funzione negativa semidefinita}
Consideriamo
\[
V(x_1,x_2)=-x_2^2.
\]
Allora
\[
V(x_1,x_2)\le 0 \quad \forall (x_1,x_2),
\]
ma \(V=0\) per tutti i punti della forma \((x_1,0)\), non solo nell'origine.  
Quindi \(V\) è negativa semidefinita, ma non negativa definita.

\begin{figure}[h!]
\centering
\begin{tikzpicture}[scale=0.95]
    \draw[->] (-2.2,0) -- (2.3,0) node[right] {$x_1$};
    \draw[->] (0,-2.2) -- (0,2.3) node[above] {$x_2$};
    \draw[very thick] (-2,0) -- (2,0);
    \node[above right] at (0.15,0) {\small \(V=0\)};
\end{tikzpicture}
\caption{Per \(V(x_1,x_2)=-x_2^2\), la funzione si annulla su tutto l'asse \(x_1\).}
\end{figure}

\subsection*{Esempio 4: funzione indefinita}
Consideriamo
\[
V(x_1,x_2)=x_1^2-x_2^2.
\]
Se \(x_2=0\), allora \(V=x_1^2\ge 0\).  
Se \(x_1=0\), allora \(V=-x_2^2\le 0\).  
Quindi vicino all'origine la funzione assume sia valori positivi sia negativi: non è definita in segno.

\begin{figure}[h!]
\centering
\begin{tikzpicture}[scale=0.95]
    \draw[->] (-2.2,0) -- (2.3,0) node[right] {$x_1$};
    \draw[->] (0,-2.2) -- (0,2.3) node[above] {$x_2$};

    \draw[red,thick] (-1.6,-1.6) -- (1.6,1.6);
    \draw[red,thick] (-1.6,1.6) -- (1.6,-1.6);

    \node at (1.45,0.55) {\small \(V>0\)};
    \node at (0.65,1.45) {\small \(V<0\)};
\end{tikzpicture}
\caption{Per \(V(x_1,x_2)=x_1^2-x_2^2\) il segno cambia a seconda della direzione.}
\end{figure}

\section{Forme quadratiche}

Un caso molto importante è quello in cui
\[
V(x)=x^\top Qx,
\]
dove \(Q\in\mathbb{R}^{n\times n}\).

\subsection{Riduzione alla parte simmetrica}
Ogni matrice \(Q\) si può decomporre come somma della sua parte simmetrica e della sua parte antisimmetrica:
\[
Q=Q_s+Q_a,
\]
dove
\[
Q_s=\frac{Q+Q^\top}{2},
\qquad
Q_a=\frac{Q-Q^\top}{2}.
\]
Vale
\[
Q_s^\top=Q_s,
\qquad
Q_a^\top=-Q_a.
\]

Ora osserviamo che, per ogni \(x\in\mathbb{R}^n\),
\[
x^\top Q_a x = 0.
\]
Infatti \(x^\top Q_a x\) è uno scalare, quindi coincide con la sua trasposta:
\[
x^\top Q_a x = (x^\top Q_a x)^\top = x^\top Q_a^\top x = x^\top(-Q_a)x = -x^\top Q_a x,
\]
da cui segue immediatamente
\[
x^\top Q_a x = 0.
\]

Pertanto
\[
V(x)=x^\top Qx=x^\top Q_s x.
\]
Quindi, nello studio del segno di una forma quadratica, conta solo la \textbf{parte simmetrica} della matrice.

\subsection{Caso simmetrico}
Se \(P=P^\top\), allora:
\[
V(x)=x^\top P x
\]
è positiva definita se e solo se tutti gli autovalori di \(P\) sono positivi.

Analogamente:
\begin{itemize}
    \item \(V\) è positiva semidefinita se gli autovalori di \(P\) sono non negativi;
    \item \(V\) è negativa definita se gli autovalori di \(P\) sono tutti negativi;
    \item \(V\) è indefinita se \(P\) ha autovalori di segno diverso.
\end{itemize}

\section{Idea del metodo diretto di Lyapunov}

Consideriamo un sistema non lineare autonomo tempo continuo
\[
\dot x = f(x),
\qquad
f(0)=0.
\]
Vogliamo capire se l'origine sia stabile, asintoticamente stabile oppure instabile, \emph{senza} risolvere esplicitamente il sistema.

L'idea di Lyapunov è costruire una funzione
\[
V:\mathbb{R}^n\to\mathbb{R}
\]
che giochi il ruolo di una specie di ``energia generalizzata''.

Se \(V\) è positiva definita e, lungo le traiettorie del sistema, il suo valore decresce, allora intuitivamente lo stato si muove verso l'origine.

\section{Derivata di \(V\) lungo le traiettorie}

\subsection{Caso continuo}
Se
\[
\dot x = f(x),
\]
allora la derivata temporale della funzione \(V(x(t))\) lungo una traiettoria \(x(t)\) è
\[
\frac{d}{dt}V(x(t))
=
\frac{\partial V}{\partial x}(x(t))\,\dot x(t)
=
\frac{\partial V}{\partial x}(x(t))\,f(x(t)).
\]
Questa quantità si indica con
\[
L_fV(x)=\nabla V(x)^\top f(x),
\]
e si chiama \textbf{derivata di Lie} di \(V\) rispetto al campo vettoriale \(f\), oppure più semplicemente \textbf{derivata di \(V\) lungo le traiettorie}.

\subsection{Caso discreto}
Se invece il sistema è tempo discreto
\[
x(k+1)=f(x(k)),
\]
si considera la variazione
\[
\Delta V(x(k)) = V(x(k+1))-V(x(k)).
\]

In entrambi i casi, il significato è lo stesso: misurare come cambia \(V\) mentre il sistema evolve.

\section{Interpretazione geometrica}

Se \(V\) è positiva definita, le sue curve di livello
\[
V(x)=c
\]
sono, intuitivamente, come ``linee di quota energetica'' attorno all'origine.

Se lungo le traiettorie vale
\[
L_fV(x)<0
\qquad\text{oppure}\qquad
\Delta V(x)<0,
\]
allora il sistema passa a livelli di energia sempre più bassi. Questo suggerisce che la traiettoria si sta avvicinando all'origine.

\begin{figure}[h!]
\centering
\begin{tikzpicture}[scale=1.0]
    % sinistra: superficie
    \begin{scope}
        \draw[->] (-1.8,0) -- (2.1,0) node[right] {$x_1$};
        \draw[->] (0,-1.6) -- (0,2.2) node[above] {$x_2$};
        \draw[->] (-1.0,-1.0) -- (-1.8,-1.8) node[left] {$V$};

        \draw (0,1.7) .. controls (0.8,1.6) and (1.2,1.0) .. (1.15,0.2);
        \draw (0,1.7) .. controls (-0.8,1.6) and (-1.2,1.0) .. (-1.15,0.2);
        \draw (-1.15,0.2) .. controls (-0.7,-0.15) and (0.7,-0.15) .. (1.15,0.2);
        \draw[dashed] (-0.75,0.35) .. controls (-0.45,0.15) and (0.45,0.15) .. (0.75,0.35);
        \node at (0.25,-0.45) {\small \(V=x_1^2+x_2^2\)};
    \end{scope}

    % destra: curve livello
    \begin{scope}[xshift=8cm]
        \draw[->] (-2.0,0) -- (2.2,0) node[right] {$x_1$};
        \draw[->] (0,-2.0) -- (0,2.2) node[above] {$x_2$};
        \draw (0,0) circle (0.45);
        \draw (0,0) circle (0.9);
        \draw (0,0) circle (1.35);
        \draw[->,thick] (1.3,0.2) -- (0.95,0.15);
        \draw[->,thick] (-1.15,-0.35) -- (-0.8,-0.25);
        \node at (0.55,0.45) {\small \(V=1\)};
        \node at (0.95,0.95) {\small \(V=2\)};
        \node at (1.35,1.35) {\small \(V=3\)};
    \end{scope}
\end{tikzpicture}
\caption{Interpretazione geometrica: se \(V\) decresce lungo le traiettorie, il sistema tende a muoversi verso curve di livello sempre più interne.}
\end{figure}

\section{Traslazione di un equilibrio non nullo}

Se il sistema ha un equilibrio in \(\bar x\neq 0\), non cambia nulla di concettuale: basta traslare le coordinate.

Poniamo
\[
z=x-\bar x.
\]
Allora l'equilibrio \(\bar x\) nel sistema originale diventa l'origine nel sistema traslato.

Perciò il metodo di Lyapunov si applica sempre, in sostanza, allo studio della stabilità dell'origine del sistema traslato.

\section{Insiemi di livello}

Per la dimostrazione del teorema diretto conviene introdurre alcuni insiemi standard.

Fissato un raggio \(\varepsilon>0\), indichiamo con
\[
B_\varepsilon=\{x\in\mathbb{R}^n:\|x\|<\varepsilon\},
\qquad
\partial B_\varepsilon=\{x\in\mathbb{R}^n:\|x\|=\varepsilon\}.
\]

Sia ora \(V\) una funzione positiva definita. Per un livello \(\rho>0\), definiamo:
\[
\Omega_\rho=\{x\in B_\varepsilon : V(x)\le \rho\},
\]
e la corrispondente curva/superficie di livello
\[
C_\rho=\partial\Omega_\rho
=
\{x\in B_\varepsilon : V(x)=\rho\}.
\]

\medskip

\noindent
L'idea è che, scegliendo \(\rho\) piccolo, l'insieme \(\Omega_\rho\) sia un intorno dell'origine contenuto in \(B_\varepsilon\).

\section{Teorema diretto di Lyapunov}

Consideriamo il sistema autonomo
\[
\dot x=f(x),
\qquad
f(0)=0,
\]
oppure, nel caso discreto,
\[
x(k+1)=f(x(k)),
\qquad
f(0)=0.
\]

Sia \(V\) una funzione positiva definita in un intorno dell'origine.

\subsection*{Enunciato}

\begin{enumerate}
    \item Se
    \[
    L_fV(x)\le 0
    \qquad
    \text{(oppure, nel discreto, } \Delta V(x)\le 0\text{)},
    \]
    allora l'origine è \textbf{stabile}.

    \item Se
    \[
    L_fV(x)<0
    \qquad
    \text{(oppure } \Delta V(x)<0\text{)}
    \]
    per ogni \(x\neq 0\), allora l'origine è \textbf{asintoticamente stabile}.

    \item Se
    \[
    L_fV(x)>0
    \qquad
    \text{(oppure } \Delta V(x)>0\text{)}
    \]
    per ogni \(x\neq 0\), allora l'origine è \textbf{instabile}.
\end{enumerate}

\medskip

\noindent
In questa lezione i passaggi di dimostrazione sono sviluppati soprattutto per il caso continuo. Il caso discreto è del tutto analogo, sostituendo l'integrale con somme finite e \(L_fV\) con \(\Delta V\).

\section{Lemma geometrico preliminare}

Prima della dimostrazione del teorema serve un fatto semplice ma fondamentale.

Sia \(V\) positiva definita e continua in \(\overline{B}_\varepsilon\). Allora:
\begin{enumerate}
    \item la funzione \(V\) ammette minimo su \(\partial B_\varepsilon\), perché \(\partial B_\varepsilon\) è compatto;
    \item tale minimo è strettamente positivo, perché su \(\partial B_\varepsilon\) non compare l'origine e \(V\) è positiva definita.
\end{enumerate}

Dunque, posto
\[
\alpha = \min_{x\in \partial B_\varepsilon} V(x),
\]
si ha
\[
\alpha>0.
\]

Ora scegliamo un numero \(\beta\) tale che
\[
0<\beta<\alpha.
\]
Allora il sottoinsieme di livello
\[
\Omega_\beta=\{x\in B_\varepsilon:V(x)\le \beta\}
\]
è contenuto strettamente dentro \(B_\varepsilon\), cioè
\[
\Omega_\beta \subset B_\varepsilon.
\]

Inoltre, essendo \(V(0)=0\) e \(V\) continua, esiste \(\delta>0\) tale che
\[
x\in B_\delta \quad \Longrightarrow \quad V(x)<\beta.
\]
Quindi
\[
B_\delta\subset \Omega_\beta.
\]

\medskip

\noindent
Questo lemma è esattamente il ponte che collega i livelli della funzione \(V\) con gli intorni metrici dell'origine.

\section{Dimostrazione del punto 1: stabilità}

Supponiamo che:
\[
V \text{ sia positiva definita},
\qquad
L_fV(x)\le 0.
\]

Vogliamo dimostrare che l'origine è stabile.

\subsection*{Passo 1: fissato un intorno finale, costruiamo un intorno iniziale}
Sia assegnato \(\varepsilon>0\).  
Per il lemma precedente, esiste \(\beta>0\) tale che
\[
B_\delta \subset \Omega_\beta \subset B_\varepsilon
\]
per un opportuno \(\delta>0\).

\subsection*{Passo 2: la traiettoria non può uscire da \(\Omega_\beta\)}
Prendiamo ora una traiettoria con condizione iniziale
\[
x(t_0)=x_0\in B_\delta.
\]
Allora, per costruzione,
\[
V(x_0)\le \beta.
\]

Poiché \(L_fV\le 0\), la funzione \(t\mapsto V(x(t))\) è non crescente. Dunque
\[
V(x(t))\le V(x_0)\le \beta
\qquad
\forall t\ge t_0.
\]
Quindi
\[
x(t)\in \Omega_\beta
\qquad
\forall t\ge t_0.
\]
Ma \(\Omega_\beta \subset B_\varepsilon\), perciò
\[
x(t)\in B_\varepsilon
\qquad
\forall t\ge t_0.
\]

\subsection*{Conclusione}
Abbiamo dimostrato che:
\[
\forall \varepsilon>0 \ \exists \delta>0
\quad \text{tale che} \quad
\|x_0\|<\delta
\Rightarrow
\|x(t)\|<\varepsilon \ \forall t\ge t_0.
\]
Questa è esattamente la definizione di stabilità.

\begin{figure}[h!]
\centering
\begin{tikzpicture}[scale=1.0]
    \draw[->] (-2.2,0) -- (2.3,0) node[right] {$x_1$};
    \draw[->] (0,-2.2) -- (0,2.3) node[above] {$x_2$};

    \draw[red,thick] (0,0) circle (1.7);
    \draw[blue,thick] (0,0) circle (1.0);
    \fill[blue!20] (0,0) circle (1.0);

    \node[red] at (1.95,0.2) {\small \(B_\varepsilon\)};
    \node[blue] at (1.15,-0.2) {\small \(\Omega_\beta\)};
\end{tikzpicture}
\caption{Se \(V\) non cresce, una traiettoria che parte in \(\Omega_\beta\) non può uscirne.}
\end{figure}

\section{Dimostrazione del punto 2: stabilità asintotica}

Supponiamo ora che:
\[
V \text{ sia positiva definita},
\qquad
L_fV(x)<0 \quad \forall x\neq 0.
\]

Dal punto 1 sappiamo già che l'origine è stabile. Resta da dimostrare che è anche attrattiva, cioè
\[
x(t)\to 0 \qquad \text{per } t\to +\infty.
\]

\subsection*{Passo 1: \(V(x(t))\) è monotona e limitata inferiormente}
Per una traiettoria iniziale sufficientemente vicina all'origine, grazie alla stabilità essa resta in un insieme compatto \(\Omega_\beta\).  
La funzione \(V(x(t))\) è decrescente e soddisfa
\[
V(x(t))\ge 0.
\]
Dunque ammette limite:
\[
\lim_{t\to +\infty}V(x(t))=c
\qquad \text{con } c\ge 0.
\]

\subsection*{Passo 2: mostriamo che \(c=0\)}
Procediamo per assurdo. Supponiamo
\[
c>0.
\]
Scegliamo allora un valore intermedio
\[
0<d<c.
\]
Consideriamo l'anello compatto
\[
K=\Omega_\beta \setminus \operatorname{int}(\Omega_d).
\]
Su questo insieme vale:
\begin{itemize}
    \item \(K\) è compatto;
    \item \(0\notin K\);
    \item \(L_fV(x)<0\) per ogni \(x\in K\).
\end{itemize}
Per continuità, \(L_fV\) raggiunge massimo su \(K\); tale massimo è negativo, quindi esiste \(\gamma>0\) tale che
\[
L_fV(x)\le -\gamma
\qquad \forall x\in K.
\]

Poiché \(V(x(t))\to c>d\), per tempi sufficientemente grandi la traiettoria resta in \(K\). Quindi, per \(t\) abbastanza grande,
\[
\frac{d}{dt}V(x(t)) \le -\gamma.
\]
Integrando,
\[
V(x(t)) \le V(x(T))-\gamma(t-T),
\]
e facendo \(t\to+\infty\) si ottiene
\[
V(x(t))\to -\infty,
\]
che è impossibile perché \(V(x)\ge 0\).

Quindi l'assurdo mostra che
\[
c=0.
\]

\subsection*{Passo 3: da \(V(x(t))\to 0\) segue \(x(t)\to 0\)}
Fissiamo un qualunque \(\eta>0\).  
Poiché \(V\) è positiva definita, sulla sfera \(\partial B_\eta\) vale
\[
m_\eta := \min_{x\in \partial B_\eta}V(x) >0.
\]
Se per infiniti tempi fosse \(\|x(t)\|\ge \eta\), allora dovremmo avere
\[
V(x(t))\ge m_\eta >0,
\]
in contraddizione con \(V(x(t))\to 0\). Dunque
\[
x(t)\to 0.
\]

\subsection*{Conclusione}
L'origine è asintoticamente stabile.

\section{Dimostrazione del punto 3: instabilità}

Supponiamo infine che:
\[
V \text{ sia positiva definita},
\qquad
L_fV(x)>0 \quad \forall x\neq 0.
\]

Vogliamo dimostrare che l'origine è instabile.

\subsection*{Idea}
Se \(V\) cresce lungo le traiettorie, allora una traiettoria che parte molto vicina all'origine tende ad allontanarsi da livelli bassi di \(V\).  
Il ragionamento è speculare a quello usato per la stabilità.

\subsection*{Costruzione}
Sia fissato un piccolo \(\varepsilon>0\).  
Scegliamo un livello \(c>0\) tale che
\[
\Omega_c \subset B_\varepsilon.
\]
Prendiamo poi un livello \(a\) con
\[
0<a<c.
\]
La superficie
\[
C_a=\{x\in B_\varepsilon:V(x)=a\}
\]
può essere scelta arbitrariamente vicina all'origine scegliendo \(a\) piccolo.

Consideriamo l'anello compatto
\[
K=\Omega_c \setminus \operatorname{int}(\Omega_a).
\]
Su \(K\), la funzione \(L_fV\) è continua e strettamente positiva. Dunque esiste \(\gamma>0\) tale che
\[
L_fV(x)\ge \gamma
\qquad \forall x\in K.
\]

Se una traiettoria parte da un punto \(x_0\in C_a\) e rimanesse per sempre in \(\Omega_c\), allora finché resta in \(K\) si avrebbe
\[
\frac{d}{dt}V(x(t))\ge \gamma,
\]
da cui
\[
V(x(t))\ge a+\gamma(t-t_0).
\]
Quindi, in tempo finito, si arriverebbe a
\[
V(x(t))>c,
\]
il che contraddice l'ipotesi di permanenza in \(\Omega_c\).

Pertanto la traiettoria esce da \(\Omega_c\), e quindi anche da \(B_\varepsilon\).

\subsection*{Conclusione}
Esistono condizioni iniziali arbitrariamente vicine all'origine le cui traiettorie escono da un intorno prefissato dell'origine.  
Quindi l'origine è instabile.

\section{Osservazione importante}

Il teorema diretto di Lyapunov è estremamente potente perché consente di concludere la stabilità \emph{senza risolvere il sistema}.  
Il problema vero, nella pratica, è trovare una buona funzione \(V\).

In sintesi:
\begin{itemize}
    \item se troviamo \(V>0\) e \(L_fV\le 0\), otteniamo stabilità;
    \item se troviamo \(V>0\) e \(L_fV<0\), otteniamo stabilità asintotica;
    \item se troviamo \(V>0\) e \(L_fV>0\), otteniamo instabilità.
\end{itemize}

\section{Interpretazione energetica}

Il metodo di Lyapunov generalizza l'idea fisica di energia.

Se il sistema possiede una funzione energia \(E(x)\) che:
\begin{itemize}
    \item è minima nell'equilibrio;
    \item non aumenta lungo le traiettorie;
\end{itemize}
allora il sistema non può allontanarsi troppo dall'equilibrio.

Se addirittura tale energia diminuisce strettamente, allora il sistema dissipa energia e tende all'equilibrio.

\medskip

\noindent
Per questo motivo, nelle applicazioni meccaniche, spesso la funzione di Lyapunov si costruisce proprio a partire da energia cinetica e potenziale.

\section{Conclusioni della lezione}

In questa lezione abbiamo introdotto i concetti che stanno alla base della teoria diretta di Lyapunov:

\begin{enumerate}
    \item funzioni positive definite, semidefinite, negative definite e indefinite;
    \item forme quadratiche e loro riduzione alla sola parte simmetrica;
    \item derivata di una funzione lungo le traiettorie:
    \[
    L_fV(x)=\nabla V(x)^\top f(x)
    \quad \text{nel continuo},
    \qquad
    \Delta V(x)=V(x^+)-V(x)
    \quad \text{nel discreto};
    \]
    \item insiemi di livello e loro uso nella dimostrazione;
    \item teorema diretto di Lyapunov:
    \begin{itemize}
        \item \(L_fV\le 0 \Rightarrow\) stabilità,
        \item \(L_fV<0 \Rightarrow\) stabilità asintotica,
        \item \(L_fV>0 \Rightarrow\) instabilità.
    \end{itemize}
\end{enumerate}

Il cuore del metodo è questo: invece di cercare la soluzione del sistema, costruiamo una quantità scalare \(V\) che ci dica se il sistema tende ad avvicinarsi o ad allontanarsi dall'equilibrio.

% =========================
% L22.tex
% =========================

\chapter{L22 -- Funzioni di Lyapunov, stima della RAS e primi esempi}
\label{chap:L22}

\section{Definizione di funzione di Lyapunov}

Nella lezione precedente abbiamo visto il teorema diretto di Lyapunov: se esiste una funzione
\[
V:\mathbb{R}^n \to \mathbb{R}
\]
positiva definita e la sua derivata lungo le traiettorie è non positiva, allora l'origine è stabile; se tale derivata è strettamente negativa fuori dall'origine, allora l'origine è asintoticamente stabile.

Da qui nasce la definizione seguente.

\medskip

\noindent
\textbf{Definizione.}
Consideriamo un sistema autonomo con equilibrio nell'origine:
\[
\dot x = f(x), \qquad f(0)=0,
\]
oppure, nel caso discreto,
\[
x^{+}=f(x), \qquad f(0)=0.
\]
Una funzione \(V\) si dice \textbf{funzione di Lyapunov} per il sistema se:
\begin{itemize}
    \item \(V(0)=0\);
    \item \(V(x)>0\) per ogni \(x\neq 0\) in un intorno dell'origine, cioè \(V\) è positiva definita;
    \item la derivata lungo le traiettorie soddisfa
    \[
    L_fV(x)\le 0
    \]
    nel caso continuo, oppure
    \[
    \Delta V(x)=V(f(x))-V(x)\le 0
    \]
    nel caso discreto.
\end{itemize}

Se addirittura vale
\[
L_fV(x)<0 \quad \forall x\neq 0
\]
oppure
\[
\Delta V(x)<0 \quad \forall x\neq 0,
\]
si parla spesso di \textbf{funzione di Lyapunov stretta}.

\medskip

\noindent
In altre parole:
\begin{itemize}
    \item \(V\) deve comportarsi come una specie di energia;
    \item lungo i movimenti del sistema questa energia non deve crescere;
    \item se decresce strettamente, il sistema tende effettivamente all'equilibrio.
\end{itemize}

\section{Primo esempio: \(\dot x=-x^3\)}

Consideriamo il sistema scalare
\[
\dot x=-x^3, \qquad x\in\mathbb{R}.
\]

\subsection{Equilibrio e intuizione qualitativa}

Gli equilibri si trovano imponendo \(\dot x=0\):
\[
-x^3=0 \quad \Longrightarrow \quad x=0.
\]
Quindi l'unico equilibrio è l'origine.

Dal segno del campo vettoriale:
\[
x>0 \Longrightarrow \dot x<0,
\qquad
x<0 \Longrightarrow \dot x>0.
\]
Dunque il moto punta sempre verso l'origine.

\begin{figure}[h!]
\centering
\begin{tikzpicture}[scale=1.0]
    \draw[->] (-4,0) -- (4,0) node[right] {$x$};
    \fill (0,0) circle (1.5pt);
    \node[below] at (0,0) {$0$};

    \draw[-{Latex[length=3mm]},thick] (3,0.35) -- (2,0.35);
    \draw[-{Latex[length=3mm]},thick] (1.8,0.35) -- (0.8,0.35);

    \draw[-{Latex[length=3mm]},thick] (-3,0.35) -- (-2,0.35);
    \draw[-{Latex[length=3mm]},thick] (-1.8,0.35) -- (-0.8,0.35);
\end{tikzpicture}
\caption{Campo di direzione sulla retta per \(\dot x=-x^3\): tutti i moti puntano verso l'origine.}
\end{figure}

Questa sola osservazione suggerisce che l'origine sia asintoticamente stabile.

\subsection{Metodo indiretto}

Linearizziamo in \(x=0\). Poiché
\[
f(x)=-x^3,
\]
si ha
\[
A=\frac{df}{dx}=-3x^2.
\]
Valutando nell'equilibrio:
\[
A_0 = -3\cdot 0^2 = 0.
\]
Quindi il sistema linearizzato è
\[
\dot z = 0.
\]

Il metodo indiretto di Lyapunov qui non conclude nulla, perché l'unico autovalore è
\[
\lambda=0,
\]
cioè ha parte reale nulla.

\subsection{Metodo diretto di Lyapunov}

Proviamo con la funzione naturale
\[
V(x)=x^2.
\]
Essa è positiva definita. La derivata lungo le traiettorie è
\[
L_fV(x)=\frac{dV}{dx}f(x)=2x(-x^3)=-2x^4.
\]
Per ogni \(x\neq 0\),
\[
-2x^4<0.
\]
Quindi \(L_fV\) è negativa definita.

Per il teorema diretto di Lyapunov, l'origine è \textbf{asintoticamente stabile}.  
Questo conferma l'intuizione qualitativa iniziale e mostra bene una situazione tipica: il metodo indiretto è inconcludente, mentre il metodo diretto funziona perfettamente.

\section{Secondo esempio: \(\dot x=x^3\)}

Consideriamo ora
\[
\dot x=x^3.
\]

L'equilibrio è ancora
\[
x=0.
\]
Ma stavolta il campo vettoriale cambia completamente significato:
\[
x>0 \Longrightarrow \dot x>0,
\qquad
x<0 \Longrightarrow \dot x<0.
\]
Quindi il moto si allontana dall'origine da entrambi i lati.

\begin{figure}[h!]
\centering
\begin{tikzpicture}[scale=1.0]
    \draw[->] (-4,0) -- (4,0) node[right] {$x$};
    \fill (0,0) circle (1.5pt);
    \node[below] at (0,0) {$0$};

    \draw[-{Latex[length=3mm]},thick] (2,0.35) -- (3,0.35);
    \draw[-{Latex[length=3mm]},thick] (0.8,0.35) -- (1.8,0.35);

    \draw[-{Latex[length=3mm]},thick] (-2,0.35) -- (-3,0.35);
    \draw[-{Latex[length=3mm]},thick] (-0.8,0.35) -- (-1.8,0.35);
\end{tikzpicture}
\caption{Campo di direzione sulla retta per \(\dot x=x^3\): i moti si allontanano dall'origine.}
\end{figure}

Se scegliamo ancora
\[
V(x)=x^2,
\]
otteniamo
\[
L_fV(x)=2x(x^3)=2x^4.
\]
Per ogni \(x\neq 0\),
\[
2x^4>0.
\]
Dunque, per il criterio di Lyapunov, l'origine è \textbf{instabile}.

Questo esempio è il duale del precedente: stessa funzione \(V\), ma segno opposto di \(L_fV\), e quindi conclusione opposta.

\section{Terzo esempio piano: una dinamica non lineare con rotazione}

Consideriamo il sistema
\begin{equation}
\label{eq:L22-sistema-rotazione}
\begin{cases}
\dot x_1 = x_1\bigl(x_1^2+x_2^2-1\bigr)-x_2,\\[4pt]
\dot x_2 = x_1+x_2\bigl(x_1^2+x_2^2-1\bigr).
\end{cases}
\end{equation}

Questo sistema è interessante perché combina:
\begin{itemize}
    \item una parte radiale, governata da \(x_1^2+x_2^2-1\),
    \item una parte rotazionale, data dai termini \(-x_2\) e \(x_1\).
\end{itemize}

\subsection{Calcolo degli equilibri}

Poniamo \(\dot x_1=\dot x_2=0\):
\[
\begin{cases}
x_1(r^2-1)-x_2=0,\\
x_1+x_2(r^2-1)=0,
\end{cases}
\qquad r^2=x_1^2+x_2^2.
\]
Dalla prima ricaviamo
\[
x_2=x_1(r^2-1).
\]
Sostituendo nella seconda:
\[
x_1 + x_1(r^2-1)^2 = 0,
\]
cioè
\[
x_1\Bigl(1+(r^2-1)^2\Bigr)=0.
\]
Il fattore
\[
1+(r^2-1)^2
\]
è sempre strettamente positivo, quindi necessariamente
\[
x_1=0.
\]
Dalla prima equazione segue allora
\[
x_2=0.
\]
L'unico equilibrio è dunque
\[
(0,0).
\]

\subsection{Studio con il linearizzato}

Scriviamo il Jacobiano:
\[
A(x)=
\begin{pmatrix}
\frac{\partial f_1}{\partial x_1} & \frac{\partial f_1}{\partial x_2}\\[4pt]
\frac{\partial f_2}{\partial x_1} & \frac{\partial f_2}{\partial x_2}
\end{pmatrix}.
\]
Con
\[
f_1(x)=x_1(x_1^2+x_2^2-1)-x_2,
\qquad
f_2(x)=x_1+x_2(x_1^2+x_2^2-1),
\]
si ottiene
\[
A(x)=
\begin{pmatrix}
(x_1^2+x_2^2-1)+2x_1^2 & 2x_1x_2-1\\[4pt]
1+2x_1x_2 & (x_1^2+x_2^2-1)+2x_2^2
\end{pmatrix}.
\]
Valutando nell'origine:
\[
A(0,0)=
\begin{pmatrix}
-1 & -1\\
1 & -1
\end{pmatrix}.
\]
Il polinomio caratteristico è
\[
p(\lambda)=\det(\lambda I-A)=
\det\begin{pmatrix}
\lambda+1 & 1\\
-1 & \lambda+1
\end{pmatrix}
=(\lambda+1)^2+1.
\]
Quindi gli autovalori sono
\[
\lambda_{1,2}=-1\pm j.
\]
Hanno parte reale negativa, dunque per il metodo indiretto l'origine è \textbf{localmente asintoticamente stabile}.

\subsection{Studio con Lyapunov}

Proviamo con la funzione naturale
\[
V(x_1,x_2)=\frac{1}{2}(x_1^2+x_2^2).
\]
È positiva definita. La derivata lungo le traiettorie vale
\[
L_fV
=
\frac{\partial V}{\partial x_1}f_1+\frac{\partial V}{\partial x_2}f_2
=
x_1f_1+x_2f_2.
\]
Sostituendo:
\begin{align*}
L_fV
&=
x_1\Bigl(x_1(x_1^2+x_2^2-1)-x_2\Bigr)
+
x_2\Bigl(x_1+x_2(x_1^2+x_2^2-1)\Bigr)\\
&=
x_1^2(x_1^2+x_2^2-1)-x_1x_2+x_1x_2+x_2^2(x_1^2+x_2^2-1)\\
&=
(x_1^2+x_2^2)\bigl(x_1^2+x_2^2-1\bigr).
\end{align*}
Quindi
\[
L_fV(x)<0
\qquad \text{se} \qquad
0<x_1^2+x_2^2<1.
\]
In particolare, nel disco unitario aperto
\[
B_1=\{x\in\mathbb{R}^2:\|x\|<1\}
\]
la derivata è negativa definita.

Per il metodo diretto di Lyapunov, l'origine è asintoticamente stabile e, inoltre, otteniamo una stima della regione di attrazione:
\[
B_1 \subseteq \mathcal{R}_{AS}(0).
\]

\begin{figure}[h!]
\centering
\begin{tikzpicture}[scale=1.0]
    \draw[->] (-2.2,0) -- (2.2,0) node[right] {$x_1$};
    \draw[->] (0,-2.2) -- (0,2.2) node[above] {$x_2$};

    \draw[thick] (0,0) circle (1.5);
    \node at (1.2,1.2) {\small \(x_1^2+x_2^2=1\)};

    \draw[red,thick] (0,0) circle (0.35);
    \draw[red,thick] (0,0) circle (0.65);
    \draw[red,thick] (0,0) circle (1.0);
    \fill (0,0) circle (1.5pt);

    \node[red] at (0.5,0.2) {\small curve di livello di \(V\)};
\end{tikzpicture}
\caption{Per il sistema \eqref{eq:L22-sistema-rotazione}, all'interno del disco unitario \(L_fV<0\): le curve di livello di \(V\) forniscono una stima della regione di attrazione.}
\end{figure}

\subsection{Osservazione importante}

In questo esempio il teorema di Lyapunov fornisce una \emph{stima} della regione di attrazione.  
Non sempre la stima coincide con la regione di attrazione esatta.

In realtà, per questo sistema si può andare oltre introducendo coordinate polari:
\[
x_1=r\cos\theta,
\qquad
x_2=r\sin\theta.
\]
Si ottiene
\[
\dot r = r(r^2-1),
\qquad
\dot\theta = 1.
\]
Da qui si vede subito che:
\begin{itemize}
    \item se \(0<r<1\), allora \(\dot r<0\), quindi \(r(t)\to 0\);
    \item se \(r>1\), allora \(\dot r>0\), quindi il raggio cresce.
\end{itemize}
Quindi in questo caso
\[
\mathcal{R}_{AS}(0)=B_1.
\]
Il metodo di Lyapunov aveva già individuato correttamente tale regione.

\section{Come si stima la regione di attrazione}

L'esempio precedente suggerisce un principio generale molto utile.

Supponiamo di avere:
\begin{itemize}
    \item una funzione \(V\) positiva definita;
    \item un insieme definito da una curva di livello
    \[
    C_c=\{x:V(x)=c\};
    \]
    \item la proprietà
    \[
    L_fV(x)<0
    \quad \text{all'interno di } C_c.
    \]
\end{itemize}

Allora l'insieme
\[
\Omega_c=\{x:V(x)\le c\}
\]
è positivamente invariante e, se contiene l'origine ed è contenuto in una zona in cui valgono le ipotesi del teorema, fornisce una stima della regione di attrazione.

\medskip

\noindent
In pratica si procede così:
\begin{enumerate}
    \item si sceglie una funzione \(V\) semplice;
    \item si calcola \(L_fV\);
    \item si individua una curva di livello di \(V\) all'interno della quale \(L_fV<0\);
    \item l'interno di quella curva è una stima della RAS.
\end{enumerate}

\section{Condizione di globalità: crescita radiale di \(V\)}

Per ottenere \textbf{stabilità asintotica globale} non basta che \(V\) sia positiva definita e che \(L_fV<0\).  
Occorre anche che le curve di livello di \(V\) chiudano tutto lo spazio, cioè che \(V\) cresca indefinitamente quando \(\|x\|\to\infty\).

\medskip

\noindent
\textbf{Condizione tipica di globalità.}
Se:
\begin{itemize}
    \item \(V\) è positiva definita in tutto \(\mathbb{R}^n\),
    \item \(L_fV\) è negativa definita in tutto \(\mathbb{R}^n\),
    \item
    \[
    \lim_{\|x\|\to\infty} V(x)=+\infty,
    \]
\end{itemize}
allora l'origine è \textbf{globalmente asintoticamente stabile}.

L'ultima condizione si chiama spesso \textbf{radialmente illimitata} o \textbf{proper}.

\subsection{Esempio di funzione non radialmente illimitata}

Consideriamo
\[
V(x_1,x_2)=\frac{x_1^2}{1+x_1^2}+x_2^2.
\]
Si ha:
\[
V(0,0)=0,
\qquad
V(x_1,x_2)>0 \ \text{per } (x_1,x_2)\neq (0,0),
\]
quindi \(V\) è positiva definita.

Però non è radialmente illimitata. Infatti, lungo l'asse \(x_2=0\),
\[
V(x_1,0)=\frac{x_1^2}{1+x_1^2}\xrightarrow[|x_1|\to\infty]{}1.
\]
Quindi
\[
\lim_{\|x\|\to\infty}V(x)\neq +\infty.
\]

Questo significa che le curve di livello possono non chiudersi a infinito e quindi \(V\) non è adatta, da sola, a concludere la stabilità globale.

\begin{figure}[h!]
\centering
\begin{tikzpicture}[scale=0.95]
    \draw[->] (-3.5,0) -- (3.5,0) node[right] {$x_1$};
    \draw[->] (0,-3.0) -- (0,3.0) node[above] {$x_2$};

    \draw[blue,thick]
        plot[smooth cycle,tension=0.9]
        coordinates {(0,1.3) (0.8,1.05) (1.6,0.55) (2.8,0.2) (1.6,-0.55) (0.8,-1.05) (0,-1.3) (-0.8,-1.05) (-1.6,-0.55) (-2.8,0.2) (-1.6,0.55) (-0.8,1.05)};
    \draw[dashed,blue]
        plot[smooth cycle,tension=0.9]
        coordinates {(0,2.0) (1.0,1.55) (2.2,1.0) (3.2,0.55) (2.2,-1.0) (1.0,-1.55) (0,-2.0) (-1.0,-1.55) (-2.2,-1.0) (-3.2,0.55) (-2.2,1.0) (-1.0,1.55)};
\end{tikzpicture}
\caption{Schema qualitativo di curve di livello di una funzione positiva definita ma non radialmente illimitata.}
\end{figure}

\section{Esempio parametrico e limiti del metodo indiretto}

Consideriamo il sistema
\begin{equation}
\label{eq:L22-esempio-parametrico}
\begin{cases}
\dot x_1 = -x_1-(x_2-1)^2,\\[4pt]
\dot x_2 = k(x_2-1)+4x_1(x_2-1),
\end{cases}
\qquad k\in\mathbb{R}.
\end{equation}

\subsection{Traslazione delle coordinate}

Poiché compare naturalmente il termine \(x_2-1\), è comodo traslare il sistema ponendo
\[
z_1=x_1,
\qquad
z_2=x_2-1.
\]
Allora il sistema diventa
\begin{equation}
\label{eq:L22-z}
\begin{cases}
\dot z_1 = -z_1-z_2^2,\\[4pt]
\dot z_2 = kz_2+4z_1z_2.
\end{cases}
\end{equation}

\subsection{Equilibri}

Gli equilibri soddisfano
\[
\begin{cases}
0=-z_1-z_2^2,\\
0=z_2(k+4z_1).
\end{cases}
\]
Dalla seconda equazione distinguiamo due casi.

\paragraph{Caso 1: \(z_2=0\).}
Allora la prima equazione dà
\[
z_1=0.
\]
Quindi esiste sempre l'equilibrio
\[
(0,0).
\]

\paragraph{Caso 2: \(z_2\neq 0\).}
Allora deve valere
\[
k+4z_1=0
\quad \Longrightarrow \quad
z_1=-\frac{k}{4}.
\]
Sostituendo nella prima:
\[
0=-\left(-\frac{k}{4}\right)-z_2^2
=
\frac{k}{4}-z_2^2,
\]
cioè
\[
z_2^2=\frac{k}{4}.
\]
Questa equazione ammette soluzioni reali se e solo se \(k\ge 0\). In tal caso:
\[
z_2=\pm \frac{\sqrt{k}}{2}.
\]
Quindi, per \(k\ge 0\), compaiono anche due ulteriori equilibri:
\[
\left(-\frac{k}{4},\frac{\sqrt{k}}{2}\right),
\qquad
\left(-\frac{k}{4},-\frac{\sqrt{k}}{2}\right).
\]

Tornando alle coordinate originali:
\[
x_1=z_1,
\qquad
x_2=z_2+1,
\]
gli equilibri sono:
\[
(0,1)
\]
per ogni \(k\), e inoltre, se \(k\ge 0\),
\[
\left(-\frac{k}{4},\,1+\frac{\sqrt{k}}{2}\right),
\qquad
\left(-\frac{k}{4},\,1-\frac{\sqrt{k}}{2}\right).
\]

\subsection{Stabilità dell'equilibrio \((0,0)\) in coordinate \(z\)}

Il Jacobiano del sistema \eqref{eq:L22-z} è
\[
A(z)=
\begin{pmatrix}
-1 & -2z_2\\
4z_2 & k+4z_1
\end{pmatrix}.
\]
Valutando nell'origine:
\[
A(0,0)=
\begin{pmatrix}
-1 & 0\\
0 & k
\end{pmatrix}.
\]
Gli autovalori sono quindi
\[
\lambda_1=-1,
\qquad
\lambda_2=k.
\]

Da qui segue subito:
\begin{itemize}
    \item se \(k<0\), allora entrambi gli autovalori hanno parte reale negativa, quindi l'origine è \textbf{asintoticamente stabile};
    \item se \(k>0\), allora c'è un autovalore positivo, quindi l'origine è \textbf{instabile};
    \item se \(k=0\), il metodo indiretto è \textbf{inconcludente}.
\end{itemize}

\subsection{Caso critico \(k=0\): uso di Lyapunov}

Per \(k=0\) il sistema diventa
\[
\begin{cases}
\dot z_1 = -z_1-z_2^2,\\[4pt]
\dot z_2 = 4z_1z_2.
\end{cases}
\]
Il linearizzato in \((0,0)\) è
\[
A_0=
\begin{pmatrix}
-1 & 0\\
0 & 0
\end{pmatrix},
\]
quindi il metodo indiretto non basta.

Scegliamo la funzione
\[
V(z_1,z_2)=\frac{1}{2}\bigl(4z_1^2+z_2^2\bigr).
\]
Essa è positiva definita. La derivata lungo le traiettorie è
\[
L_fV = 4z_1\dot z_1 + z_2\dot z_2.
\]
Sostituendo:
\begin{align*}
L_fV
&=
4z_1(-z_1-z_2^2)+z_2(4z_1z_2)\\
&=
-4z_1^2-4z_1z_2^2+4z_1z_2^2\\
&=
-4z_1^2\le 0.
\end{align*}
Quindi \(L_fV\) è negativa semidefinita.

Con il solo teorema diretto classico concludiamo subito che l'origine è \textbf{stabile}.  
Per concludere l'asintotica stabilità, però, serve un passo in più.

\section{Richiamo al principio di invarianza di LaSalle}

Quando \(L_fV\le 0\) ma non è negativa definita, il teorema diretto di Lyapunov garantisce la stabilità, ma non sempre l'asintotica stabilità. In questi casi si guarda il sottoinsieme
\[
E=\{x : L_fV(x)=0\}.
\]

Nel caso in esame,
\[
L_fV(z_1,z_2)=-4z_1^2,
\]
quindi
\[
E=\{(z_1,z_2): z_1=0\}.
\]
Studiamo la dinamica su \(E\). Se \(z_1=0\), il sistema diventa
\[
\begin{cases}
\dot z_1=-z_2^2,\\
\dot z_2=0.
\end{cases}
\]
Affinché una traiettoria resti per sempre in \(E\), deve valere anche \(\dot z_1=0\), cioè
\[
-z_2^2=0 \quad \Longrightarrow \quad z_2=0.
\]
Dunque l'unico insieme invariante contenuto in \(E\) è il solo punto
\[
(0,0).
\]

Per il principio di LaSalle si conclude che, per \(k=0\), l'origine è \textbf{asintoticamente stabile}.

\medskip

\noindent
Riassumendo per l'equilibrio \((0,0)\) del sistema \eqref{eq:L22-z}:
\[
\boxed{
\begin{array}{ll}
k<0 & \Rightarrow \text{origine asintoticamente stabile},\\[4pt]
k=0 & \Rightarrow \text{origine asintoticamente stabile (via Lyapunov + LaSalle)},\\[4pt]
k>0 & \Rightarrow \text{origine instabile}.
\end{array}
}
\]

\section{Osservazione sul ritratto di fase nel caso \(k=0\)}

Nel caso \(k=0\), il sistema
\[
\begin{cases}
\dot z_1=-z_1-z_2^2,\\
\dot z_2=4z_1z_2
\end{cases}
\]
mostra bene perché la sola semidefinità negativa di \(L_fV\) non impedisce comunque la convergenza all'origine.

Infatti:
\begin{itemize}
    \item se \(z_1>0\), il termine \(-z_1\) tende a spingere \(z_1\) verso sinistra;
    \item il termine \(-z_2^2\) rende \(\dot z_1\) ancora più negativo;
    \item la dinamica di \(z_2\) dipende dal prodotto \(z_1z_2\), quindi cambia segno a seconda dei quadranti.
\end{itemize}

Le traiettorie vengono comunque progressivamente riportate verso l'origine.  
Il ritratto di fase suggerisce la stessa conclusione ottenuta rigorosamente con LaSalle.

\begin{figure}[h!]
\centering
\begin{tikzpicture}[scale=1.0]
    \draw[->] (-3.2,0) -- (3.2,0) node[right] {$z_1$};
    \draw[->] (0,-3.0) -- (0,3.0) node[above] {$z_2$};

    % nullcline z1 = -z2^2
    \draw[orange,thick,domain=-1.6:1.6,samples=100]
        plot ({-(\x*\x)}, {\x});
    \node[orange] at (-2.2,2.2) {\small \(z_1=-z_2^2\)};

    % z2=0
    \draw[blue,thick] (-3,0) -- (3,0);
    \node[blue] at (2.3,0.25) {\small \(z_2=0\)};

    \fill (0,0) circle (1.8pt);

    % qualitative arrows
    \draw[-{Latex[length=2.5mm]},red] (1.8,2.0) -- (1.25,1.8);
    \draw[-{Latex[length=2.5mm]},red] (1.8,-2.0) -- (1.25,-1.8);
    \draw[-{Latex[length=2.5mm]},red] (-2.0,2.0) -- (-2.2,1.4);
    \draw[-{Latex[length=2.5mm]},red] (-2.0,-2.0) -- (-2.2,-1.4);
    \draw[-{Latex[length=2.5mm]},red] (0.9,0.7) -- (0.25,0.45);
    \draw[-{Latex[length=2.5mm]},red] (0.9,-0.7) -- (0.25,-0.45);
\end{tikzpicture}
\caption{Schema qualitativo del ritratto di fase per il caso \(k=0\).}
\end{figure}

\section{Conclusioni operative della lezione}

Questa lezione chiarisce alcuni punti molto importanti.

\subsection*{1. Il metodo diretto può riuscire quando il linearizzato fallisce}
Negli esempi
\[
\dot x=-x^3
\qquad \text{e} \qquad
\dot x=x^3
\]
il linearizzato nell'origine è nullo e quindi non conclude nulla; il metodo diretto di Lyapunov invece distingue perfettamente il caso stabile da quello instabile.

\subsection*{2. Le curve di livello di \(V\) servono anche a stimare la RAS}
Non servono solo a dimostrare la stabilità locale, ma anche a costruire insiemi invarianti contenuti nella regione di attrazione.

\subsection*{3. Per la stabilità globale serve una condizione di crescita a infinito}
La proprietà
\[
\lim_{\|x\|\to\infty}V(x)=+\infty
\]
è fondamentale se si vuole passare da una conclusione locale a una conclusione globale.

\subsection*{4. Se \(L_fV\) è solo semidefinita, spesso bisogna usare LaSalle}
Il teorema diretto classico dà la stabilità.  
Per l'asintotica stabilità bisogna studiare l'insieme
\[
\{x : L_fV(x)=0\}
\]
e cercare il suo massimo insieme invariante.

\section{Schema riassuntivo}

Per un sistema autonomo con equilibrio nell'origine:
\[
\dot x=f(x),
\qquad
f(0)=0,
\]
si cerca una funzione \(V\) tale che:
\begin{itemize}
    \item \(V(x)>0\) per \(x\neq 0\);
    \item \(V(0)=0\);
    \item \(L_fV(x)\le 0\) oppure \(<0\).
\end{itemize}

Allora:
\[
\boxed{
\begin{array}{ll}
V \text{ p.d.},\ L_fV\le 0
&\Rightarrow \text{origine stabile},\\[4pt]
V \text{ p.d.},\ L_fV<0
&\Rightarrow \text{origine asintoticamente stabile},\\[4pt]
V \text{ p.d.},\ L_fV>0
&\Rightarrow \text{origine instabile}.
\end{array}
}
\]

Se inoltre
\[
\lim_{\|x\|\to\infty}V(x)=+\infty,
\]
si può spesso concludere la \textbf{stabilità asintotica globale}.

\medskip

\noindent
Questa è l'idea di fondo che useremo anche nelle lezioni successive: costruire una funzione scalare semplice che catturi il comportamento qualitativo di un sistema non lineare, anche quando risolverlo esplicitamente è impossibile.
% =========================
% L23.tex
% =========================

\chapter{L23 -- Principio di invarianza di LaSalle, esercizi ed elementi di sintesi con Lyapunov}
\label{chap:L23}

\section{Insiemi invarianti}

Prima di enunciare il principio di LaSalle conviene richiamare con precisione il concetto di insieme invariante.

\medskip

\noindent
\textbf{Definizione.}
Sia dato un sistema dinamico autonomo.

\begin{itemize}
    \item Nel caso continuo:
    \[
    \dot x = f(x), \qquad x\in\mathbb{R}^n.
    \]
    \item Nel caso discreto:
    \[
    x(k+1)=f(x(k)).
    \]
\end{itemize}

Un insieme \(M\subseteq \mathbb{R}^n\) si dice \textbf{positivamente invariante} se ogni traiettoria che parte in \(M\) rimane in \(M\) per tutti i tempi futuri.

Nel caso continuo questo significa:
\[
x(t_0)\in M
\quad \Longrightarrow \quad
\varphi(t,t_0,x(t_0))\in M
\qquad \forall t\ge t_0.
\]

Nel caso discreto:
\[
x(0)\in M
\quad \Longrightarrow \quad
x(k)\in M
\qquad \forall k\ge 0.
\]

\medskip

\noindent
In modo intuitivo: un insieme è invariante se, una volta entrati dentro, non se ne esce più.

\begin{figure}[h!]
\centering
\begin{tikzpicture}[scale=1]
    \draw[thick] (0,0) ellipse (3 and 2);
    \node at (0,2.35) {$M$};

    \fill (-1,0.3) circle (1.5pt);
    \fill (1.2,-0.5) circle (1.5pt);

    \draw[-{Latex[length=3mm]},thick] (-1,0.3) .. controls (-0.2,0.8) and (0.4,0.6) .. (1.2,-0.5);
    \draw[-{Latex[length=3mm]},thick] (1.2,-0.5) .. controls (1.8,-0.3) and (2.0,0.4) .. (1.4,1.0);

    \draw[-{Latex[length=3mm]},dashed] (-2.7,-0.7) -- (-2.2,-0.1);
    \node at (-3.2,-0.9) {\small traiettoria};
\end{tikzpicture}
\caption{Schema qualitativo: una traiettoria che parte in \(M\) rimane in \(M\).}
\end{figure}

\section{Perché serve LaSalle}

Il teorema diretto di Lyapunov classico dice:

\begin{itemize}
    \item se esiste \(V\) positiva definita e \(L_fV\le 0\), allora l'equilibrio è stabile;
    \item se esiste \(V\) positiva definita e \(L_fV<0\) per ogni \(x\neq 0\), allora l'equilibrio è asintoticamente stabile.
\end{itemize}

Il punto delicato è il caso intermedio:
\[
L_fV(x)\le 0
\quad \text{ma non strettamente } <0.
\]
In questa situazione il teorema diretto, da solo, spesso garantisce soltanto la stabilità, non la convergenza all'equilibrio.

Il principio di invarianza di LaSalle serve esattamente a colmare questo vuoto: ci dice \emph{dove vanno a finire}, nel lungo periodo, le traiettorie di un sistema quando la funzione di Lyapunov non cresce ma può rimanere costante su qualche sottoinsieme.

\section{Enunciato del principio di invarianza di LaSalle}

\subsection{Versione per sistemi a tempo continuo}

Consideriamo il sistema
\[
\dot x = f(x), \qquad x\in\mathbb{R}^n,
\]
con \(f\) sufficientemente regolare da garantire esistenza e unicità delle soluzioni.

Sia \(\Omega\subseteq \mathbb{R}^n\) un insieme chiuso, limitato e positivamente invariante.  
Sia inoltre
\[
V:\Omega\to\mathbb{R}
\]
una funzione di classe \(C^1\) tale che
\[
L_fV(x)=\nabla V(x)^\top f(x)\le 0
\qquad \forall x\in\Omega.
\]

Definiamo l'insieme
\[
E=\{x\in\Omega : L_fV(x)=0\}.
\]
Sia poi \(M\) il \textbf{massimo insieme invariante} contenuto in \(E\).

\medskip

\noindent
\textbf{Teorema di LaSalle.}
Per ogni traiettoria che parte in \(\Omega\),
\[
x(0)\in\Omega,
\]
si ha
\[
\operatorname{dist}(x(t),M)\to 0
\qquad \text{per } t\to+\infty.
\]

\medskip

\noindent
In parole: la traiettoria converge al massimo insieme invariante contenuto nella zona in cui \(V\) smette di diminuire.

\subsection{Versione per sistemi a tempo discreto}

Per un sistema discreto
\[
x(k+1)=f(x(k)),
\]
si ragiona in modo del tutto analogo.  
Se esiste una funzione \(V\) tale che
\[
\Delta V(x)=V(f(x))-V(x)\le 0,
\]
allora si considera l'insieme
\[
E=\{x\in\Omega : \Delta V(x)=0\},
\]
e il massimo insieme invariante \(M\subseteq E\).  
Anche in questo caso le traiettorie convergono a \(M\).

\section{Interpretazione geometrica}

Il principio di LaSalle si capisce bene così.

\begin{itemize}
    \item La funzione \(V(x)\) si comporta come una ``energia''.
    \item Se \(L_fV\le 0\), questa energia non aumenta.
    \item Quindi la traiettoria non può vagare in modo arbitrario: è costretta a scendere, o al massimo a restare sullo stesso livello di energia.
    \item Nel lungo periodo, allora, il sistema deve finire nella regione in cui la discesa si arresta, cioè in
    \[
    E=\{x:L_fV(x)=0\}.
    \]
    \item Ma non tutto \(E\) è necessariamente compatibile con la dinamica: la traiettoria può fermarsi solo nel \emph{massimo insieme invariante} contenuto in \(E\).
\end{itemize}

\begin{figure}[h!]
\centering
\begin{tikzpicture}[scale=1]
    \draw[->] (-3.2,0) -- (3.2,0) node[right] {$x_1$};
    \draw[->] (0,-3.2) -- (0,3.2) node[above] {$x_2$};

    \draw[thick] (0,0) circle (2.4);
    \draw[thick] (0,0) circle (1.7);
    \draw[thick] (0,0) circle (1.0);

    \draw[red,thick] (-2.8,0.8) .. controls (-1.6,1.1) and (-0.5,0.8) .. (0.1,0.2);
    \draw[-{Latex[length=3mm]},red,thick] (-0.05,0.25) -- (0.1,0.2);

    \draw[blue,thick] (-1.3,-1.3) -- (1.3,1.3);
    \node[blue] at (1.7,1.6) {\small \(E=\{L_fV=0\}\)};

    \fill (0,0) circle (1.7pt);
    \node at (0.35,-0.3) {\small \(M\)};
\end{tikzpicture}
\caption{Idea qualitativa: la traiettoria scende lungo le curve di livello di \(V\) e converge al massimo insieme invariante contenuto in \(E\).}
\end{figure}

\section{Corollario fondamentale per la stabilità asintotica}

Il caso più importante, in pratica, è questo.

\medskip

\noindent
\textbf{Corollario.}
Se l'origine è un equilibrio, esiste una funzione \(V\) positiva definita e
\[
L_fV(x)\le 0,
\]
e inoltre il massimo insieme invariante contenuto in
\[
E=\{x:L_fV(x)=0\}
\]
è il solo punto
\[
M=\{0\},
\]
allora l'origine è \textbf{asintoticamente stabile}.

\medskip

\noindent
Questo risultato è estremamente utile perché permette di concludere l'asintotica stabilità anche quando \(L_fV\) è solo semidefinita negativa.

\section{Mini-esempio: convergenza a una retta di equilibri}

Consideriamo il sistema
\[
\begin{cases}
\dot x_1 = -x_1,\\[4pt]
\dot x_2 = 0.
\end{cases}
\]
Scegliamo
\[
V(x)=\frac{1}{2}x_1^2.
\]
Allora
\[
L_fV=x_1\dot x_1=-x_1^2\le 0.
\]
Quindi
\[
E=\{x:L_fV=0\}=\{x_1=0\}.
\]
Ma l'insieme \(\{x_1=0\}\) è esso stesso invariante, perché se \(x_1(0)=0\), allora \(x_1(t)=0\) per ogni \(t\), mentre \(x_2\) resta costante.

Dunque il massimo insieme invariante contenuto in \(E\) è
\[
M=\{(0,x_2):x_2\in\mathbb{R}\}.
\]
Per LaSalle, le traiettorie convergono a questa retta, non necessariamente all'origine.

Questo esempio è importante perché mostra bene che:
\begin{center}
LaSalle non dice automaticamente che si converge a un punto; dice che si converge al massimo insieme invariante contenuto in \(E\).
\end{center}

\section{Collegamento con l'esercizio della lezione precedente}

Riprendiamo il sistema già studiato in precedenza, nel caso \(k=0\):
\[
\begin{cases}
\dot z_1 = -z_1-z_2^2,\\[4pt]
\dot z_2 = 4z_1z_2.
\end{cases}
\]
Scegliamo
\[
V(z)=\frac{1}{2}\left(4z_1^2+z_2^2\right).
\]
Allora
\[
L_fV = -4z_1^2\le 0.
\]
Quindi
\[
E=\{z:L_fV(z)=0\}=\{z_1=0\}.
\]

Per capire il massimo insieme invariante contenuto in \(E\), imponiamo \(z_1(t)=0\) per tutti i tempi.  
Ma allora dalla prima equazione segue
\[
\dot z_1 = -z_2^2.
\]
Affinché \(z_1\) resti identicamente nullo, deve valere necessariamente
\[
-z_2^2=0 \quad \Longrightarrow \quad z_2=0.
\]
Quindi l'unico insieme invariante contenuto in \(E\) è il solo punto
\[
M=\{(0,0)\}.
\]
Per il corollario di LaSalle segue allora che l'origine è \textbf{asintoticamente stabile}.

Questo è esattamente il tipo di situazione in cui LaSalle è più potente del teorema diretto classico.

\section{Esercizio 1: studio di un sistema non lineare}
\label{sec:L23-es1}

Studiamo il sistema
\begin{equation}
\label{eq:L23-es1}
\begin{cases}
\dot x_1 = -x_1^3 + x_2,\\[4pt]
\dot x_2 = -x_1^3 - x_2^3.
\end{cases}
\end{equation}

\subsection{Equilibri}

Gli equilibri si trovano imponendo
\[
\begin{cases}
0 = -x_1^3 + x_2,\\[4pt]
0 = -x_1^3 - x_2^3.
\end{cases}
\]
Dalla prima equazione:
\[
x_2=x_1^3.
\]
Sostituendo nella seconda:
\[
0=-x_1^3-(x_1^3)^3=-x_1^3-x_1^9=-x_1^3(1+x_1^6).
\]
Poiché
\[
1+x_1^6>0 \qquad \forall x_1\in\mathbb{R},
\]
segue
\[
x_1=0.
\]
Di conseguenza, anche
\[
x_2=0.
\]
L'unico equilibrio è quindi
\[
(0,0).
\]

\subsection{Metodo indiretto}

Il Jacobiano è
\[
A(x)=
\begin{pmatrix}
-3x_1^2 & 1\\[4pt]
-3x_1^2 & -3x_2^2
\end{pmatrix}.
\]
Valutando nell'origine:
\[
A(0,0)=
\begin{pmatrix}
0 & 1\\
0 & 0
\end{pmatrix}.
\]
Il polinomio caratteristico è
\[
\lambda^2=0,
\]
quindi gli autovalori sono entrambi nulli. Il metodo indiretto è dunque inconcludente.

\subsection{Scelta di una funzione di Lyapunov}

Una scelta quadratica standard,
\[
V=\frac{1}{2}(x_1^2+x_2^2),
\]
non è conveniente perché genera termini misti difficili da controllare:
\[
L_fV=x_1(-x_1^3+x_2)+x_2(-x_1^3-x_2^3).
\]
Per eliminare i termini misti e adattarsi alle potenze del sistema conviene scegliere una funzione di Lyapunov \emph{non quadratica}.

Una scelta efficace è
\[
V(x)=\frac{1}{4}x_1^4+\frac{1}{2}x_2^2.
\]
Essa è positiva definita e radialmente illimitata.

Calcoliamo la derivata lungo le traiettorie:
\[
L_fV = x_1^3\dot x_1 + x_2\dot x_2.
\]
Sostituendo \eqref{eq:L23-es1}:
\begin{align*}
L_fV
&= x_1^3(-x_1^3+x_2)+x_2(-x_1^3-x_2^3)\\
&= -x_1^6 + x_1^3x_2 - x_1^3x_2 - x_2^4\\
&= -x_1^6 - x_2^4.
\end{align*}
Quindi
\[
L_fV(x)<0
\qquad \forall x\neq 0.
\]
Essendo \(V\) positiva definita e radialmente illimitata, e \(L_fV\) negativa definita, concludiamo che l'origine è \textbf{globalmente asintoticamente stabile}.

\medskip

\noindent
\textbf{Osservazione metodologica.}
Questo esercizio è istruttivo perché mostra che, nei sistemi non lineari, la funzione di Lyapunov non deve per forza essere quadratica: spesso bisogna scegliere le potenze giuste per adattarsi alla struttura del campo vettoriale.

\section{Esercizio 2: pendolo smorzato}
\label{sec:L23-pendolo}

Consideriamo il pendolo smorzato
\begin{equation}
\label{eq:L23-pendolo}
\begin{cases}
\dot x_1 = x_2,\\[4pt]
\dot x_2 = -\dfrac{c}{m\ell^2}x_2 - \dfrac{g}{\ell}\sin x_1,
\end{cases}
\end{equation}
dove:
\begin{itemize}
    \item \(x_1=\theta\) è l'angolo,
    \item \(x_2=\dot\theta\) è la velocità angolare,
    \item \(m\) è la massa,
    \item \(\ell\) la lunghezza,
    \item \(c\ge 0\) il coefficiente di smorzamento.
\end{itemize}

\subsection{Linearizzazione nell'origine}

L'origine \((0,0)\) è un equilibrio perché
\[
x_2=0, \qquad \sin x_1 = 0 \text{ per } x_1=0.
\]
Il Jacobiano nell'origine è
\[
A(0,0)=
\begin{pmatrix}
0 & 1\\[4pt]
-\dfrac{g}{\ell} & -\dfrac{c}{m\ell^2}
\end{pmatrix}.
\]

\begin{itemize}
    \item Se \(c>0\), il sistema linearizzato è asintoticamente stabile.
    \item Se \(c=0\), il linearizzato ha autovalori puramente immaginari e il metodo indiretto non conclude.
\end{itemize}

\subsection{Funzione energia}

La funzione naturale da provare è l'energia meccanica:
\[
V(x)=mg\ell\bigl(1-\cos x_1\bigr)+\frac{m\ell^2}{2}x_2^2.
\]
Vicino all'origine questa funzione è positiva definita, perché:
\[
1-\cos x_1 \sim \frac{x_1^2}{2}
\qquad \text{per } x_1\to 0.
\]

La derivata lungo le traiettorie è
\[
L_fV
=
mg\ell\sin x_1\,\dot x_1 + m\ell^2 x_2\dot x_2.
\]
Sostituendo \eqref{eq:L23-pendolo}:
\begin{align*}
L_fV
&=
mg\ell\sin x_1\,x_2
+
m\ell^2x_2\left(
-\frac{c}{m\ell^2}x_2 - \frac{g}{\ell}\sin x_1
\right)\\
&=
mg\ell\sin x_1\,x_2
-cx_2^2
-mg\ell x_2\sin x_1\\
&=
-cx_2^2.
\end{align*}
Quindi:
\[
L_fV = -cx_2^2 \le 0.
\]

\subsection{Caso \(c>0\): applicazione di LaSalle}

Se \(c>0\), allora
\[
L_fV=0 \iff x_2=0.
\]
Quindi
\[
E=\{(x_1,x_2):x_2=0\}.
\]

Per applicare bene LaSalle conviene limitarsi a un sottolivello di energia che non includa il punto alto del pendolo, cioè scegliere
\[
\Omega_\rho = \{x: V(x)\le \rho\}
\qquad \text{con } 0<\rho<2mg\ell.
\]
In questo modo restiamo in una regione attorno all'origine e non includiamo l'equilibrio instabile in \(x_1=\pi\).

Se una traiettoria resta in \(E\), allora \(x_2(t)=0\) per ogni \(t\). Ma dalla seconda equazione di \eqref{eq:L23-pendolo} segue
\[
\dot x_2 = -\frac{g}{\ell}\sin x_1.
\]
Affinché \(x_2\) resti identicamente nullo, deve essere
\[
\sin x_1 = 0.
\]
Nel sottolivello scelto, l'unico punto che soddisfa questa condizione è
\[
x_1=0.
\]
Dunque il massimo insieme invariante contenuto in \(E\cap \Omega_\rho\) è il solo punto
\[
M=\{(0,0)\}.
\]
Per LaSalle, l'origine è \textbf{asintoticamente stabile}.

\subsection{Caso \(c=0\): solo stabilità}

Se \(c=0\), allora
\[
L_fV=0
\qquad \forall x.
\]
Quindi il teorema di LaSalle non può dare convergenza all'origine, perché tutto l'insieme considerato appartiene a \(E\). In effetti il sistema è conservativo: l'energia si conserva e le traiettorie, in generale, non convergono all'equilibrio.

Quindi:
\[
c=0
\quad \Longrightarrow \quad
(0,0) \text{ è stabile, ma non asintoticamente stabile.}
\]

\section{Esercizio 3: sistema discreto}
\label{sec:L23-TD}

Consideriamo il sistema discreto
\begin{equation}
\label{eq:L23-TD}
\begin{cases}
x_1(k+1)=\alpha x_2(k)-x_1^2(k)x_2(k),\\[4pt]
x_2(k+1)=-\alpha x_1(k)+x_1(k)x_2^2(k).
\end{cases}
\end{equation}

\subsection{Linearizzazione nell'origine}

L'origine è un equilibrio.  
Il linearizzato nell'origine è
\[
x(k+1)=Ax(k),
\qquad
A=
\begin{pmatrix}
0 & \alpha\\
-\alpha & 0
\end{pmatrix}.
\]
Gli autovalori sono
\[
\lambda_{1,2}=\pm j\alpha,
\]
quindi il loro modulo è
\[
|\lambda_{1,2}|=|\alpha|.
\]

Da qui segue:
\[
|\alpha|<1 \Rightarrow (0,0)\text{ asintoticamente stabile},
\]
\[
|\alpha|>1 \Rightarrow (0,0)\text{ instabile},
\]
\[
|\alpha|=1 \Rightarrow \text{metodo indiretto inconcludente}.
\]

\subsection{Caso critico \(\alpha=1\)}

Poniamo ora \(\alpha=1\).  
Proviamo con
\[
V(x)=x_1^2+x_2^2.
\]
Calcoliamo
\[
\Delta V(x)=V(x(k+1))-V(x(k)).
\]
Sostituendo la dinamica:
\begin{align*}
\Delta V
&=
\bigl(x_2-x_1^2x_2\bigr)^2+\bigl(-x_1+x_1x_2^2\bigr)^2-x_1^2-x_2^2\\
&=
x_1^2x_2^2\bigl(x_1^2+x_2^2-4\bigr).
\end{align*}
Quindi, in un intorno sufficientemente piccolo dell'origine,
\[
\Delta V \le 0.
\]
Per esempio, nel disco
\[
x_1^2+x_2^2<4
\]
si ha effettivamente \(\Delta V\le 0\).

Tuttavia
\[
\Delta V=0
\]
non si verifica solo nell'origine, ma anche sugli assi \(x_1=0\) e \(x_2=0\). Bisogna quindi studiare il massimo insieme invariante contenuto in tale insieme.

Se partiamo da un punto dell'asse \(x_1\), per esempio
\[
x(0)=(a,0),
\]
si ottiene:
\[
x(1)=(0,-a),
\qquad
x(2)=(-a,0),
\qquad
x(3)=(0,a),
\qquad
x(4)=(a,0).
\]
Quindi compare un'orbita periodica di periodo \(4\).

Dunque il massimo insieme invariante contenuto in \(\{\Delta V=0\}\) non è il solo punto \((0,0)\).  
Perciò LaSalle non dà asintotica stabilità. In effetti l'origine risulta soltanto \textbf{stabile}, non asintoticamente stabile.

\medskip

\noindent
Questo esempio è importante perché mostra una differenza sostanziale tra:
\begin{itemize}
    \item ``\(\Delta V\le 0\)'',
    \item ``\(\Delta V<0\)''.
\end{itemize}
Nel primo caso, senza studiare l'insieme invariante, non si può concludere la convergenza all'equilibrio.

\section{Lyapunov e sintesi del controllo}

Una delle idee più importanti della lezione è che i metodi di Lyapunov non servono solo per \emph{analizzare} la stabilità, ma anche per \emph{progettare} controlli stabilizzanti.

\subsection{Idea di fondo}

Supponiamo di avere un sistema non lineare controllato:
\[
\dot x = f(x,u).
\]
Se vogliamo stabilizzare un equilibrio \((\bar x,\bar u)\), conviene traslare il sistema:
\[
z=x-\bar x,
\qquad
v=u-\bar u.
\]
In queste coordinate l'equilibrio da stabilizzare diventa l'origine:
\[
z=0,\qquad v=0.
\]

L'idea è allora cercare una funzione \(V(z)\) positiva definita e scegliere una legge di controllo
\[
v=\alpha(z)
\]
in modo che
\[
L_fV(z)\le 0
\quad \text{o meglio} \quad
L_fV(z)<0.
\]

In questo modo il controllo è \emph{stabilizzante per costruzione}.

\subsection{Perché il controllo non deve spostare l'equilibrio}

Se il controllo serve a stabilizzare l'equilibrio, è essenziale che nel punto di equilibrio esso si annulli in coordinate traslate:
\[
v(0)=0.
\]
Equivalentemente, in coordinate originali deve valere
\[
u(\bar x)=\bar u.
\]
Se così non fosse, il controllo introdurrebbe una forza costante anche nel punto che vogliamo mantenere come equilibrio, spostandolo.

\section{Controllare prima e linearizzare dopo, o viceversa}

Per un sistema non lineare
\[
\dot x=f(x,u),
\]
sia \((\bar x,\bar u)\) un equilibrio, cioè
\[
f(\bar x,\bar u)=0.
\]
Linearizzando attorno a \((\bar x,\bar u)\), otteniamo
\[
\dot z = Az+Bv,
\]
dove
\[
A=\left.\frac{\partial f}{\partial x}\right|_{(\bar x,\bar u)},
\qquad
B=\left.\frac{\partial f}{\partial u}\right|_{(\bar x,\bar u)}.
\]

Se imponiamo una retroazione lineare locale
\[
v=Kz,
\]
il linearizzato in anello chiuso diventa
\[
\dot z = (A+BK)z.
\]

Questa osservazione è fondamentale: localmente, progettare il feedback sul sistema linearizzato equivale a linearizzare il sistema già controllato.  
Perciò, se \(A+BK\) è Hurwitz, il sistema non lineare controllato è localmente asintoticamente stabile attorno all'equilibrio, per il metodo indiretto di Lyapunov.

\section{Esempio di sintesi con Lyapunov}

Consideriamo il sistema
\begin{equation}
\label{eq:L23-controllo}
\begin{cases}
\dot x_1 = (x_2+1)^3,\\[4pt]
\dot x_2 = (x_1-1)^3 + u + 2.
\end{cases}
\end{equation}

Vogliamo scegliere \(u(t)\) in modo da rendere stabile l'equilibrio
\[
(\bar x_1,\bar x_2)=(1,-1).
\]

\subsection{Traslazione}

Poniamo
\[
z_1=x_1-1,\qquad z_2=x_2+1,
\qquad
v=u+2.
\]
In questo modo l'equilibrio desiderato diventa l'origine:
\[
z=0,\qquad v=0.
\]
Il sistema si riscrive come
\[
\begin{cases}
\dot z_1 = z_2^3,\\[4pt]
\dot z_2 = z_1^3 + v.
\end{cases}
\]

\subsection{Scelta della funzione di Lyapunov}

Proviamo con
\[
V(z)=\frac12(z_1^2+z_2^2).
\]
Allora
\[
L_fV=z_1\dot z_1+z_2\dot z_2
= z_1 z_2^3 + z_2(z_1^3+v).
\]
Sviluppando:
\[
L_fV = z_1z_2^3 + z_1^3z_2 + z_2v.
\]
I primi due termini hanno segno indefinito, quindi da soli non consentono di concludere nulla.  
L'idea è usare il controllo per cancellarli e aggiungere un termine dissipativo.

Scegliamo
\[
v = -z_1z_2^2 - z_1^3 - z_2.
\]
Allora
\begin{align*}
L_fV
&= z_1z_2^3 + z_1^3z_2 + z_2\bigl(-z_1z_2^2-z_1^3-z_2\bigr)\\
&= -z_2^2 \le 0.
\end{align*}
Quindi il sistema controllato è stabile.

Per vedere l'asintotica stabilità, osserviamo che
\[
L_fV=0 \iff z_2=0.
\]
Se \(z_2=0\), allora
\[
\dot z_1=0,
\qquad
\dot z_2=z_1^3.
\]
Affinché la traiettoria resti in \(\{z_2=0\}\), deve essere
\[
z_1^3=0 \quad \Longrightarrow \quad z_1=0.
\]
Quindi il massimo insieme invariante contenuto in \(\{L_fV=0\}\) è il solo punto
\[
M=\{0\}.
\]
Per LaSalle, l'origine del sistema controllato è asintoticamente stabile.

\subsection{Legge di controllo in coordinate originali}

Ricordando che
\[
z_1=x_1-1,\qquad z_2=x_2+1,\qquad v=u+2,
\]
si ottiene
\[
u = v-2,
\]
quindi una legge di controllo stabilizzante è
\[
\boxed{
u(x)= -2 - (x_1-1)(x_2+1)^2 - (x_1-1)^3 - (x_2+1).
}
\]

\section{Cicli limite: cenno qualitativo}

Nella parte finale della lezione compare un'osservazione importante sui sistemi non lineari: oltre agli equilibri, possono esistere anche \textbf{cicli limite}.

\medskip

\noindent
\textbf{Definizione informale.}
Un ciclo limite è una traiettoria periodica isolata nello spazio delle fasi.  
In altre parole, il sistema torna periodicamente sugli stessi stati e la traiettoria descrive una curva chiusa.

\medskip

\noindent
Questo fenomeno è tipico dei sistemi non lineari.  
Nei sistemi lineari autonomi non banali, infatti, le traiettorie periodiche non sono isolate nello stesso senso: se esiste una traiettoria chiusa, tipicamente ne esiste un'intera famiglia.

\medskip

\noindent
Un ciclo limite può essere:
\begin{itemize}
    \item attrattivo,
    \item repulsivo,
    \item semistabile.
\end{itemize}

Talvolta, per studiarne l'attrattività, si costruisce una funzione del tipo
\[
V(x)=C(x)^2,
\]
dove la curva
\[
C(x)=0
\]
descrive il ciclo o una curva invariante candidata.  
L'idea è che \(V\) sia nulla sulla curva e positiva fuori, così da misurare la distanza energetica dalla traiettoria chiusa.

\section{Schema operativo per gli esercizi di fine lezione}

Gli esercizi di questa lezione si risolvono in genere seguendo questo schema.

\subsection*{A. Trovare gli equilibri}
Si impone:
\[
f(x)=0
\]
oppure, nel caso controllato,
\[
f(x,u)=0
\]
con l'ingresso assegnato o da scegliere.

\subsection*{B. Provare il linearizzato}
Si calcola il Jacobiano:
\[
A=\frac{\partial f}{\partial x}.
\]
Se tutti gli autovalori hanno parte reale negativa (o modulo minore di \(1\) nel discreto), allora si conclude subito la stabilità asintotica locale.

\subsection*{C. Se il linearizzato è inconcludente, usare Lyapunov}
Si prova una funzione:
\begin{itemize}
    \item quadratica, se il sistema lo suggerisce;
    \item energetica, se c'è una struttura meccanica;
    \item con potenze adattate alla non linearità, se compaiono monomi di grado alto.
\end{itemize}

\subsection*{D. Se \(L_fV\) o \(\Delta V\) è solo semidefinita, usare LaSalle}
Si calcola
\[
E=\{x:L_fV(x)=0\}
\qquad \text{oppure} \qquad
E=\{x:\Delta V(x)=0\},
\]
poi si cerca il massimo insieme invariante contenuto in \(E\).

\subsection*{E. Nei problemi di controllo, traslare prima l'equilibrio}
Se l'equilibrio desiderato non è l'origine:
\[
z=x-\bar x,\qquad v=u-\bar u.
\]
Poi si progetta il controllo in \((z,v)\), imponendo
\[
v(0)=0.
\]

\section{Riassunto finale}

La lezione porta a quattro idee chiave.

\begin{enumerate}
    \item \textbf{LaSalle completa Lyapunov:} se \(V\) non cresce, la traiettoria converge al massimo insieme invariante contenuto in \(\{L_fV=0\}\).
    \item \textbf{Per concludere l'asintotica stabilità} basta spesso mostrare che tale massimo insieme invariante è il solo equilibrio.
    \item \textbf{Le funzioni di Lyapunov si possono progettare}, non solo usare a posteriori: questo apre direttamente alla sintesi del controllo.
    \item \textbf{Nei sistemi non lineari} compaiono fenomeni più ricchi degli equilibri, come i cicli limite, che richiedono strumenti qualitativi più raffinati.
\end{enumerate}

\medskip

\noindent
In forma compatta, il messaggio da ricordare è:

\begin{center}
\emph{Se trovi una funzione \(V\) che non aumenta, non hai ancora finito: devi chiederti dove il sistema può restare quando \(V\) smette di diminuire. La risposta la dà il massimo insieme invariante di LaSalle.}
\end{center}
% =========================
% L24.tex
% =========================

\chapter{L24 -- Sintesi via Lyapunov, cicli limite, teoremi di instabilit\`a e metodo indiretto}
\label{chap:L24}

\section{Richiamo: progetto di un controllo tramite Lyapunov}

Riprendiamo l'idea vista nella parte finale della lezione precedente: la teoria di Lyapunov non serve solo per \emph{analizzare} la stabilit\`a di un sistema, ma pu\`o essere usata anche per \emph{progettare} un controllo stabilizzante.

Consideriamo il sistema non lineare controllato
\begin{equation}
\label{eq:L24-sistema-controllato}
\begin{cases}
\dot x_1 = (x_2+1)^3,\\[4pt]
\dot x_2 = (x_1-1)^3 + u + 2.
\end{cases}
\end{equation}

Vogliamo costruire un controllo \(u(t)\) tale che il punto
\[
(\bar x_1,\bar x_2)=(1,-1)
\]
sia un equilibrio stabile.

\subsection{Scelta dell'equilibrio e traslazione}

Per prima cosa osserviamo che, se fissiamo
\[
\bar u=-2,
\]
allora il punto \((1,-1)\) \`e effettivamente un equilibrio del sistema, perch\'e
\[
(x_2+1)^3 = (-1+1)^3 = 0,
\]
e inoltre
\[
(x_1-1)^3 + \bar u + 2 = (1-1)^3 -2 + 2 = 0.
\]

Per lavorare pi\`u comodamente trasliamo il sistema nell'origine:
\[
z_1 = x_1-1,
\qquad
z_2 = x_2+1,
\qquad
v = u-\bar u = u+2.
\]
Equivalentemente,
\[
u = v-2.
\]

In queste nuove coordinate il sistema diventa
\begin{equation}
\label{eq:L24-sistema-traslato}
\begin{cases}
\dot z_1 = z_2^3,\\[4pt]
\dot z_2 = z_1^3 + v.
\end{cases}
\end{equation}
Ora l'equilibrio desiderato coincide con l'origine \(z=0\).

\subsection{Scelta della funzione di Lyapunov}

Proviamo con la candidata naturale
\[
V(z)=\frac{1}{2}(z_1^2+z_2^2),
\]
che \`e positiva definita.

La derivata lungo le traiettorie di \eqref{eq:L24-sistema-traslato} \`e
\[
L_fV(z)=\frac{\partial V}{\partial z}\dot z
= z_1\dot z_1 + z_2\dot z_2.
\]
Sostituendo:
\begin{align*}
L_fV(z)
&= z_1 z_2^3 + z_2(z_1^3+v)\\
&= z_1 z_2^3 + z_1^3 z_2 + z_2 v.
\end{align*}

I primi due termini hanno segno indefinito. L'idea allora \`e usare il controllo per cancellarli e aggiungere un termine dissipativo.

Scegliamo
\[
v = - z_1 z_2^2 - z_1^3 - z_2.
\]
Allora
\begin{align*}
L_fV(z)
&= z_1 z_2^3 + z_1^3 z_2 + z_2\bigl(- z_1 z_2^2 - z_1^3 - z_2\bigr)\\
&= z_1 z_2^3 + z_1^3 z_2 - z_1 z_2^3 - z_1^3 z_2 - z_2^2\\
&= -z_2^2 \le 0.
\end{align*}
Quindi il sistema in anello chiuso \`e stabile.

\subsection{Legge di controllo in coordinate originali}

Ricordando che \(u=v-2\), otteniamo la legge di controllo
\[
\boxed{
u(x)=-(x_1-1)(x_2+1)^2-(x_1-1)^3-(x_2+1)-2.
}
\]

\medskip

\noindent
\textbf{Osservazione importante.}
Questo esempio mostra bene l'idea fondamentale:
\begin{center}
\emph{si pu\`o progettare il controllo imponendo direttamente il segno della derivata di una funzione di Lyapunov.}
\end{center}

\section{Controllo affine in \(x\)}

Una legge di controllo si dice \textbf{affine in \(x\)} se ha la forma
\[
u = Kx + h,
\]
oppure, in componenti,
\[
u = k_1 x_1 + k_2 x_2 + \cdots + k_n x_n + h.
\]

Si tratta quindi di un controllo lineare nelle variabili di stato, pi\`u una costante.

\medskip

\noindent
\textbf{Perch\'e compare la costante?}  
Perch\'e in generale l'equilibrio da stabilizzare non coincide con l'origine. Se vogliamo stabilizzare un punto \(\bar x\), allora il controllo deve soddisfare
\[
u(\bar x)=\bar u,
\]
cio\`e deve assumere il valore corretto nel punto di equilibrio. In coordinate traslate, questo equivale a chiedere che il controllo si annulli nell'origine:
\[
v(0)=0.
\]

\medskip

\noindent
In altre parole:
\begin{center}
\emph{un controllo stabilizzante non deve spostare il punto di equilibrio che vogliamo mantenere.}
\end{center}

\section{Sistema non lineare controllato vs linearizzato controllato}

Questa parte \`e concettualmente molto importante.

Sia dato il sistema non lineare controllato
\[
\dot x = f(x,u),
\]
e sia \((\bar x,\bar u)\) una coppia di equilibrio, cio\`e
\[
f(\bar x,\bar u)=0.
\]

Trasliamo le variabili:
\[
z = x-\bar x,
\qquad
v = u-\bar u.
\]

Sviluppando \(f(x,u)\) al primo ordine attorno a \((\bar x,\bar u)\), otteniamo
\[
f(x,u)\approx f(\bar x,\bar u)
+ \left.\frac{\partial f}{\partial x}\right|_{(\bar x,\bar u)} (x-\bar x)
+ \left.\frac{\partial f}{\partial u}\right|_{(\bar x,\bar u)} (u-\bar u).
\]
Poich\'e \(f(\bar x,\bar u)=0\), segue
\[
\dot z \approx Az + Bv,
\]
dove
\[
A=\left.\frac{\partial f}{\partial x}\right|_{(\bar x,\bar u)},
\qquad
B=\left.\frac{\partial f}{\partial u}\right|_{(\bar x,\bar u)}.
\]

Supponiamo ora di applicare una retroazione statica affine
\[
u = Kx + h.
\]
Affinch\'e \((\bar x,\bar u)\) resti equilibrio deve valere
\[
\bar u = K\bar x + h.
\]
Dunque, nelle variabili traslate,
\[
v = u-\bar u = K(x-\bar x)=Kz.
\]
Sostituendo nel linearizzato:
\[
\dot z \approx Az + BKz = (A+BK)z.
\]

\subsection{Conclusione concettuale}

Il \textbf{linearizzato del sistema non lineare gi\`a controllato} coincide, al primo ordine, con il \textbf{sistema linearizzato non controllato a cui si applica poi la stessa retroazione lineare}.

\medskip

\noindent
In simboli:
\[
\text{linearizza}\bigl(f(x,Kx+h)\bigr)
\quad \Longleftrightarrow \quad
\dot z = (A+BK)z.
\]

Questa osservazione spiega perch\'e, localmente, il progetto del controllo sul modello linearizzato \`e significativo anche per il sistema non lineare.

\medskip

\noindent
\textbf{Messaggio da ricordare.}
Se il linearizzato in anello chiuso
\[
\dot z=(A+BK)z
\]
\`e asintoticamente stabile, allora il sistema non lineare controllato \`e localmente asintoticamente stabile attorno all'equilibrio, per il metodo indiretto di Lyapunov.

\section{Ciclo limite}

Nei sistemi non lineari possono comparire fenomeni che nei sistemi lineari non si osservano in questo modo. Il pi\`u importante \`e il \textbf{ciclo limite}.

\subsection{Definizione qualitativa}

Un ciclo limite \`e una traiettoria periodica isolata nello spazio delle fasi.  
Questo significa che:

\begin{itemize}
    \item la traiettoria \`e chiusa;
    \item il sistema vi ritorna periodicamente;
    \item le traiettorie vicine possono essere attratte oppure respinte da tale curva.
\end{itemize}

\medskip

\noindent
\textbf{Attenzione:} un ciclo limite \emph{non} \`e un equilibrio.  
Un equilibrio \`e un punto fermo; un ciclo limite \`e un moto periodico.

\subsection{Idea per studiarne l'attrattivit\`a}

Se una curva chiusa \`e descritta implicitamente da
\[
C(x)=0,
\]
si pu\`o provare a costruire
\[
V(x)=\bigl(C(x)\bigr)^2.
\]
Questa funzione:

\begin{itemize}
    \item vale zero sulla curva;
    \item \`e positiva fuori dalla curva;
    \item misura, in un certo senso, la distanza energetica dalla curva.
\end{itemize}

La derivata lungo le traiettorie \`e
\[
L_fV(x)=2\,C(x)\,\frac{\partial C}{\partial x}f(x).
\]
Il segno di \(L_fV\) ci dice se le traiettorie si avvicinano alla curva oppure se ne allontanano.

\begin{itemize}
    \item Se \(L_fV<0\) vicino alla curva, il ciclo \`e attrattivo.
    \item Se \(L_fV>0\) vicino alla curva, il ciclo \`e repulsivo.
\end{itemize}

\subsection{Esempio classico}

Consideriamo il sistema
\begin{equation}
\label{eq:L24-ciclo}
\begin{cases}
\dot x_1 = -x_2 + x_1(1-x_1^2-x_2^2),\\[4pt]
\dot x_2 = x_1 + x_2(1-x_1^2-x_2^2).
\end{cases}
\end{equation}

\subsubsection{Studio dell'origine}

Con
\[
V(x)=\frac12(x_1^2+x_2^2),
\]
si ottiene
\begin{align*}
L_fV
&= x_1\dot x_1 + x_2\dot x_2\\
&= x_1\bigl(-x_2 + x_1(1-r^2)\bigr)
+ x_2\bigl(x_1 + x_2(1-r^2)\bigr),
\end{align*}
dove
\[
r^2=x_1^2+x_2^2.
\]
I termini misti si cancellano:
\[
L_fV = (x_1^2+x_2^2)(1-x_1^2-x_2^2)=r^2(1-r^2).
\]

Vicino all'origine vale \(r<1\), quindi
\[
L_fV>0
\qquad \text{per } x\neq 0 \text{ sufficientemente vicino all'origine}.
\]
Dunque l'origine \`e instabile.

\subsubsection{Studio della circonferenza \(x_1^2+x_2^2=1\)}

Poniamo
\[
C(x)=1-x_1^2-x_2^2 = 1-r^2,
\qquad
V_C(x)=\bigl(C(x)\bigr)^2.
\]
Allora
\[
\dot C = -2x_1\dot x_1 -2x_2\dot x_2.
\]
Sostituendo \eqref{eq:L24-ciclo}:
\begin{align*}
\dot C
&=
-2x_1\bigl(-x_2 + x_1(1-r^2)\bigr)
-2x_2\bigl(x_1 + x_2(1-r^2)\bigr)\\
&=
-2(x_1^2+x_2^2)(1-r^2)\\
&=
-2r^2 C.
\end{align*}
Di conseguenza
\[
L_fV_C = \frac{d}{dt}(C^2)=2C\dot C = -4r^2 C^2 \le 0.
\]
Fuori dall'origine e fuori dalla circonferenza unitaria, la derivata \`e strettamente negativa. Dunque la curva
\[
x_1^2+x_2^2=1
\]
\`e un ciclo limite attrattivo.

\begin{figure}[h!]
\centering
\begin{tikzpicture}[scale=1.2]
    \draw[->] (-2,0) -- (2,0) node[right] {$x_1$};
    \draw[->] (0,-2) -- (0,2) node[above] {$x_2$};

    \draw[thick,red] (0,0) circle (1);
    \node[red] at (1.35,1.2) {\small ciclo limite};

    \draw[orange,dashed,thick] (0,0) circle (0.55);
    \draw[orange,dashed,thick] (0,0) circle (1.45);

    \draw[-{Latex[length=3mm]},orange,thick] (0.55,0) -- (0.8,0);
    \draw[-{Latex[length=3mm]},orange,thick] (-0.55,0) -- (-0.8,0);
    \draw[-{Latex[length=3mm]},orange,thick] (0,0.55) -- (0,0.8);
    \draw[-{Latex[length=3mm]},orange,thick] (0,-0.55) -- (0,-0.8);

    \draw[-{Latex[length=3mm]},orange,thick] (1.45,0) -- (1.15,0);
    \draw[-{Latex[length=3mm]},orange,thick] (-1.45,0) -- (-1.15,0);
    \draw[-{Latex[length=3mm]},orange,thick] (0,1.45) -- (0,1.15);
    \draw[-{Latex[length=3mm]},orange,thick] (0,-1.45) -- (0,-1.15);

    \fill (0,0) circle (1.4pt);
    \node at (-0.45,-0.2) {\small origine instabile};
\end{tikzpicture}
\caption{Schema qualitativo: l'origine \`e instabile e la circonferenza unitaria \`e un ciclo limite attrattivo.}
\end{figure}

\section{Teoremi inversi di Lyapunov: idea generale}

I teoremi diretti di Lyapunov dicono:
\begin{center}
\emph{se trovi una funzione \(V\) con certe propriet\`a, allora puoi concludere qualcosa sulla stabilit\`a.}
\end{center}

I \textbf{teoremi inversi} affermano, in sostanza, il viceversa:

\begin{itemize}
    \item se un equilibrio \`e stabile, allora esiste una funzione di Lyapunov positiva definita con derivata non positiva;
    \item se un equilibrio \`e asintoticamente stabile, allora esiste una funzione di Lyapunov positiva definita con derivata negativa definita.
\end{itemize}

\medskip

\noindent
Questi risultati sono molto profondi, ma in pratica sono poco costruttivi:
\begin{center}
\emph{garantiscono che la funzione \(V\) esiste, ma non ti dicono come trovarla.}
\end{center}

Per questo, negli esercizi, il problema resta sempre lo stesso: bisogna saper scegliere una buona candidata.

\section{Teorema di instabilit\`a di Chetaev}
\label{sec:L24-chetaev}

Quando vogliamo dimostrare che un equilibrio \`e instabile, esiste un criterio molto utile: il teorema di Chetaev.

\subsection{Enunciato}

Consideriamo il sistema
\[
\dot x=f(x),
\]
con equilibrio nell'origine.  
Sia \(W\) un intorno dell'origine e sia \(V\in C^1(W)\). Supponiamo che esista un aperto \(U\subset W\) tale che:

\begin{enumerate}
    \item l'origine appartiene al bordo di \(U\);
    \item \(V(x)>0\) per ogni \(x\in U\);
    \item \(V(x)=0\) per ogni \(x\in \partial U \cap W\);
    \item \(L_fV(x)>0\) per ogni \(x\in U\).
\end{enumerate}

Allora l'origine \`e \textbf{instabile}.

\medskip

\noindent
L'idea geometrica \`e molto semplice: se in una regione \(U\) arbitrariamente vicina all'origine la funzione \(V\) cresce lungo le traiettorie, allora una traiettoria che parte in \(U\) non pu\`o rimanere confinata vicino all'origine.

\begin{figure}[h!]
\centering
\begin{tikzpicture}[scale=1.1]
    \draw[->] (-2.2,0) -- (2.2,0) node[right] {$x_1$};
    \draw[->] (0,-2.2) -- (0,2.2) node[above] {$x_2$};

    \draw[gray,thick] (0,0) circle (1.4);
    \node[gray] at (-1.6,1.6) {\small \(W\)};

    \fill[pattern=north east lines,pattern color=blue!60] (0,0) -- (1.4,0) arc (0:60:1.4) -- cycle;
    \draw[red,thick] (0,0) -- (1.4,0);
    \draw[red,thick] (0,0) -- ({1.4*cos(60)},{1.4*sin(60)});
    \draw[red,thick] (1.4,0) arc (0:60:1.4);
    \node at (0.8,0.45) {\small \(U\)};

    \fill (0,0) circle (1.5pt);
    \node at (-0.2,-0.2) {\small \(0\)};
\end{tikzpicture}
\caption{Schema qualitativo del teorema di Chetaev: l'origine sta sul bordo della regione \(U\).}
\end{figure}

\subsection{Esempio}

Consideriamo il sistema
\[
\begin{cases}
\dot x_1 = x_2^3,\\[4pt]
\dot x_2 = x_1^3.
\end{cases}
\]
L'origine \`e un equilibrio.

Scegliamo
\[
V(x)=x_1x_2.
\]
Allora
\[
V(0)=0,
\]
e nella regione del primo quadrante
\[
U=\{(x_1,x_2):x_1>0,\ x_2>0,\ x_1^2+x_2^2<\rho^2\}
\]
si ha
\[
V(x)>0.
\]
La derivata lungo le traiettorie \`e
\[
L_fV = x_2 \dot x_1 + x_1 \dot x_2 = x_2 x_2^3 + x_1 x_1^3 = x_1^4 + x_2^4 > 0
\]
per ogni \(x\in U\setminus\{0\}\). Quindi, per Chetaev, l'origine \`e instabile.

\section{Metodo indiretto di Lyapunov per sistemi lineari e non lineari}

Questa parte \`e centrale perch\'e spiega come il linearizzato permette di dedurre la stabilit\`a locale del sistema non lineare.

\subsection{Caso lineare continuo: equazione di Lyapunov}

Consideriamo il sistema lineare tempo continuo
\[
\dot x = Ax.
\]
Cerchiamo una candidata quadratica:
\[
V(x)=x^\top P x,
\]
con \(P=P^\top\).

La derivata lungo le traiettorie \`e
\begin{align*}
\dot V(x)
&= \dot x^\top P x + x^\top P \dot x\\
&= x^\top A^\top P x + x^\top P A x\\
&= x^\top (A^\top P + PA)x.
\end{align*}

Se imponiamo
\[
A^\top P + PA = -Q,
\]
con \(Q=Q^\top>0\), otteniamo
\[
\dot V(x) = -x^\top Q x <0 \qquad \forall x\neq 0.
\]

Questa \`e la \textbf{equazione algebrica di Lyapunov}.

\subsection{Teorema classico}

Dato \(Q=Q^\top>0\), esiste un'unica soluzione simmetrica \(P=P^\top>0\) dell'equazione
\[
A^\top P + PA = -Q
\]
se e solo se tutti gli autovalori di \(A\) hanno parte reale strettamente negativa.

Quindi:
\[
A \text{ Hurwitz}
\quad \Longleftrightarrow \quad
\exists\, P>0 \text{ t.c. } A^\top P + PA = -Q.
\]

\medskip

\noindent
Una rappresentazione utile della soluzione \`e
\[
P=\int_0^{+\infty} e^{A^\top t} Q e^{At}\,dt.
\]

\subsection{Caso lineare discreto}

Per il sistema
\[
x(k+1)=Ax(k),
\]
poniamo ancora
\[
V(x)=x^\top P x.
\]
Allora
\begin{align*}
\Delta V(x)
&=V(x(k+1))-V(x(k))\\
&=x(k)^\top A^\top P A x(k) - x(k)^\top P x(k)\\
&=x(k)^\top (A^\top P A - P)x(k).
\end{align*}
Se imponiamo
\[
A^\top P A - P = -Q,
\qquad Q>0,
\]
si ottiene
\[
\Delta V(x) = -x^\top Q x <0.
\]

Anche qui vale il teorema analogo:
\[
\exists\, P>0 \text{ t.c. } A^\top P A - P = -Q
\quad \Longleftrightarrow \quad
|\lambda_i(A)|<1 \ \forall i.
\]

Una formula utile per \(P\) \`e
\[
P = \sum_{k=0}^{+\infty} (A^\top)^k Q A^k.
\]

\section{Metodo indiretto per sistemi non lineari}

Consideriamo ora un sistema non lineare
\[
\dot x = f(x),
\qquad
f(0)=0.
\]
Supponiamo che \(f\) sia regolare e sviluppabile attorno all'origine:
\[
f(x)=Ax+h(x),
\qquad
A=\left.\frac{\partial f}{\partial x}\right|_{x=0},
\]
dove il resto \(h(x)\) soddisfa
\[
\lim_{x\to 0}\frac{\|h(x)\|}{\|x\|}=0.
\]

\subsection{Caso stabile del linearizzato}

Se \(A\) \`e Hurwitz, scegliamo \(Q=I\) e risolviamo
\[
A^\top P + PA = -I,
\]
ottenendo \(P>0\). Consideriamo allora
\[
V(x)=x^\top P x.
\]
Lungo le traiettorie del sistema non lineare:
\begin{align*}
L_fV(x)
&= x^\top(PA+A^\top P)x + 2x^\top P h(x)\\
&= -\|x\|^2 + 2x^\top P h(x).
\end{align*}
Per \(x\) sufficientemente piccolo, il termine \(h(x)\) \`e di ordine superiore e soddisfa una stima del tipo
\[
\|h(x)\| \le \frac{1}{4\|P\|}\|x\|.
\]
Di conseguenza
\[
2|x^\top P h(x)| \le 2\|x\|\,\|P\|\,\|h(x)\|
\le \frac12 \|x\|^2,
\]
e quindi
\[
L_fV(x)\le -\frac12\|x\|^2<0
\]
per \(x\neq 0\) sufficientemente vicino all'origine.

Allora l'origine del sistema non lineare \`e \textbf{localmente asintoticamente stabile}.

\subsection{Caso instabile del linearizzato}

Se \(A\) possiede almeno un autovalore con parte reale positiva, allora l'origine del sistema non lineare \`e instabile.

L'idea della dimostrazione \`e separare, tramite cambio di base, la parte instabile e quella stabile del linearizzato:
\[
A \sim
\begin{pmatrix}
A_1 & 0\\
0 & A_2
\end{pmatrix},
\]
dove:
\begin{itemize}
    \item \(A_1\) contiene gli autovalori con parte reale positiva;
    \item \(A_2\) contiene quelli con parte reale negativa.
\end{itemize}

Si costruisce quindi una funzione di Chetaev usando due soluzioni di Lyapunov sui due sottoblocchi e si dimostra che esiste una regione in cui la funzione cresce lungo le traiettorie. Da qui segue l'instabilit\`a dell'origine.

\subsection{Caso inconcludente}

Se tutti gli autovalori hanno parte reale \(\le 0\), ma qualcuno ha parte reale nulla, il metodo indiretto \textbf{non conclude}.  
In questo caso servono strumenti ulteriori: Lyapunov diretto, LaSalle, studio qualitativo, centro-manifold, eccetera.

\section{Sistemi non stazionari: perch\'e il criterio sugli autovalori non vale pi\`u}
\label{sec:L24-nonstazionari}

Questa parte \`e soprattutto concettuale.

Consideriamo un sistema lineare tempo-variante
\[
\dot x(t)=A(t)x(t).
\]
A differenza del caso stazionario, qui gli autovalori istantanei di \(A(t)\) \emph{non bastano} per dedurre la stabilit\`a del sistema.

\subsection{Esempio 1: autovalori sempre negativi ma sistema instabile}

Consideriamo
\[
A(t)=
\begin{pmatrix}
-1 & e^{2t}\\
0 & -1
\end{pmatrix}.
\]
Gli autovalori di \(A(t)\) sono sempre
\[
\lambda_1(t)=\lambda_2)=\lambda_2)=\lambda_2(t)=-1,
\]
quindi sembrerebbe un sistema stabile. Ma il sistema \`e
\[
\begin{cases}
\dot x_1 = -x_1 + e^{2t}x_2,\\[4pt]
\dot x_2 = -x_2.
\end{cases}
\]
Dalla seconda equazione:
\[
x_2(t)=e^{-t}x_2(0).
\]
Sostituendo nella prima:
\[
\dot x_1 + x_1 = e^{2t}e^{-t}x_2(0)=e^t x_2(0).
\]
Moltiplicando per \(e^t\):
\[
\frac{d}{dt}\bigl(e^t x_1(t)\bigr)=e^{2t}x_2(0).
\]
Integrando:
\[
e^t x_1(t)=x_1(0)+\frac{x_2(0)}{2}(e^{2t}-1),
\]
quindi
\[
x_1(t)=e^{-t}x_1(0)+\frac{x_2(0)}{2}(e^t-e^{-t}).
\]
Se \(x_2(0)\neq 0\), allora \(x_1(t)\) diverge come \(e^t\). Quindi il sistema \`e instabile, pur avendo autovalori istantanei sempre uguali a \(-1\).

\subsection{Messaggio da ricordare}

Per i sistemi non stazionari:
\begin{center}
\emph{gli autovalori ``congelati'' di \(A(t)\) non determinano da soli la stabilit\`a del sistema.}
\end{center}

Ci\`o che conta \`e l'evoluzione complessiva nel tempo, non il solo spettro istantaneo.

\section{Schema finale della lezione}

La lezione raccoglie diversi messaggi importanti.

\begin{enumerate}
    \item \textbf{Lyapunov pu\`o essere usato per la sintesi del controllo}: scegli una funzione \(V\) e costruisci \(u\) in modo che \(L_fV\le 0\) o \(L_fV<0\).
    
    \item \textbf{Il linearizzato controllato} del sistema non lineare coincide, al primo ordine, con il sistema \((A+BK)\) ottenuto applicando la retroazione al linearizzato.

    \item \textbf{I cicli limite} sono traiettorie periodiche isolate tipiche dei sistemi non lineari: non sono equilibri, ma possono essere attrattivi o repulsivi.

    \item \textbf{Il teorema di Chetaev} \`e uno strumento utile per dimostrare instabilit\`a quando si riesce a costruire una funzione che cresce in una regione arbitrariamente vicina all'origine.

    \item \textbf{L'equazione di Lyapunov} fornisce una procedura concreta per costruire candidate quadratiche nei sistemi lineari stabili.

    \item \textbf{Il metodo indiretto di Lyapunov} collega il segno degli autovalori del linearizzato con la stabilit\`a locale del sistema non lineare.

    \item \textbf{Nei sistemi tempo-varianti}, invece, il criterio sugli autovalori istantanei non \`e pi\`u affidabile.
\end{enumerate}

\medskip

\noindent
\textbf{Frase da ricordare all'orale.}
\begin{center}
\emph{Nel caso non lineare, il linearizzato descrive il comportamento locale solo vicino all'equilibrio; se il linearizzato \`e asintoticamente stabile, allora anche il sistema non lineare lo \`e localmente. Ma se il linearizzato \`e inconcludente, servono Lyapunov diretto, LaSalle o strumenti qualitativi pi\`u fini.}
\end{center}
% =========================
% L25.tex
% =========================

\chapter{L25 -- Sistemi non stazionari, stabilit\`a uniforme e Lyapunov tempo-variante}
\label{chap:L25}

\section{Chiusura del metodo indiretto e passaggio ai sistemi non stazionari}

Concludiamo anzitutto l'idea vista alla fine della lezione precedente sul \emph{metodo indiretto di Lyapunov}.

Quando il linearizzato di un sistema non lineare attorno a un equilibrio ha una parte stabile e una parte instabile, tramite un opportuno cambio di base si pu\`o scrivere, localmente, una dinamica del tipo
\[
\dot x = Ax + h(x),
\qquad
A=
\begin{pmatrix}
A_1 & 0\\
0 & A_2
\end{pmatrix},
\]
dove:
\begin{itemize}
    \item \(A_1\) raccoglie i modi con parte reale positiva,
    \item \(A_2\) raccoglie i modi con parte reale negativa,
    \item \(h(x)\) contiene i termini non lineari di ordine superiore.
\end{itemize}

Per studiare l'instabilit\`a si costruisce una funzione del tipo
\[
V(x)=x^\top P x,
\qquad
P=
\begin{pmatrix}
P_1 & 0\\
0 & -P_2
\end{pmatrix},
\]
con \(P_1>0\), \(P_2>0\).  
Il segno meno sul blocco \(P_2\) serve proprio a costruire una funzione che non sia positiva definita su tutto lo spazio, ma che sia utile per applicare un criterio di instabilit\`a alla Chetaev.

Lungo le traiettorie del sistema non lineare, per \(x\) piccolo, si ottiene una relazione del tipo
\[
L_fV(x)=x^\top x + 2x^\top P h(x),
\]
e il termine \(2x^\top P h(x)\) \`e di ordine superiore.  
Questo consente di concludere che, se nel linearizzato c'\`e almeno un modo esponenzialmente divergente, allora l'equilibrio del sistema non lineare \`e instabile.

\medskip

\noindent
A questo punto passiamo a una nuova classe di sistemi.

\section{Sistemi non stazionari}

Un sistema \`e detto \textbf{non stazionario} o \textbf{tempo-variante} se la dinamica dipende esplicitamente dal tempo:
\begin{equation}
\label{eq:L25-nonstat}
\dot x = f(x,t).
\end{equation}

Nel caso lineare si scrive
\begin{equation}
\label{eq:L25-ltv}
\dot x(t)=A(t)x(t).
\end{equation}

La differenza concettuale rispetto ai sistemi stazionari \`e fondamentale:
\begin{center}
\emph{per i sistemi tempo-varianti, conoscere gli autovalori istantanei di \(A(t)\) non basta pi\`u per concludere sulla stabilit\`a del sistema.}
\end{center}

\subsection{Primo controesempio: autovalori sempre negativi ma soluzione divergente}

Consideriamo il sistema
\[
\dot x(t)=A(t)x(t),
\qquad
A(t)=
\begin{pmatrix}
-1 & a(t)\\
0 & -1
\end{pmatrix},
\qquad
a(t)=e^{2t}.
\]
Gli autovalori di \(A(t)\) sono sempre
\[
\lambda_1(t)=\lambda_2(t)=-1.
\]
Se ragionassimo come nel caso stazionario, saremmo portati a dire che il sistema \`e stabile. Ma non \`e cos\`\i.

Infatti il sistema vale
\[
\begin{cases}
\dot x_1 = -x_1 + e^{2t}x_2,\\[4pt]
\dot x_2 = -x_2.
\end{cases}
\]
Dalla seconda equazione otteniamo immediatamente
\[
x_2(t)=e^{-t}x_2(0).
\]
Sostituendo nella prima:
\[
\dot x_1 + x_1 = e^{2t}e^{-t}x_2(0)=e^t x_2(0).
\]
Moltiplichiamo per il fattore integrante \(e^t\):
\[
\frac{d}{dt}\bigl(e^t x_1(t)\bigr)=e^{2t}x_2(0).
\]
Integrando:
\[
e^t x_1(t)=x_1(0)+\frac{x_2(0)}{2}(e^{2t}-1),
\]
quindi
\[
x_1(t)=e^{-t}x_1(0)+\frac{x_2(0)}{2}\bigl(e^t-e^{-t}\bigr).
\]
Se \(x_2(0)\neq 0\), allora \(x_1(t)\) cresce come \(e^t\), dunque diverge.

\medskip

\noindent
\textbf{Conclusione:}
nonostante gli autovalori istantanei siano sempre uguali a \(-1\), il sistema non \`e stabile.

\subsection{Messaggio da ricordare}

Per i sistemi non stazionari il criterio ``autovalori con parte reale negativa'' non vale pi\`u in generale.  
Conta l'evoluzione complessiva nel tempo, non lo spettro istantaneo della matrice \(A(t)\).

\IfFileExists{25_page-0001.jpg}{
\begin{figure}[h!]
\centering
\includegraphics[width=.82\textwidth]{25_page-0001.jpg}
\caption{Appunti originali: introduzione ai sistemi non stazionari e controesempi sullo spettro istantaneo.}
\end{figure}
}{}

\section{Stabilit\`a per sistemi dipendenti dal tempo}

Nel caso non stazionario la definizione di stabilit\`a deve tener conto del fatto che il comportamento pu\`o dipendere anche dall'istante iniziale \(t_0\).

Consideriamo il sistema \eqref{eq:L25-nonstat}, con equilibrio nell'origine.

\subsection{Definizione di stabilit\`a}

L'origine \(\bar x=0\) si dice \textbf{stabile} se
\[
\forall \varepsilon>0,\ \forall t_0\ge 0,\ \exists \delta=\delta(\varepsilon,t_0)>0
\]
tale che
\[
\|x(t_0)\|<\delta
\quad\Longrightarrow\quad
\|x(t)\|<\varepsilon
\qquad \forall t\ge t_0.
\]

La differenza rispetto al caso stazionario \`e che qui \(\delta\) pu\`o dipendere da \(t_0\).

\subsection{Definizione di stabilit\`a uniforme}

L'origine si dice \textbf{uniformemente stabile} se
\[
\forall \varepsilon>0,\ \exists \delta=\delta(\varepsilon)>0
\]
indipendente da \(t_0\), tale che
\[
\|x(t_0)\|<\delta
\quad\Longrightarrow\quad
\|x(t)\|<\varepsilon
\qquad \forall t\ge t_0.
\]

\medskip

\noindent
\textbf{Interpretazione.}
Nel caso uniforme, la soglia di sicurezza \(\delta\) dipende solo da quanto piccolo vogliamo l'errore finale (\(\varepsilon\)), ma non dipende dal tempo in cui facciamo partire il sistema.

\section{Funzioni di classe \(\mathcal K\) e \(\mathcal K_\infty\)}

Quando si lavora con Lyapunov nei sistemi non stazionari, si usano molto le cosiddette funzioni di classe \(\mathcal K\).

\subsection{Definizione}

Una funzione
\[
\alpha:[0,a)\to [0,+\infty)
\]
si dice di \textbf{classe \(\mathcal K\)} se:
\begin{itemize}
    \item \(\alpha\) \`e continua;
    \item \(\alpha\) \`e strettamente crescente;
    \item \(\alpha(0)=0\).
\end{itemize}

Se inoltre \(a=+\infty\) e
\[
\alpha(r)\to +\infty
\qquad \text{per } r\to+\infty,
\]
allora \(\alpha\) si dice di \textbf{classe \(\mathcal K_\infty\)}.

\medskip

\noindent
Queste funzioni servono per ``incastrare'' una funzione di Lyapunov tra due stime radiali dipendenti solo dalla norma dello stato.

\begin{figure}[h!]
\centering
\begin{tikzpicture}[scale=1]
\draw[->] (-0.2,0) -- (4.3,0) node[right] {$r$};
\draw[->] (0,-0.2) -- (0,3.4) node[above] {$\alpha(r)$};
\draw[thick,blue] (0,0) .. controls (0.7,0.05) and (1.2,0.8) .. (1.8,1.4)
.. controls (2.4,2.0) and (3.2,2.4) .. (4,3.0);
\node[blue] at (3.6,2.5) {\small classe \(\mathcal K\)};
\end{tikzpicture}
\caption{Andamento qualitativo di una funzione di classe \(\mathcal K\).}
\end{figure}

\section{Lemma fondamentale: una funzione positiva definita si stima con funzioni di classe \(\mathcal K\)}

Sia
\[
V:D\to\mathbb R,
\qquad
0\in D\subseteq \mathbb R^n,
\]
continua e positiva definita.  
Se \(B_r\subset D\), allora esistono due funzioni \(\alpha_1,\alpha_2\in\mathcal K\) tali che
\begin{equation}
\label{eq:L25-K-bounds}
\alpha_1(\|x\|)\le V(x)\le \alpha_2(\|x\|)
\qquad \forall x\in B_r.
\end{equation}

Se inoltre \(D=\mathbb R^n\) e \(V\) \`e radialmente illimitata, allora si possono scegliere
\[
\alpha_1,\alpha_2\in\mathcal K_\infty.
\]

\medskip

\noindent
\textbf{Interpretazione.}
La disuguaglianza \eqref{eq:L25-K-bounds} dice che una funzione di Lyapunov positiva definita si comporta come una funzione crescente della distanza dall'origine.

\section{Caratterizzazione della stabilit\`a uniforme}

Un lemma molto utile afferma che l'origine \`e uniformemente stabile se e solo se esistono:
\begin{itemize}
    \item una funzione \(\alpha\in\mathcal K\),
    \item una costante \(c>0\),
\end{itemize}
tali che
\[
\|x(t)\|\le \alpha\bigl(\|x(t_0)\|\bigr)
\qquad \forall t\ge t_0,
\qquad \forall \|x(t_0)\|<c.
\]

\medskip

\noindent
Questa formulazione \`e comoda perch\'e permette di tradurre la stabilit\`a uniforme in una stima quantitativa esplicita.

\IfFileExists{25_page-0002.jpg}{
\begin{figure}[h!]
\centering
\includegraphics[width=.82\textwidth]{25_page-0002.jpg}
\caption{Appunti originali: definizioni di stabilit\`a uniforme e funzioni di classe \(\mathcal K\).}
\end{figure}
}{}

\section{Teorema di Lyapunov per sistemi non stazionari}

Passiamo ora al corrispondente criterio di Lyapunov per i sistemi tempo-varianti.

Consideriamo
\[
\dot x = f(x,t),
\qquad
f(0,t)=0 \ \ \forall t\ge 0.
\]
Sia
\[
V:[0,+\infty)\times D \to \mathbb R
\]
una funzione continuamente differenziabile.

Supponiamo che esistano due funzioni positive definite \(W_1,W_2\) tali che
\begin{equation}
\label{eq:L25-V-bounds}
W_1(x)\le V(t,x)\le W_2(x)
\qquad \forall t\ge 0,\ \forall x\in D,
\end{equation}
e che inoltre
\begin{equation}
\label{eq:L25-Vdot-timevarying}
\frac{\partial V}{\partial t}(t,x)
+
\frac{\partial V}{\partial x}(t,x)\,f(x,t)
\le 0.
\end{equation}

Allora l'origine \`e \textbf{uniformemente stabile}.

\medskip

\noindent
\textbf{Commento.}
La differenza rispetto al caso autonomo \`e che adesso la derivata totale lungo le traiettorie diventa
\[
\frac{d}{dt}V(t,x(t))
=
\frac{\partial V}{\partial t}(t,x(t))
+
\frac{\partial V}{\partial x}(t,x(t))\,\dot x(t),
\]
perch\'e \(V\) dipende esplicitamente anche dal tempo.

Se invece vale la stima pi\`u forte
\[
\frac{\partial V}{\partial t}(t,x)
+
\frac{\partial V}{\partial x}(t,x)\,f(x,t)
\le -W_3(x),
\]
con \(W_3\) positiva definita, allora si ottiene una forma di stabilit\`a pi\`u forte; nel caso lineare, addirittura stabilit\`a uniforme esponenziale.

\IfFileExists{25_page-0003.jpg}{
\begin{figure}[h!]
\centering
\includegraphics[width=.82\textwidth]{25_page-0003.jpg}
\caption{Appunti originali: teorema di Lyapunov per sistemi non stazionari e interpretazione geometrica tramite insiemi di livello.}
\end{figure}
}{}

\section{Esempio: funzione di Lyapunov che dipende dal tempo}

Consideriamo il sistema
\begin{equation}
\label{eq:L25-esempio-g}
\begin{cases}
\dot x_1 = -x_1 - g(t)x_2,\\[4pt]
\dot x_2 = x_1 - x_2,
\end{cases}
\end{equation}
dove \(g(t)\) \`e continua, derivabile e soddisfa
\[
0\le g(t)\le k,
\qquad
\dot g(t)\le g(t)
\qquad \forall t\ge 0.
\]

Vogliamo mostrare che l'origine \`e uniformemente stabile.

\subsection{Scelta della candidata}

Una candidata naturale \`e
\[
V(t,x)=x_1^2 + (1+g(t))x_2^2.
\]
Poich\'e \(0\le g(t)\le k\), otteniamo le stime
\[
x_1^2+x_2^2 \le V(t,x)\le x_1^2+(1+k)x_2^2.
\]
Dunque possiamo porre
\[
W_1(x)=x_1^2+x_2^2,
\qquad
W_2(x)=x_1^2+(1+k)x_2^2.
\]

\subsection{Derivata lungo le traiettorie}

Calcoliamo
\[
\frac{d}{dt}V(t,x(t))
=
\frac{\partial V}{\partial t}
+
\frac{\partial V}{\partial x}\dot x.
\]
Si ha
\[
\frac{\partial V}{\partial t}=\dot g(t)x_2^2,
\qquad
\frac{\partial V}{\partial x}=(2x_1,\ 2(1+g(t))x_2).
\]
Quindi
\begin{align*}
\dot V
&=
\dot g(t)x_2^2
+2x_1(-x_1-gx_2)
+2(1+g)x_2(x_1-x_2)\\
&=
\dot g\,x_2^2 -2x_1^2 -2g x_1x_2 +2x_1x_2 +2g x_1x_2 -2x_2^2 -2g x_2^2\\
&=
-2x_1^2 +2x_1x_2 -(2+2g-\dot g)x_2^2.
\end{align*}
Poich\'e \(\dot g \le g\), segue
\[
2+2g-\dot g \ge 2+g \ge 1,
\]
e dunque
\[
\dot V \le -2x_1^2 + 2x_1x_2 - x_2^2.
\]
Ma
\[
-2x_1^2 + 2x_1x_2 - x_2^2
=
-(x_1^2 + (x_1-x_2)^2)\le 0.
\]
Quindi \(\dot V \le 0\) e, per il teorema precedente, l'origine \`e uniformemente stabile.

\medskip

\noindent
\textbf{Osservazione.}
Qui non basta una funzione \(V(x)\) dipendente solo dallo stato: la scelta giusta \`e una funzione \(V(t,x)\) che incorpora esplicitamente il coefficiente tempo-variante \(g(t)\).

\IfFileExists{25_page-0004.jpg}{
\begin{figure}[h!]
\centering
\includegraphics[width=.72\textwidth]{25_page-0004.jpg}
\caption{Calcolo svolto a mano per l'esempio con \(g(t)\).}
\end{figure}
}{}

\section{Caso lineare non stazionario e equazione differenziale di Lyapunov}

Consideriamo ora il caso lineare
\[
\dot x = A(t)x,
\]
con \(A(t)\) continua.

Supponiamo che esista una matrice simmetrica \(P(t)\), continua e derivabile, tale che
\begin{equation}
\label{eq:L25-P-bounds}
0<c_1 I \le P(t)\le c_2 I
\qquad \forall t\ge t_0,
\end{equation}
e che soddisfi l'equazione differenziale di Lyapunov
\begin{equation}
\label{eq:L25-lyap-diff}
-\dot P(t)=P(t)A(t)+A^\top(t)P(t)+Q(t),
\end{equation}
dove \(Q(t)=Q^\top(t)\) \`e continua e
\begin{equation}
\label{eq:L25-Q-lower}
Q(t)\ge c_3 I>0.
\end{equation}

Poniamo
\[
V(t,x)=x^\top P(t)x.
\]
Dalle stime \eqref{eq:L25-P-bounds} segue
\[
c_1\|x\|^2 \le V(t,x)\le c_2\|x\|^2.
\]
Inoltre
\begin{align*}
\dot V
&=
x^\top \dot P x + x^\top P \dot x + \dot x^\top P x\\
&=
x^\top\bigl(\dot P + PA + A^\top P\bigr)x\\
&=
-x^\top Q(t)x\\
&\le -c_3\|x\|^2.
\end{align*}
Usando \(V\le c_2\|x\|^2\), otteniamo
\[
\dot V \le -\frac{c_3}{c_2}V.
\]
Per confronto scalare,
\[
V(t)\le e^{-\frac{c_3}{c_2}(t-t_0)}V(t_0),
\]
e quindi
\[
\|x(t)\|
\le
\sqrt{\frac{c_2}{c_1}}
\,e^{-\frac{c_3}{2c_2}(t-t_0)}
\|x(t_0)\|.
\]

\medskip

\noindent
\textbf{Conclusione.}
L'origine \`e \textbf{uniformemente esponenzialmente stabile}.

\section{Richiami: LaSalle--Yoshizawa e lemma di Barb\u{a}lat}

Questi risultati vengono spesso usati per rafforzare i criteri di Lyapunov nei sistemi tempo-varianti.

\subsection{Principio di LaSalle--Yoshizawa}

Se esiste una funzione \(V(t,x)\) tale che
\[
\alpha_1(\|x\|)\le V(t,x)\le \alpha_2(\|x\|),
\]
e
\[
\frac{\partial V}{\partial t} + \frac{\partial V}{\partial x}f(x,t)\le -W(x)\le 0,
\]
con \(W\) semidefinita positiva, allora le traiettorie restano limitate e
\[
W(x(t))\to 0 \qquad \text{per } t\to+\infty.
\]
Se si riesce poi a capire quali traiettorie possono restare nell'insieme
\[
S=\{x:W(x)=0\},
\]
si pu\`o concludere verso quale insieme tende la soluzione.

\subsection{Lemma di Barb\u{a}lat}

Una forma pratica molto usata \`e la seguente: se una funzione \(r(t)\) \`e uniformemente continua e il suo integrale converge, allora
\[
r(t)\to 0
\qquad \text{per } t\to+\infty.
\]

In applicazioni Lyapunov, spesso si prende
\[
r(t)=\dot V(t,x(t)),
\]
si mostra che \(\dot V\le 0\), che \(V\) \`e inferiormente limitata e che \(\dot V\) \`e uniformemente continua; da qui si conclude
\[
\dot V(t)\to 0.
\]

\IfFileExists{25_page-0005.jpg}{
\begin{figure}[h!]
\centering
\includegraphics[width=.82\textwidth]{25_page-0005.jpg}
\caption{Appunti originali: richiamo a LaSalle--Yoshizawa e al lemma di Barb\u{a}lat.}
\end{figure}
}{}

\section{Esercizio d'esame: 13/04/2021}
\label{sec:L25-esercizio-130421}

Consideriamo il sistema
\begin{equation}
\label{eq:L25-ex-130421}
\begin{cases}
\dot x_1 = x_1^2 - x_1x_2,\\[4pt]
\dot x_2 = -x_2^2 + a x_1x_2,
\end{cases}
\qquad a\in\mathbb R.
\end{equation}

\subsection{Punto 1: equilibri e metodo indiretto}

Gli equilibri si ottengono imponendo
\[
\begin{cases}
0 = x_1^2 - x_1x_2 = x_1(x_1-x_2),\\[4pt]
0 = -x_2^2 + a x_1x_2 = x_2(ax_1-x_2).
\end{cases}
\]

Dalla prima equazione:
\[
x_1=0
\qquad \text{oppure} \qquad
x_1=x_2.
\]

\paragraph{Caso A: \(x_1=0\).}
La seconda equazione diventa
\[
0=-x_2^2,
\]
quindi necessariamente \(x_2=0\).  
Otteniamo l'equilibrio
\[
(0,0).
\]

\paragraph{Caso B: \(x_1=x_2\).}
La seconda equazione diventa
\[
0=(a-1)x_2^2.
\]
Quindi:
\begin{itemize}
    \item se \(a\neq 1\), allora \(x_2=0\) e si torna a \((0,0)\);
    \item se \(a=1\), allora ogni punto della retta
    \[
    x_1=x_2
    \]
    \`e equilibrio.
\end{itemize}

\medskip

\noindent
\textbf{Conclusione sugli equilibri.}
\begin{itemize}
    \item Se \(a\neq 1\), l'unico equilibrio \`e \((0,0)\).
    \item Se \(a=1\), c'\`e una retta di punti di equilibrio:
    \[
    \{(s,s):s\in\mathbb R\}.
    \]
\end{itemize}

\subsection{Linearizzazione}

La matrice jacobiana \`e
\[
A(x)=
\begin{pmatrix}
2x_1-x_2 & -x_1\\
a x_2 & -2x_2 + a x_1
\end{pmatrix}.
\]

\paragraph{All'origine.}
\[
A(0,0)=
\begin{pmatrix}
0 & 0\\
0 & 0
\end{pmatrix}.
\]
Tutti gli autovalori sono nulli, quindi il metodo indiretto non conclude nulla.

\paragraph{Sulla retta di equilibrio, quando \(a=1\).}
Per \(x_1=x_2=s\),
\[
A(s,s)=
\begin{pmatrix}
s & -s\\
s & -s
\end{pmatrix}.
\]
Si ha
\[
\operatorname{tr}A=0,
\qquad
\det A=0,
\]
quindi gli autovalori sono ancora entrambi nulli. Anche qui il metodo indiretto \`e inconcludente.

\subsection{Punto 2: si pu\`o stabilizzare con un solo ingresso?}

Supponiamo di aggiungere un solo ingresso in una sola equazione. Ad esempio:
\[
\begin{cases}
\dot x_1 = x_1^2 - x_1x_2 + u,\\[4pt]
\dot x_2 = -x_2^2 + a x_1x_2,
\end{cases}
\]
oppure, analogamente, nella seconda.

Se linearizziamo all'origine, otteniamo sempre
\[
A(0,0)=0.
\]
Nel primo caso
\[
B=
\begin{pmatrix}
1\\
0
\end{pmatrix},
\]
nel secondo
\[
B=
\begin{pmatrix}
0\\
1
\end{pmatrix}.
\]
In entrambi i casi la matrice di raggiungibilit\`a \`e
\[
\mathcal R=[B\ AB]=[B\ 0],
\]
e quindi ha rango \(1<2\).  
Dunque la coppia \((A,B)\) non \`e completamente raggiungibile e il sistema non \`e completamente stabilizzabile con un solo ingresso locale all'origine.

\subsection{Punto 3: studio dell'instabilit\`a con Chetaev nel caso \(a=1\)}

Nel caso \(a=1\), proviamo la candidata
\[
V(x)=x_1^2-x_2^2.
\]
Questa funzione \emph{non} \`e definita positiva, quindi non serve per dimostrare stabilit\`a.  
Pu\`o per\`o essere usata in senso ``inverso'' per dimostrare instabilit\`a.

Calcoliamo la derivata lungo le traiettorie:
\begin{align*}
L_fV
&= 2x_1(x_1^2-x_1x_2) -2x_2(-x_2^2+x_1x_2)\\
&=2x_1^3-2x_1^2x_2+2x_2^3-2x_1x_2^2\\
&=2(x_1-x_2)^2(x_1+x_2).
\end{align*}

Consideriamo la regione
\[
U=\left\{x:\ x_1^2+x_2^2<\rho^2,\ \ x_1^2-x_2^2>0,\ \ x_1+x_2>0\right\}.
\]
In \(U\):
\begin{itemize}
    \item \(V(x)>0\),
    \item \(L_fV(x)>0\),
    \item l'origine sta sul bordo di \(U\).
\end{itemize}
Per il teorema di Chetaev, l'origine \`e instabile.

\medskip

\noindent
\textbf{Osservazione importante.}
Questo \`e coerente con il fatto che per \(a=1\) l'origine non \`e un equilibrio isolato, ma appartiene a una retta di equilibri.

\IfFileExists{25_page-0006.jpg}{
\begin{figure}[h!]
\centering
\includegraphics[width=.76\textwidth]{25_page-0006.jpg}
\caption{Appunti originali sull'esercizio del 13/04/2021: equilibri, linearizzazione e commenti sulla stabilizzabilit\`a.}
\end{figure}
}{}

\IfFileExists{25_page-0007.jpg}{
\begin{figure}[h!]
\centering
\includegraphics[width=.70\textwidth]{25_page-0007.jpg}
\caption{Appunti originali: uso di una candidata non definita in segno per dimostrare instabilit\`a alla Chetaev.}
\end{figure}
}{}

\section{Esercizio d'esame: 07/04/2022}
\label{sec:L25-esercizio-070422}

Consideriamo il sistema discreto
\begin{equation}
\label{eq:L25-ex-070422}
\begin{cases}
x_1(k+1)= -\dfrac13 x_1(k)+\alpha x_1(k)^2 + 2u(k)x_2(k),\\[8pt]
x_2(k+1)= 2x_1(k)x_2(k)+\dfrac13 x_2(k).
\end{cases}
\end{equation}

\subsection{Punto 1: equilibri del sistema con \(u=0\)}

Poniamo \(u=0\). Gli equilibri \((\bar x_1,\bar x_2)\) soddisfano
\[
\begin{cases}
\bar x_1 = -\dfrac13 \bar x_1 + \alpha \bar x_1^2,\\[8pt]
\bar x_2 = 2\bar x_1\bar x_2 + \dfrac13 \bar x_2.
\end{cases}
\]

La prima equazione diventa
\[
\frac43 \bar x_1 - \alpha \bar x_1^2 = 0
\quad\Longleftrightarrow\quad
\bar x_1\left(\frac43 - \alpha \bar x_1\right)=0.
\]
Quindi
\[
\bar x_1=0
\qquad \text{oppure} \qquad
\bar x_1=\frac{4}{3\alpha}\quad (\alpha\neq 0).
\]

La seconda diventa
\[
\bar x_2 - 2\bar x_1\bar x_2 - \frac13 \bar x_2 = 0
\quad\Longleftrightarrow\quad
\bar x_2\left(\frac23 - 2\bar x_1\right)=0,
\]
cio\`e
\[
\bar x_2=0
\qquad \text{oppure} \qquad
\bar x_1=\frac13.
\]

\paragraph{Combinando i casi:}
\begin{itemize}
    \item \((0,0)\) \`e equilibrio per ogni \(\alpha\in\mathbb R\);
    \item se \(\alpha\neq 0\), esiste anche l'equilibrio
    \[
    \left(\frac{4}{3\alpha},\,0\right);
    \]
    \item se \(\alpha=4\), allora \(\bar x_1=\frac13\) e quindi tutti i punti
    \[
    \left(\frac13,\bar x_2\right),\qquad \bar x_2\in\mathbb R,
    \]
    sono equilibri.
\end{itemize}

\subsection{Stabilit\`a locale degli equilibri}

Con \(u=0\), il jacobiano rispetto allo stato \`e
\[
A(x)=
\begin{pmatrix}
-\dfrac13 + 2\alpha x_1 & 0\\[8pt]
2x_2 & 2x_1+\dfrac13
\end{pmatrix}.
\]

\paragraph{All'origine.}
\[
A(0,0)=
\begin{pmatrix}
-\dfrac13 & 0\\[6pt]
0 & \dfrac13
\end{pmatrix}.
\]
Attenzione: qui bisogna essere coerenti con la forma del sistema. Se la seconda equazione \`e davvero
\[
x_2(k+1)=2x_1(k)x_2(k)+\frac13 x_2(k),
\]
allora gli autovalori all'origine sono
\[
\lambda_1=-\frac13,
\qquad
\lambda_2=\frac13.
\]
Poich\'e entrambi hanno modulo minore di \(1\), l'origine \`e localmente asintoticamente stabile.

\paragraph{All'equilibrio \(\left(\frac{4}{3\alpha},0\right)\), \(\alpha\neq 0\).}
Si ottiene
\[
A\!\left(\frac{4}{3\alpha},0\right)=
\begin{pmatrix}
\frac73 & 0\\[6pt]
0 & \frac{8+\alpha}{3\alpha}
\end{pmatrix}.
\]
Il primo autovalore vale sempre
\[
\lambda_1=\frac73,
\]
quindi ha modulo maggiore di \(1\). Ne segue che tale equilibrio \`e sempre instabile.

\paragraph{Sulla retta \(\left(\frac13,\bar x_2\right)\), quando \(\alpha=4\).}
In questo caso
\[
A\!\left(\frac13,\bar x_2\right)=
\begin{pmatrix}
\frac73 & 0\\[6pt]
2\bar x_2 & 1
\end{pmatrix}.
\]
Gli autovalori sono
\[
\lambda_1=\frac73,
\qquad
\lambda_2=1.
\]
Anche qui l'equilibrio non \`e asintoticamente stabile; in effetti \`e instabile perch\'e \(|\lambda_1|>1\).

\subsection{Punto 2: raggiungibilit\`a del linearizzato}

Il vettore d'ingresso \`e
\[
B(x)=
\begin{pmatrix}
\frac{\partial f_1}{\partial u}\\[4pt]
\frac{\partial f_2}{\partial u}
\end{pmatrix}
=
\begin{pmatrix}
2x_2\\
0
\end{pmatrix}.
\]

\paragraph{All'origine.}
\[
B(0,0)=
\begin{pmatrix}
0\\
0
\end{pmatrix},
\]
quindi il linearizzato non \`e controllabile.

\paragraph{All'equilibrio \(\left(\frac{4}{3\alpha},0\right)\).}
Anche qui
\[
B\!\left(\frac{4}{3\alpha},0\right)=
\begin{pmatrix}
0\\
0
\end{pmatrix},
\]
quindi ancora nessuna controllabilit\`a.

\paragraph{Sulla retta \(\left(\frac13,\bar x_2\right)\), con \(\alpha=4\).}
Se \(\bar x_2\neq 0\),
\[
B\!\left(\frac13,\bar x_2\right)=
\begin{pmatrix}
2\bar x_2\\
0
\end{pmatrix}.
\]
La matrice di raggiungibilit\`a \`e
\[
\mathcal R=
\begin{bmatrix}
B & AB
\end{bmatrix}
=
\begin{bmatrix}
2\bar x_2 & \dfrac{14}{3}\bar x_2\\[6pt]
0 & 4\bar x_2^2
\end{bmatrix}.
\]
Se \(\bar x_2\neq 0\), allora
\[
\operatorname{rank}\mathcal R = 2.
\]
Quindi il linearizzato \`e completamente raggiungibile e si pu\`o progettare un controllo locale stabilizzante.

\subsection{Forma di un controllo affine locale}

Se \((\bar x,\bar u)\) \`e il punto di equilibrio scelto, un controllo affine locale si scrive come
\[
u = K(x-\bar x)+\bar u.
\]
Esplicitamente, in dimensione \(2\),
\[
u = k_1(x_1-\bar x_1)+k_2(x_2-\bar x_2)+\bar u.
\]
Equivalentemente,
\[
u = k_1x_1+k_2x_2+k_3,
\]
dove \(k_3\) \`e scelto in modo da non spostare l'equilibrio.

\medskip

\noindent
Il progetto si fa come sempre sul sistema delle variazioni
\[
z(k+1)=Az(k)+Bv(k),
\qquad
v(k)=Kz(k),
\]
scegliendo \(K\) in modo che gli autovalori di \(A+BK\) cadano dove desiderato.

\IfFileExists{25_page-0008.jpg}{
\begin{figure}[h!]
\centering
\includegraphics[width=.76\textwidth]{25_page-0008.jpg}
\caption{Appunti originali sull'esercizio del 07/04/2022: equilibri e studio del linearizzato.}
\end{figure}
}{}

\IfFileExists{25_page-0009.jpg}{
\begin{figure}[h!]
\centering
\includegraphics[width=.76\textwidth]{25_page-0009.jpg}
\caption{Appunti originali: matrice di ingresso, raggiungibilit\`a e forma del controllo affine.}
\end{figure}
}{}

\section{Schema riassuntivo della lezione}

La lezione contiene quattro idee fondamentali.

\subsection*{1. Nei sistemi non stazionari la stabilit\`a dipende anche da \(t_0\)}

Per questo introduciamo:
\begin{itemize}
    \item stabilit\`a,
    \item stabilit\`a uniforme,
    \item funzioni di classe \(\mathcal K\).
\end{itemize}

\subsection*{2. Lyapunov si estende ai sistemi tempo-varianti}

La funzione candidata pu\`o dipendere da \(t\):
\[
V=V(t,x).
\]
La derivata corretta lungo le traiettorie \`e
\[
\frac{d}{dt}V(t,x(t))
=
\frac{\partial V}{\partial t}
+
\frac{\partial V}{\partial x}f(x,t).
\]

\subsection*{3. Nel caso lineare tempo-variante si usa una matrice \(P(t)\)}

Si ottiene l'equazione differenziale di Lyapunov
\[
-\dot P = PA + A^\top P + Q,
\]
che gioca lo stesso ruolo dell'equazione algebrica nel caso stazionario.

\subsection*{4. Gli esercizi d'esame combinano sempre le stesse idee}

In pratica, nei compiti bisogna saper fare bene queste quattro mosse:
\begin{enumerate}
    \item trovare gli equilibri;
    \item linearizzare correttamente;
    \item capire quando il metodo indiretto conclude e quando no;
    \item se serve, usare Chetaev, Lyapunov diretto o un controllo affine locale.
\end{enumerate}

\medskip

\noindent
\textbf{Frase da ricordare all'orale.}
\begin{center}
\emph{Per i sistemi tempo-varianti il criterio spettrale istantaneo non basta. La stabilit\`a si studia con strumenti pi\`u robusti, in particolare con funzioni di Lyapunov dipendenti dal tempo e con stime uniformi rispetto all'istante iniziale.}
\end{center}

\end{document}